\documentclass[numbers=enddot,12pt,final,onecolumn,notitlepage]{scrartcl}%
\usepackage[headsepline,footsepline,manualmark]{scrlayer-scrpage}
\usepackage[all,cmtip]{xy}
\usepackage{amssymb}
\usepackage{framed}
\usepackage{amsmath}
\usepackage{comment}
\usepackage{color}
\usepackage{hyperref}
\usepackage[sc]{mathpazo}
\usepackage[T1]{fontenc}
\usepackage{amsthm}
%TCIDATA{OutputFilter=latex2.dll}
%TCIDATA{Version=5.50.0.2960}
%TCIDATA{LastRevised=Tuesday, May 26, 2026 02:11:50}
%TCIDATA{SuppressPackageManagement}
%TCIDATA{<META NAME="GraphicsSave" CONTENT="32">}
%TCIDATA{<META NAME="SaveForMode" CONTENT="1">}
%TCIDATA{BibliographyScheme=Manual}
%BeginMSIPreambleData
\providecommand{\U}[1]{\protect\rule{.1in}{.1in}}
%EndMSIPreambleData
\newenvironment{statement}{\begin{quote}}{\end{quote}}
\theoremstyle{definition}
\newtheorem{theo}{Theorem}[section]
\newenvironment{theorem}[1][]
{\begin{theo}[#1]\begin{leftbar}}
{\end{leftbar}\end{theo}}
\newtheorem{lem}[theo]{Lemma}
\newenvironment{lemma}[1][]
{\begin{lem}[#1]\begin{leftbar}}
{\end{leftbar}\end{lem}}
\newenvironment{lemmain}[1][]
{\begin{statement}\begin{lem}[#1]}
{\end{lem}\end{statement}}
\newtheorem{prop}[theo]{Proposition}
\newenvironment{proposition}[1][]
{\begin{prop}[#1]\begin{leftbar}}
{\end{leftbar}\end{prop}}
\newtheorem{defi}[theo]{Definition}
\newenvironment{definition}[1][]
{\begin{defi}[#1]\begin{leftbar}}
{\end{leftbar}\end{defi}}
\newtheorem{remk}[theo]{Remark}
\newenvironment{remark}[1][]
{\begin{remk}[#1]\begin{leftbar}}
{\end{leftbar}\end{remk}}
\newtheorem{coro}[theo]{Corollary}
\newenvironment{corollary}[1][]
{\begin{coro}[#1]\begin{leftbar}}
{\end{leftbar}\end{coro}}
\newtheorem{conv}[theo]{Convention}
\newenvironment{condition}[1][]
{\begin{conv}[#1]\begin{leftbar}}
{\end{leftbar}\end{conv}}
\newtheorem{quest}[theo]{Question}
\newenvironment{algorithm}[1][]
{\begin{quest}[#1]\begin{leftbar}}
{\end{leftbar}\end{quest}}
\newtheorem{warn}[theo]{Warning}
\newenvironment{conclusion}[1][]
{\begin{warn}[#1]\begin{leftbar}}
{\end{leftbar}\end{warn}}
\newtheorem{conj}[theo]{Conjecture}
\newenvironment{conjecture}[1][]
{\begin{conj}[#1]\begin{leftbar}}
{\end{leftbar}\end{conj}}
\newtheorem{exmp}[theo]{Example}
\newenvironment{example}[1][]
{\begin{exmp}[#1]\begin{leftbar}}
{\end{leftbar}\end{exmp}}
\iffalse
\newenvironment{proof}[1][Proof]{\noindent\textbf{#1.} }{\ \rule{0.5em}{0.5em}}
\newenvironment{lemmain}[1][Proof]{\noindent\textbf{#1.} }{\ \rule{0.5em}{0.5em}}
\fi
\newenvironment{verlong}{}{}
\newenvironment{vershort}{}{}
\newenvironment{noncompile}{}{}
\excludecomment{verlong}
\includecomment{vershort}
\excludecomment{noncompile}
\newcommand{\kk}{\mathbf{k}}
\newcommand{\id}{\operatorname{id}}
\newcommand{\ev}{\operatorname{ev}}
\newcommand{\Comp}{\operatorname{Comp}}
\newcommand{\bk}{\mathbf{k}}
\newcommand{\Nplus}{\mathbb{N}_{+}}
\newcommand{\NN}{\mathbb{N}}
\let\sumnonlimits\sum
\let\prodnonlimits\prod
\renewcommand{\sum}{\sumnonlimits\limits}
\renewcommand{\prod}{\prodnonlimits\limits}
\renewcommand{\theequation}{\thesection.\arabic{equation}}
\setlength\textheight{22.5cm}
\setlength\textwidth{15cm}
\ihead{Errata to ``Schur functions: Theme and variations''}
\ohead{\today}
\begin{document}

\begin{center}
\textbf{Schur functions: Theme and variations}

\textit{I. G. Macdonald}

\href{https://www.mat.univie.ac.at/~slc/opapers/s28macdonald.html}{Publ.
I.R.M.A. Strasbourg, 1992, 498/S--27,}

\href{https://www.mat.univie.ac.at/~slc/opapers/s28macdonald.html}{Actes
28$^{\text{e}}$\ S\'{e}minaire Lotharingien, p. 5-39.}

\textbf{Errata and addenda by Darij Grinberg}

\bigskip
\end{center}

I will refer to the results appearing in the paper \textquotedblleft Schur
functions: Theme and variations\textquotedblright\ by the numbers under which
they appear in this paper. Page numbering goes from 5 to 39.

I have read the \textquotedblleft6th Variation\textquotedblright,
\textquotedblleft7th Variation\textquotedblright\ and \textquotedblleft9th
Variation\textquotedblright\ sections of the paper, with the exception of the
Frobenius-notation parts (i.e., (6.11), the analogous formulas for
$S_{\lambda}\left(  V\right)  $, and the 9th Variation from (9.7) onwards) and
of the conjectural formulas (7.24 ?) and (7.25 ?).

The list below contains both actual corrections (of what few typos and minor
errors exist in the paper) and clarifications (mostly elaborations of details
that have been omitted from Macdonald's notoriously terse proofs), along with
the occasional comment on related work. The clarifications, in fact, take up
the major part of the list.

\appendix \setcounter{section}{1}

\section{Errata and addenda}

\subsection{Notations}

I shall use the following notations:

\begin{itemize}
\item If $p$ and $q$ are two integers such that $p\leq q+1$, then $\left[
p,q\right]  $ shall denote the set of all integers $m$ such that $p\leq m\leq
q$. We call this set an \emph{integer interval}. It has size $q-p+1$ (so it is
empty if $p=q+1$).

\item For a given $r\in\mathbb{N}$, we shall denote the integer interval
$\left[  1,r\right]  =\left\{  1,2,\ldots,r\right\}  $ by $\left[  r\right]  $.

\item If $A=\left(  a_{i,j}\right)  _{i\in S,\ j\in T}$ is an arbitrary matrix
(where the sets $S$ and $T$ may be finite or infinite, but must be totally
ordered), and if $P=\left\{  p_{1}<p_{2}<\cdots<p_{\alpha}\right\}  $ is a
finite subset of $S$, and if $Q=\left\{  q_{1}<q_{2}<\cdots<q_{\beta}\right\}
$ is a finite subset of $T$, then $\operatorname*{sub}\nolimits_{P}^{Q}A$
shall denote the submatrix $\left(  a_{p_{i},q_{j}}\right)  _{i\in\left[
\alpha\right]  ,\ j\in\left[  \beta\right]  }$ of $A$. For instance, if
$A=\left(  a_{i,j}\right)  _{i,j\in\left[  4\right]  }$ is a $4\times
4$-matrix, then $\operatorname*{sub}\nolimits_{\left\{  2,3\right\}
}^{\left\{  2,4\right\}  }A=\left(
\begin{array}
[c]{cc}%
a_{2,2} & a_{2,4}\\
a_{3,2} & a_{3,4}%
\end{array}
\right)  $.

\item The length (i.e., the number of nonzero entries) of a partition
$\lambda$ will be denoted by $\ell\left(  \lambda\right)  $. (Macdonald writes
$l\left(  \lambda\right)  $ instead.)

\item If $\lambda=\left(  \lambda_{1},\lambda_{2},\lambda_{3},\ldots\right)  $
and $\mu=\left(  \mu_{1},\mu_{2},\mu_{3},\ldots\right)  $ are two partitions,
then the notation \textquotedblleft$\lambda\subseteq\mu$\textquotedblright%
\ shall mean that the Young diagram of $\lambda$ is a subset of the Young
diagram of $\mu$, or, equivalently, that each $i\geq1$ satisfies $\lambda
_{i}\leq\mu_{i}$. This is denoted by \textquotedblleft$\lambda\subset\mu
$\textquotedblright\ in Macdonald's paper, but I shall use the notation
\textquotedblleft$\lambda\subseteq\mu$\textquotedblright\ instead, since it
corresponds better to my use of \textquotedblleft$\subseteq$\textquotedblright%
\ for subsets. (So the symbols \textquotedblleft$\subset$\textquotedblright%
\ and \textquotedblleft$\subseteq$\textquotedblright\ are synonymous as far as
partitions are concerned.)

\item The partition $0=\left(  0,0,0,\ldots\right)  $ will be denoted by
$\varnothing$.
\end{itemize}

Now, the actual corrections:

\subsection{Corrections to the Introduction}

\begin{enumerate}
\item The \emph{Frobenius notation} for partitions is defined as follows: If
$\lambda=\left(  \lambda_{1},\lambda_{2},\lambda_{3},\ldots\right)  $ is a
partition with conjugate partition $\lambda^{\prime}=\left(  \lambda
_{1}^{\prime},\lambda_{2}^{\prime},\lambda_{3}^{\prime},\ldots\right)  $, then
the \emph{Frobenius notation} for $\lambda$ is the $2r$-tuple $\left(
\alpha_{1},\alpha_{2},\ldots,\alpha_{r}\mid\beta_{1},\beta_{2},\ldots
,\beta_{r}\right)  $ (the bar in the middle separates the first $r$ from the
last $r$ entries), where

\begin{itemize}
\item the number $r$ is the largest integer $i\geq1$ such that the cell
$\left(  i,i\right)  $ lies in the Young diagram of $\lambda$ (or $0$ if no
such $i$ exists);

\item the numbers $\alpha_{i}$ for $i\in\left[  r\right]  $ are defined by
$\alpha_{i}:=\lambda_{i}-i$;

\item the numbers $\beta_{i}$ for $i\in\left[  r\right]  $ are defined by
$\beta_{i}:=\lambda_{i}^{\prime}-i$.
\end{itemize}

For instance, the Frobenius notation for the partition $\left(
4,4,2,1\right)  $ is $\left(  3,2\mid3,1\right)  $.
\end{enumerate}

\subsection{Corrections to the 6th Variation}

\begin{enumerate}
\item \textbf{page 15, line 2:} \textquotedblleft for each $\alpha=\left(
\alpha_{1},\ldots,\alpha_{r}\right)  \in\mathbb{N}^{n}$\textquotedblright%
\ should be \textquotedblleft for each $\alpha=\left(  \alpha_{1}%
,\ldots,\alpha_{n}\right)  \in\mathbb{N}^{n}$\textquotedblright.

\item \textbf{page 15, after (6.3):} For a detailed proof of the fact that the
quotient
\[
s_{\lambda}\left(  x\mid a\right)  =A_{\lambda+\delta}\left(  x\mid a\right)
\diagup A_{\delta}\left(  x\mid a\right)
\]
is a symmetric polynomial in $x_{1},\ldots,x_{n}$ with coefficients in $R$,
see \cite[Corollary 9.14]{regpol}. (Apply this corollary to $X_{i}=x_{i}$ and
$P_{j}\left(  T\right)  =\left(  T\mid a\right)  ^{\lambda_{j}+n-j}$.)

\item \textbf{page 15, definition of }$h_{r}\left(  x\mid a\right)
$\textbf{:} It should also be said that $h_{r}\left(  x\mid a\right)  $ is
defined to be $0$ when $n=0$ and $r>0$. (In this case, $s_{\left(  r\right)
}\left(  x\mid a\right)  $ is undefined, since $\left(  r\right)  $ is not a
partition of length $\leq n$, although it would make perfect sense to just
define $s_{\lambda}\left(  x\mid a\right)  $ to be $0$ when $\lambda$ is a
partition of length $>n$.)

\item \textbf{page 16, proof of (6.6):} This argument only shows that
$\mathbf{E}\left(  x\mid a\right)  \cdot\mathbf{H}\left(  x\mid a\right)
=I_{\mathbb{Z}}$ (where $I_{\mathbb{Z}}$ is the identity matrix with rows and
columns indexed by all integers). In order to prove that the two matrices
$\mathbf{E}\left(  x\mid a\right)  $ and $\mathbf{H}\left(  x\mid a\right)  $
are inverse to each other (i.e., in order to prove (6.6)), it must also be
shown that $\mathbf{H}\left(  x\mid a\right)  \cdot\mathbf{E}\left(  x\mid
a\right)  =I_{\mathbb{Z}}$.

Fortunately, there is a simple shortcut for this: Let $\operatorname*{UT}%
\nolimits_{\mathbb{Z}}$ be the set of all upper unitriangular matrices with
rows and columns indexed by all integers (and with entries in a given base
ring, which in our case is the polynomial ring over $\mathbb{Z}$ in the
variables $x_{i}$ for $i\in\left[  n\right]  $ and $a_{j}$ for $j\in
\mathbb{Z}$). This set $\operatorname*{UT}\nolimits_{\mathbb{Z}}$ is closed
under matrix multiplication, thus is a monoid. Moreover, each matrix
$A\in\operatorname*{UT}\nolimits_{\mathbb{Z}}$ can be written as $I+M$ for
some strictly upper-triangular matrix $M$, and thus has an inverse
$A^{-1}=\left(  I+M\right)  ^{-1}$, which can be computed by the formula
$\left(  I+M\right)  ^{-1}=I-M+M^{2}-M^{3}+M^{4}\pm\cdots$ (this infinite sum
makes sense, since the nonzero entries of each power $M^{i}$ start no earlier
than $i$ steps above the main diagonal\footnote{That is: For any
$u,v\in\mathbb{Z}$ and any $i\in\mathbb{N}$, the $\left(  u,v\right)  $-th
entry of $M^{i}$ is $0$ whenever $v-u<i$, and therefore the $\left(
u,v\right)  $-th entry of the infinite sum $I-M+M^{2}-M^{3}+M^{4}\pm\cdots$ is
only affected by the first $v-u+1$ addends of this sum.}). This inverse
$A^{-1}=I-M+M^{2}-M^{3}+M^{4}\pm\cdots$ again belongs to $\operatorname*{UT}%
\nolimits_{\mathbb{Z}}$ (since all the matrices $M,M^{2},M^{3},M^{4},\ldots$
are strictly upper-triangular). Hence, each element of the monoid
$\operatorname*{UT}\nolimits_{\mathbb{Z}}$ has an inverse. Thus,
$\operatorname*{UT}\nolimits_{\mathbb{Z}}$ is a group with respect to matrix
multiplication. Since both matrices $\mathbf{E}\left(  x\mid a\right)  $ and
$\mathbf{H}\left(  x\mid a\right)  $ belong to this group $\operatorname*{UT}%
\nolimits_{\mathbb{Z}}$, we can thus conclude $\mathbf{E}\left(  x\mid
a\right)  =\mathbf{H}\left(  x\mid a\right)  ^{-1}$ from $\mathbf{E}\left(
x\mid a\right)  \cdot\mathbf{H}\left(  x\mid a\right)  =I_{\mathbb{Z}}$.

\item \textbf{page 17, proof of (6.7):} It is worth saying that all the
matrices that appear in the proof of the first of the formulas (6.7) are
understood to be $n\times n$-matrices.

\item \textbf{page 17, proof of (6.7):} Let me give a proof of the second of
the formulas (6.7) along with the more general formula (6.9). We will need
some notations and some lemmas.

Let $\operatorname*{UT}\nolimits_{\mathbb{Z}}$ be the set of all upper
unitriangular matrices with rows and columns indexed by all integers (and with
entries in a given commutative ring). For any $m\in\mathbb{N}$, we let
$\operatorname*{UT}\nolimits_{m}$ be the set of all upper unitriangular
$m\times m$-matrices (again, with entries in our given commutative ring). Both
of these sets $\operatorname*{UT}\nolimits_{\mathbb{Z}}$ and
$\operatorname*{UT}\nolimits_{m}$ are groups (under matrix multiplication).
The following fact is easy:

\begin{lemmain}
\label{lem.A1}Let $T=\left[  p,q\right]  $ be an integer interval of size
$m=q-p+1$ (so that $q=p+m-1$). Then, the map%
\begin{align}
\operatorname*{UT}\nolimits_{\mathbb{Z}}  &  \rightarrow\operatorname*{UT}%
\nolimits_{m},\nonumber\\
A  &  \mapsto\operatorname*{sub}\nolimits_{T}^{T}A \label{p17.lemA1.map}%
\end{align}
is a group morphism.
\end{lemmain}

\begin{proof}
[Proof of Lemma \ref{lem.A1}.]It is easy to see that this map is well-defined
(i.e., that $\operatorname*{sub}\nolimits_{T}^{T}A\in\operatorname*{UT}%
\nolimits_{m}$ for each $A\in\operatorname*{UT}\nolimits_{\mathbb{Z}}$).
(Indeed, more generally, any principal submatrix of an upper unitriangular
matrix is again upper unitriangular.) It is clear that $\operatorname*{sub}%
\nolimits_{T}^{T}\left(  I_{\mathbb{Z}}\right)  =I_{m}$. It remains to show
that $\operatorname*{sub}\nolimits_{T}^{T}\left(  AB\right)
=\operatorname*{sub}\nolimits_{T}^{T}A\cdot\operatorname*{sub}\nolimits_{T}%
^{T}B$ for all $A,B\in\operatorname*{UT}\nolimits_{\mathbb{Z}}$.

So let $A,B\in\operatorname*{UT}\nolimits_{\mathbb{Z}}$ be arbitrary. Write
these matrices $A$ and $B$ as $A=\left(  a_{i,j}\right)  _{i,j\in\mathbb{Z}}$
and $B=\left(  b_{i,j}\right)  _{i,j\in\mathbb{Z}}$. Since $T=\left[
p,q\right]  =\left[  p,p+m-1\right]  $ (because $q=p+m-1$), we thus have%
\[
\operatorname*{sub}\nolimits_{T}^{T}A=\left(  a_{p+i-1,p+j-1}\right)
_{i,j\in\left[  m\right]  }\ \ \ \ \ \ \ \ \ \ \text{and}%
\ \ \ \ \ \ \ \ \ \ \operatorname*{sub}\nolimits_{T}^{T}B=\left(
b_{p+i-1,p+j-1}\right)  _{i,j\in\left[  m\right]  }%
\]
and therefore%
\begin{align}
\operatorname*{sub}\nolimits_{T}^{T}A\cdot\operatorname*{sub}\nolimits_{T}%
^{T}B  &  =\left(  a_{p+i-1,p+j-1}\right)  _{i,j\in\left[  m\right]  }%
\cdot\left(  b_{p+i-1,p+j-1}\right)  _{i,j\in\left[  m\right]  }\nonumber\\
&  =\left(  \sum_{k=1}^{m}a_{p+i-1,p+k-1}b_{p+k-1,p+j-1}\right)
_{i,j\in\left[  m\right]  } \label{p17.lemA1.pf.0}%
\end{align}
(by the definition of matrix multiplication).

The matrix $A$ is upper-triangular (since $A\in\operatorname*{UT}%
\nolimits_{\mathbb{Z}}$), so we have%
\begin{equation}
a_{i,k}=0\ \ \ \ \ \ \ \ \ \ \text{for all }i>k. \label{p17.lemA1.pf.1}%
\end{equation}
Likewise,%
\begin{equation}
b_{k,j}=0\ \ \ \ \ \ \ \ \ \ \text{for all }k>j. \label{p17.lemA1.pf.2}%
\end{equation}
Thus, we can easily see that if $i,j\in T$, then%
\begin{equation}
a_{i,k}b_{k,j}=0\ \ \ \ \ \ \ \ \ \ \text{for all integers }k\notin T.
\label{p17.lemA1.pf.3}%
\end{equation}
(\textit{Proof:} Let $i,j\in T$. Let $k$ be an integer such that $k\notin T$.
Thus, $k\notin T=\left[  p,q\right]  $. Hence, either $k<p$ or $k>q$. In the
former case, we have $k<p\leq i$ (since $i\in T=\left[  p,q\right]  $) and
therefore $i>k$, so that $a_{i,k}=0$ (by (\ref{p17.lemA1.pf.1})), whence
$\underbrace{a_{i,k}}_{=0}b_{k,j}=0$. In the latter case, we have $k>q\geq j$
(since $j\in T=\left[  p,q\right]  $ entails $j\leq q$) and therefore
$b_{k,j}=0$ (by (\ref{p17.lemA1.pf.2})), whence $a_{i,k}\underbrace{b_{k,j}%
}_{=0}=0$. Hence, we have proved $a_{i,k}b_{k,j}=0$ in both cases. Thus,
(\ref{p17.lemA1.pf.3}) is proved.)

Now, for any $i,j\in T$, we have%
\begin{align}
&  \left(  \text{the }\left(  i,j\right)  \text{-th entry of the matrix
}AB\right) \nonumber\\
&  =\sum_{k\in\mathbb{Z}}a_{i,k}b_{k,j}\ \ \ \ \ \ \ \ \ \ \left(  \text{by
the definition of the matrix product }AB\right) \nonumber\\
&  =\sum_{k\in T}a_{i,k}b_{k,j} \label{p17.lemA1.pf.5}%
\end{align}
(here, we have restricted the sum $\sum_{k\in\mathbb{Z}}a_{i,k}b_{k,j}$ to
only range over the $k\in T$, since (\ref{p17.lemA1.pf.3}) shows that any
addend $a_{i,k}b_{k,j}$ equals $0$ unless $k\in T$). However, from $T=\left[
p,q\right]  $, we obtain%
\begin{align}
\operatorname*{sub}\nolimits_{T}^{T}\left(  AB\right)   &  =\left(
\underbrace{\text{the }\left(  p+i-1,p+j-1\right)  \text{-th entry of the
matrix }AB}_{\substack{=\sum_{k\in T}a_{p+i-1,k}b_{k,p+j-1}\\\text{(by
(\ref{p17.lemA1.pf.5}))}}}\right)  _{i,j\in\left[  m\right]  }\nonumber\\
&  =\left(  \sum_{k\in T}a_{p+i-1,k}b_{k,p+j-1}\right)  _{i,j\in\left[
m\right]  }. \label{p17.lemA1.pf.6}%
\end{align}
However, each $i,j\in\left[  m\right]  $ satisfy%
\begin{align*}
&  \sum_{k=1}^{m}a_{p+i-1,p+k-1}b_{p+k-1,p+j-1}\\
&  =\underbrace{\sum_{k=p}^{p+m-1}}_{\substack{=\sum_{k\in T}\\\text{(since
}T=\left[  p,p+m-1\right]  \text{)}}}a_{p+i-1,k}b_{k,p+j-1}\\
&  \ \ \ \ \ \ \ \ \ \ \ \ \ \ \ \ \ \ \ \ \left(  \text{here, we have
substituted }k\text{ for }p+k-1\text{ in the sum}\right) \\
&  =\sum_{k\in T}a_{p+i-1,k}b_{k,p+j-1}.
\end{align*}
Thus, the right hand sides of the equalities (\ref{p17.lemA1.pf.0}) and
(\ref{p17.lemA1.pf.6}) are equal. Hence, their left hand sides are also equal.
In other words, $\operatorname*{sub}\nolimits_{T}^{T}\left(  AB\right)
=\operatorname*{sub}\nolimits_{T}^{T}A\cdot\operatorname*{sub}\nolimits_{T}%
^{T}B$. This completes the proof of Lemma \ref{lem.A1}.
\end{proof}

\begin{lemmain}
\label{lem.A2}Let $A$ be an invertible $m\times m$-matrix. Let $P$ and $Q$ be
two subsets of $\left[  m\right]  $ such that $\left\vert P\right\vert
=\left\vert Q\right\vert $. Let $\widetilde{P}:=\left[  m\right]  \setminus P$
and $\widetilde{Q}:=\left[  m\right]  \setminus Q$ be their complements. Let
$\sum P$ be the sum of all elements of $P$, and let $\sum Q$ be the sum of all
elements of $Q$. Then,%
\[
\det\left(  \operatorname*{sub}\nolimits_{P}^{Q}A\right)  =\left(  -1\right)
^{\sum P+\sum Q}\det A\cdot\det\left(  \operatorname*{sub}%
\nolimits_{\widetilde{Q}}^{\widetilde{P}}\left(  A^{-1}\right)  \right)  .
\]

\end{lemmain}

\begin{proof}
[Proof of Lemma \ref{lem.A2}.]Lemma \ref{lem.A2} is \cite[Exercise
6.56]{detnotes} (with slightly different notations: $\operatorname*{sub}%
\nolimits_{P}^{Q}A$ is denoted $\operatorname*{sub}\nolimits_{w\left(
P\right)  }^{w\left(  Q\right)  }A$ there). Alternatively, it can be easily
derived from \cite[(APP.1.5.2)]{LLPT95} (since the adjugate matrix
$\operatorname*{adj}A$ of $A$, which is denoted by $^{\dag}A$ in
\cite{LLPT95}, is known to equal $\left(  \det A\right)  \cdot A^{-1}$).
\end{proof}

\begin{lemmain}
\label{lem.A3}Let $B$ be an infinite matrix in $\operatorname*{UT}%
\nolimits_{\mathbb{Z}}$. Let $T=\left[  p,q\right]  $ be an integer interval.
Let $U$ and $V$ be two subsets of $T$ such that $\left\vert U\right\vert
=\left\vert V\right\vert $. Then,%
\[
\det\left(  \operatorname*{sub}\nolimits_{U}^{V}B\right)  =\left(  -1\right)
^{\sum U+\sum V}\det\left(  \operatorname*{sub}\nolimits_{T\setminus
V}^{T\setminus U}\left(  B^{-1}\right)  \right)  .
\]
(Note that $B^{-1}$ exists: Indeed, the matrix $B$ belongs to the group
$\operatorname*{UT}\nolimits_{\mathbb{Z}}$ and thus has an inverse.)
\end{lemmain}

\begin{proof}
[Proof of Lemma \ref{lem.A3}.]Let $m$ be the size $q-p+1$ of the interval
$T=\left[  p,q\right]  $. Then, $q=m+p-1$. Furthermore, Lemma \ref{lem.A1}
says that the map (\ref{p17.lemA1.map}) is a group morphism. Hence, $\left(
\operatorname*{sub}\nolimits_{T}^{T}B\right)  ^{-1}=\operatorname*{sub}%
\nolimits_{T}^{T}\left(  B^{-1}\right)  $. Moreover, $\operatorname*{sub}%
\nolimits_{T}^{T}B\in\operatorname*{UT}\nolimits_{m}$ (since the map
(\ref{p17.lemA1.map}) has target $\operatorname*{UT}\nolimits_{m}$), so that
the matrix $\operatorname*{sub}\nolimits_{T}^{T}B$ is upper unitriangular.
Thus, $\det\left(  \operatorname*{sub}\nolimits_{T}^{T}B\right)  =1$ (since
the determinant of a triangular matrix is the product of its diagonal entries,
and therefore the determinant of a unitriangular matrix is $1$).

Let $A$ be the submatrix $\operatorname*{sub}\nolimits_{T}^{T}B$ of $B$.
Then,
\begin{align}
A^{-1}  &  =\left(  \operatorname*{sub}\nolimits_{T}^{T}B\right)
^{-1}=\operatorname*{sub}\nolimits_{T}^{T}\left(  B^{-1}\right)
\ \ \ \ \ \ \ \ \ \ \text{and}\label{p17.lemA3.pf.A1}\\
\det A  &  =\det\left(  \operatorname*{sub}\nolimits_{T}^{T}B\right)  =1.
\label{p17.lemA3.pf.A2}%
\end{align}


For any subset $S$ of $\mathbb{Z}$ and any integer $z$, we let $S+z$ denote
the set $\left\{  s+z\ \mid\ s\in S\right\}  $. Visually speaking, this is
simply the set $S$ shifted by $z$ units to the right along the number line.
Clearly, $\left\vert S+z\right\vert =\left\vert S\right\vert $ and $\left(
S+z\right)  +\left(  -z\right)  =S$. Moreover, any two subsets $S$ and $T$ of
$\mathbb{Z}$ and any integer $z$ satisfy%
\begin{equation}
\left(  S\setminus T\right)  +z=\left(  S+z\right)  \setminus\left(
T+z\right)  , \label{p17.lemA3.pf.diff}%
\end{equation}
since the operation of adding $z$ to each integer is a bijection.

We have%
\begin{align}
T &  =\left[  p,q\right]  \nonumber\\
&  =\left[  1+p-1,\ m+p-1\right]  \ \ \ \ \ \ \ \ \ \ \left(  \text{since
}p=1+p-1\text{ and }q=m+p-1\right)  \nonumber\\
&  =\left[  1,m\right]  +\left(  p-1\right)  \nonumber\\
&  =\left[  m\right]  +\left(  p-1\right)  \ \ \ \ \ \ \ \ \ \ \left(
\text{since }\left[  1,m\right]  =\left[  m\right]  \right)
.\label{p17.lemA3.pf.T}%
\end{align}
Therefore, the map $i\mapsto i+\left(  p-1\right)  $ is a bijection from
$\left[  m\right]  $ to $T$. This bijection induces a bijection $J\mapsto
J+\left(  p-1\right)  $ from the set of all subsets of $\left[  m\right]  $ to
the set of all subsets of $T$. Hence, any subset $S$ of $T$ has the form
$S=S^{\prime}+\left(  p-1\right)  $ for a unique subset $S^{\prime}%
\subseteq\left[  m\right]  $. In particular, the subsets $U$ and $V$ of $T$
thus have the forms $U=P+\left(  p-1\right)  $ and $V=Q+\left(  p-1\right)  $
for unique subsets $P,Q\subseteq\left[  m\right]  $. Consider these $P,Q$.
Clearly, $\left\vert U\right\vert =\left\vert P+\left(  p-1\right)
\right\vert =\left\vert P\right\vert $ and $\left\vert V\right\vert
=\left\vert Q+\left(  p-1\right)  \right\vert =\left\vert Q\right\vert $, so
that $\left\vert P\right\vert =\left\vert U\right\vert =\left\vert
V\right\vert =\left\vert Q\right\vert $.

Let $\widetilde{P}:=\left[  m\right]  \setminus P$ and $\widetilde{Q}:=\left[
m\right]  \setminus Q$ be the complements of $P$ and $Q$ within $\left[
m\right]  $. Let $\sum P$ be the sum of all elements of $P$, and let $\sum Q$
be the sum of all elements of $Q$. From $U=P+\left(  p-1\right)  $, we obtain
\[
\sum U=\sum\left(  P+\left(  p-1\right)  \right)  =\sum P+\left\vert
P\right\vert \cdot\left(  p-1\right)
\]
(since the elements of $P+\left(  p-1\right)  $ are simply the $\left\vert
P\right\vert $ elements of $P$ with $p-1$ added to each). Likewise, $\sum
V=\sum Q+\left\vert Q\right\vert \cdot\left(  p-1\right)  $. Adding these two
equalities together, we find%
\begin{align*}
\sum U+\sum V  &  =\left(  \sum P+\left\vert P\right\vert \cdot\left(
p-1\right)  \right)  +\left(  \sum Q+\left\vert Q\right\vert \cdot\left(
p-1\right)  \right) \\
&  =\sum P+\sum Q+\underbrace{\left(  \left\vert P\right\vert +\left\vert
Q\right\vert \right)  }_{\substack{=\left\vert Q\right\vert +\left\vert
Q\right\vert \\\text{(since }\left\vert P\right\vert =\left\vert Q\right\vert
\text{)}}}\cdot\left(  p-1\right) \\
&  =\sum P+\sum Q+\underbrace{\left(  \left\vert Q\right\vert +\left\vert
Q\right\vert \right)  }_{=2\left\vert Q\right\vert \equiv0\operatorname{mod}%
2}\cdot\left(  p-1\right) \\
&  \equiv\sum P+\sum Q\operatorname{mod}2.
\end{align*}
Hence,%
\begin{equation}
\left(  -1\right)  ^{\sum U+\sum V}=\left(  -1\right)  ^{\sum P+\sum Q}.
\label{p17.lemA3.pf.sign}%
\end{equation}


Next, we recall that $A=\operatorname*{sub}\nolimits_{T}^{T}%
B=\operatorname*{sub}\nolimits_{\left[  p,q\right]  }^{\left[  p,q\right]  }B$
(since $T=\left[  p,q\right]  $). Hence, for all $i,j\in\left[  m\right]  $,
the $\left(  i,j\right)  $-th entry of $A$ is the $\left(  i+p-1,j+p-1\right)
$-th entry of $B$. Consequently, for any $X,Y\subseteq\left[  m\right]  $, we
have%
\[
\operatorname*{sub}\nolimits_{X}^{Y}A=\operatorname*{sub}\nolimits_{X+\left(
p-1\right)  }^{Y+\left(  p-1\right)  }B.
\]
Applying this to $X=P$ and $Y=Q$, we obtain%
\begin{equation}
\operatorname*{sub}\nolimits_{P}^{Q}A=\operatorname*{sub}\nolimits_{P+\left(
p-1\right)  }^{Q+\left(  p-1\right)  }B=\operatorname*{sub}\nolimits_{U}^{V}B
\label{p17.lemA3.pf.4}%
\end{equation}
(since $P+\left(  p-1\right)  =U$ and $Q+\left(  p-1\right)  =V$).

Furthermore, from $\widetilde{P}=\left[  m\right]  \setminus P$, we obtain
\begin{align*}
\widetilde{P}+\left(  p-1\right)   &  =\left(  \left[  m\right]  \setminus
P\right)  +\left(  p-1\right) \\
&  =\underbrace{\left(  \left[  m\right]  +\left(  p-1\right)  \right)  }%
_{=T}\setminus\underbrace{\left(  P+\left(  p-1\right)  \right)  }%
_{=U}\ \ \ \ \ \ \ \ \ \ \left(  \text{by (\ref{p17.lemA3.pf.diff})}\right) \\
&  =T\setminus U.
\end{align*}
Similarly, $\widetilde{Q}+\left(  p-1\right)  =T\setminus V$.

However, (\ref{p17.lemA3.pf.A1}) says that $A^{-1}=\operatorname*{sub}%
\nolimits_{T}^{T}\left(  B^{-1}\right)  =\operatorname*{sub}\nolimits_{\left[
p,q\right]  }^{\left[  p,q\right]  }\left(  B^{-1}\right)  $ (since $T=\left[
p,q\right]  $). Thus, for all $i,j\in\left[  m\right]  $, the $\left(
i,j\right)  $-th entry of $A^{-1}$ is the $\left(  i+p-1,j+p-1\right)  $-th
entry of $B^{-1}$. Consequently, for any $X,Y\subseteq\left[  m\right]  $, we
have%
\[
\operatorname*{sub}\nolimits_{X}^{Y}\left(  A^{-1}\right)
=\operatorname*{sub}\nolimits_{X+\left(  p-1\right)  }^{Y+\left(  p-1\right)
}\left(  B^{-1}\right)  .
\]
Applying this to $X=\widetilde{Q}$ and $Y=\widetilde{P}$, we obtain%
\begin{equation}
\operatorname*{sub}\nolimits_{\widetilde{Q}}^{\widetilde{P}}\left(
A^{-1}\right)  =\operatorname*{sub}\nolimits_{\widetilde{Q}+\left(
p-1\right)  }^{\widetilde{P}+\left(  p-1\right)  }\left(  B^{-1}\right)
=\operatorname*{sub}\nolimits_{T\setminus V}^{T\setminus U}\left(
B^{-1}\right)  \label{p17.lemA3.pf.5}%
\end{equation}
(since $\widetilde{P}+\left(  p-1\right)  =T\setminus U$ and $\widetilde{Q}%
+\left(  p-1\right)  =T\setminus V$).

Now, Lemma \ref{lem.A2} yields%
\begin{align*}
\det\left(  \operatorname*{sub}\nolimits_{P}^{Q}A\right)   &
=\underbrace{\left(  -1\right)  ^{\sum P+\sum Q}}_{\substack{=\left(
-1\right)  ^{\sum U+\sum V}\\\text{(by (\ref{p17.lemA3.pf.sign}))}%
}}\ \ \underbrace{\det A}_{\substack{=1\\\text{(by (\ref{p17.lemA3.pf.A2}))}%
}}\cdot\det\left(  \operatorname*{sub}\nolimits_{\widetilde{Q}}^{\widetilde{P}%
}\left(  A^{-1}\right)  \right) \\
&  =\left(  -1\right)  ^{\sum U+\sum V}\cdot\det\left(  \operatorname*{sub}%
\nolimits_{\widetilde{Q}}^{\widetilde{P}}\left(  A^{-1}\right)  \right)  .
\end{align*}
In view of (\ref{p17.lemA3.pf.4}) and (\ref{p17.lemA3.pf.5}), we can rewrite
this as%
\[
\det\left(  \operatorname*{sub}\nolimits_{U}^{V}B\right)  =\left(  -1\right)
^{\sum U+\sum V}\det\left(  \operatorname*{sub}\nolimits_{T\setminus
V}^{T\setminus U}\left(  B^{-1}\right)  \right)  .
\]
Thus, Lemma \ref{lem.A3} is proved.
\end{proof}

\begin{lemmain}
\label{lem.A4}Let $\lambda=\left(  \lambda_{1},\lambda_{2},\ldots,\lambda
_{q}\right)  $ be a partition, and let $\lambda^{\prime}=\left(  \lambda
_{1}^{\prime},\lambda_{2}^{\prime},\ldots,\lambda_{p}^{\prime}\right)  $ be
its conjugate partition. (The entries $\lambda_{i}$ and $\lambda_{j}^{\prime}$
are allowed to be $0$.) For each $i\in\left[  q\right]  $, let us set
$\alpha_{i}:=\lambda_{i}-i$. For each $j\in\left[  p\right]  $, let us set
$\beta_{j}:=\lambda_{j}^{\prime}-j$ and $\eta_{j}:=-1-\beta_{j}$. Then, the
two sets $\left\{  \alpha_{1},\alpha_{2},\ldots,\alpha_{q}\right\}  $ and
$\left\{  \eta_{1},\eta_{2},\ldots,\eta_{p}\right\}  $ are disjoint, and their
union is the integer interval $\left[  -q,p-1\right]  $.
\end{lemmain}

\begin{proof}
[Proof of Lemma \ref{lem.A4}.]Lemma \ref{lem.A4} is Proposition 3.18
\textbf{(f)} in the detailed version of the paper \cite{petriesym} (this
detailed version is downloadable from the arXiv as
\href{https://arxiv.org/src/2004.11194v3/anc/petriesym-long.pdf}{an ancillary
file}). (Note that the $\mu$ in the paper corresponds to our $\lambda^{\prime
}$, and that the conditions $p\geq\lambda_{1}$ and $q\geq\mu_{1}$ in the paper
follow from our assumptions $\lambda^{\prime}=\left(  \lambda_{1}^{\prime
},\lambda_{2}^{\prime},\ldots,\lambda_{p}^{\prime}\right)  $ and
$\lambda=\left(  \lambda_{1},\lambda_{2},\ldots,\lambda_{q}\right)  $, respectively.)
\end{proof}

\begin{lemmain}
\label{lem.A5}Let $\lambda=\left(  \lambda_{1},\lambda_{2},\ldots,\lambda
_{q}\right)  $ be a partition, and let $\lambda^{\prime}=\left(  \lambda
_{1}^{\prime},\lambda_{2}^{\prime},\ldots,\lambda_{p}^{\prime}\right)  $ be
its conjugate partition. (The entries $\lambda_{i}$ and $\lambda_{j}^{\prime}$
are allowed to be $0$.)

Let $\mu=\left(  \mu_{1},\mu_{2},\ldots,\mu_{q}\right)  $ be a partition, and
let $\mu^{\prime}=\left(  \mu_{1}^{\prime},\mu_{2}^{\prime},\ldots,\mu
_{p}^{\prime}\right)  $ be its conjugate partition. (The entries $\mu_{i}$ and
$\mu_{j}^{\prime}$ are allowed to be $0$.)

Let $B=\left(  b_{i,j}\right)  _{i,j\in\mathbb{Z}}\in\operatorname*{UT}%
\nolimits_{\mathbb{Z}}$ be an upper unitriangular matrix. Let $\left(
c_{i,j}\right)  _{i,j\in\mathbb{Z}}=B^{-1}$ be its inverse matrix. Then,%
\begin{align*}
&  \det\left(  b_{\mu_{i}-i,\ \lambda_{j}-j}\right)  _{i,j\in\left[  q\right]
}\\
&  =\left(  -1\right)  ^{\left(  \lambda_{1}+\lambda_{2}+\cdots+\lambda
_{q}\right)  +\left(  \mu_{1}+\mu_{2}+\cdots+\mu_{q}\right)  }\det\left(
c_{i-\lambda_{i}^{\prime}-1,\ j-\mu_{j}^{\prime}-1}\right)  _{i,j\in\left[
p\right]  }.
\end{align*}

\end{lemmain}

\begin{proof}
[Proof of Lemma \ref{lem.A5}.]Let $T$ be the integer interval $\left[
-q,p-1\right]  $. For each $i\in\left[  q\right]  $, let us set $\alpha
_{i}:=\lambda_{i}-i$ and $\gamma_{i}:=\mu_{i}-i$. For each $j\in\left[
p\right]  $, let us set $\beta_{j}:=\lambda_{j}^{\prime}-j$ and $\eta
_{j}:=-1-\beta_{j}$ and $\delta_{j}:=\mu_{j}^{\prime}-j$ and $\omega
_{j}:=-1-\delta_{j}$.

Then, Lemma \ref{lem.A4} says that the two sets $\left\{  \alpha_{1}%
,\alpha_{2},\ldots,\alpha_{q}\right\}  $ and $\left\{  \eta_{1},\eta
_{2},\ldots,\eta_{p}\right\}  $ are disjoint, and their union is the integer
interval $\left[  -q,p-1\right]  $. Therefore, the set $\left\{  \eta_{1}%
,\eta_{2},\ldots,\eta_{p}\right\}  $ is the complement of the set $\left\{
\alpha_{1},\alpha_{2},\ldots,\alpha_{q}\right\}  $ in the interval $\left[
-q,p-1\right]  $. In other words,%
\begin{align}
\left\{  \eta_{1},\eta_{2},\ldots,\eta_{p}\right\}   &  =\underbrace{\left[
-q,p-1\right]  }_{=T}\setminus\left\{  \alpha_{1},\alpha_{2},\ldots,\alpha
_{q}\right\} \nonumber\\
&  =T\setminus\left\{  \alpha_{1},\alpha_{2},\ldots,\alpha_{q}\right\}  .
\label{p17.lemA5.pf.1}%
\end{align}
The same argument (applied to $\mu$, $\mu^{\prime}$, $\gamma_{i}$, $\delta
_{j}$ and $\omega_{j}$ instead of $\lambda$, $\lambda^{\prime}$, $\alpha_{i}$,
$\beta_{j}$ and $\eta_{j}$) yields%
\begin{equation}
\left\{  \omega_{1},\omega_{2},\ldots,\omega_{p}\right\}  =T\setminus\left\{
\gamma_{1},\gamma_{2},\ldots,\gamma_{q}\right\}  . \label{p17.lemA5.pf.2}%
\end{equation}


Since $\lambda$ is a partition, we have $\lambda_{1}\geq\lambda_{2}\geq
\cdots\geq\lambda_{q}$ and thus $\lambda_{1}-1>\lambda_{2}-2>\cdots
>\lambda_{q}-q$. In other words, $\alpha_{1}>\alpha_{2}>\cdots>\alpha_{q}$
(since $\alpha_{i}=\lambda_{i}-i$ for each $i$). Hence,
\begin{equation}
\left\{  \alpha_{1},\alpha_{2},\ldots,\alpha_{q}\right\}  =\left\{  \alpha
_{q}<\alpha_{q-1}<\cdots<\alpha_{1}\right\}  . \label{p17.lemA5.pf.3a}%
\end{equation}
Similarly,%
\begin{equation}
\left\{  \gamma_{1},\gamma_{2},\ldots,\gamma_{q}\right\}  =\left\{  \gamma
_{q}<\gamma_{q-1}<\cdots<\gamma_{1}\right\}  . \label{p17.lemA5.pf.3b}%
\end{equation}
These two equalities show that both sets $\left\{  \gamma_{1},\gamma
_{2},\ldots,\gamma_{q}\right\}  $ and $\left\{  \alpha_{1},\alpha_{2}%
,\ldots,\alpha_{q}\right\}  $ have size $q$, so that they have the same size.
In other words, $\left\vert \left\{  \gamma_{1},\gamma_{2},\ldots,\gamma
_{q}\right\}  \right\vert =\left\vert \left\{  \alpha_{1},\alpha_{2}%
,\ldots,\alpha_{q}\right\}  \right\vert $.

Furthermore, recall that the two sets $\left\{  \alpha_{1},\alpha_{2}%
,\ldots,\alpha_{q}\right\}  $ and $\left\{  \eta_{1},\eta_{2},\ldots,\eta
_{p}\right\}  $ are disjoint, and their union is the integer interval $\left[
-q,p-1\right]  $. Hence, in particular, $\left\{  \alpha_{1},\alpha_{2}%
,\ldots,\alpha_{q}\right\}  $ is a subset of $\left[  -q,p-1\right]  =T$.
Similarly, $\left\{  \gamma_{1},\gamma_{2},\ldots,\gamma_{q}\right\}  $ is a
subset of $T$ as well.

Thus, Lemma \ref{lem.A3} (applied to $-q$, $p-1$, $\left\{  \gamma_{1}%
,\gamma_{2},\ldots,\gamma_{q}\right\}  $ and $\left\{  \alpha_{1},\alpha
_{2},\ldots,\alpha_{q}\right\}  $ instead of $p$, $q$, $U$ and $V$) yields%
\begin{align}
&  \det\left(  \operatorname*{sub}\nolimits_{\left\{  \gamma_{1},\gamma
_{2},\ldots,\gamma_{q}\right\}  }^{\left\{  \alpha_{1},\alpha_{2}%
,\ldots,\alpha_{q}\right\}  }B\right) \nonumber\\
&  =\left(  -1\right)  ^{\sum\left\{  \gamma_{1},\gamma_{2},\ldots,\gamma
_{q}\right\}  +\sum\left\{  \alpha_{1},\alpha_{2},\ldots,\alpha_{q}\right\}
}\nonumber\\
&  \ \ \ \ \ \ \ \ \ \ \det\left(  \operatorname*{sub}\nolimits_{T\setminus
\left\{  \alpha_{1},\alpha_{2},\ldots,\alpha_{q}\right\}  }^{T\setminus
\left\{  \gamma_{1},\gamma_{2},\ldots,\gamma_{q}\right\}  }\left(
B^{-1}\right)  \right)  . \label{p17.lemA5.pf.4}%
\end{align}


In view of (\ref{p17.lemA5.pf.3a}) and (\ref{p17.lemA5.pf.3b}), we have%
\[
\operatorname*{sub}\nolimits_{\left\{  \gamma_{1},\gamma_{2},\ldots,\gamma
_{q}\right\}  }^{\left\{  \alpha_{1},\alpha_{2},\ldots,\alpha_{q}\right\}
}B=\operatorname*{sub}\nolimits_{\left\{  \gamma_{q}<\gamma_{q-1}%
<\cdots<\gamma_{1}\right\}  }^{\left\{  \alpha_{q}<\alpha_{q-1}<\cdots
<\alpha_{1}\right\}  }B=\left(  b_{\gamma_{q+1-i},\alpha_{q+1-j}}\right)
_{i,j\in\left[  q\right]  }%
\]
(since $B=\left(  b_{i,j}\right)  _{i,j\in\mathbb{Z}}$). Hence,%
\begin{align}
\det\left(  \operatorname*{sub}\nolimits_{\left\{  \gamma_{1},\gamma
_{2},\ldots,\gamma_{q}\right\}  }^{\left\{  \alpha_{1},\alpha_{2}%
,\ldots,\alpha_{q}\right\}  }B\right)   &  =\det\left(  b_{\gamma
_{q+1-i},\alpha_{q+1-j}}\right)  _{i,j\in\left[  q\right]  }\nonumber\\
&  =\det\left(  b_{\gamma_{i},\alpha_{j}}\right)  _{i,j\in\left[  q\right]  }
\label{p17.lemA5.pf.5}%
\end{align}
(here, we have substituted $q+1-i$ and $q+1-j$ for $i$ and $j$ in the matrix,
which effectively rotates the matrix by $180^{\circ}$; this rotation does not
change the determinant, because it is a composition of a row permutation and a
column permutation with the same sign).

On the other hand, let us recall that $\left\{  \eta_{1},\eta_{2},\ldots
,\eta_{p}\right\}  =T\setminus\left\{  \alpha_{1},\alpha_{2},\ldots,\alpha
_{q}\right\}  $. But $\lambda^{\prime}$ is a partition; thus, $\lambda
_{1}^{\prime}\geq\lambda_{2}^{\prime}\geq\cdots\geq\lambda_{p}^{\prime}$ and
therefore $\lambda_{1}^{\prime}-1>\lambda_{2}^{\prime}-2>\cdots>\lambda
_{p}^{\prime}-p$. In other words, $\beta_{1}>\beta_{2}>\cdots>\beta_{p}$
(since $\beta_{j}=\lambda_{j}^{\prime}-j$ for each $j$). Hence, $-1-\beta
_{1}<-1-\beta_{2}<\cdots<-1-\beta_{p}$. In other words, $\eta_{1}<\eta
_{2}<\cdots<\eta_{p}$ (since $\eta_{j}=-1-\beta_{j}$ for each $j$). Thus,
$\left\{  \eta_{1},\eta_{2},\ldots,\eta_{p}\right\}  =\left\{  \eta_{1}%
<\eta_{2}<\cdots<\eta_{p}\right\}  $. Comparing this with
(\ref{p17.lemA5.pf.1}), we obtain%
\[
T\setminus\left\{  \alpha_{1},\alpha_{2},\ldots,\alpha_{q}\right\}  =\left\{
\eta_{1}<\eta_{2}<\cdots<\eta_{p}\right\}  .
\]
The same argument (applied to $\mu$, $\mu^{\prime}$, $\gamma_{i}$, $\delta
_{j}$ and $\omega_{j}$ instead of $\lambda$, $\lambda^{\prime}$, $\alpha_{i}$,
$\beta_{j}$ and $\eta_{j}$) yields%
\[
T\setminus\left\{  \gamma_{1},\gamma_{2},\ldots,\gamma_{q}\right\}  =\left\{
\omega_{1}<\omega_{2}<\cdots<\omega_{p}\right\}  .
\]
In view of these two equalities, we have%
\[
\operatorname*{sub}\nolimits_{T\setminus\left\{  \alpha_{1},\alpha_{2}%
,\ldots,\alpha_{q}\right\}  }^{T\setminus\left\{  \gamma_{1},\gamma_{2}%
,\ldots,\gamma_{q}\right\}  }\left(  B^{-1}\right)  =\operatorname*{sub}%
\nolimits_{\left\{  \eta_{1}<\eta_{2}<\cdots<\eta_{p}\right\}  }^{\left\{
\omega_{1}<\omega_{2}<\cdots<\omega_{p}\right\}  }\left(  B^{-1}\right)
=\left(  c_{\eta_{i},\omega_{j}}\right)  _{i,j\in\left[  p\right]  }%
\]
(since $B^{-1}=\left(  c_{i,j}\right)  _{i,j\in\mathbb{Z}}$). Thus,%
\begin{equation}
\det\left(  \operatorname*{sub}\nolimits_{T\setminus\left\{  \alpha_{1}%
,\alpha_{2},\ldots,\alpha_{q}\right\}  }^{T\setminus\left\{  \gamma_{1}%
,\gamma_{2},\ldots,\gamma_{q}\right\}  }\left(  B^{-1}\right)  \right)
=\det\left(  c_{\eta_{i},\omega_{j}}\right)  _{i,j\in\left[  p\right]  }.
\label{p17.lemA5.pf.6}%
\end{equation}


Furthermore, $\alpha_{1}>\alpha_{2}>\cdots>\alpha_{q}$ shows that the numbers
$\alpha_{1},\alpha_{2},\ldots,\alpha_{q}$ are distinct; hence,%
\begin{align*}
\sum\left\{  \alpha_{1},\alpha_{2},\ldots,\alpha_{q}\right\}   &  =\alpha
_{1}+\alpha_{2}+\cdots+\alpha_{q}\\
&  =\left(  \lambda_{1}-1\right)  +\left(  \lambda_{2}-2\right)
+\cdots+\left(  \lambda_{q}-q\right) \\
&  \ \ \ \ \ \ \ \ \ \ \ \ \ \ \ \ \ \ \ \ \left(  \text{since }\alpha
_{i}=\lambda_{i}-i\text{ for all }i\right) \\
&  =\left(  \lambda_{1}+\lambda_{2}+\cdots+\lambda_{q}\right)  -\left(
1+2+\cdots+q\right)  .
\end{align*}
Similarly,%
\[
\sum\left\{  \gamma_{1},\gamma_{2},\ldots,\gamma_{q}\right\}  =\left(  \mu
_{1}+\mu_{2}+\cdots+\mu_{q}\right)  -\left(  1+2+\cdots+q\right)  .
\]
Adding these two equalities together, we find%
\begin{align*}
&  \sum\left\{  \alpha_{1},\alpha_{2},\ldots,\alpha_{q}\right\}  +\sum\left\{
\gamma_{1},\gamma_{2},\ldots,\gamma_{q}\right\} \\
&  =\left(  \lambda_{1}+\lambda_{2}+\cdots+\lambda_{q}\right)  -\left(
1+2+\cdots+q\right) \\
&  \ \ \ \ \ \ \ \ \ \ +\left(  \mu_{1}+\mu_{2}+\cdots+\mu_{q}\right)
-\left(  1+2+\cdots+q\right) \\
&  =\left(  \lambda_{1}+\lambda_{2}+\cdots+\lambda_{q}\right)  +\left(
\mu_{1}+\mu_{2}+\cdots+\mu_{q}\right)  -\underbrace{2\left(  1+2+\cdots
+q\right)  }_{\equiv0\operatorname{mod}2}\\
&  \equiv\left(  \lambda_{1}+\lambda_{2}+\cdots+\lambda_{q}\right)  +\left(
\mu_{1}+\mu_{2}+\cdots+\mu_{q}\right)  \operatorname{mod}2.
\end{align*}
Hence,%
\begin{align}
&  \left(  -1\right)  ^{\sum\left\{  \alpha_{1},\alpha_{2},\ldots,\alpha
_{q}\right\}  +\sum\left\{  \gamma_{1},\gamma_{2},\ldots,\gamma_{q}\right\}
}\nonumber\\
&  =\left(  -1\right)  ^{\left(  \lambda_{1}+\lambda_{2}+\cdots+\lambda
_{q}\right)  +\left(  \mu_{1}+\mu_{2}+\cdots+\mu_{q}\right)  }.
\label{p17.lemA5.pf.7}%
\end{align}


Finally, using (\ref{p17.lemA5.pf.5}), we can rewrite (\ref{p17.lemA5.pf.4})
as%
\begin{align*}
&  \det\left(  b_{\gamma_{i},\alpha_{j}}\right)  _{i,j\in\left[  q\right]  }\\
&  =\underbrace{\left(  -1\right)  ^{\sum\left\{  \gamma_{1},\gamma_{2}%
,\ldots,\gamma_{q}\right\}  +\sum\left\{  \alpha_{1},\alpha_{2},\ldots
,\alpha_{q}\right\}  }}_{\substack{=\left(  -1\right)  ^{\sum\left\{
\alpha_{1},\alpha_{2},\ldots,\alpha_{q}\right\}  +\sum\left\{  \gamma
_{1},\gamma_{2},\ldots,\gamma_{q}\right\}  }\\=\left(  -1\right)  ^{\left(
\lambda_{1}+\lambda_{2}+\cdots+\lambda_{q}\right)  +\left(  \mu_{1}+\mu
_{2}+\cdots+\mu_{q}\right)  }\\\text{(by (\ref{p17.lemA5.pf.7}))}%
}}\underbrace{\det\left(  \operatorname*{sub}\nolimits_{T\setminus\left\{
\alpha_{1},\alpha_{2},\ldots,\alpha_{q}\right\}  }^{T\setminus\left\{
\gamma_{1},\gamma_{2},\ldots,\gamma_{q}\right\}  }\left(  B^{-1}\right)
\right)  }_{\substack{=\det\left(  c_{\eta_{i},\omega_{j}}\right)
_{i,j\in\left[  p\right]  }\\\text{(by (\ref{p17.lemA5.pf.6}))}}}\\
&  =\left(  -1\right)  ^{\left(  \lambda_{1}+\lambda_{2}+\cdots+\lambda
_{q}\right)  +\left(  \mu_{1}+\mu_{2}+\cdots+\mu_{q}\right)  }\det\left(
c_{\eta_{i},\omega_{j}}\right)  _{i,j\in\left[  p\right]  }.
\end{align*}
In view of $\gamma_{i}=\mu_{i}-i$ and $\alpha_{j}=\lambda_{j}-j$ and $\eta
_{i}=-1-\underbrace{\beta_{i}}_{=\lambda_{i}^{\prime}-i}=-1-\left(
\lambda_{i}^{\prime}-i\right)  =i-\lambda_{i}^{\prime}-1$ and $\omega
_{j}=-1-\underbrace{\delta_{j}}_{=\mu_{j}^{\prime}-j}=-1-\left(  \mu
_{j}^{\prime}-j\right)  =j-\mu_{j}^{\prime}-1$, we can rewrite this as%
\begin{align*}
&  \det\left(  b_{\mu_{i}-i,\ \lambda_{j}-j}\right)  _{i,j\in\left[  q\right]
}\\
&  =\left(  -1\right)  ^{\left(  \lambda_{1}+\lambda_{2}+\cdots+\lambda
_{q}\right)  +\left(  \mu_{1}+\mu_{2}+\cdots+\mu_{q}\right)  }\det\left(
c_{i-\lambda_{i}^{\prime}-1,\ j-\mu_{j}^{\prime}-1}\right)  _{i,j\in\left[
p\right]  }.
\end{align*}
This proves Lemma \ref{lem.A5}.
\end{proof}

\begin{lemmain}
\label{lem.A6}Let $\left(  u_{i,j}\right)  _{i,j\in\left[  p\right]  }$ be a
$p\times p$-matrix. Furthermore, let $\alpha_{1},\alpha_{2},\ldots,\alpha
_{p},\beta_{1},\beta_{2},\ldots,\beta_{p}$ be any $2p$ scalars. Then,%
\[
\det\left(  \alpha_{i}\beta_{j}u_{i,j}\right)  _{i,j\in\left[  p\right]
}=\left(  \prod_{i=1}^{p}\alpha_{i}\right)  \left(  \prod_{j=1}^{p}\beta
_{j}\right)  \cdot\det\left(  u_{i,j}\right)  _{i,j\in\left[  p\right]  }.
\]

\end{lemmain}

\begin{proof}
[Proof of Lemma \ref{lem.A6}.]The matrix $\left(  \alpha_{i}\beta_{j}%
u_{i,j}\right)  _{i,j\in\left[  p\right]  }$ is obtained from the matrix
$\left(  u_{i,j}\right)  _{i,j\in\left[  p\right]  }$ by

\begin{itemize}
\item multiplying the $i$-th row by $\alpha_{i}$ for each $i\in\left[
p\right]  $, and then

\item multiplying the $j$-th column by $\beta_{j}$ for each $j\in\left[
p\right]  $.
\end{itemize}

Each of these operations multiplies the determinant of the matrix by the
corresponding factor $\alpha_{i}$ or $\beta_{j}$. Thus, in total, the
determinant gets multiplied by $\left(  \prod_{i=1}^{p}\alpha_{i}\right)
\left(  \prod_{j=1}^{p}\beta_{j}\right)  $. This proves Lemma \ref{lem.A6}.
\end{proof}

\begin{lemmain}
\label{lem.A6prods}Let $\lambda$ and $\mu$ be two partitions. Write the
partitions $\lambda$, $\lambda^{\prime}$, $\mu$ and $\mu^{\prime}$ in the
forms%
\begin{align*}
\lambda &  =\left(  \lambda_{1},\lambda_{2},\ldots,\lambda_{q}\right)
,\ \ \ \ \ \ \ \ \ \ \lambda^{\prime}=\left(  \lambda_{1}^{\prime},\lambda
_{2}^{\prime},\ldots,\lambda_{p}^{\prime}\right)  ,\\
\mu &  =\left(  \mu_{1},\mu_{2},\ldots,\mu_{q}\right)
,\ \ \ \ \ \ \ \ \ \ \mu^{\prime}=\left(  \mu_{1}^{\prime},\mu_{2}^{\prime
},\ldots,\mu_{p}^{\prime}\right)
\end{align*}
for some $p,q\in\mathbb{N}$ (where, of course, the entries $\lambda_{i}$,
$\lambda_{j}^{\prime}$, $\mu_{i}$ and $\mu_{j}^{\prime}$ are allowed to be
$0$). Then,%
\[
\left(  -1\right)  ^{\left(  \lambda_{1}+\lambda_{2}+\cdots+\lambda
_{q}\right)  +\left(  \mu_{1}+\mu_{2}+\cdots+\mu_{q}\right)  }\left(
\prod_{i=1}^{p}\left(  -1\right)  ^{\lambda_{i}^{\prime}-i}\right)  \left(
\prod_{j=1}^{p}\left(  -1\right)  ^{\mu_{j}^{\prime}-j}\right)  =1.
\]

\end{lemmain}

\begin{proof}
[Proof of Lemma \ref{lem.A6prods}.]We have $\left\vert \lambda\right\vert
=\lambda_{1}+\lambda_{2}+\cdots+\lambda_{q}$ and $\left\vert \mu\right\vert
=\mu_{1}+\mu_{2}+\cdots+\mu_{q}$ and $\left\vert \lambda^{\prime}\right\vert
=\lambda_{1}^{\prime}+\lambda_{2}^{\prime}+\cdots+\lambda_{p}^{\prime}%
=\sum_{i=1}^{p}\lambda_{i}^{\prime}$ and $\left\vert \mu^{\prime}\right\vert
=\mu_{1}^{\prime}+\mu_{2}^{\prime}+\cdots+\mu_{p}^{\prime}=\sum_{j=1}^{p}%
\mu_{j}^{\prime}$. But we also have $\left\vert \lambda^{\prime}\right\vert
=\left\vert \lambda\right\vert $ (since the transpose of a partition has the
same size as the original partition) and $\left\vert \mu^{\prime}\right\vert
=\left\vert \mu\right\vert $ (for the same reason). Now, the integer%
\begin{align*}
&  \underbrace{\left(  \lambda_{1}+\lambda_{2}+\cdots+\lambda_{q}\right)
}_{=\left\vert \lambda\right\vert }+\underbrace{\left(  \mu_{1}+\mu_{2}%
+\cdots+\mu_{q}\right)  }_{=\left\vert \mu\right\vert }+\underbrace{\sum
_{i=1}^{p}\left(  \lambda_{i}^{\prime}-i\right)  }_{=\sum_{i=1}^{p}\lambda
_{i}^{\prime}-\sum_{i=1}^{p}i}+\underbrace{\sum_{j=1}^{p}\left(  \mu
_{j}^{\prime}-j\right)  }_{=\sum_{j=1}^{p}\mu_{j}^{\prime}-\sum_{j=1}^{p}j}\\
&  =\left\vert \lambda\right\vert +\left\vert \mu\right\vert +\underbrace{\sum
_{i=1}^{p}\lambda_{i}^{\prime}}_{=\left\vert \lambda^{\prime}\right\vert
=\left\vert \lambda\right\vert }-\underbrace{\sum_{i=1}^{p}i}_{=1+2+\cdots
+p}+\underbrace{\sum_{j=1}^{p}\mu_{j}^{\prime}}_{=\left\vert \mu^{\prime
}\right\vert =\left\vert \mu\right\vert }-\underbrace{\sum_{j=1}^{p}%
j}_{=1+2+\cdots+p}\\
&  =\left\vert \lambda\right\vert +\left\vert \mu\right\vert +\left\vert
\lambda\right\vert -\left(  1+2+\cdots+p\right)  +\left\vert \mu\right\vert
-\left(  1+2+\cdots+p\right)  \\
&  =2\left(  \left\vert \lambda\right\vert +\left\vert \mu\right\vert -\left(
1+2+\cdots+p\right)  \right)
\end{align*}
is even. Thus,%
\[
\left(  -1\right)  ^{\left(  \lambda_{1}+\lambda_{2}+\cdots+\lambda
_{q}\right)  +\left(  \mu_{1}+\mu_{2}+\cdots+\mu_{q}\right)  +\sum_{i=1}%
^{p}\left(  \lambda_{i}^{\prime}-i\right)  +\sum_{j=1}^{p}\left(  \mu
_{j}^{\prime}-j\right)  }=1.
\]
Finally,%
\begin{align*}
&  \left(  -1\right)  ^{\left(  \lambda_{1}+\lambda_{2}+\cdots+\lambda
_{q}\right)  +\left(  \mu_{1}+\mu_{2}+\cdots+\mu_{q}\right)  }\left(
\prod_{i=1}^{p}\left(  -1\right)  ^{\lambda_{i}^{\prime}-i}\right)  \left(
\prod_{j=1}^{p}\left(  -1\right)  ^{\mu_{j}^{\prime}-j}\right)  \\
&  =\left(  -1\right)  ^{\left(  \lambda_{1}+\lambda_{2}+\cdots+\lambda
_{q}\right)  +\left(  \mu_{1}+\mu_{2}+\cdots+\mu_{q}\right)  +\sum_{i=1}%
^{p}\left(  \lambda_{i}^{\prime}-i\right)  +\sum_{j=1}^{p}\left(  \mu
_{j}^{\prime}-j\right)  }\\
&  =1.
\end{align*}
This proves Lemma \ref{lem.A6prods}.
\end{proof}

Now, we can prove (6.9):

\begin{proof}
[Proof of (6.9).]Write the partitions $\lambda$, $\lambda^{\prime}$, $\mu$ and
$\mu^{\prime}$ in the forms%
\begin{align*}
\lambda &  =\left(  \lambda_{1},\lambda_{2},\ldots,\lambda_{q}\right)
,\ \ \ \ \ \ \ \ \ \ \lambda^{\prime}=\left(  \lambda_{1}^{\prime},\lambda
_{2}^{\prime},\ldots,\lambda_{p}^{\prime}\right)  ,\\
\mu &  =\left(  \mu_{1},\mu_{2},\ldots,\mu_{q}\right)
,\ \ \ \ \ \ \ \ \ \ \mu^{\prime}=\left(  \mu_{1}^{\prime},\mu_{2}^{\prime
},\ldots,\mu_{p}^{\prime}\right)
\end{align*}
for some $p,q\in\mathbb{N}$ (where, of course, the entries $\lambda_{i}$,
$\lambda_{j}^{\prime}$, $\mu_{i}$ and $\mu_{j}^{\prime}$ are allowed to be $0$).

Recall that%
\begin{align*}
\mathbf{H}\left(  x\mid a\right)   &  =\left(  h_{j-i}\left(  x\mid\tau
^{i+1}a\right)  \right)  _{i,j\in\mathbb{Z}}\ \ \ \ \ \ \ \ \ \ \text{and}\\
\mathbf{E}\left(  x\mid a\right)   &  =\left(  \left(  -1\right)
^{j-i}e_{j-i}\left(  x\mid\tau^{j}a\right)  \right)  _{i,j\in\mathbb{Z}}%
\end{align*}
are two upper unitriangular matrices in $\operatorname*{UT}%
\nolimits_{\mathbb{Z}}$ (since $h_{0}\left(  x\mid a\right)  =e_{0}\left(
x\mid a\right)  =1$ and $h_{k}\left(  x\mid a\right)  =e_{k}\left(  x\mid
a\right)  =0$ whenever $k<0$). From (6.6), we know that the matrix
$\mathbf{E}\left(  x\mid a\right)  $ is the inverse matrix of $\mathbf{H}%
\left(  x\mid a\right)  $. Hence, Lemma \ref{lem.A5} (applied to
$B=\mathbf{H}\left(  x\mid a\right)  $ and $b_{i,j}=h_{j-i}\left(  x\mid
\tau^{i+1}a\right)  $ and $c_{i,j}=\left(  -1\right)  ^{j-i}e_{j-i}\left(
x\mid\tau^{j}a\right)  $) yields%
\begin{align*}
&  \det\left(  h_{\left(  \lambda_{j}-j\right)  -\left(  \mu_{i}-i\right)
}\left(  x\mid\tau^{\mu_{i}-i+1}a\right)  \right)  _{i,j\in\left[  q\right]
}\\
&  =\left(  -1\right)  ^{\left(  \lambda_{1}+\lambda_{2}+\cdots+\lambda
_{q}\right)  +\left(  \mu_{1}+\mu_{2}+\cdots+\mu_{q}\right)  }\\
&  \ \ \ \ \ \ \ \ \ \ \det\left(  \left(  -1\right)  ^{\left(  j-\mu
_{j}^{\prime}-1\right)  -\left(  i-\lambda_{i}^{\prime}-1\right)  }e_{\left(
j-\mu_{j}^{\prime}-1\right)  -\left(  i-\lambda_{i}^{\prime}-1\right)
}\left(  x\mid\tau^{j-\mu_{j}^{\prime}-1}a\right)  \right)  _{i,j\in\left[
p\right]  }.
\end{align*}
In view of $\left(  \lambda_{j}-j\right)  -\left(  \mu_{i}-i\right)
=\lambda_{j}-\mu_{i}-j+i$ and $\left(  j-\mu_{j}^{\prime}-1\right)  -\left(
i-\lambda_{i}^{\prime}-1\right)  =\lambda_{i}^{\prime}-\mu_{j}^{\prime}-i+j$,
we can rewrite this as%
\begin{align*}
&  \det\left(  h_{\lambda_{j}-\mu_{i}-j+i}\left(  x\mid\tau^{\mu_{i}%
-i+1}a\right)  \right)  _{i,j\in\left[  q\right]  }\\
&  =\left(  -1\right)  ^{\left(  \lambda_{1}+\lambda_{2}+\cdots+\lambda
_{q}\right)  +\left(  \mu_{1}+\mu_{2}+\cdots+\mu_{q}\right)  }\\
&  \ \ \ \ \ \ \ \ \ \ \det\left(  \left(  -1\right)  ^{\lambda_{i}^{\prime
}-\mu_{j}^{\prime}-i+j}e_{\lambda_{i}^{\prime}-\mu_{j}^{\prime}-i+j}\left(
x\mid\tau^{j-\mu_{j}^{\prime}-1}a\right)  \right)  _{i,j\in\left[  p\right]
}.
\end{align*}
In view of%
\[
\det\left(  h_{\lambda_{j}-\mu_{i}-j+i}\left(  x\mid\tau^{\mu_{i}%
-i+1}a\right)  \right)  _{i,j\in\left[  q\right]  }=\det\left(  h_{\lambda
_{i}-\mu_{j}-i+j}\left(  x\mid\tau^{\mu_{j}-j+1}a\right)  \right)
_{i,j\in\left[  q\right]  }%
\]
(since the determinant of a matrix does not change when we transpose it) and%
\begin{align*}
&  \det\left(  \underbrace{\left(  -1\right)  ^{\lambda_{i}^{\prime}-\mu
_{j}^{\prime}-i+j}}_{=\left(  -1\right)  ^{\lambda_{i}^{\prime}-i}\left(
-1\right)  ^{\mu_{j}^{\prime}-j}}e_{\lambda_{i}^{\prime}-\mu_{j}^{\prime}%
-i+j}\left(  x\mid\tau^{j-\mu_{j}^{\prime}-1}a\right)  \right)  _{i,j\in
\left[  p\right]  }\\
&  =\det\left(  \left(  -1\right)  ^{\lambda_{i}^{\prime}-i}\left(  -1\right)
^{\mu_{j}^{\prime}-j}e_{\lambda_{i}^{\prime}-\mu_{j}^{\prime}-i+j}\left(
x\mid\tau^{j-\mu_{j}^{\prime}-1}a\right)  \right)  _{i,j\in\left[  p\right]
}\\
&  =\left(  \prod_{i=1}^{p}\left(  -1\right)  ^{\lambda_{i}^{\prime}%
-i}\right)  \left(  \prod_{j=1}^{p}\left(  -1\right)  ^{\mu_{j}^{\prime}%
-j}\right)  \cdot\det\left(  e_{\lambda_{i}^{\prime}-\mu_{j}^{\prime}%
-i+j}\left(  x\mid\tau^{j-\mu_{j}^{\prime}-1}a\right)  \right)  _{i,j\in
\left[  p\right]  }%
\end{align*}
(by Lemma \ref{lem.A6}, applied to $\alpha_{i}=\left(  -1\right)
^{\lambda_{i}^{\prime}-i}$ and $\beta_{j}=\left(  -1\right)  ^{\mu_{j}%
^{\prime}-j}$ and $u_{i,j}=e_{\lambda_{i}^{\prime}-\mu_{j}^{\prime}%
-i+j}\left(  x\mid\tau^{j-\mu_{j}^{\prime}-1}a\right)  $), we can rewrite this
as%
\begin{align*}
&  \det\left(  h_{\lambda_{i}-\mu_{j}-i+j}\left(  x\mid\tau^{\mu_{j}%
-j+1}a\right)  \right)  _{i,j\in\left[  q\right]  }\\
&  =\left(  -1\right)  ^{\left(  \lambda_{1}+\lambda_{2}+\cdots+\lambda
_{q}\right)  +\left(  \mu_{1}+\mu_{2}+\cdots+\mu_{q}\right)  }\\
&  \ \ \ \ \ \ \ \ \ \ \left(  \prod_{i=1}^{p}\left(  -1\right)  ^{\lambda
_{i}^{\prime}-i}\right)  \left(  \prod_{j=1}^{p}\left(  -1\right)  ^{\mu
_{j}^{\prime}-j}\right)  \cdot\det\left(  e_{\lambda_{i}^{\prime}-\mu
_{j}^{\prime}-i+j}\left(  x\mid\tau^{j-\mu_{j}^{\prime}-1}a\right)  \right)
_{i,j\in\left[  p\right]  }\\
&  =\underbrace{\left(  -1\right)  ^{\left(  \lambda_{1}+\lambda_{2}%
+\cdots+\lambda_{q}\right)  +\left(  \mu_{1}+\mu_{2}+\cdots+\mu_{q}\right)
}\left(  \prod_{i=1}^{p}\left(  -1\right)  ^{\lambda_{i}^{\prime}-i}\right)
\left(  \prod_{j=1}^{p}\left(  -1\right)  ^{\mu_{j}^{\prime}-j}\right)
}_{\substack{=1\\\text{(by Lemma \ref{lem.A6prods})}}}\\
&  \ \ \ \ \ \ \ \ \ \ \cdot\det\left(  e_{\lambda_{i}^{\prime}-\mu
_{j}^{\prime}-i+j}\left(  x\mid\tau^{j-\mu_{j}^{\prime}-1}a\right)  \right)
_{i,j\in\left[  p\right]  }\\
&  =\det\left(  e_{\lambda_{i}^{\prime}-\mu_{j}^{\prime}-i+j}\left(  x\mid
\tau^{j-\mu_{j}^{\prime}-1}a\right)  \right)  _{i,j\in\left[  p\right]  }\\
&  =\det\left(  e_{\lambda_{i}^{\prime}-\mu_{j}^{\prime}-i+j}\left(  x\mid
\tau^{-\mu_{j}^{\prime}+j-1}a\right)  \right)  _{i,j\in\left[  p\right]
}\ \ \ \ \ \ \ \ \ \ \left(  \text{since }j-\mu_{j}^{\prime}-1=-\mu
_{j}^{\prime}+j-1\right)  .
\end{align*}
In view of (6.8), this can be rewritten as%
\[
s_{\lambda/\mu}\left(  x\mid a\right)  =\det\left(  e_{\lambda_{i}^{\prime
}-\mu_{j}^{\prime}-i+j}\left(  x\mid\tau^{-\mu_{j}^{\prime}+j-1}a\right)
\right)  _{i,j\in\left[  p\right]  }.
\]
This proves (6.9).
\end{proof}

Applying (6.9) to $\mu=\varnothing$, we obtain the second equality in (6.7).

\item \textbf{page 18, (6.8):} The matrix on the right hand side of (6.8)
should be understood as%
\[
\left(  h_{\lambda_{i}-\mu_{j}-i+j}\left(  x\mid\tau^{\mu_{j}-j+1}a\right)
\right)  _{i,j\in\left[  m\right]  },
\]
where $m$ is an integer satisfying $\ell\left(  \lambda\right)  \leq m$ and
$\ell\left(  \mu\right)  \leq m$. Note that different choices of $m$ are
possible, and they lead to different matrices, but all these matrices have the
same determinant (because if $\lambda_{m}=0$ and $\mu_{m}=0$, then the last
row of the matrix $\left(  h_{\lambda_{i}-\mu_{j}-i+j}\left(  x\mid\tau
^{\mu_{j}-j+1}a\right)  \right)  _{i,j\in\left[  m\right]  }$ is $\left(
0,0,\ldots,0,1\right)  $, and therefore Laplace expansion along this row
reduces the determinant to the determinant of the submatrix $\left(
h_{\lambda_{i}-\mu_{j}-i+j}\left(  x\mid\tau^{\mu_{j}-j+1}a\right)  \right)
_{i,j\in\left[  m-1\right]  }$), so that the definition (6.8) of
$s_{\lambda/\mu}\left(  x\mid a\right)  $ is still unambiguous.

\item \textbf{page 18, (6.9):} See one of the above comments for a proof of (6.9).

\item \textbf{page 18, after (6.9):} It is worth pointing out that (unlike the
polynomial $s_{\lambda}\left(  x\mid a\right)  $) the polynomial
$s_{\lambda/\mu}\left(  x\mid a\right)  $ for a skew partition $\lambda/\mu$
can depend not only on the \textquotedblleft positive\textquotedblright%
\ $a$-variables $a_{1},a_{2},a_{3},\ldots$ but also on some of the
\textquotedblleft nonpositive\textquotedblright\ $a$-variables $a_{0}%
,a_{-1},a_{-2},\ldots$.

Indeed, if we shift the skew Young diagram $\lambda/\mu$ by one step downwards
(i.e., if we replace the partitions $\lambda$ and $\mu$ by $\left(
N,\lambda_{1},\lambda_{2},\lambda_{3},\ldots\right)  $ and $\left(  N,\mu
_{1},\mu_{2},\mu_{3},\ldots\right)  $, where $N\geq\lambda_{1}$ is an
integer), then the polynomial $s_{\lambda/\mu}\left(  x\mid a\right)  $
undergoes the substitution $\tau^{-1}$ (that is, each $a$-variable $a_{i}$ is
replaced by $a_{i-1}$). This is easy to see from the formula (6.9) (indeed,
our shift is simply replacing each entry $e_{\lambda_{i}^{\prime}-\mu
_{j}^{\prime}-i+j}\left(  x\mid\tau^{-\mu_{j}^{\prime}+j-1}a\right)  $ of the
matrix by
\[
e_{\left(  \lambda_{i}^{\prime}+1\right)  -\left(  \mu_{j}^{\prime}+1\right)
-i+j}\left(  x\mid\tau^{-\left(  \mu_{j}^{\prime}+1\right)  +j-1}a\right)
=e_{\lambda_{i}^{\prime}-\mu_{j}^{\prime}-i+j}\left(  x\mid\tau^{-\mu
_{j}^{\prime}-1+j-1}a\right)  ,
\]
which is the same entry but with $a$ replaced by $\tau^{-1}a$). Thus, by
repeatedly performing this shift, we can \textquotedblleft drive all the
$a$-variables into the negative region\textquotedblright.

Conversely, if we shift $\lambda/\mu$ by $k$ steps to the right (i.e., if we
replace the partitions $\lambda$ and $\mu$ by $\left(  \lambda_{1}%
+k,\ \lambda_{2}+k,\ \ldots,\ \lambda_{m}+k\right)  $ and $\left(  \mu
_{1}+k,\ \mu_{2}+k,\ \ldots,\ \mu_{m}+k\right)  $, where $m\in\mathbb{N}$ is
such that $\ell\left(  \lambda\right)  \leq m$ and $\ell\left(  \mu\right)
\leq m$), then the polynomial $s_{\lambda/\mu}\left(  x\mid a\right)  $
undergoes the substitution $\tau^{k}$ (that is, each $a$-variable $a_{i}$ is
replaced by $a_{i+k}$). This is easiest to see using (6.8).

\item \textbf{page 18, (6.10):} Let me prove (6.10) here. We will need a
simple lemma about determinants.

If $p$ and $q$ are two integers such that $p\leq q+1$, then $\left[
p,q\right]  $ shall denote the set of all integers $m$ such that $p\leq m\leq
q$. We call this set an \emph{integer interval}. It has size $q-p+1$ (so it is
empty if $p=q+1$).

For a given $r\in\mathbb{N}$, we shall denote the integer interval $\left[
1,r\right]  =\left\{  1,2,\ldots,r\right\}  $ by $\left[  r\right]  $.

Our lemma says the following:

\begin{lemmain}
\label{lem.A7}Let $\left(  u_{i,j}\right)  _{i,j\in\left[  p\right]  }$ be a
$p\times p$-matrix for some $p\in\mathbb{N}$. Let $U$ and $V$ be two subsets
of $\left[  p\right]  $ satisfying $\left\vert U\right\vert +\left\vert
V\right\vert >p$. Assume that%
\begin{equation}
u_{i,j}=0\ \ \ \ \ \ \ \ \ \ \text{for all }i\in U\text{ and }j\in V.
\label{p18.lemA7.ass}%
\end{equation}
Then, $\det\left(  u_{i,j}\right)  _{i,j\in\left[  p\right]  }=0$.
\end{lemmain}

\begin{proof}
[Proof of Lemma \ref{lem.A7}.]By the definition of a determinant, we have
\begin{equation}
\det\left(  u_{i,j}\right)  _{i,j\in\left[  p\right]  }=\sum_{\sigma\in S_{p}%
}\operatorname*{sign}\sigma\cdot\prod_{i=1}^{p}u_{i,\sigma\left(  i\right)  }
\label{p18.lemA7.pf.0}%
\end{equation}
(where $S_{p}$ is the symmetric group of all permutations of $\left[
p\right]  $). Now we shall show that each $\sigma\in S_{p}$ satisfies%
\begin{equation}
\prod_{i=1}^{p}u_{i,\sigma\left(  i\right)  }=0. \label{p18.lemA7.pf.1}%
\end{equation}


(\textit{Proof:} Let $\sigma\in S_{p}$. Then, the map $\sigma$ is a
permutation, so that $\left\vert \sigma\left(  U\right)  \right\vert
=\left\vert U\right\vert >p-\left\vert V\right\vert $ (since $\left\vert
U\right\vert +\left\vert V\right\vert >p$), so that $\sigma\left(  U\right)
\not \subseteq \left[  p\right]  \setminus V$ (since $\sigma\left(  U\right)
\subseteq\left[  p\right]  \setminus V$ would entail $\left\vert \sigma\left(
U\right)  \right\vert \leq\left\vert \left[  p\right]  \setminus V\right\vert
=p-\left\vert V\right\vert $, contradicting $\left\vert \sigma\left(
U\right)  \right\vert >p-\left\vert V\right\vert $). In other words, there
exists some $j\in\sigma\left(  U\right)  $ such that $j\notin\left[  p\right]
\setminus V$. Consider this $j$. From $j\in\sigma\left(  U\right)  $, we
obtain $j=\sigma\left(  k\right)  $ for some $k\in U$. Consider this $k$. Now,
$k\in U$ and $j\in V$ (since $j\in\left[  p\right]  $ but $j\notin\left[
p\right]  \setminus V$). Hence, (\ref{p18.lemA7.ass}) (applied to $i=k$)
yields $u_{k,j}=0$. In other words, $u_{k,\sigma\left(  k\right)  }=0$ (since
$j=\sigma\left(  k\right)  $). Thus, one of the factors of the product
$\prod_{i=1}^{p}u_{i,\sigma\left(  i\right)  }$ is $0$ (namely, the $k$-th
factor). Hence, this whole product is $0$. This proves (\ref{p18.lemA7.pf.1}).)

Now, (\ref{p18.lemA7.pf.0}) becomes%
\[
\det\left(  u_{i,j}\right)  _{i,j\in\left[  p\right]  }=\sum_{\sigma\in S_{p}%
}\operatorname*{sign}\sigma\cdot\underbrace{\prod_{i=1}^{p}u_{i,\sigma\left(
i\right)  }}_{\substack{=0\\\text{(by (\ref{p18.lemA7.pf.1}))}}}=0.
\]
This proves Lemma \ref{lem.A7}.
\end{proof}

We can now prove (6.10):

\begin{proof}
[Proof of (6.10).]Assume that we don't have $0\leq\lambda_{i}^{\prime}-\mu
_{i}^{\prime}\leq n$ for all $i$. We must prove that $s_{\lambda/\mu}\left(
x\mid a\right)  =0$.

Write the partitions $\lambda^{\prime}$ and $\mu^{\prime}$ as $\lambda
^{\prime}=\left(  \lambda_{1}^{\prime},\lambda_{2}^{\prime},\ldots,\lambda
_{p}^{\prime}\right)  $ and $\mu^{\prime}=\left(  \mu_{1}^{\prime},\mu
_{2}^{\prime},\ldots,\mu_{p}^{\prime}\right)  $. Then, all $i>p$ satisfy
$\lambda_{i}^{\prime}=0$ and $\mu_{i}^{\prime}=0$ and therefore $\lambda
_{i}^{\prime}-\mu_{i}^{\prime}=0-0=0$. Note also that $\lambda_{1}^{\prime
}\geq\lambda_{2}^{\prime}\geq\lambda_{3}^{\prime}\geq\cdots$ (since
$\lambda^{\prime}$ is a partition) and $\mu_{1}^{\prime}\geq\mu_{2}^{\prime
}\geq\mu_{3}^{\prime}\geq\cdots$ (similarly). From (6.9), we obtain%
\begin{equation}
s_{\lambda/\mu}\left(  x\mid a\right)  =\det\left(  e_{\lambda_{i}^{\prime
}-\mu_{j}^{\prime}-i+j}\left(  x\mid\tau^{-\mu_{j}^{\prime}+j-1}a\right)
\right)  _{i,j\in\left[  p\right]  }. \label{p18.6.10.pf.1}%
\end{equation}


We have assumed that we don't have $0\leq\lambda_{i}^{\prime}-\mu_{i}^{\prime
}\leq n$ for all $i$. In other words, there exists some $i\geq1$ such that we
don't have $0\leq\lambda_{i}^{\prime}-\mu_{i}^{\prime}\leq n$. Consider this
$i$, and denote it by $k$. Thus, we don't have $0\leq\lambda_{k}^{\prime}%
-\mu_{k}^{\prime}\leq n$. Hence, we have either $\lambda_{k}^{\prime}-\mu
_{k}^{\prime}<0$ or $\lambda_{k}^{\prime}-\mu_{k}^{\prime}>n$. We are thus in
one of the following two cases:

\textit{Case 1:} We have $\lambda_{k}^{\prime}-\mu_{k}^{\prime}<0$.

\textit{Case 2:} We have $\lambda_{k}^{\prime}-\mu_{k}^{\prime}>n$.

Consider Case 1 first. In this case, we have $\lambda_{k}^{\prime}-\mu
_{k}^{\prime}<0$. Hence, we cannot have $k>p$ (since all $i>p$ satisfy
$\lambda_{i}^{\prime}-\mu_{i}^{\prime}=0$, which would yield $\lambda
_{k}^{\prime}-\mu_{k}^{\prime}=0$ if we had $k>p$). Thus, $k\leq p$, so that
$k\in\left[  p\right]  $.

Now, let $i\in\left[  k,p\right]  $ and $j\in\left[  k\right]  $. Then, $i\geq
k$ (since $i\in\left[  k,p\right]  $), so that $k\leq i$ and thus $\lambda
_{k}^{\prime}\geq\lambda_{i}^{\prime}$ (since $\lambda_{1}^{\prime}\geq
\lambda_{2}^{\prime}\geq\lambda_{3}^{\prime}\geq\cdots$). Thus, $\lambda
_{i}^{\prime}\leq\lambda_{k}^{\prime}$. Furthermore, $j\leq k$ (since
$j\in\left[  k\right]  $), so that $\mu_{j}^{\prime}\geq\mu_{k}^{\prime}$
(since $\mu_{1}^{\prime}\geq\mu_{2}^{\prime}\geq\mu_{3}^{\prime}\geq\cdots$).
Thus,%
\[
\underbrace{\lambda_{i}^{\prime}}_{\leq\lambda_{k}^{\prime}}-\underbrace{\mu
_{j}^{\prime}}_{\geq\mu_{k}^{\prime}}-\underbrace{i}_{\geq k}+\underbrace{j}%
_{\leq k}\leq\lambda_{k}^{\prime}-\mu_{k}^{\prime}-k+k=\lambda_{k}^{\prime
}-\mu_{k}^{\prime}<0,
\]
and therefore $e_{\lambda_{i}^{\prime}-\mu_{j}^{\prime}-i+j}\left(  x\mid
\tau^{-\mu_{j}^{\prime}+j-1}a\right)  =0$ (since $e_{\ell}\left(  x\mid
\tau^{-\mu_{j}^{\prime}+j-1}a\right)  =0$ for any $\ell<0$).

Forget that we fixed $i$ and $j$. We thus have shown that
\[
e_{\lambda_{i}^{\prime}-\mu_{j}^{\prime}-i+j}\left(  x\mid\tau^{-\mu
_{j}^{\prime}+j-1}a\right)  =0\ \ \ \ \ \ \ \ \ \ \text{for all }i\in\left[
k,p\right]  \text{ and }j\in\left[  k\right]  .
\]
Hence, Lemma \ref{lem.A7} (applied to $u_{i,j}=e_{\lambda_{i}^{\prime}-\mu
_{j}^{\prime}-i+j}\left(  x\mid\tau^{-\mu_{j}^{\prime}+j-1}a\right)  $) yields
that%
\[
\det\left(  e_{\lambda_{i}^{\prime}-\mu_{j}^{\prime}-i+j}\left(  x\mid
\tau^{-\mu_{j}^{\prime}+j-1}a\right)  \right)  _{i,j\in\left[  p\right]  }=0
\]
(since $\underbrace{\left\vert \left[  k,p\right]  \right\vert }%
_{=p-k+1}+\underbrace{\left\vert \left[  k\right]  \right\vert }_{=k}=\left(
p-k+1\right)  +k=p+1>p$). Hence, (\ref{p18.6.10.pf.1}) rewrites as
$s_{\lambda/\mu}\left(  x\mid a\right)  =0$. This proves (6.10) in Case 1.

Let us now consider Case 2. In this case, we have $\lambda_{k}^{\prime}%
-\mu_{k}^{\prime}>n$. Hence, we cannot have $k>p$ (since all $i>p$ satisfy
$\lambda_{i}^{\prime}-\mu_{i}^{\prime}=0$, which would yield $\lambda
_{k}^{\prime}-\mu_{k}^{\prime}=0\leq n$ if we had $k>p$). Thus, $k\leq p$, so
that $k\in\left[  p\right]  $.

Now, let $i\in\left[  k\right]  $ and $j\in\left[  k,p\right]  $. Then, $j\geq
k$ (since $j\in\left[  k,p\right]  $), so that $k\leq j$ and thus $\mu
_{k}^{\prime}\geq\mu_{j}^{\prime}$ (since $\mu_{1}^{\prime}\geq\mu_{2}%
^{\prime}\geq\mu_{3}^{\prime}\geq\cdots$). Hence, $\mu_{j}^{\prime}\leq\mu
_{k}^{\prime}$. Furthermore, $i\leq k$ (since $i\in\left[  k\right]  $), so
that $\lambda_{i}^{\prime}\geq\lambda_{k}^{\prime}$ (since $\lambda
_{1}^{\prime}\geq\lambda_{2}^{\prime}\geq\lambda_{3}^{\prime}\geq\cdots$).
Thus,%
\[
\underbrace{\lambda_{i}^{\prime}}_{\geq\lambda_{k}^{\prime}}-\underbrace{\mu
_{j}^{\prime}}_{\leq\mu_{k}^{\prime}}-\underbrace{i}_{\leq k}+\underbrace{j}%
_{\geq k}\geq\lambda_{k}^{\prime}-\mu_{k}^{\prime}-k+k=\lambda_{k}^{\prime
}-\mu_{k}^{\prime}>n,
\]
and therefore $e_{\lambda_{i}^{\prime}-\mu_{j}^{\prime}-i+j}\left(  x\mid
\tau^{-\mu_{j}^{\prime}+j-1}a\right)  =0$ (since $e_{\ell}\left(  x\mid
\tau^{-\mu_{j}^{\prime}+j-1}a\right)  =0$ for any $\ell>n$).

Forget that we fixed $i$ and $j$. We thus have shown that
\[
e_{\lambda_{i}^{\prime}-\mu_{j}^{\prime}-i+j}\left(  x\mid\tau^{-\mu
_{j}^{\prime}+j-1}a\right)  =0\ \ \ \ \ \ \ \ \ \ \text{for all }i\in\left[
k\right]  \text{ and }j\in\left[  k,p\right]  .
\]
Hence, Lemma \ref{lem.A7} (applied to $u_{i,j}=e_{\lambda_{i}^{\prime}-\mu
_{j}^{\prime}-i+j}\left(  x\mid\tau^{-\mu_{j}^{\prime}+j-1}a\right)  $) yields
that%
\[
\det\left(  e_{\lambda_{i}^{\prime}-\mu_{j}^{\prime}-i+j}\left(  x\mid
\tau^{-\mu_{j}^{\prime}+j-1}a\right)  \right)  _{i,j\in\left[  p\right]  }=0
\]
(since $\underbrace{\left\vert \left[  k\right]  \right\vert }_{=k}%
+\underbrace{\left\vert \left[  k,p\right]  \right\vert }_{=p-k+1}=k+\left(
p-k+1\right)  =p+1>p$). Hence, (\ref{p18.6.10.pf.1}) rewrites as
$s_{\lambda/\mu}\left(  x\mid a\right)  =0$. This proves (6.10) in Case 2.

We have now proved (6.10) in both Cases 1 and 2; this completes the proof of (6.10).
\end{proof}

\item \textbf{page 19, proof of (6.12):} In the last displayed equation of the
proof, replace \textquotedblleft$e_{j-k}\left(  y\mid\tau^{n+j}a\right)
$\textquotedblright\ by \textquotedblleft$e_{j-i}\left(  y\mid\tau
^{n+j}a\right)  $\textquotedblright.

\item \textbf{page 19, proof of (6.13):} Let me explain in some more detail
how the equality
\[
\bigwedge\nolimits^{r}\mathbf{H}\left(  x,y\mid a\right)  =\bigwedge
\nolimits^{r}\mathbf{H}\left(  x\mid a\right)  .\bigwedge\nolimits^{r}%
\mathbf{H}\left(  y\mid\tau^{n}a\right)
\]
yields the claim (6.13).

Indeed, a well-known corollary of the Cauchy--Binet theorem (specifically,
\cite[Corollary 7.182]{detnotes}, or rather its version for infinite
matrices\footnote{The version for infinite matrices is proved in the same way
as the version for finite matrices, as long as (formal) convergence is taken
care of (and that is easy when the matrices in question are
upper-triangular).}) yields%
\begin{align*}
&  \det\left(  \operatorname*{sub}\nolimits_{\left\{  \mu_{1}-1,\mu
_{2}-2,\ldots,\mu_{r}-r\right\}  }^{\left\{  \lambda_{1}-1,\lambda
_{2}-2,\ldots,\lambda_{r}-r\right\}  }\left(  \mathbf{H}\left(  x\mid
a\right)  \cdot\mathbf{H}\left(  y\mid\tau^{n}a\right)  \right)  \right) \\
&  =\sum_{g_{1}<g_{2}<\cdots<g_{r}}\det\left(  \operatorname*{sub}%
\nolimits_{\left\{  \mu_{1}-1,\mu_{2}-2,\ldots,\mu_{r}-r\right\}  }^{\left\{
g_{1},g_{2},\ldots,g_{r}\right\}  }\left(  \mathbf{H}\left(  x\mid a\right)
\right)  \right) \\
&  \ \ \ \ \ \ \ \ \ \ \ \ \ \ \ \ \ \ \ \ \cdot\det\left(
\operatorname*{sub}\nolimits_{\left\{  g_{1},g_{2},\ldots,g_{r}\right\}
}^{\left\{  \lambda_{1}-1,\lambda_{2}-2,\ldots,\lambda_{r}-r\right\}  }\left(
\mathbf{H}\left(  y\mid\tau^{n}a\right)  \right)  \right)  ,
\end{align*}
where the sum ranges over all strictly increasing $r$-tuples $\left(
g_{1}<g_{2}<\cdots<g_{r}\right)  $ of integers. Using (6.12)(ii), we can
rewrite this as%
\begin{align}
&  \det\left(  \operatorname*{sub}\nolimits_{\left\{  \mu_{1}-1,\mu
_{2}-2,\ldots,\mu_{r}-r\right\}  }^{\left\{  \lambda_{1}-1,\lambda
_{2}-2,\ldots,\lambda_{r}-r\right\}  }\left(  \mathbf{H}\left(  xy\mid
a\right)  \right)  \right) \nonumber\\
&  =\sum_{g_{1}<g_{2}<\cdots<g_{r}}\det\left(  \operatorname*{sub}%
\nolimits_{\left\{  \mu_{1}-1,\mu_{2}-2,\ldots,\mu_{r}-r\right\}  }^{\left\{
g_{1},g_{2},\ldots,g_{r}\right\}  }\left(  \mathbf{H}\left(  x\mid a\right)
\right)  \right) \nonumber\\
&  \ \ \ \ \ \ \ \ \ \ \ \ \ \ \ \ \ \ \ \ \cdot\det\left(
\operatorname*{sub}\nolimits_{\left\{  g_{1},g_{2},\ldots,g_{r}\right\}
}^{\left\{  \lambda_{1}-1,\lambda_{2}-2,\ldots,\lambda_{r}-r\right\}  }\left(
\mathbf{H}\left(  y\mid\tau^{n}a\right)  \right)  \right) \nonumber\\
&  =\sum_{\substack{S\text{ is an }r\text{-element}\\\text{set of integers}%
}}\det\left(  \operatorname*{sub}\nolimits_{\left\{  \mu_{1}-1,\mu
_{2}-2,\ldots,\mu_{r}-r\right\}  }^{S}\left(  \mathbf{H}\left(  x\mid
a\right)  \right)  \right) \nonumber\\
&  \ \ \ \ \ \ \ \ \ \ \ \ \ \ \ \ \ \ \ \ \cdot\det\left(
\operatorname*{sub}\nolimits_{S}^{\left\{  \lambda_{1}-1,\lambda_{2}%
-2,\ldots,\lambda_{r}-r\right\}  }\left(  \mathbf{H}\left(  y\mid\tau
^{n}a\right)  \right)  \right)
\ \ \ \ \ \ \ \ \ \ \ \ \ \ \ \ \ \ \label{p19.6.13.pf.2}%
\end{align}
(here, we have substituted $S$ for the set $\left\{  g_{1},g_{2},\ldots
,g_{r}\right\}  $).

Now, let $\mathbb{Z}_{\geq-r}$ be the set of all integers $\geq-r$. We observe
the following:

\begin{itemize}
\item If $S$ is an $r$-element set of integers that does not satisfy
$S\subseteq\mathbb{Z}_{\geq-r}$, then%
\begin{equation}
\det\left(  \operatorname*{sub}\nolimits_{\left\{  \mu_{1}-1,\mu_{2}%
-2,\ldots,\mu_{r}-r\right\}  }^{S}\left(  \mathbf{H}\left(  x\mid a\right)
\right)  \right)  =0. \label{p19.6.13.pf.3a}%
\end{equation}


(\textit{Proof:} Let $S$ be an $r$-element set of integers that does not
satisfy $S\subseteq\mathbb{Z}_{\geq-r}$. Write this set $S$ as $S=\left\{
s_{1}<s_{2}<\cdots<s_{r}\right\}  $. Then, $s_{1}<-r$ (since otherwise, we
would have $s_{1}\geq-r$ and therefore $-r\leq s_{1}<s_{2}<\cdots<s_{r}$, so
that all the elements $s_{1},s_{2},\ldots,s_{r}$ of $S$ would be $\geq-r$, and
thus we would have $S\subseteq\mathbb{Z}_{\geq-r}$, which would contradict the
assumption that $S$ does not satisfy $S\subseteq\mathbb{Z}_{\geq-r}$).
However, since $\mu$ is a partition, we have $\mu_{1}\geq\mu_{2}\geq\cdots
\geq\mu_{r}$ and thus $\mu_{1}-1>\mu_{2}-2>\cdots>\mu_{r}-r$. Thus,
\[
\left\{  \mu_{1}-1,\mu_{2}-2,\ldots,\mu_{r}-r\right\}  =\left\{  \mu_{r}%
-r<\mu_{r-1}-\left(  r-1\right)  <\cdots<\mu_{1}-1\right\}  .
\]
Using this equality as well as $S=\left\{  s_{1}<s_{2}<\cdots<s_{r}\right\}
$, we obtain%
\begin{align}
&  \operatorname*{sub}\nolimits_{\left\{  \mu_{1}-1,\mu_{2}-2,\ldots,\mu
_{r}-r\right\}  }^{S}\left(  \mathbf{H}\left(  x\mid a\right)  \right)
\nonumber\\
&  =\operatorname*{sub}\nolimits_{\left\{  \mu_{r}-r<\mu_{r-1}-\left(
r-1\right)  <\cdots<\mu_{1}-1\right\}  }^{\left\{  s_{1}<s_{2}<\cdots
<s_{r}\right\}  }\left(  \mathbf{H}\left(  x\mid a\right)  \right) \nonumber\\
&  =\left(  \left(  \mathbf{H}\left(  x\mid a\right)  \right)  _{\mu
_{r+1-i}-\left(  r+1-i\right)  ,\ s_{j}}\right)  _{i,j\in\left[  r\right]  },
\label{p19.6.13.pf.3a.pf.1}%
\end{align}
where $\left(  \mathbf{H}\left(  x\mid a\right)  \right)  _{p,q}$ denotes the
$\left(  p,q\right)  $-th entry of the matrix $\mathbf{H}\left(  x\mid
a\right)  $. However, each $i\in\left[  r\right]  $ satisfies
\[
\underbrace{\mu_{r+1-i}}_{\geq0}-\left(  r+1-i\right)  \geq-\left(
r+1-i\right)  =\underbrace{i}_{\geq1}-\,r-1\geq1-r-1=-r>s_{1}%
\]
(since $s_{1}<-r$) and therefore $\left(  \mathbf{H}\left(  x\mid a\right)
\right)  _{\mu_{r+1-i}-\left(  r+1-i\right)  ,\ s_{1}}=0$ (since the matrix
$\mathbf{H}\left(  x\mid a\right)  $ is upper-triangular and thus satisfies
$\left(  \mathbf{H}\left(  x\mid a\right)  \right)  _{p,q}=0$ whenever $p>q$).
In other words, all entries in the $1$-st row of the matrix $\left(  \left(
\mathbf{H}\left(  x\mid a\right)  \right)  _{\mu_{r+1-i}-\left(  r+1-i\right)
,\ s_{j}}\right)  _{i,j\in\left[  r\right]  }$ are $0$. In view of
(\ref{p19.6.13.pf.3a.pf.1}), we can restate this as follows: All entries in
the $1$-st row of the matrix $\operatorname*{sub}\nolimits_{\left\{  \mu
_{1}-1,\mu_{2}-2,\ldots,\mu_{r}-r\right\}  }^{S}\left(  \mathbf{H}\left(
x\mid a\right)  \right)  $ are $0$. Thus, the determinant of this matrix is
$0$ as well. This proves (\ref{p19.6.13.pf.3a}).)
\end{itemize}

The observation (\ref{p19.6.13.pf.3a}) allows us to rewrite the sum on the
right hand side of (\ref{p19.6.13.pf.2}), by removing all the addends
corresponding to sets $S$ that do not satisfy $S\subseteq\mathbb{Z}_{\geq-r}$
(since (\ref{p19.6.13.pf.3a}) shows that all these addends are $0$). The
surviving addends are the ones corresponding to sets $S$ that do satisfy
$S\subseteq\mathbb{Z}_{\geq-r}$. In other words, they are the ones
corresponding to the $r$-element subsets $S$ of $\mathbb{Z}_{\geq-r}$. Hence,
(\ref{p19.6.13.pf.2}) can be rewritten as%
\begin{align}
&  \det\left(  \operatorname*{sub}\nolimits_{\left\{  \mu_{1}-1,\mu
_{2}-2,\ldots,\mu_{r}-r\right\}  }^{\left\{  \lambda_{1}-1,\lambda
_{2}-2,\ldots,\lambda_{r}-r\right\}  }\left(  \mathbf{H}\left(  xy\mid
a\right)  \right)  \right) \nonumber\\
&  =\sum_{\substack{S\text{ is an }r\text{-element}\\\text{subset of
}\mathbb{Z}_{\geq-r}}}\det\left(  \operatorname*{sub}\nolimits_{\left\{
\mu_{1}-1,\mu_{2}-2,\ldots,\mu_{r}-r\right\}  }^{S}\left(  \mathbf{H}\left(
x\mid a\right)  \right)  \right) \nonumber\\
&  \ \ \ \ \ \ \ \ \ \ \ \ \ \ \ \ \ \ \ \ \cdot\det\left(
\operatorname*{sub}\nolimits_{S}^{\left\{  \lambda_{1}-1,\lambda_{2}%
-2,\ldots,\lambda_{r}-r\right\}  }\left(  \mathbf{H}\left(  y\mid\tau
^{n}a\right)  \right)  \right)
.\ \ \ \ \ \ \ \ \ \ \ \ \ \label{p19.6.13.pf.3b}%
\end{align}


Now, we observe the following:

\begin{itemize}
\item If $\nu=\left(  \nu_{1},\nu_{2},\ldots,\nu_{r}\right)  $ is a partition
of length $\leq r$, then we have $\nu_{1}\geq\nu_{2}\geq\cdots\geq\nu_{r}%
\geq0$ and therefore $\nu_{1}-1>\nu_{2}-2>\cdots>\nu_{r}-r\geq0-r=-r$, so that
$\left\{  \nu_{1}-1>\nu_{2}-2>\cdots>\nu_{r}-r\right\}  $ is an $r$-element
subset of $\mathbb{Z}_{\geq-r}$. Thus, we can define a map%
\begin{align}
\left\{  \text{partitions of length }\leq r\right\}   &  \rightarrow\left\{
r\text{-element subsets of }\mathbb{Z}_{\geq-r}\right\}  ,\nonumber\\
\nu &  \mapsto\left\{  \nu_{1}-1>\nu_{2}-2>\cdots>\nu_{r}-r\right\}  .
\label{p19.6.13.pf.bij}%
\end{align}
In the opposite direction, we can define a map%
\begin{align*}
\left\{  r\text{-element subsets of }\mathbb{Z}_{\geq-r}\right\}   &
\rightarrow\left\{  \text{partitions of length }\leq r\right\}  ,\\
\left\{  h_{1}>h_{2}>\cdots>h_{r}\right\}   &  \mapsto\left(  h_{1}%
+1,\ h_{2}+2,\ \ldots,\ h_{r}+r\right)
\end{align*}
(this is well-defined because if $\left\{  h_{1}>h_{2}>\cdots>h_{r}\right\}  $
is an $r$-element subset of $\mathbb{Z}_{\geq-r}$, then $h_{1}>h_{2}%
>\cdots>h_{r}\geq-r$ and therefore $h_{1}+1\geq h_{2}+2\geq\cdots\geq
h_{r}+r\geq\left(  -r\right)  +r=0$, so that $\left(  h_{1}+1,\ h_{2}%
+2,\ \ldots,\ h_{r}+r\right)  $ is a partition of length $\leq r$). It is easy
to see that these two maps are mutually inverse, and thus are bijections. In
particular, the map (\ref{p19.6.13.pf.bij}) is thus a bijection. In other
words, the map%
\begin{align}
\left\{  \text{partitions of length }\leq r\right\}   &  \rightarrow\left\{
r\text{-element subsets of }\mathbb{Z}_{\geq-r}\right\}  ,\nonumber\\
\nu &  \mapsto\left\{  \nu_{1}-1,\ \nu_{2}-2,\ \ldots,\ \nu_{r}-r\right\}
\ \ \ \ \ \ \ \ \ \label{p19.6.13.pf.bij2}%
\end{align}
is a bijection.
\end{itemize}

Now, (\ref{p19.6.13.pf.3b}) becomes
\begin{align}
&  \det\left(  \operatorname*{sub}\nolimits_{\left\{  \mu_{1}-1,\mu
_{2}-2,\ldots,\mu_{r}-r\right\}  }^{\left\{  \lambda_{1}-1,\lambda
_{2}-2,\ldots,\lambda_{r}-r\right\}  }\left(  \mathbf{H}\left(  xy\mid
a\right)  \right)  \right) \nonumber\\
&  =\sum_{\substack{S\text{ is an }r\text{-element}\\\text{subset of
}\mathbb{Z}_{\geq-r}}}\det\left(  \operatorname*{sub}\nolimits_{\left\{
\mu_{1}-1,\mu_{2}-2,\ldots,\mu_{r}-r\right\}  }^{S}\left(  \mathbf{H}\left(
x\mid a\right)  \right)  \right) \nonumber\\
&  \ \ \ \ \ \ \ \ \ \ \ \ \ \ \ \ \ \ \ \ \cdot\det\left(
\operatorname*{sub}\nolimits_{S}^{\left\{  \lambda_{1}-1,\lambda_{2}%
-2,\ldots,\lambda_{r}-r\right\}  }\left(  \mathbf{H}\left(  y\mid\tau
^{n}a\right)  \right)  \right) \nonumber\\
&  =\sum_{\substack{\nu\text{ is a partition}\\\text{of length }\leq r}%
}\det\left(  \operatorname*{sub}\nolimits_{\left\{  \mu_{1}-1,\mu_{2}%
-2,\ldots,\mu_{r}-r\right\}  }^{\left\{  \nu_{1}-1,\nu_{2}-2,\ldots,\nu
_{r}-r\right\}  }\left(  \mathbf{H}\left(  x\mid a\right)  \right)  \right)
\nonumber\\
&  \ \ \ \ \ \ \ \ \ \ \ \ \ \ \ \ \ \ \ \ \cdot\det\left(
\operatorname*{sub}\nolimits_{\left\{  \nu_{1}-1,\nu_{2}-2,\ldots,\nu
_{r}-r\right\}  }^{\left\{  \lambda_{1}-1,\lambda_{2}-2,\ldots,\lambda
_{r}-r\right\}  }\left(  \mathbf{H}\left(  y\mid\tau^{n}a\right)  \right)
\right)  \ \ \ \ \ \ \ \ \ \ \ \ \ \label{p19.6.13.pf.4}%
\end{align}
(here, we have substituted $\left\{  \nu_{1}-1,\nu_{2}-2,\ldots,\nu
_{r}-r\right\}  $ for $S$ in the sum, since the map (\ref{p19.6.13.pf.bij2})
is a bijection).

However, in the second sentence of the proof of (6.13), it was said that%
\[
s_{\lambda/\mu}\left(  x,y\mid a\right)  =\det\left(  \operatorname*{sub}%
\nolimits_{\left\{  \mu_{1}-1,\mu_{2}-2,\ldots,\mu_{r}-r\right\}  }^{\left\{
\lambda_{1}-1,\lambda_{2}-2,\ldots,\lambda_{r}-r\right\}  }\left(
\mathbf{H}\left(  xy\mid a\right)  \right)  \right)  .
\]
Similarly, for any partition $\nu$ of length $\leq r$, we have%
\[
s_{\nu/\mu}\left(  x\mid a\right)  =\det\left(  \operatorname*{sub}%
\nolimits_{\left\{  \mu_{1}-1,\mu_{2}-2,\ldots,\mu_{r}-r\right\}  }^{\left\{
\nu_{1}-1,\nu_{2}-2,\ldots,\nu_{r}-r\right\}  }\left(  \mathbf{H}\left(  x\mid
a\right)  \right)  \right)
\]
and%
\[
s_{\lambda/\nu}\left(  y\mid\tau^{n}a\right)  =\det\left(  \operatorname*{sub}%
\nolimits_{\left\{  \nu_{1}-1,\nu_{2}-2,\ldots,\nu_{r}-r\right\}  }^{\left\{
\lambda_{1}-1,\lambda_{2}-2,\ldots,\lambda_{r}-r\right\}  }\left(
\mathbf{H}\left(  y\mid\tau^{n}a\right)  \right)  \right)  .
\]
In view of these three equalities, we can rewrite (\ref{p19.6.13.pf.4}) as%
\begin{align*}
s_{\lambda/\mu}\left(  x,y\mid a\right)   &  =\sum_{\substack{\nu\text{ is a
partition}\\\text{of length }\leq r}}s_{\nu/\mu}\left(  x\mid a\right)
s_{\lambda/\nu}\left(  y\mid\tau^{n}a\right) \\
&  =\sum_{\nu\text{ is a partition}}s_{\nu/\mu}\left(  x\mid a\right)
s_{\lambda/\nu}\left(  y\mid\tau^{n}a\right)
\end{align*}
(here, we have removed the condition \textquotedblleft of length $\leq
r$\textquotedblright\ from the sum; this does not change the sum, since all
newly introduced addends are zero\footnote{\textit{Proof.} Let $\nu$ is a
partition of length $>r$. We must show that $s_{\nu/\mu}\left(  x\mid
a\right)  s_{\lambda/\nu}\left(  y\mid\tau^{n}a\right)  =0$.
\par
Since $\nu$ has length $>r$, we have $l\left(  \nu\right)  >r$. Thus, $\nu
_{1}^{\prime}=l\left(  \nu\right)  >r\geq\max\left(  l\left(  \lambda\right)
,l\left(  \mu\right)  \right)  \geq l\left(  \lambda\right)  =\lambda
_{1}^{\prime}$ and therefore $0>\lambda_{1}^{\prime}-\nu_{1}^{\prime}$. Hence,
we don't have $0\leq\lambda_{1}^{\prime}-\nu_{1}^{\prime}\leq n$. Thus, we
don't have $0\leq\lambda_{i}^{\prime}-\nu_{i}^{\prime}\leq n$ for all $i$
(since this inequality fails for $i=1$). Thus, (6.10) (applied to $\nu$,
$\tau^{n}a$ and $y$ instead of $\mu$, $a$ and $x$) yields $s_{\lambda/\nu
}\left(  y\mid\tau^{n}a\right)  =0$. Therefore, $s_{\nu/\mu}\left(  x\mid
a\right)  \underbrace{s_{\lambda/\nu}\left(  y\mid\tau^{n}a\right)  }_{=0}=0$
as well.}). This proves (6.13).

\item \textbf{page 19, the paragraph containing (6.14):} Replace
\textquotedblleft Let $x^{\left(  i\right)  },\ \ldots,\ x^{\left(  n\right)
}$ be\textquotedblright\ by \textquotedblleft Let $x^{\left(  1\right)
},\ \ldots,\ x^{\left(  n\right)  }$ be\textquotedblright.

\item \textbf{page 19, the paragraph containing (6.14):} Replace
\textquotedblleft where $x^{\left(  i\right)  }=\left(  x_{1}^{\left(
1\right)  },\ldots,x_{r_{i}}^{\left(  i\right)  }\right)  $\textquotedblright%
\ by \textquotedblleft where $x^{\left(  i\right)  }=\left(  x_{1}^{\left(
i\right)  },\ldots,x_{r_{i}}^{\left(  i\right)  }\right)  $\textquotedblright.

\item \textbf{page 19, (6.14):} On the left hand side of (6.14), replace
\textquotedblleft$x^{\left(  i\right)  },\ \ldots,\ x^{\left(  n\right)  }%
$\textquotedblright\ by \textquotedblleft$x^{\left(  1\right)  }%
,\ \ldots,\ x^{\left(  n\right)  }$\textquotedblright.

\item \textbf{page 19, last paragraph:} Let me explain why
\begin{equation}
s_{\lambda/\mu}\left(  x\mid a\right)  =\prod_{i\geq1}h_{\lambda_{i}-\mu_{i}%
}\left(  x\mid\tau^{\mu_{i}-i+1}a\right)  \label{p19.last.diag}%
\end{equation}
when $\lambda-\mu$ is a horizontal strip:

\begin{proof}
[Proof of (\ref{p19.last.diag}).]Assume that $\lambda-\mu$ is a horizontal
strip. Write the partitions $\lambda$ and $\mu$ as $\lambda=\left(
\lambda_{1},\lambda_{2},\ldots,\lambda_{q}\right)  $ and $\mu=\left(  \mu
_{1},\mu_{2},\ldots,\mu_{q}\right)  $, where $q\in\mathbb{N}$ is sufficiently
large (namely, $q\geq\max\left(  \ell\left(  \lambda\right)  ,\ell\left(
\mu\right)  \right)  $). Then, for any $i,j\in\left[  q\right]  $ satisfying
$i>j$, we have%
\begin{equation}
\lambda_{i}-\mu_{j}-i+j<0. \label{p19.last.pf.1}%
\end{equation}


(\textit{Proof:} Let $i,j\in\left[  q\right]  $ satisfy $i>j$. Thus, $j<i$, so
that $\lambda_{j}\geq\lambda_{i}$ (since $\lambda_{1}\geq\lambda_{2}%
\geq\lambda_{3}\geq\cdots$) and $\mu_{j}\geq\mu_{i}$ (similarly). In other
words, $\lambda_{i}\leq\lambda_{j}$ and $\mu_{i}\leq\mu_{j}$.

The skew partition $\lambda-\mu$ is a horizontal strip, i.e., contains no two
cells in the same column. However, if we had $\mu_{j}<\lambda_{i}$, then the
two distinct cells $\left(  i,\lambda_{i}\right)  $ and $\left(  j,\lambda
_{i}\right)  $ would both belong to $\lambda-\mu$ (indeed, we would have
$\left(  i,\lambda_{i}\right)  \in\lambda-\mu$ because of $\mu_{i}\leq\mu
_{j}<\lambda_{i}\leq\lambda_{i}$, and we would have $\left(  j,\lambda
_{i}\right)  \in\lambda-\mu$ because of $\mu_{j}<\lambda_{i}\leq\lambda_{j}$),
which would contradict the preceding sentence (since these two cells clearly
lie in the same column). Thus, we cannot have $\mu_{j}<\lambda_{i}$. In other
words, we have $\mu_{j}\geq\lambda_{i}$. Hence, $\lambda_{i}-\underbrace{\mu
_{j}}_{\geq\lambda_{i}}-\underbrace{i}_{>j}+\,j<\lambda_{i}-\lambda_{i}%
-j+j=0$. This proves (\ref{p19.last.pf.1}).)

Now, for any $i,j\in\left[  q\right]  $ satisfying $i>j$, we have $\lambda
_{i}-\mu_{j}-i+j<0$ (by (\ref{p19.last.pf.1})) and therefore
\[
h_{\lambda_{i}-\mu_{j}-i+j}\left(  x\mid\tau^{\mu_{j}-j+1}a\right)  =0
\]
(since $h_{k}\left(  x\mid\tau^{\mu_{j}-j+1}a\right)  =0$ for all $k<0$). In
other words, the matrix $\left(  h_{\lambda_{i}-\mu_{j}-i+j}\left(  x\mid
\tau^{\mu_{j}-j+1}a\right)  \right)  _{i,j\in\left[  q\right]  }$ is
upper-triangular. Hence, its determinant is the product of its diagonal
entries. In other words,%
\begin{align*}
&  \det\left(  h_{\lambda_{i}-\mu_{j}-i+j}\left(  x\mid\tau^{\mu_{j}%
-j+1}a\right)  \right)  _{i,j\in\left[  q\right]  }\\
&  =\prod_{i=1}^{q}\underbrace{h_{\lambda_{i}-\mu_{i}-i+i}\left(  x\mid
\tau^{\mu_{i}-i+1}a\right)  }_{\substack{=h_{\lambda_{i}-\mu_{i}}\left(
x\mid\tau^{\mu_{i}-i+1}a\right)  \\\text{(since }\lambda_{i}-\mu
_{i}-i+i=\lambda_{i}-\mu_{i}\text{)}}}\\
&  =\prod_{i=1}^{q}h_{\lambda_{i}-\mu_{i}}\left(  x\mid\tau^{\mu_{i}%
-i+1}a\right)  .
\end{align*}
In view of (6.8), we can rewrite this as%
\begin{equation}
s_{\lambda/\mu}\left(  x\mid a\right)  =\prod_{i=1}^{q}h_{\lambda_{i}-\mu_{i}%
}\left(  x\mid\tau^{\mu_{i}-i+1}a\right)  . \label{p19.last.diag-1}%
\end{equation}
This is a finite product, but we can extend it to an infinite product over all
$i\geq1$; this will not change the value of the product, since all the newly
inserted factors $h_{\lambda_{i}-\mu_{i}}\left(  x\mid\tau^{\mu_{i}%
-i+1}a\right)  $ for $i>q$ will equal $1$ (because if $i>q$, then $\lambda
_{i}=0$ and $\mu_{i}=0$ and therefore $\lambda_{i}-\mu_{i}=0-0=0$, so that
$h_{\lambda_{i}-\mu_{i}}\left(  x\mid\tau^{\mu_{i}-i+1}a\right)  =h_{0}\left(
x\mid\tau^{\mu_{i}-i+1}a\right)  =1$). Hence, (\ref{p19.last.diag-1}) can be
rewritten as%
\[
s_{\lambda/\mu}\left(  x\mid a\right)  =\prod_{i\geq1}h_{\lambda_{i}-\mu_{i}%
}\left(  x\mid\tau^{\mu_{i}-i+1}a\right)  .
\]
This proves (\ref{p19.last.diag}).
\end{proof}

\item \textbf{page 19, last paragraph:} At the end of the last display on this
page, there is a period. This period should be a comma.

\item \textbf{page 20:} \textquotedblleft\textit{column strict}%
\textquotedblright\ should be \textquotedblleft\textit{column-strict}%
\textquotedblright.

\item \textbf{page 20, proof of (6.17):} Replace \textquotedblleft%
$A_{\delta_{m+n}}\left(  x,y\right)  $\textquotedblright\ by \textquotedblleft%
$A_{\delta_{m+n}}\left(  x,y\mid a\right)  $\textquotedblright. Likewise,
replace \textquotedblleft$A_{\delta_{n}}\left(  x\right)  $\textquotedblright%
\ by \textquotedblleft$A_{\delta_{n}}\left(  x\mid a\right)  $%
\textquotedblright. Likewise, replace \textquotedblleft$A_{\delta_{m}}\left(
y\right)  $\textquotedblright\ by \textquotedblleft$A_{\delta_{m}}\left(
y\mid a\right)  $\textquotedblright.

\item \textbf{page 20, proof of (6.17):} The equality%
\[
A_{\delta_{m+n}}\left(  x,y\right)  =\sum_{\lambda\subset\left(  m^{n}\right)
}\left(  -1\right)  ^{\left\vert \widehat{\lambda}\right\vert }A_{\lambda
+\delta_{n}}\left(  x\right)  A_{\widehat{\lambda}^{\prime}+\delta_{m}}\left(
y\right)
\]
should be
\begin{equation}
A_{\delta_{m+n}}\left(  x,y\mid a\right)  =\sum_{\lambda\subset\left(
m^{n}\right)  }\left(  -1\right)  ^{\left\vert \widehat{\lambda}\right\vert
}A_{\lambda+\delta_{n}}\left(  x\mid a\right)  A_{\widehat{\lambda}^{\prime
}+\delta_{m}}\left(  y\mid a\right)  \ \ \ \ \ \ \ \ \ \ \label{p20.claim}%
\end{equation}
instead.

\item \textbf{page 20, proof of (6.17):} Let me explain how the equality
(\ref{p20.claim}) is proved. Some lemmas will be needed for the proof. The
first lemma is known as the Laplace expansion formula along multiple rows:

\begin{lemmain}
\label{lem.lap-rows}Let $k\in\mathbb{N}$. Let $A$ be any $k\times k$-matrix.
For any subset $I$ of $\left[  k\right]  $, we let $\sum I$ denote the sum of
all elements of $I$, and we let $\widetilde{I}$ denote the complement $\left[
k\right]  \setminus I$ of $I$. (For instance, if $k=4$ and $I=\left\{
1,4\right\}  $, then $\sum I=1+4=5$ and $\widetilde{I}=\left\{  2,3\right\}  $.)

Let $P$ be a subset of $\left[  k\right]  $. Then,%
\[
\det A=\sum_{\substack{Q\subseteq\left[  k\right]  ;\\\left\vert Q\right\vert
=\left\vert P\right\vert }}\left(  -1\right)  ^{\sum P+\sum Q}\det\left(
\operatorname*{sub}\nolimits_{P}^{Q}A\right)  \det\left(  \operatorname*{sub}%
\nolimits_{\widetilde{P}}^{\widetilde{Q}}A\right)  .
\]

\end{lemmain}

\begin{proof}
[Proof of Lemma \ref{lem.lap-rows}.]Lemma \ref{lem.lap-rows} is \cite[Theorem
6.156]{detnotes}, applied to $n=k$. (Note that our $\operatorname*{sub}%
\nolimits_{P}^{Q}A$ is denoted $\operatorname*{sub}\nolimits_{w\left(
P\right)  }^{w\left(  Q\right)  }A$ in \cite{detnotes}.)
\end{proof}

Next, we state two combinatorial lemmas. Recall that the partitions
$\lambda=\left(  \lambda_{1},\lambda_{2},\ldots,\lambda_{n}\right)  $ that
satisfy $\lambda_{1}\leq m$ are precisely the partitions $\lambda$ that
satisfy $\lambda\subseteq\left(  m^{n}\right)  $ (where \textquotedblleft%
$\left(  m^{n}\right)  $ is to be read in exponential notation -- i.e., it
means an $n$-tuple of $m$'s). For each such partition $\lambda$, we can define
the set%
\[
Q_{\lambda}^{\left(  n,m\right)  }:=\left\{  m-\lambda_{i}+i\ \mid
\ i\in\left[  n\right]  \right\}  =\left\{  m-\lambda_{1}+1,\ m-\lambda
_{2}+2,\ \ldots,\ m-\lambda_{n}+n\right\}
\]
of integers. We first claim the following:

\begin{lemmain}
\label{lem.Qlam1}Let $n,m\in\mathbb{N}$. Then:

\begin{enumerate}
\item[\textbf{(a)}] For any partition $\lambda\subseteq\left(  m^{n}\right)  $
(that is, for any partition $\lambda$ that satisfies $\lambda\subseteq\left(
m^{n}\right)  $), we have
\[
Q_{\lambda}^{\left(  n,m\right)  }=\left\{  m-\lambda_{1}+1<m-\lambda
_{2}+2<\cdots<m-\lambda_{n}+n\right\}  .
\]


\item[\textbf{(b)}] For any partition $\lambda\subseteq\left(  m^{n}\right)
$, the set $Q_{\lambda}^{\left(  n,m\right)  }$ is an $n$-element subset of
$\left[  n+m\right]  $.

\item[\textbf{(c)}] The map%
\begin{align*}
\left\{  \text{partitions }\lambda\subseteq\left(  m^{n}\right)  \right\}   &
\rightarrow\left\{  n\text{-element subsets of }\left[  n+m\right]  \right\}
,\\
\lambda &  \mapsto Q_{\lambda}^{\left(  n,m\right)  }%
\end{align*}
is a bijection.
\end{enumerate}
\end{lemmain}

\begin{proof}
\textbf{(a)} Let $\lambda\subseteq\left(  m^{n}\right)  $ be a partition. We
must show that%
\[
Q_{\lambda}^{\left(  n,m\right)  }=\left\{  m-\lambda_{1}+1<m-\lambda
_{2}+2<\cdots<m-\lambda_{n}+n\right\}  .
\]
But the set $Q_{\lambda}^{\left(  n,m\right)  }$ is defined as $\left\{
m-\lambda_{1}+1,\ m-\lambda_{2}+2,\ \ldots,\ m-\lambda_{n}+n\right\}  $. Thus,
it remains to prove that
\[
m-\lambda_{1}+1<m-\lambda_{2}+2<\cdots<m-\lambda_{n}+n.
\]
In other words, it remains to prove that $m-\lambda_{i}+i<m-\lambda
_{i+1}+\left(  i+1\right)  $ for each $i\in\left[  n-1\right]  $.

So let us prove this. Let $i\in\left[  n-1\right]  $. Since $\lambda$ is a
partition, we have $\lambda_{i}\geq\lambda_{i+1}$. Hence,%
\[
m-\underbrace{\lambda_{i}}_{\geq\lambda_{i+1}}+\underbrace{i}_{<i+1}%
<m-\lambda_{i+1}+\left(  i+1\right)  .
\]
This is exactly what we needed to prove. Thus, Lemma \ref{lem.Qlam1}
\textbf{(a)} is proved. \medskip

\textbf{(b)} Let $\lambda\subseteq\left(  m^{n}\right)  $ be a partition.
Lemma \ref{lem.Qlam1} \textbf{(a)} shows that%
\[
Q_{\lambda}^{\left(  n,m\right)  }=\left\{  m-\lambda_{1}+1<m-\lambda
_{2}+2<\cdots<m-\lambda_{n}+n\right\}  .
\]
Hence, $Q_{\lambda}^{\left(  n,m\right)  }$ is an $n$-element set.

Let $i\in\left[  n\right]  $. Then, $\lambda_{i}\leq m$ (since $\lambda
\subseteq\left(  m^{n}\right)  $) and $\lambda_{i}\geq0$ (obviously) and
$i\geq1$ and $i\leq n$ (since $i\in\left[  n\right]  $). Now,
$m-\underbrace{\lambda_{i}}_{\geq0}+\underbrace{i}_{\leq n}\leq m+n$ and
$m-\underbrace{\lambda_{i}}_{\leq m}+\underbrace{i}_{\geq1}\geq m-m+1=1$.
Combining these two inequalities, we obtain $m-\lambda_{i}+i\in\left[
m+n\right]  $.

Forget that we fixed $i$. We thus have shown that $m-\lambda_{i}+i\in\left[
m+n\right]  $ for each $i\in\left[  n\right]  $. In other words, $\left\{
m-\lambda_{i}+i\ \mid\ i\in\left[  n\right]  \right\}  \subseteq\left[
m+n\right]  $. Thus,
\[
Q_{\lambda}^{\left(  n,m\right)  }=\left\{  m-\lambda_{i}+i\ \mid\ i\in\left[
n\right]  \right\}  \subseteq\left[  m+n\right]  =\left[  n+m\right]  .
\]
Since we know that $Q_{\lambda}^{\left(  n,m\right)  }$ is an $n$-element set,
we thus conclude that $Q_{\lambda}^{\left(  n,m\right)  }$ is an $n$-element
subset of $\left[  n+m\right]  $. This proves Lemma \ref{lem.Qlam1}
\textbf{(b)}. \medskip

\textbf{(c)} Lemma \ref{lem.Qlam1} \textbf{(b)} shows that for any partition
$\lambda\subseteq\left(  m^{n}\right)  $, the set $Q_{\lambda}^{\left(
n,m\right)  }$ is an $n$-element subset of $\left[  n+m\right]  $. Hence, the
map%
\begin{align}
\left\{  \text{partitions }\lambda\subseteq\left(  m^{n}\right)  \right\}   &
\rightarrow\left\{  n\text{-element subsets of }\left[  n+m\right]  \right\}
,\nonumber\\
\lambda &  \mapsto Q_{\lambda}^{\left(  n,m\right)  }
\label{pf.lem.Qlam1.c.map}%
\end{align}
is well-defined. It remains to show that this map is bijective. We will prove
its injectivity and its surjectivity separately:

\begin{itemize}
\item \textit{Injectivity:} Let us show that the map (\ref{pf.lem.Qlam1.c.map}%
) is injective. For this purpose, we must show that any partition
$\lambda\subseteq\left(  m^{n}\right)  $ can be uniquely reconstructed from
the set $Q_{\lambda}^{\left(  n,m\right)  }$. But this is easy: Lemma
\ref{lem.Qlam1} \textbf{(a)} shows that
\[
Q_{\lambda}^{\left(  n,m\right)  }=\left\{  m-\lambda_{1}+1<m-\lambda
_{2}+2<\cdots<m-\lambda_{n}+n\right\}  .
\]
Hence, the $n$ numbers $m-\lambda_{1}+1,\ m-\lambda_{2}+2,\ \ldots
,\ m-\lambda_{n}+n$ are the elements of the set $Q_{\lambda}^{\left(
n,m\right)  }$ in increasing order. Thus, knowing the set $Q_{\lambda
}^{\left(  n,m\right)  }$, we can reconstruct these $n$ numbers, and therefore
also the entries $\lambda_{1},\lambda_{2},\ldots,\lambda_{n}$ of $\lambda$;
hence, the whole partition $\lambda$ can be reconstructed from $Q_{\lambda
}^{\left(  n,m\right)  }$. This proves that the map (\ref{pf.lem.Qlam1.c.map})
is injective.

\item \textit{Surjectivity:} Let us now show that the map
(\ref{pf.lem.Qlam1.c.map}) is surjective. For this purpose, we fix some
$n$-element subset $R$ of $\left[  n+m\right]  $. We must show that there is a
partition $\lambda\subseteq\left(  m^{n}\right)  $ such that $Q_{\lambda
}^{\left(  n,m\right)  }=R$.

We know that $R$ is an $n$-element subset of $\left[  n+m\right]  $; thus, it
can be written as $R=\left\{  r_{1}<r_{2}<\cdots<r_{n}\right\}  $ for some
elements $r_{1}<r_{2}<\cdots<r_{n}$ of $\left[  n+m\right]  $. Consider these
elements. Define an $n$-tuple $\lambda:=\left(  \lambda_{1},\lambda_{2}%
,\ldots,\lambda_{n}\right)  \in\mathbb{Z}^{n}$ by setting
\[
\lambda_{i}:=m-r_{i}+i\ \ \ \ \ \ \ \ \ \ \text{for each }i\in\left[
n\right]  .
\]
Then, we have $\lambda_{1}\geq\lambda_{2}\geq\cdots\geq\lambda_{n}%
$\ \ \ \ \footnote{\textit{Proof.} We must show that $\lambda_{i}\geq
\lambda_{i+1}$ for each $i\in\left[  n-1\right]  $.
\par
So let $i\in\left[  n-1\right]  $. Then, $r_{i}<r_{i+1}$ (since $r_{1}%
<r_{2}<\cdots<r_{n}$) and thus $r_{i}\leq r_{i+1}-1$ (since $r_{i}$ and
$r_{i+1}$ are integers). But the definition of $\lambda_{i}$ shows that
$\lambda_{i}=m-r_{i}+i$, whereas the definition of $\lambda_{i+1}$ shows that
$\lambda_{i+1}=m-r_{i+1}+\left(  i+1\right)  $. Hence,%
\[
\lambda_{i}=m-\underbrace{r_{i}}_{\leq r_{i+1}-1}+\,i\geq m-\left(
r_{i+1}-1\right)  +i=m-r_{i+1}+\left(  i+1\right)  =\lambda_{i+1}.
\]
Thus, we have proved that $\lambda_{i}\geq\lambda_{i+1}$ for each $i\in\left[
n-1\right]  $, qed.} and $m\geq\lambda_{1}$\ \ \ \ \footnote{\textit{Proof.}
By the definition of $\lambda_{1}$, we have $\lambda_{1}=m-r_{1}+1$. But
$r_{1}\geq1$, since $r_{1}\in\left\{  r_{1}<r_{2}<\cdots<r_{n}\right\}
=R\subseteq\left[  n+m\right]  $. Thus, $\lambda_{1}=m-\underbrace{r_{1}%
}_{\geq1}+\,1\leq m-1+1=m$, so that $m\geq\lambda_{1}$.} and $\lambda_{n}%
\geq0$\ \ \ \ \footnote{\textit{Proof.} By the definition of $\lambda_{n}$, we
have $\lambda_{n}=m-r_{n}+n$. But $r_{n}\leq n+m$, since $r_{n}\in\left\{
r_{1}<r_{2}<\cdots<r_{n}\right\}  =R\subseteq\left[  n+m\right]  $. Thus,
$\lambda_{n}=m-\underbrace{r_{n}}_{\leq n+m}+\,n\geq m-\left(  n+m\right)
+n=0$.}. Combining these inequalities, we obtain $m\geq\lambda_{1}\geq
\lambda_{2}\geq\cdots\geq\lambda_{n}\geq0$. Hence, $\lambda=\left(
\lambda_{1},\lambda_{2},\ldots,\lambda_{n}\right)  $ is a partition satisfying
$\lambda\subseteq\left(  m^{n}\right)  $. Furthermore, each $i\in\left[
n\right]  $ satisfies $m-\lambda_{i}+i=r_{i}$ (since the definition of
$\lambda$ yields $\lambda_{i}=m-r_{i}+i$). But the definition of $Q_{\lambda
}^{\left(  n,m\right)  }$ says that
\begin{align*}
Q_{\lambda}^{\left(  n,m\right)  }  &  =\left\{  m-\lambda_{i}+i\ \mid
\ i\in\left[  n\right]  \right\} \\
&  =\left\{  r_{i}\ \mid\ i\in\left[  n\right]  \right\}
\ \ \ \ \ \ \ \ \ \ \left(  \text{since each }i\in\left[  n\right]  \text{
satisfies }m-\lambda_{i}+i=r_{i}\right) \\
&  =R\ \ \ \ \ \ \ \ \ \ \left(  \text{since }R=\left\{  r_{1}<r_{2}%
<\cdots<r_{n}\right\}  =\left\{  r_{i}\ \mid\ i\in\left[  n\right]  \right\}
\right)  .
\end{align*}
Thus, we have found a partition $\lambda\subseteq\left(  m^{n}\right)  $ such
that $Q_{\lambda}^{\left(  n,m\right)  }=R$. This completes our proof that the
map (\ref{pf.lem.Qlam1.c.map}) is surjective.
\end{itemize}

We have now shown that the map (\ref{pf.lem.Qlam1.c.map}) is both injective
and surjective. Hence, this map is bijective, i.e., is a bijection. This
proves Lemma \ref{lem.Qlam1} \textbf{(c)}.
\end{proof}

Next, we shall connect the $Q_{\lambda}^{\left(  n,m\right)  }$ construction
to the partition $\widehat{\lambda}^{\prime}$ that appears on the right hand
sides of (0.11') and of (6.17). First, we recall some notations:

\begin{itemize}
\item If $\lambda=\left(  \lambda_{1},\lambda_{2},\ldots,\lambda_{n}\right)
\subseteq\left(  m^{n}\right)  $ is a partition, then its complementary
partition $\widehat{\lambda}=\left(  \widehat{\lambda}_{1},\widehat{\lambda
}_{2},\ldots,\widehat{\lambda}_{n}\right)  $ is defined by $\widehat{\lambda
}_{i}=m-\lambda_{n+1-i}$ for all $i\in\left[  n\right]  $. Note that this
depends on $m$ and $n$. The complementary partition $\widehat{\lambda}$ again
satisfies $\widehat{\lambda}\subseteq\left(  m^{n}\right)  $, and can be
described visually as follows: Its Young diagram is the complement of the
Young diagram of $\lambda$ in the rectangle $\left(  m^{n}\right)  $, rotated
by $180^{\circ}$. It is furthermore easy to see that
$\widehat{\widehat{\lambda}}=\lambda$ and%
\begin{equation}
\left\vert \widehat{\lambda}\right\vert =nm-\left\vert \lambda\right\vert .
\label{p20.complement.size}%
\end{equation}


\item If $\lambda$ is any partition, then $\lambda^{\prime}$ is the conjugate
partition of $\lambda$, obtained by reflecting the Young diagram of $\lambda$
across the main diagonal (so that rows become columns and vice versa).
Clearly, if $\lambda\subseteq\left(  m^{n}\right)  $, then $\lambda^{\prime
}\subseteq\left(  n^{m}\right)  $.
\end{itemize}

Now our next lemma is the following:

\begin{lemmain}
\label{lem.Qlam2}Let $\lambda\subseteq\left(  m^{n}\right)  $ be a partition.
Then, $\lambda^{\prime}\subseteq\left(  n^{m}\right)  $ and
\[
Q_{\lambda^{\prime}}^{\left(  m,n\right)  }=\left[  n+m\right]  \setminus
Q_{\widehat{\lambda}}^{\left(  n,m\right)  }.
\]

\end{lemmain}

\begin{proof}
From $\lambda\subseteq\left(  m^{n}\right)  $, we obtain $\lambda^{\prime
}\subseteq\left(  n^{m}\right)  $. Thus, $\lambda=\left(  \lambda_{1}%
,\lambda_{2},\ldots,\lambda_{n}\right)  $ (since $\lambda\subseteq\left(
m^{n}\right)  $) and $\lambda^{\prime}=\left(  \lambda_{1}^{\prime}%
,\lambda_{2}^{\prime},\ldots,\lambda_{m}^{\prime}\right)  $ (since
$\lambda^{\prime}\subseteq\left(  n^{m}\right)  $).

For each $i\in\left[  n\right]  $, let us set $\alpha_{i}:=\lambda_{i}-i$. For
each $j\in\left[  m\right]  $, let us set $\beta_{j}:=\lambda_{j}^{\prime}-j$
and $\eta_{j}:=-1-\beta_{j}$. Then, Lemma \ref{lem.A4} (applied to $q=n$ and
$p=m$) yields that the two sets $\left\{  \alpha_{1},\alpha_{2},\ldots
,\alpha_{n}\right\}  $ and $\left\{  \eta_{1},\eta_{2},\ldots,\eta
_{m}\right\}  $ are disjoint, and their union is the integer interval $\left[
-n,m-1\right]  $. Hence,%
\begin{equation}
\left\{  \eta_{1},\eta_{2},\ldots,\eta_{m}\right\}  =\left[  -n,m-1\right]
\setminus\left\{  \alpha_{1},\alpha_{2},\ldots,\alpha_{n}\right\}  .
\label{pf.lem.Qlam2.comp1}%
\end{equation}


For any subset $S$ of $\mathbb{Z}$ and any integer $z$, we let $S+z$ denote
the set $\left\{  s+z\ \mid\ s\in S\right\}  $. Visually speaking, this is
simply the set $S$ shifted by $z$ units to the right along the number line.
Clearly, any two subsets $S$ and $T$ of $\mathbb{Z}$ and any integer $z$
satisfy%
\begin{equation}
\left(  S\setminus T\right)  +z=\left(  S+z\right)  \setminus\left(
T+z\right)  , \label{pf.lem.Qlam2.comp+z}%
\end{equation}
since the operation of adding $z$ to each integer is a bijection.

Now, from (\ref{pf.lem.Qlam2.comp1}), we obtain%
\begin{align}
&  \left\{  \eta_{1},\eta_{2},\ldots,\eta_{m}\right\}  +\left(  n+1\right)
\nonumber\\
&  =\left(  \left[  -n,m-1\right]  \setminus\left\{  \alpha_{1},\alpha
_{2},\ldots,\alpha_{n}\right\}  \right)  +\left(  n+1\right) \nonumber\\
&  =\underbrace{\left(  \left[  -n,m-1\right]  +\left(  n+1\right)  \right)
}_{\substack{=\left[  1,\left(  m-1\right)  +\left(  n+1\right)  \right]
\\=\left[  1,n+m\right]  =\left[  n+m\right]  }}\setminus\underbrace{\left(
\left\{  \alpha_{1},\alpha_{2},\ldots,\alpha_{n}\right\}  +\left(  n+1\right)
\right)  }_{\substack{=\left\{  \alpha_{1}+\left(  n+1\right)  ,\alpha
_{2}+\left(  n+1\right)  ,\ldots,\alpha_{n}+\left(  n+1\right)  \right\}
\\=\left\{  \alpha_{i}+\left(  n+1\right)  \ \mid\ i\in\left[  n\right]
\right\}  \\=\left\{  \alpha_{n+1-i}+\left(  n+1\right)  \ \mid\ i\in\left[
n\right]  \right\}  \\\text{(here, we have substituted }n+1-i\text{ for
}i\text{ in the set,}\\\text{since the map }\left[  n\right]  \rightarrow
\left[  n\right]  ,\ i\mapsto n+1-i\text{ is a bijection)}}%
}\ \ \ \ \ \ \ \ \ \ \left(  \text{by (\ref{pf.lem.Qlam2.comp+z})}\right)
\nonumber\\
&  =\left[  n+m\right]  \setminus\left\{  \alpha_{n+1-i}+\left(  n+1\right)
\ \mid\ i\in\left[  n\right]  \right\}  . \label{pf.lem.Qlam2.4}%
\end{align}


However, each $i\in\left[  n\right]  $ satisfies $\alpha_{n+1-i}%
=\lambda_{n+1-i}-\left(  n+1-i\right)  $ (by the definition of $\alpha
_{n+1-i}$) and thus%
\begin{align*}
\alpha_{n+1-i}+\left(  n+1\right)   &  =\lambda_{n+1-i}-\left(  n+1-i\right)
+\left(  n+1\right) \\
&  =\underbrace{\lambda_{n+1-i}}_{\substack{=m-\widehat{\lambda}%
_{i}\\\text{(since the definition of }\widehat{\lambda}\\\text{yields
}\widehat{\lambda}_{i}=m-\lambda_{n+1-i}\text{)}}}+\,i=m-\widehat{\lambda}%
_{i}+i.
\end{align*}
Hence,%
\[
\left\{  \underbrace{\alpha_{n+1-i}+\left(  n+1\right)  }_{=m-\widehat{\lambda
}_{i}+i}\ \mid\ i\in\left[  n\right]  \right\}  =\left\{  m-\widehat{\lambda
}_{i}+i\ \mid\ i\in\left[  n\right]  \right\}  =Q_{\widehat{\lambda}}^{\left(
n,m\right)  }%
\]
(since the definition of $Q_{\widehat{\lambda}}^{\left(  n,m\right)  }$ yields
$Q_{\widehat{\lambda}}^{\left(  n,m\right)  }=\left\{  m-\widehat{\lambda}%
_{i}+i\ \mid\ i\in\left[  n\right]  \right\}  $). Thus, we can rewrite
(\ref{pf.lem.Qlam2.4}) as%
\begin{equation}
\left\{  \eta_{1},\eta_{2},\ldots,\eta_{m}\right\}  +\left(  n+1\right)
=\left[  n+m\right]  \setminus Q_{\widehat{\lambda}}^{\left(  n,m\right)  }.
\label{pf.lem.Qlam2.6}%
\end{equation}


On the other hand, each $i\in\left[  m\right]  $ satisfies%
\begin{align}
\underbrace{\eta_{i}}_{\substack{=-1-\beta_{i}\\\text{(by the definition of
}\eta_{i}\text{)}}}+\left(  n+1\right)   &  =-1-\beta_{i}+\left(  n+1\right)
=n-\underbrace{\beta_{i}}_{\substack{=\lambda_{i}^{\prime}-i\\\text{(by the
definition of }\beta_{i}\text{)}}}\nonumber\\
&  =n-\left(  \lambda_{i}^{\prime}-i\right)  =n-\lambda_{i}^{\prime}+i.
\label{pf.lem.Qlam2.7}%
\end{align}
Now,%
\begin{align*}
&  \left\{  \eta_{1},\eta_{2},\ldots,\eta_{m}\right\}  +\left(  n+1\right) \\
&  =\left\{  \eta_{1}+\left(  n+1\right)  ,\ \eta_{2}+\left(  n+1\right)
,\ \ldots,\ \eta_{m}+\left(  n+1\right)  \right\} \\
&  =\left\{  \eta_{i}+\left(  n+1\right)  \ \mid\ i\in\left[  m\right]
\right\} \\
&  =\left\{  n-\lambda_{i}^{\prime}+i\ \mid\ i\in\left[  m\right]  \right\}
\ \ \ \ \ \ \ \ \ \ \left(  \text{by (\ref{pf.lem.Qlam2.7})}\right) \\
&  =Q_{\lambda^{\prime}}^{\left(  m,n\right)  }%
\end{align*}
(since the definition of $Q_{\lambda^{\prime}}^{\left(  m,n\right)  }$ yields
$Q_{\lambda^{\prime}}^{\left(  m,n\right)  }=\left\{  n-\lambda_{i}^{\prime
}+i\ \mid\ i\in\left[  m\right]  \right\}  $). Comparing this with
(\ref{pf.lem.Qlam2.6}), we obtain $Q_{\lambda^{\prime}}^{\left(  m,n\right)
}=\left[  n+m\right]  \setminus Q_{\widehat{\lambda}}^{\left(  n,m\right)  }$.
This proves Lemma \ref{lem.Qlam2}.
\end{proof}

We can easily restate Lemma \ref{lem.Qlam2} as follows:

\begin{lemmain}
\label{lem.Qlam3}Let $\lambda\subseteq\left(  m^{n}\right)  $ be a partition.
Then, $\widehat{\lambda}^{\prime}\subseteq\left(  n^{m}\right)  $ and
\[
Q_{\widehat{\lambda}^{\prime}}^{\left(  m,n\right)  }=\left[  n+m\right]
\setminus Q_{\lambda}^{\left(  n,m\right)  }.
\]

\end{lemmain}

\begin{proof}
From $\lambda\subseteq\left(  m^{n}\right)  $, we obtain $\widehat{\lambda
}\subseteq\left(  m^{n}\right)  $ and $\widehat{\widehat{\lambda}}=\lambda$
(by the basic properties of $\widehat{\lambda}$). Thus, Lemma \ref{lem.Qlam2}
(applied to $\widehat{\lambda}$ instead of $\lambda$) yields $\widehat{\lambda
}^{\prime}\subseteq\left(  n^{m}\right)  $ and
\[
Q_{\widehat{\lambda}^{\prime}}^{\left(  m,n\right)  }=\left[  n+m\right]
\setminus Q_{\widehat{\widehat{\lambda}}}^{\left(  n,m\right)  }=\left[
n+m\right]  \setminus Q_{\lambda}^{\left(  n,m\right)  }%
\]
(since $\widehat{\widehat{\lambda}}=\lambda$). This proves Lemma
\ref{lem.Qlam3}.
\end{proof}

We are now ready to prove (\ref{p20.claim}):

\begin{proof}
[Proof of (\ref{p20.claim}).]From the definition (6.2) of $A_{\delta_{m+n}%
}\left(  x,y\mid a\right)  $, we know that
\begin{equation}
A_{\delta_{m+n}}\left(  x,y\mid a\right)  =\det B, \label{p20.claim.pf.1}%
\end{equation}
where $B$ is the $\left(  n+m\right)  \times\left(  n+m\right)  $-matrix whose
$\left(  i,j\right)  $-th entry (for all $i,j\in\left[  n+m\right]  $) is%
\[%
\begin{cases}
\left(  x_{i}\mid a\right)  ^{m+n-j}, & \text{if }i\leq n;\\
\left(  y_{i-n}\mid a\right)  ^{m+n-j}, & \text{if }i>n.
\end{cases}
\]
Consider this matrix $B$.

Now, set $k:=n+m$, so that $B$ is a $k\times k$-matrix. Let furthermore
$P:=\left[  n\right]  $. This is a subset of $\left[  k\right]  $, since
$n\leq n+m=k$. Moreover, $\left\vert P\right\vert =\left\vert \left[
n\right]  \right\vert =n$.

We shall use the notations of Lemma \ref{lem.lap-rows}, so that in particular
we have
\begin{align*}
\widetilde{P}  &  =\left[  \underbrace{k}_{=n+m}\right]  \setminus
\underbrace{P}_{=\left[  n\right]  }=\left[  n+m\right]  \setminus\left[
n\right]  =\left\{  n+1,n+2,\ldots,n+m\right\} \\
&  =\left\{  n+1<n+2<\cdots<n+m\right\}  .
\end{align*}
Now, Lemma \ref{lem.lap-rows} (applied to $A=B$) yields
\begin{align*}
\det B  &  =\sum_{\substack{Q\subseteq\left[  k\right]  ;\\\left\vert
Q\right\vert =\left\vert P\right\vert }}\left(  -1\right)  ^{\sum P+\sum
Q}\det\left(  \operatorname*{sub}\nolimits_{P}^{Q}B\right)  \det\left(
\operatorname*{sub}\nolimits_{\widetilde{P}}^{\widetilde{Q}}B\right) \\
&  =\sum_{\substack{Q\subseteq\left[  n+m\right]  ;\\\left\vert Q\right\vert
=n}}\left(  -1\right)  ^{\sum P+\sum Q}\det\left(  \operatorname*{sub}%
\nolimits_{P}^{Q}B\right)  \det\left(  \operatorname*{sub}%
\nolimits_{\widetilde{P}}^{\widetilde{Q}}B\right) \\
&  \ \ \ \ \ \ \ \ \ \ \ \ \ \ \ \ \ \ \ \ \left(  \text{since }k=n+m\text{
and }\left\vert P\right\vert =n\right) \\
&  =\sum_{\lambda\subseteq\left(  m^{n}\right)  }\left(  -1\right)  ^{\sum
P+\sum Q_{\lambda}^{\left(  n,m\right)  }}\det\left(  \operatorname*{sub}%
\nolimits_{P}^{Q_{\lambda}^{\left(  n,m\right)  }}B\right)  \det\left(
\operatorname*{sub}\nolimits_{\widetilde{P}}^{\widetilde{Q_{\lambda}^{\left(
n,m\right)  }}}B\right)
\end{align*}
(here, we have substituted $Q_{\lambda}^{\left(  n,m\right)  }$ for $Q$,
because Lemma \ref{lem.Qlam1} \textbf{(c)} says that the map%
\begin{align*}
\left\{  \text{partitions }\lambda\subseteq\left(  m^{n}\right)  \right\}   &
\rightarrow\left\{  n\text{-element subsets of }\left[  n+m\right]  \right\}
,\\
\lambda &  \mapsto Q_{\lambda}^{\left(  n,m\right)  }%
\end{align*}
is a bijection).

Now, we shall simplify the right hand side of this equality. Fix a partition
$\lambda\subseteq\left(  m^{n}\right)  $. Then, by the definition of
$\widetilde{Q_{\lambda}^{\left(  n,m\right)  }}$, we have%
\begin{align}
\widetilde{Q_{\lambda}^{\left(  n,m\right)  }}  &  =\left[  \underbrace{k}%
_{=n+m}\right]  \setminus Q_{\lambda}^{\left(  n,m\right)  }=\left[
n+m\right]  \setminus Q_{\lambda}^{\left(  n,m\right)  }\nonumber\\
&  =Q_{\widehat{\lambda}^{\prime}}^{\left(  m,n\right)  }
\label{p20.claim.pf.l1}%
\end{align}
(by Lemma \ref{lem.Qlam3}). Note that Lemma \ref{lem.Qlam3} also yields
$\widehat{\lambda}^{\prime}\subseteq\left(  n^{m}\right)  $.

Furthermore, recall that $P=\left[  n\right]  =\left\{  1<2<\cdots<n\right\}
$ and \newline$Q_{\lambda}^{\left(  n,m\right)  }=\left\{  m-\lambda
_{1}+1<m-\lambda_{2}+2<\cdots<m-\lambda_{n}+n\right\}  $ (by Lemma
\ref{lem.Qlam1} \textbf{(a)}). Hence,%
\begin{align*}
\operatorname*{sub}\nolimits_{P}^{Q_{\lambda}^{\left(  n,m\right)  }}B  &
=\operatorname*{sub}\nolimits_{\left\{  1<2<\cdots<n\right\}  }^{\left\{
m-\lambda_{1}+1<m-\lambda_{2}+2<\cdots<m-\lambda_{n}+n\right\}  }B\\
&  =\left(  \left(  x_{i}\mid a\right)  ^{m+n-\left(  m-\lambda_{j}+j\right)
}\right)  _{1\leq i,j\leq n}\\
&  \ \ \ \ \ \ \ \ \ \ \ \ \ \ \ \ \ \ \ \ \left(
\begin{array}
[c]{c}%
\text{by the definition of }B\text{, since}\\
\text{all }i\in\left\{  1<2<\cdots<n\right\}  \text{ satisfy }i\leq n
\end{array}
\right) \\
&  =\left(  \left(  x_{i}\mid a\right)  ^{\left(  \lambda+\delta_{n}\right)
_{j}}\right)  _{1\leq i,j\leq n}%
\end{align*}
(since $m+n-\left(  m-\lambda_{j}+j\right)  =\lambda_{j}+n-j=\left(
\lambda+\delta_{n}\right)  _{j}$ for all $j\in\left[  n\right]  $) and
therefore%
\begin{align}
\det\left(  \operatorname*{sub}\nolimits_{P}^{Q_{\lambda}^{\left(  n,m\right)
}}B\right)   &  =\det\left(  \left(  x_{i}\mid a\right)  ^{\left(
\lambda+\delta_{n}\right)  _{j}}\right)  _{1\leq i,j\leq n}\nonumber\\
&  =A_{\lambda+\delta_{n}}\left(  x\mid a\right)  \label{p20.claim.pf.l2}%
\end{align}
(by the definition of $A_{\lambda+\delta_{n}}\left(  x\mid a\right)  $).
Furthermore, we have \newline$\widetilde{P}=\left\{  n+1<n+2<\cdots
<n+m\right\}  $ and%
\begin{align*}
\widetilde{Q_{\lambda}^{\left(  n,m\right)  }}  &  =Q_{\widehat{\lambda
}^{\prime}}^{\left(  m,n\right)  }\ \ \ \ \ \ \ \ \ \ \left(  \text{by
(\ref{p20.claim.pf.l1})}\right) \\
&  =\left\{  n-\widehat{\lambda}_{1}^{\prime}+1<n-\widehat{\lambda}%
_{2}^{\prime}+2<\cdots<n-\widehat{\lambda}_{m}^{\prime}+m\right\}
\end{align*}
(by Lemma \ref{lem.Qlam1} \textbf{(a)}, applied to $m$, $n$ and
$\widehat{\lambda}^{\prime}$ instead of $n$, $m$ and $\lambda$). Hence,%
\begin{align*}
\operatorname*{sub}\nolimits_{\widetilde{P}}^{\widetilde{Q_{\lambda}^{\left(
n,m\right)  }}}B  &  =\operatorname*{sub}\nolimits_{\left\{  n+1<n+2<\cdots
<n+m\right\}  }^{\left\{  n-\widehat{\lambda}_{1}^{\prime}%
+1<n-\widehat{\lambda}_{2}^{\prime}+2<\cdots<n-\widehat{\lambda}_{m}^{\prime
}+m\right\}  }B\\
&  =\left(  \left(  y_{i}\mid a\right)  ^{m+n-\left(  n-\widehat{\lambda}%
_{j}^{\prime}+j\right)  }\right)  _{1\leq i,j\leq m}\\
&  \ \ \ \ \ \ \ \ \ \ \ \ \ \ \ \ \ \ \ \ \left(
\begin{array}
[c]{c}%
\text{by the definition of }B\text{, since}\\
\text{all }i\in\left\{  n+1<n+2<\cdots<n+m\right\}  \text{ satisfy }i>n
\end{array}
\right) \\
&  =\left(  \left(  y_{i}\mid a\right)  ^{\left(  \widehat{\lambda}^{\prime
}+\delta_{m}\right)  _{j}}\right)  _{1\leq i,j\leq m}%
\end{align*}
(since $m+n-\left(  n-\widehat{\lambda}_{j}^{\prime}+j\right)
=\widehat{\lambda}_{j}^{\prime}+m-j=\left(  \widehat{\lambda}^{\prime}%
+\delta_{m}\right)  _{j}$ for all $j\in\left[  m\right]  $) and therefore%
\begin{align}
\det\left(  \operatorname*{sub}\nolimits_{\widetilde{P}}%
^{\widetilde{Q_{\lambda}^{\left(  n,m\right)  }}}B\right)   &  =\det\left(
\left(  y_{i}\mid a\right)  ^{\left(  \widehat{\lambda}^{\prime}+\delta
_{m}\right)  _{j}}\right)  _{1\leq i,j\leq m}\nonumber\\
&  =A_{\widehat{\lambda}^{\prime}+\delta_{m}}\left(  y\mid a\right)
\label{p20.claim.pf.l3}%
\end{align}
(by the definition of $A_{\widehat{\lambda}^{\prime}+\delta_{m}}\left(  y\mid
a\right)  $). Finally, we have%
\begin{align*}
&  \sum\underbrace{P}_{=\left[  n\right]  }+\sum\underbrace{Q_{\lambda
}^{\left(  n,m\right)  }}_{\substack{=\left\{  m-\lambda_{1}+1<m-\lambda
_{2}+2<\cdots<m-\lambda_{n}+n\right\}  \\\text{(by Lemma \ref{lem.Qlam1}
\textbf{(a)})}}}\\
&  =\underbrace{\sum\left[  n\right]  }_{=1+2+\cdots+n}+\underbrace{\sum
\left\{  m-\lambda_{1}+1<m-\lambda_{2}+2<\cdots<m-\lambda_{n}+n\right\}
}_{=\left(  m-\lambda_{1}+1\right)  +\left(  m-\lambda_{2}+2\right)
+\cdots+\left(  m-\lambda_{n}+n\right)  }\\
&  =\left(  1+2+\cdots+n\right)  +\underbrace{\left(  m-\lambda_{1}+1\right)
+\left(  m-\lambda_{2}+2\right)  +\cdots+\left(  m-\lambda_{n}+n\right)
}_{=\underbrace{\left(  m+m+\cdots+m\right)  }_{n\text{ times}}-\left(
\lambda_{1}+\lambda_{2}+\cdots+\lambda_{n}\right)  +\left(  1+2+\cdots
+n\right)  }\\
&  =\left(  1+2+\cdots+n\right)  +\underbrace{\underbrace{\left(
m+m+\cdots+m\right)  }_{n\text{ times}}}_{=nm}-\underbrace{\left(  \lambda
_{1}+\lambda_{2}+\cdots+\lambda_{n}\right)  }_{\substack{=\left\vert
\lambda\right\vert \\\text{(since }\lambda=\left(  \lambda_{1},\lambda
_{2},\ldots,\lambda_{n}\right)  \text{)}}}+\left(  1+2+\cdots+n\right) \\
&  =\left(  1+2+\cdots+n\right)  +nm-\left\vert \lambda\right\vert +\left(
1+2+\cdots+n\right) \\
&  =\underbrace{2\cdot\left(  1+2+\cdots+n\right)  }_{\equiv
0\operatorname{mod}2}+\,nm-\left\vert \lambda\right\vert \equiv nm-\left\vert
\lambda\right\vert =\left\vert \widehat{\lambda}\right\vert \operatorname{mod}%
2\ \ \ \ \ \ \ \ \ \ \left(  \text{by (\ref{p20.complement.size})}\right)
\end{align*}
and thus
\begin{equation}
\left(  -1\right)  ^{\sum P+\sum Q_{\lambda}^{\left(  n,m\right)  }}=\left(
-1\right)  ^{\left\vert \widehat{\lambda}\right\vert }.
\label{p20.claim.pf.l4}%
\end{equation}


Forget that we fixed $\lambda$. We thus have proved (\ref{p20.claim.pf.l1}),
(\ref{p20.claim.pf.l2}), (\ref{p20.claim.pf.l3}) and (\ref{p20.claim.pf.l4})
for each partition $\lambda\subseteq\left(  m^{n}\right)  $.

Hence, our above computation of $\det B$ can be continued as follows:%
\begin{align*}
\det B  &  =\sum_{\lambda\subseteq\left(  m^{n}\right)  }\underbrace{\left(
-1\right)  ^{\sum P+\sum Q_{\lambda}^{\left(  n,m\right)  }}}%
_{\substack{=\left(  -1\right)  ^{\left\vert \widehat{\lambda}\right\vert
}\\\text{(by (\ref{p20.claim.pf.l4}))}}}\underbrace{\det\left(
\operatorname*{sub}\nolimits_{P}^{Q_{\lambda}^{\left(  n,m\right)  }}B\right)
}_{\substack{=A_{\lambda+\delta_{n}}\left(  x\mid a\right)  \\\text{(by
(\ref{p20.claim.pf.l2}))}}}\underbrace{\det\left(  \operatorname*{sub}%
\nolimits_{\widetilde{P}}^{\widetilde{Q_{\lambda}^{\left(  n,m\right)  }}%
}B\right)  }_{\substack{=A_{\widehat{\lambda}^{\prime}+\delta_{m}}\left(
y\mid a\right)  \\\text{(by (\ref{p20.claim.pf.l3}))}}}\\
&  =\sum_{\lambda\subseteq\left(  m^{n}\right)  }\left(  -1\right)
^{\left\vert \widehat{\lambda}\right\vert }A_{\lambda+\delta_{n}}\left(  x\mid
a\right)  A_{\widehat{\lambda}^{\prime}+\delta_{m}}\left(  y\mid a\right)  .
\end{align*}
In view of (\ref{p20.claim.pf.1}), we can rewrite this as%
\[
A_{\delta_{m+n}}\left(  x,y\mid a\right)  =\sum_{\lambda\subseteq\left(
m^{n}\right)  }\left(  -1\right)  ^{\left\vert \widehat{\lambda}\right\vert
}A_{\lambda+\delta_{n}}\left(  x\mid a\right)  A_{\widehat{\lambda}^{\prime
}+\delta_{m}}\left(  y\mid a\right)  .
\]
This proves (\ref{p20.claim}).
\end{proof}

\item \textbf{page 21, Remark:} The chain of equalities%
\begin{align*}
A_{\alpha}\left(  x\mid a\right)   &  =\det\left(  \sum_{\beta_{k}\geq0}%
x_{i}^{\beta_{k}}e_{\beta_{k}-\alpha_{j}}\left(  a^{\left(  \alpha_{j}\right)
}\right)  \right) \\
&  =\sum_{\beta}\det\left(  x_{i}^{\beta_{k}}\right)  \det\left(  e_{\beta
_{k}-\alpha_{j}}\left(  a^{\left(  \alpha_{j}\right)  }\right)  \right)
\end{align*}
should be replaced by%
\begin{align}
A_{\alpha}\left(  x\mid a\right)   &  =\det\left(  \sum_{k\geq0}x_{i}%
^{k}e_{\alpha_{j}-k}\left(  a^{\left(  \alpha_{j}\right)  }\right)  \right)
_{1\leq i,j\leq n}\nonumber\\
&  =\sum_{\beta}\det\left(  x_{i}^{\beta_{k}}\right)  _{1\leq i,k\leq n}%
\det\left(  e_{\alpha_{j}-\beta_{k}}\left(  a^{\left(  \alpha_{j}\right)
}\right)  \right)  _{1\leq k,j\leq n}.\ \ \ \ \ \ \ \ \ \ \label{p20.eq2}%
\end{align}


Let me also explain in more detail how the last equality sign here is proved
(the first one is clear from the definition of $A_{\alpha}\left(  x\mid
a\right)  $ and the formula for $\left(  x\mid a\right)  ^{r}$ stated at the
beginning of the Remark). We shall use the Cauchy--Binet formula for infinite
matrices, but this time the infinite matrices are infinite leftwards and
upwards, respectively. Namely, we let $-\mathbb{N}$ be the set $\left\{
0,-1,-2,\ldots\right\}  $ of all nonpositive integers, and we define the
$\left[  n\right]  \times\left(  -\mathbb{N}\right)  $-matrix
\[
U:=\left(  x_{i}^{-j}\right)  _{i\in\left[  n\right]  ,\ j\in-\mathbb{N}%
}=\left(
\begin{array}
[c]{cccc}%
\cdots & x_{1}^{2} & x_{1}^{1} & x_{1}^{0}\\
\cdots & x_{2}^{2} & x_{2}^{1} & x_{2}^{0}\\
\cdots & \vdots & \vdots & \vdots\\
\cdots & x_{n}^{2} & x_{n}^{1} & x_{n}^{0}%
\end{array}
\right)
\]
and the $\left(  -\mathbb{N}\right)  \times\left[  n\right]  $-matrix%
\[
V:=\left(  e_{\alpha_{j}+i}\left(  a^{\left(  \alpha_{j}\right)  }\right)
\right)  _{i\in-\mathbb{N},\ j\in\left[  n\right]  }=\left(
\begin{array}
[c]{cccc}%
\vdots & \vdots & \vdots & \vdots\\
e_{\alpha_{1}-2}\left(  a^{\left(  \alpha_{1}\right)  }\right)  &
e_{\alpha_{2}-2}\left(  a^{\left(  \alpha_{2}\right)  }\right)  & \cdots &
e_{\alpha_{n}-2}\left(  a^{\left(  \alpha_{n}\right)  }\right) \\
e_{\alpha_{1}-1}\left(  a^{\left(  \alpha_{1}\right)  }\right)  &
e_{\alpha_{2}-1}\left(  a^{\left(  \alpha_{2}\right)  }\right)  & \cdots &
e_{\alpha_{n}-1}\left(  a^{\left(  \alpha_{n}\right)  }\right) \\
e_{\alpha_{1}-0}\left(  a^{\left(  \alpha_{1}\right)  }\right)  &
e_{\alpha_{2}-0}\left(  a^{\left(  \alpha_{2}\right)  }\right)  & \cdots &
e_{\alpha_{n}-0}\left(  a^{\left(  \alpha_{n}\right)  }\right)
\end{array}
\right)  .
\]
Then, $UV$ is the $\left[  n\right]  \times\left[  n\right]  $-matrix whose
$\left(  i,j\right)  $-th entry (for all $i,j\in\left[  n\right]  $) is%
\begin{align*}
\sum_{k\in-\mathbb{N}}x_{i}^{-k}e_{\alpha_{j}+k}\left(  a^{\left(  \alpha
_{j}\right)  }\right)   &  =\sum_{k\in\mathbb{N}}x_{i}^{k}e_{\alpha_{j}%
-k}\left(  a^{\left(  \alpha_{j}\right)  }\right) \\
&  \ \ \ \ \ \ \ \ \ \ \ \ \ \ \ \ \ \ \ \ \left(
\begin{array}
[c]{c}%
\text{here, we have substituted }-k\text{ for }k\text{,}\\
\text{since the map }\mathbb{N}\rightarrow-\mathbb{N},\ k\mapsto-k\\
\text{is a bijection}%
\end{array}
\right) \\
&  =\sum_{k\geq0}x_{i}^{k}e_{\alpha_{j}-k}\left(  a^{\left(  \alpha
_{j}\right)  }\right)  .
\end{align*}
In other words,
\begin{equation}
UV=\left(  \sum_{k\geq0}x_{i}^{k}e_{\alpha_{j}-k}\left(  a^{\left(  \alpha
_{j}\right)  }\right)  \right)  _{1\leq i,j\leq n}. \label{p21.UV=}%
\end{equation}
On the other hand, the Cauchy--Binet formula (applied to the matrices $U$ and
$V$) yields%
\begin{align*}
\det\left(  UV\right)   &  =\sum_{b_{1}<b_{2}<\cdots<b_{n}\leq0}%
\det\underbrace{\left(  \operatorname*{sub}\nolimits_{\left\{  1,2,\ldots
,n\right\}  }^{\left\{  b_{1},b_{2},\ldots,b_{n}\right\}  }U\right)
}_{\substack{=\left(  x_{i}^{-b_{j}}\right)  _{i,j\in\left[  n\right]
}\\\text{(by the definition of }U\text{)}}}\cdot\det\underbrace{\left(
\operatorname*{sub}\nolimits_{\left\{  b_{1},b_{2},\ldots,b_{n}\right\}
}^{\left\{  1,2,\ldots,n\right\}  }V\right)  }_{\substack{=\left(
e_{\alpha_{j}+b_{i}}\left(  a^{\left(  \alpha_{j}\right)  }\right)  \right)
_{i,j\in\left[  n\right]  }\\\text{(by the definition of }V\text{)}}}\\
&  =\sum_{b_{1}<b_{2}<\cdots<b_{n}\leq0}\det\left(  x_{i}^{-b_{j}}\right)
_{i,j\in\left[  n\right]  }\cdot\det\left(  e_{\alpha_{j}+b_{i}}\left(
a^{\left(  \alpha_{j}\right)  }\right)  \right)  _{i,j\in\left[  n\right]  }\\
&  =\sum_{\beta_{1}>\beta_{2}>\cdots>\beta_{n}\geq0}\det\left(  x_{i}%
^{\beta_{j}}\right)  _{i,j\in\left[  n\right]  }\cdot\det\left(  e_{\alpha
_{j}-\beta_{i}}\left(  a^{\left(  \alpha_{j}\right)  }\right)  \right)
_{i,j\in\left[  n\right]  }%
\end{align*}
(here, we have substituted $-\beta_{1},-\beta_{2},\ldots,-\beta_{n}$ for
$b_{1},b_{2},\ldots,b_{n}$ in the sum). Thus,%
\begin{align*}
\det\left(  UV\right)   &  =\sum_{\beta_{1}>\beta_{2}>\cdots>\beta_{n}\geq
0}\det\underbrace{\left(  x_{i}^{\beta_{j}}\right)  _{i,j\in\left[  n\right]
}}_{\substack{=\left(  x_{i}^{\beta_{k}}\right)  _{i,k\in\left[  n\right]
}\\=\left(  x_{i}^{\beta_{k}}\right)  _{1\leq i,k\leq n}}}\cdot\det
\underbrace{\left(  e_{\alpha_{j}-\beta_{i}}\left(  a^{\left(  \alpha
_{j}\right)  }\right)  \right)  _{i,j\in\left[  n\right]  }}%
_{\substack{=\left(  e_{\alpha_{j}-\beta_{k}}\left(  a^{\left(  \alpha
_{j}\right)  }\right)  \right)  _{k,j\in\left[  n\right]  }\\=\left(
e_{\alpha_{j}-\beta_{k}}\left(  a^{\left(  \alpha_{j}\right)  }\right)
\right)  _{1\leq k,j\leq n}}}\\
&  =\sum_{\beta_{1}>\beta_{2}>\cdots>\beta_{n}\geq0}\det\left(  x_{i}%
^{\beta_{k}}\right)  _{1\leq i,k\leq n}\det\left(  e_{\alpha_{j}-\beta_{k}%
}\left(  a^{\left(  \alpha_{j}\right)  }\right)  \right)  _{1\leq k,j\leq n}.
\end{align*}
In view of (\ref{p21.UV=}), we can rewrite this as%
\begin{align*}
&  \det\left(  \sum_{k\geq0}x_{i}^{k}e_{\alpha_{j}-k}\left(  a^{\left(
\alpha_{j}\right)  }\right)  \right)  _{1\leq i,j\leq n}\\
&  =\sum_{\beta_{1}>\beta_{2}>\cdots>\beta_{n}\geq0}\det\left(  x_{i}%
^{\beta_{k}}\right)  _{1\leq i,k\leq n}\det\left(  e_{\alpha_{j}-\beta_{k}%
}\left(  a^{\left(  \alpha_{j}\right)  }\right)  \right)  _{1\leq k,j\leq n}.
\end{align*}
This proves the second equality sign in (\ref{p20.eq2}).

\item \textbf{page 21, (6.18):} Replace \textquotedblleft$a^{\left(
\lambda_{j}+n-j\right)  }$\textquotedblright\ by \textquotedblleft$a^{\left(
\lambda_{i}+n-i\right)  }$\textquotedblright\ on the right hand side.

\item \textbf{page 21, (6.18):} Let me briefly explain why the sum here is
over the partitions $\mu\subset\lambda$ rather than over all partitions
$\mu\in\mathbb{N}^{n}$.

Indeed, of course, the sum that is initially obtained from (\ref{p20.eq2}) by
doing what is said here (i.e., dividing both sides by $\Delta\left(  x\right)
$ and replacing $\alpha,\beta$ by $\lambda+\delta,\mu+\delta$) is a sum
ranging over all partitions $\mu\in\mathbb{N}^{n}$. However, the partitions
$\mu$ that don't satisfy $\mu\subset\lambda$ do not actually contribute
anything to the sum, since the corresponding addends are all $0$ (because an
argument similar to the proof of (6.10) shows that $\det\left(  e_{\lambda
_{i}-\mu_{j}-i+j}\left(  a^{\left(  \lambda_{i}+n-i\right)  }\right)  \right)
_{i,j\in\left[  n\right]  }=0$ for all such $\mu$). Thus, the sum can be
restricted (without changing its value) to range only over the partitions
$\mu\subset\lambda$.

\item \textbf{page 21, (6.19):} Let me explain how this is proved.

We define $h_{r}\left(  x\parallel a\right)  $ (for any $r\in\mathbb{Z}$) to
be the coefficient of $t^{r}$ in the power series
\[
\prod_{i\geq1}\left(  1-tx_{i}\right)  ^{-1}\prod_{j\geq1}\left(
1+ta_{j}\right)  \in\left(  \mathbb{Z}\left[  \left[  x_{i},a_{i}%
\ \mid\ i>0\right]  \right]  \right)  \left[  \left[  t\right]  \right]  .
\]
(So we have $h_{r}\left(  x\parallel a\right)  =0$ for all $r<0$.) We define
$s_{\lambda}\left(  x\parallel a\right)  $ to be the determinant
\[
\det\left(  h_{\lambda_{i}-i+j}\left(  x\parallel a\right)  \right)
_{i,j\in\left[  m\right]  }%
\]
whenever $\lambda=\left(  \lambda_{1},\lambda_{2},\ldots,\lambda_{m}\right)  $
is a partition. (This does not depend on the choice of $m$, since one can
easily see that if $\lambda_{m}=0$, then the last row of the matrix $\left(
h_{\lambda_{i}-i+j}\left(  x\parallel a\right)  \right)  _{i,j\in\left[
m\right]  }$ is $\left(  0,0,\ldots,0,1\right)  $, and therefore Laplace
expansion along this row reduces the determinant to the determinant of the
submatrix $\left(  h_{\lambda_{i}-i+j}\left(  x\parallel a\right)  \right)
_{i,j\in\left[  m-1\right]  }$.)

We want to prove the equality (6.19). We simplify it using a shorthand: We let
$x^{\left(  n\right)  }$ denote the $n$-tuple $\left(  x_{1},x_{2}%
,\ldots,x_{n}\right)  $ for all $n\in\mathbb{N}$ (so that $s_{\lambda}\left(
x^{\left(  n\right)  }\mid a\right)  $ means $s_{\lambda}\left(  x_{1}%
,x_{2},\ldots,x_{n}\mid a\right)  $). Thus, (6.19) rewrites as
\begin{equation}
\lim\limits_{n\rightarrow\infty}s_{\lambda}\left(  x^{\left(  n\right)  }\mid
a\right)  =s_{\lambda}\left(  x\parallel a\right)  . \label{p21.6.19.rewr}%
\end{equation}
This is the equality we shall be proving.

First, we show a few lemmas:

\begin{lemmain}
\label{lem.p21.TO}Let $\lambda$ be any partition. Then,%
\[
\lim\limits_{n\rightarrow\infty}s_{\lambda}\left(  x^{\left(  n\right)  }\mid
a\right)  =\sum_{\mu\subseteq\lambda}s_{\mu}\left(  x\right)  s_{\lambda
^{\prime}/\mu^{\prime}}\left(  a\right)  .
\]

\end{lemmain}

\begin{proof}
[Proof of Lemma \ref{lem.p21.TO}.]In the paper (between (6.18) and (6.19)), it
was said that \textquotedblleft we can let $n\rightarrow\infty$%
\textquotedblright\ and that \textquotedblleft In the limit the right hand
side of (6.18) becomes, by virtue of (0.8),
\[
\sum_{\mu\subset\lambda}s_{\mu}\left(  x\right)  s_{\lambda^{\prime}%
/\mu^{\prime}}\left(  a\right)
\]
where $x=\left(  x_{1},x_{2},\ldots\right)  $ and $a=\left(  a_{1}%
,a_{2},\ldots\right)  $\textquotedblright. But the left hand side of (6.18)
obviously becomes $\lim\limits_{n\rightarrow\infty}s_{\lambda}\left(
x^{\left(  n\right)  }\mid a\right)  $ in this limit. Thus, letting
$n\rightarrow\infty$ in (6.18), we obtain%
\[
\lim\limits_{n\rightarrow\infty}s_{\lambda}\left(  x^{\left(  n\right)  }\mid
a\right)  =\sum_{\mu\subset\lambda}s_{\mu}\left(  x\right)  s_{\lambda
^{\prime}/\mu^{\prime}}\left(  a\right)  =\sum_{\mu\subseteq\lambda}s_{\mu
}\left(  x\right)  s_{\lambda^{\prime}/\mu^{\prime}}\left(  a\right)
\]
(recalling that our symbol \textquotedblleft$\subseteq$\textquotedblright\ is
equivalent with Macdonald's \textquotedblleft$\subset$\textquotedblright).
This proves Lemma \ref{lem.p21.TO}.
\end{proof}

\begin{lemmain}
\label{lem.p21.TA}Let $r\in\mathbb{Z}$. Then,%
\[
h_{r}\left(  x\parallel a\right)  =\lim\limits_{n\rightarrow\infty}%
h_{r}\left(  x^{\left(  n\right)  }\mid a\right)  .
\]

\end{lemmain}

\begin{proof}
[Proof of Lemma \ref{lem.p21.TA}.]If $r<0$, then both $h_{r}\left(  x\parallel
a\right)  $ and $h_{r}\left(  x^{\left(  n\right)  }\mid a\right)  $ equal
$0$, and thus the lemma is obvious in this case. Thus, we WLOG assume that
$r\geq0$. We also WLOG assume that $n\geq1$ (since the $n=0$ case is trivial).

Let $\left[  t^{i}\right]  f$ denote the coefficient of a monomial $t^{i}$ in
a power series $f$. Then, the definition of $h_{r}\left(  x\parallel a\right)
$ rewrites as%
\begin{align}
&  h_{r}\left(  x\parallel a\right) \nonumber\\
&  =\left[  t^{r}\right]  \left(  \prod_{i\geq1}\left(  1-tx_{i}\right)
^{-1}\prod_{j\geq1}\left(  1+ta_{j}\right)  \right) \nonumber\\
&  =\sum_{j=0}^{r}\left[  t^{j}\right]  \left(  \prod_{i\geq1}\left(
1-tx_{i}\right)  ^{-1}\right)  \cdot\left[  t^{r-j}\right]  \left(
\prod_{j\geq1}\left(  1+ta_{j}\right)  \right)  \label{pf.lem.p21.TA.1}%
\end{align}
(by the formula $\left[  t^{r}\right]  \left(  fg\right)  =\sum_{j=0}%
^{r}\left(  \left[  t^{j}\right]  f\right)  \cdot\left(  \left[
t^{r-j}\right]  g\right)  $ for the product of two power series $f$ and $g$).
But each $j\in\left\{  0,1,\ldots,r\right\}  $ satisfies%
\[
\left[  t^{j}\right]  \left(  \prod_{i\geq1}\left(  1-tx_{i}\right)
^{-1}\right)  =h_{j}\left(  x\right)  \ \ \ \ \ \ \ \ \ \ \left(  \text{since
}\prod_{i\geq1}\left(  1-tx_{i}\right)  ^{-1}=\sum_{i\in\mathbb{N}}%
h_{i}\left(  x\right)  t^{i}\right)
\]
and%
\[
\left[  t^{r-j}\right]  \left(  \prod_{j\geq1}\left(  1+ta_{j}\right)
\right)  =e_{r-j}\left(  a\right)  \ \ \ \ \ \ \ \ \ \ \left(  \text{since
}\prod_{j\geq1}\left(  1+ta_{j}\right)  =\sum_{i\in\mathbb{N}}e_{i}\left(
a\right)  t^{i}\right)  .
\]
Thus, we can rewrite (\ref{pf.lem.p21.TA.1}) as%
\begin{align*}
&  h_{r}\left(  x\parallel a\right) \\
&  =\sum_{j=0}^{r}\underbrace{\left[  t^{j}\right]  \left(  \prod_{i\geq
1}\left(  1-tx_{i}\right)  ^{-1}\right)  }_{\substack{=h_{j}\left(  x\right)
\\=s_{\left(  j\right)  }\left(  x\right)  }}\cdot\underbrace{\left[
t^{r-j}\right]  \left(  \prod_{j\geq1}\left(  1+ta_{j}\right)  \right)
}_{\substack{=e_{r-j}\left(  a\right)  \\=s_{\left(  1^{r-j}\right)  }\left(
a\right)  \\=s_{\left(  r\right)  ^{\prime}/\left(  j\right)  ^{\prime}%
}\left(  a\right)  \\\text{(since the partition }\left(  1^{r-j}\right)
\text{ and the skew partition }\left(  r\right)  ^{\prime}/\left(  j\right)
^{\prime}\\\text{have the same shape, just positioned differently)}}}\\
&  =\sum_{j=0}^{r}s_{\left(  j\right)  }\left(  x\right)  s_{\left(  r\right)
^{\prime}/\left(  j\right)  ^{\prime}}\left(  a\right)  .
\end{align*}
Comparing this with%
\begin{align*}
&  \lim\limits_{n\rightarrow\infty}\underbrace{h_{r}\left(  x^{\left(
n\right)  }\mid a\right)  }_{\substack{=s_{\left(  r\right)  }\left(
x^{\left(  n\right)  }\mid a\right)  }}\\
&  =\lim\limits_{n\rightarrow\infty}s_{\left(  r\right)  }\left(  x^{\left(
n\right)  }\mid a\right) \\
&  =\sum_{\mu\subseteq\left(  r\right)  }s_{\mu}\left(  x\right)  s_{\left(
r\right)  ^{\prime}/\mu^{\prime}}\left(  a\right)  \ \ \ \ \ \ \ \ \ \ \left(
\text{by Lemma \ref{lem.p21.TO}, applied to }\lambda=\left(  r\right)  \right)
\\
&  =\sum_{j=0}^{r}s_{\left(  j\right)  }\left(  x\right)  s_{\left(  r\right)
^{\prime}/\left(  j\right)  ^{\prime}}\left(  a\right)
\end{align*}
(since the partitions $\mu$ that satisfy $\mu\subseteq\left(  r\right)  $ are
precisely the partitions of the form $\left(  j\right)  $ for $j\in\left\{
0,1,\ldots,r\right\}  $), we obtain%
\[
h_{r}\left(  x\parallel a\right)  =\lim\limits_{n\rightarrow\infty}%
h_{r}\left(  x^{\left(  n\right)  }\mid a\right)  .
\]
This proves Lemma \ref{lem.p21.TA}.
\end{proof}

\begin{lemmain}
\label{lem.p21.TB}Let $k\geq1$ and $r$ be two integers. If we set all the
variables $a_{0},a_{-1},a_{-2},\ldots$ to $0$, then $h_{r}\left(
x\parallel\tau^{1-k}a\right)  $ becomes $h_{r}\left(  x\parallel a\right)  $.
\end{lemmain}

\begin{proof}
[Proof of Lemma \ref{lem.p21.TB}.]Recall that $h_{r}\left(  x\parallel
a\right)  $ is defined as the coefficient of $t^{r}$ in the power series
\[
\prod_{i\geq1}\left(  1-tx_{i}\right)  ^{-1}\prod_{j\geq1}\left(
1+ta_{j}\right)  \in\left(  \mathbb{Z}\left[  \left[  x_{i},a_{i}%
\ \mid\ i>0\right]  \right]  \right)  \left[  \left[  t\right]  \right]  .
\]
Hence, $h_{r}\left(  x\parallel\tau^{1-k}a\right)  $ is the coefficient of
$t^{r}$ in the power series
\[
\prod_{i\geq1}\left(  1-tx_{i}\right)  ^{-1}\prod_{j\geq1}\left(
1+ta_{j+1-k}\right)
\]
(since the $j$-th entry of the family $\tau^{1-k}a$ is $a_{j+1-k}$). If we set
all the variables $a_{0},a_{-1},a_{-2},\ldots$ to $0$, then the product
\begin{align*}
\prod_{j\geq1}\left(  1+ta_{j+1-k}\right)   &  =\left(  1+ta_{2-k}\right)
\left(  1+ta_{3-k}\right)  \left(  1+ta_{4-k}\right)  \cdots\\
&  =\left(  1+ta_{2-k}\right)  \left(  1+ta_{3-k}\right)  \cdots\left(
1+ta_{0}\right)  \cdot\prod_{j\geq1}\left(  1+ta_{j}\right)
\end{align*}
is turned into
\[
\underbrace{\left(  1+t\cdot0\right)  \left(  1+t\cdot0\right)  \cdots\left(
1+t\cdot0\right)  }_{=1\cdot1\cdot\cdots\cdot1=1}\cdot\prod_{j\geq1}\left(
1+ta_{j}\right)  =\prod_{j\geq1}\left(  1+ta_{j}\right)  ,
\]
and therefore the whole power series%
\[
\prod_{i\geq1}\left(  1-tx_{i}\right)  ^{-1}\prod_{j\geq1}\left(
1+ta_{j+1-k}\right)
\]
is turned into the power series%
\[
\prod_{i\geq1}\left(  1-tx_{i}\right)  ^{-1}\prod_{j\geq1}\left(
1+ta_{j}\right)  .
\]
Thus, $h_{r}\left(  x\parallel\tau^{1-k}a\right)  $ is turned into
$h_{r}\left(  x\parallel a\right)  $ (since $h_{r}\left(  x\parallel\tau
^{1-k}a\right)  $ is the coefficient of $t^{r}$ in the former power series,
while $h_{r}\left(  x\parallel a\right)  $ is the coefficient of $t^{r}$ in
the latter power series). This proves Lemma \ref{lem.p21.TB}.
\end{proof}

\begin{lemmain}
\label{lem.p21.slam/0}Let $\lambda$ be a partition. Let $n\in\mathbb{N}$ be
such that $\ell\left(  \lambda\right)  \leq n$. Then,%
\[
s_{\lambda}\left(  x^{\left(  n\right)  }\mid a\right)  =s_{\lambda
/\varnothing}\left(  x^{\left(  n\right)  }\mid a\right)  .
\]

\end{lemmain}

\begin{proof}
[Proof of Lemma \ref{lem.p21.slam/0}.]Write $\lambda$ as $\lambda=\left(
\lambda_{1},\lambda_{2},\ldots,\lambda_{m}\right)  $, where $m=\ell\left(
\lambda\right)  $. Then, $m\geq\ell\left(  \lambda\right)  $ and $m\geq
0=\ell\left(  \varnothing\right)  $. Hence, (6.8) (applied to $\mu
=\varnothing$) yields\footnote{Keep in mind that our $x^{\left(  n\right)  }$
is what was called $x$ in (6.8).}
\begin{align*}
s_{\lambda/\varnothing}\left(  x^{\left(  n\right)  }\mid a\right)   &
=\det\left(  h_{\lambda_{i}-\varnothing_{j}-i+j}\left(  x^{\left(  n\right)
}\mid\tau^{\varnothing_{j}-j+1}a\right)  \right)  _{i,j\in\left[  n\right]
}\\
&  =\det\left(  h_{\lambda_{i}-i+j}\left(  x^{\left(  n\right)  }\mid
\tau^{1-j}a\right)  \right)  _{i,j\in\left[  n\right]  }%
\end{align*}
(since $\lambda_{i}-\underbrace{\varnothing_{j}}_{=0}=\lambda_{i}$ and
$\underbrace{\varnothing_{j}}_{=0}-\,j+1=-j+1=1-j$). Comparing this with%
\[
s_{\lambda}\left(  x^{\left(  n\right)  }\mid a\right)  =\det\left(
h_{\lambda_{i}-i+j}\left(  x^{\left(  n\right)  }\mid\tau^{1-j}a\right)
\right)  _{i,j\in\left[  n\right]  }\ \ \ \ \ \ \ \ \ \ \left(  \text{by
(6.7)}\right)  ,
\]
we obtain $s_{\lambda}\left(  x^{\left(  n\right)  }\mid a\right)
=s_{\lambda/\varnothing}\left(  x^{\left(  n\right)  }\mid a\right)  $. This
proves Lemma \ref{lem.p21.slam/0}.
\end{proof}

\begin{lemmain}
\label{lem.p21.slam/mu=det}Let $\lambda=\left(  \lambda_{1},\lambda_{2}%
,\ldots,\lambda_{m}\right)  $ and $\mu=\left(  \mu_{1},\mu_{2},\ldots,\mu
_{m}\right)  $ be two partitions. Then,%
\[
\lim\limits_{n\rightarrow\infty}s_{\lambda/\mu}\left(  x^{\left(  n\right)
}\mid a\right)  =\det\left(  h_{\lambda_{i}-\mu_{j}-i+j}\left(  x\parallel
\tau^{\mu_{j}-j+1}a\right)  \right)  _{i,j\in\left[  m\right]  }.
\]

\end{lemmain}

\begin{proof}
[Proof of Lemma \ref{lem.p21.slam/mu=det}.]For any integer $n\geq m$, we
have\footnote{Keep in mind that our $x^{\left(  n\right)  }$ is what was
called $x$ in (6.8).}%
\[
s_{\lambda/\mu}\left(  x^{\left(  n\right)  }\mid a\right)  =\det\left(
h_{\lambda_{i}-\mu_{j}-i+j}\left(  x^{\left(  n\right)  }\mid\tau^{\mu
_{j}-j+1}a\right)  \right)  _{i,j\in\left[  m\right]  }%
\ \ \ \ \ \ \ \ \ \ \left(  \text{by (6.8)}\right)  .
\]
Taking the limit as $n\rightarrow\infty$ on both sides of this equality, we
obtain\footnote{All the limits here are well-defined, due to the equality
(\ref{pf.lem.p21.slam/mu=det.limh}) below.}%
\begin{align}
&  \lim\limits_{n\rightarrow\infty}s_{\lambda/\mu}\left(  x^{\left(  n\right)
}\mid a\right) \nonumber\\
&  =\lim\limits_{n\rightarrow\infty}\det\left(  h_{\lambda_{i}-\mu_{j}%
-i+j}\left(  x^{\left(  n\right)  }\mid\tau^{\mu_{j}-j+1}a\right)  \right)
_{i,j\in\left[  m\right]  }\nonumber\\
&  =\det\left(  \lim\limits_{n\rightarrow\infty}h_{\lambda_{i}-\mu_{j}%
-i+j}\left(  x^{\left(  n\right)  }\mid\tau^{\mu_{j}-j+1}a\right)  \right)
_{i,j\in\left[  m\right]  }. \label{pf.lem.p21.slam/mu=det.3}%
\end{align}
However, for all $i,j\in\left[  m\right]  $, we have%
\begin{align}
&  h_{\lambda_{i}-\mu_{j}-i+j}\left(  x\parallel\tau^{\mu_{j}-j+1}a\right)
\nonumber\\
&  =\lim\limits_{n\rightarrow\infty}h_{\lambda_{i}-\mu_{j}-i+j}\left(
x^{\left(  n\right)  }\mid\tau^{\mu_{j}-j+1}a\right)
\label{pf.lem.p21.slam/mu=det.limh}%
\end{align}
(by Lemma \ref{lem.p21.TA}, applied to $\tau^{\mu_{j}-j+1}a$ and $\lambda
_{i}-\mu_{j}-i+j$ instead of $a$ and $r$). Thus,
\begin{align*}
&  \det\left(  h_{\lambda_{i}-\mu_{j}-i+j}\left(  x\parallel\tau^{\mu_{j}%
-j+1}a\right)  \right)  _{i,j\in\left[  m\right]  }\\
&  =\det\left(  \lim\limits_{n\rightarrow\infty}h_{\lambda_{i}-\mu_{j}%
-i+j}\left(  x^{\left(  n\right)  }\mid\tau^{\mu_{j}-j+1}a\right)  \right)
_{i,j\in\left[  m\right]  }.
\end{align*}
Comparing this with (\ref{pf.lem.p21.slam/mu=det.3}), we obtain%
\[
\lim\limits_{n\rightarrow\infty}s_{\lambda/\mu}\left(  x^{\left(  n\right)
}\mid a\right)  =\det\left(  h_{\lambda_{i}-\mu_{j}-i+j}\left(  x\parallel
\tau^{\mu_{j}-j+1}a\right)  \right)  _{i,j\in\left[  m\right]  }.
\]
This proves Lemma \ref{lem.p21.slam/mu=det}.
\end{proof}

We can now prove (\ref{p21.6.19.rewr}):

\begin{lemmain}
\label{lem.p21.6.19}Let $\lambda$ be a partition. Then,%
\[
\lim\limits_{n\rightarrow\infty}s_{\lambda}\left(  x^{\left(  n\right)  }\mid
a\right)  =s_{\lambda}\left(  x\parallel a\right)  .
\]

\end{lemmain}

\begin{proof}
[Proof of Lemma \ref{lem.p21.6.19}.]Write $\lambda$ as $\lambda=\left(
\lambda_{1},\lambda_{2},\ldots,\lambda_{m}\right)  $, where $m=\ell\left(
\lambda\right)  $. Then, the definition of $s_{\lambda}\left(  x\parallel
a\right)  $ yields%
\begin{equation}
s_{\lambda}\left(  x\parallel a\right)  =\det\left(  h_{\lambda_{i}%
-i+j}\left(  x\parallel a\right)  \right)  _{i,j\in\left[  m\right]  }.
\label{pf.lem.p21.6.19.0}%
\end{equation}


But%
\begin{align*}
\lim\limits_{n\rightarrow\infty}\underbrace{s_{\lambda}\left(  x^{\left(
n\right)  }\mid a\right)  }_{\substack{=s_{\lambda/\varnothing}\left(
x^{\left(  n\right)  }\mid a\right)  \text{ whenever }n\geq\ell\left(
\lambda\right)  \\\text{(by Lemma \ref{lem.p21.slam/0})}}}  &  =\lim
\limits_{n\rightarrow\infty}s_{\lambda/\varnothing}\left(  x^{\left(
n\right)  }\mid a\right) \\
&  =\det\left(  h_{\lambda_{i}-\varnothing_{j}-i+j}\left(  x\parallel
\tau^{\varnothing_{j}-j+1}a\right)  \right)  _{i,j\in\left[  m\right]  }%
\end{align*}
(by Lemma \ref{lem.p21.slam/mu=det}, applied to $\mu=\varnothing$). Since all
$i,j\in\left[  m\right]  $ satisfy $\lambda_{i}-\underbrace{\varnothing_{j}%
}_{=0}=\lambda_{i}$ and $\varnothing_{j}-j+1=1-j+\underbrace{\varnothing_{j}%
}_{=0}=1-j$, we can rewrite this as%
\begin{equation}
\lim\limits_{n\rightarrow\infty}s_{\lambda}\left(  x^{\left(  n\right)  }\mid
a\right)  =\det\left(  h_{\lambda_{i}-i+j}\left(  x\parallel\tau
^{1-j}a\right)  \right)  _{i,j\in\left[  m\right]  }.
\label{pf.lem.p21.6.19.5}%
\end{equation}


Now, let us set all the variables $a_{0},a_{-1},a_{-2},\ldots$ to $0$. Then,
for any $i,j\in\left[  m\right]  $, the polynomial $h_{\lambda_{i}-i+j}\left(
x\parallel\tau^{1-j}a\right)  $ on the right hand side of
(\ref{pf.lem.p21.6.19.5}) becomes $h_{\lambda_{i}-i+j}\left(  x\parallel
a\right)  $ (by Lemma \ref{lem.p21.TB}, applied to $r=\lambda_{i}-i+j$ and
$k=j$), whereas the left hand side of (\ref{pf.lem.p21.6.19.5}) stays
unchanged (since the definition (6.4) of $s_{\lambda}\left(  x^{\left(
n\right)  }\mid a\right)  $ does not involve the variables $a_{0}%
,a_{-1},a_{-2},\ldots$ at all). Thus, the equality (\ref{pf.lem.p21.6.19.5})
turns into%
\[
\lim\limits_{n\rightarrow\infty}s_{\lambda}\left(  x^{\left(  n\right)  }\mid
a\right)  =\det\left(  h_{\lambda_{i}-i+j}\left(  x\parallel a\right)
\right)  _{i,j\in\left[  m\right]  }.
\]
Comparing this with (\ref{pf.lem.p21.6.19.0}), we obtain%
\[
\lim\limits_{n\rightarrow\infty}s_{\lambda}\left(  x^{\left(  n\right)  }\mid
a\right)  =s_{\lambda}\left(  x\parallel a\right)  .
\]
Thus, Lemma \ref{lem.p21.6.19} is proved.
\end{proof}

In other words, (\ref{p21.6.19.rewr}) and thus (6.19) is proved.
$\blacksquare$

\begin{noncompile}
We can actually go a bit further and prove a generalization of (6.19):

\begin{lemmain}
\label{lem.p21.6.19skew}Let $\lambda$ and $\mu$ be partitions such that
$\lambda\supseteq\mu$. Let $\mu^{\prime}=\left(  \mu_{1}^{\prime},\mu
_{2}^{\prime},\mu_{3}^{\prime},\ldots\right)  $ be the conjugate partition of
$\mu$. Assume that each $j\geq1$ satisfies%
\begin{equation}
\mu_{j}^{\prime}<j\ \ \ \ \ \ \ \ \ \ \text{and}\ \ \ \ \ \ \ \ \ \ \mu_{j}<j.
\label{eq.lem.p21.6.19skew.ass}%
\end{equation}
Define%
\begin{equation}
s_{\lambda/\mu}\left(  x\parallel a\right)  :=\det\left(  h_{\lambda_{i}%
-\mu_{j}-i+j}\left(  x\parallel a\right)  \right)  _{i,j\in\left[  m\right]
}, \label{eq.lem.p21.6.19skew.def}%
\end{equation}
where $m$ is an integer satisfying $\ell\left(  \lambda\right)  \leq m$ and
$\ell\left(  \mu\right)  \leq m$. Then,%
\[
\lim\limits_{n\rightarrow\infty}s_{\lambda/\mu}\left(  x^{\left(  n\right)
}\mid a\right)  =s_{\lambda/\mu}\left(  x\parallel a\right)  .
\]

\end{lemmain}

\begin{proof}
[Proof of Lemma \ref{lem.p21.6.19skew}.]Let $m$ be an integer satisfying
$\ell\left(  \lambda\right)  \leq m$ and $\ell\left(  \mu\right)  \leq m$.

Let $n\in\mathbb{N}$. If $T$ is any tableau of shape $\lambda-\mu$ with
positive integer entries, and if $\left(  i,j\right)  \in\lambda-\mu$ is a
cell of the skew diagram $\lambda-\mu$, then $T^{\ast}\left(  i,j\right)  $ is
defined as the integer $T\left(  i,j\right)  +j-i$. Then, (6.16) says that%
\begin{equation}
s_{\lambda/\mu}\left(  x^{\left(  n\right)  }\mid a\right)  =\sum
_{\substack{T:\lambda-\mu\rightarrow\left[  1,n\right]  \text{ is
a}\\\text{column-strict tableau}}}\left(  x\mid a\right)  ^{T},
\label{pf.lem.p21.6.19skew.1}%
\end{equation}
where%
\begin{equation}
\left(  x\mid a\right)  ^{T}=\prod_{s\in\lambda-\mu}\left(  x_{T\left(
s\right)  }+a_{T^{\ast}\left(  s\right)  }\right)  .
\label{pf.lem.p21.6.19skew.2}%
\end{equation}


Now, we claim that each column-strict tableau $T:\lambda-\mu\rightarrow\left[
1,n\right]  $ and each cell $s\in\lambda-\mu$ satisfy%
\begin{equation}
T^{\ast}\left(  s\right)  >0. \label{pf.lem.p21.6.19skew.pos}%
\end{equation}


[\textit{Proof:} Let $T:\lambda-\mu\rightarrow\left[  1,n\right]  $ be a
column-strict tableau, and let $s\in\lambda-\mu$ be a cell. Write $s$ as
$\left(  i,j\right)  $. Then, $\mu_{j}^{\prime}<i$ (since $\left(  i,j\right)
=s\in\lambda-\mu$), and furthermore, all the $i-\mu_{j}^{\prime}$ cells
$\left(  \mu_{j}^{\prime}+1,\ j\right)  ,\ \left(  \mu_{j}^{\prime
}+2,\ j\right)  ,\ \ldots,\ \left(  i,\ j\right)  $ belong to $\lambda-\mu$.
Since $T$ is column-strict, we have $T\left(  k,\ j\right)  <T\left(
k+1,\ j\right)  $ for each $k\in\left[  \mu_{j}^{\prime}+1,\ i-1\right]  $.
Thus, $T\left(  k,\ j\right)  \leq T\left(  k+1,\ j\right)  -1$ for each
$k\in\left[  \mu_{j}^{\prime}+1,\ i-1\right]  $ (since $T\left(  k,\ j\right)
$ and $T\left(  k+1,\ j\right)  $ are integers). Adding $j-k$ to both sides of
this inequality, we obtain%
\[
T\left(  k,\ j\right)  +j-k\leq T\left(  k+1,\ j\right)  -1+j-k=T\left(
k+1,\ j\right)  +j-\left(  k+1\right)  .
\]
So we have shown that $T\left(  k,\ j\right)  +j-k\leq T\left(
k+1,\ j\right)  +j-\left(  k+1\right)  $ for each $k\in\left[  \mu_{j}%
^{\prime}+1,\ i-1\right]  $. In other words, the function%
\begin{align*}
\left[  \mu_{j}^{\prime}+1,\ i\right]   &  \rightarrow\mathbb{Z},\\
k  &  \mapsto T\left(  k,\ j\right)  +j-k
\end{align*}
is weakly increasing. Hence, its value at $\mu_{j}^{\prime}+1$ is $\leq$ to
its value at $i$. In other words,%
\[
T\left(  \mu_{j}^{\prime}+1,\ j\right)  +j-\left(  \mu_{j}^{\prime}+1\right)
\leq T\left(  i,\ j\right)  +j-i.
\]
Hence,%
\[
T\left(  i,\ j\right)  +j-i\geq\underbrace{T\left(  \mu_{j}^{\prime
}+1,\ j\right)  }_{>0}+\underbrace{j-\left(  \mu_{j}^{\prime}+1\right)
}_{\substack{\geq0\\\text{(since }\mu_{j}^{\prime}+1\leq j\text{ (because }%
\mu_{j}^{\prime}<j\\\text{(by (\ref{eq.lem.p21.6.19skew.ass}))))}}}>0.
\]
In other words, $T^{\ast}\left(  i,j\right)  >0$ (since $T^{\ast}\left(
i,j\right)  $ is defined as the integer $T\left(  i,j\right)  +j-i$). In other
words, $T^{\ast}\left(  s\right)  >0$ (since $s=\left(  i,j\right)  $). This
proves (\ref{pf.lem.p21.6.19skew.pos}).] \medskip

Now, we draw some conclusions. If $T$ is any tableau of shape $\lambda-\mu$
with positive integer entries, then all the addends $a_{T^{\ast}\left(
s\right)  }$ on the right hand side of (\ref{pf.lem.p21.6.19skew.2}) satisfy
$T^{\ast}\left(  s\right)  >0$ (by (\ref{pf.lem.p21.6.19skew.pos})); thus, the
right hand side of (\ref{pf.lem.p21.6.19skew.2}) does not involve the
variables $a_{0},a_{-1},a_{-2},\ldots$ at all. Hence, neither does the left
hand side. In other words, if $T$ is any tableau of shape $\lambda-\mu$ with
positive integer entries, then the polynomial $\left(  x\mid a\right)  ^{T}$
does not involve the variables $a_{0},a_{-1},a_{-2},\ldots$ at all. Therefore,
(\ref{pf.lem.p21.6.19skew.1}) shows that the polynomial $s_{\lambda/\mu
}\left(  x^{\left(  n\right)  }\mid a\right)  $ does not involve the variables
$a_{0},a_{-1},a_{-2},\ldots$ at all.

On the other hand, (6.8) yields\footnote{Keep in mind that our $x^{\left(
n\right)  }$ is what was called $x$ in (6.8).}%
\begin{align*}
s_{\lambda/\mu}\left(  x^{\left(  n\right)  }\mid a\right)   &  =\det\left(
h_{\lambda_{i}-\mu_{j}-i+j}\left(  x^{\left(  n\right)  }\mid\tau^{\mu
_{j}-j+1}a\right)  \right)  _{i,j\in\left[  m\right]  }\\
&  =\det\left(  h_{\lambda_{i}-\mu_{j}-i+j}\left(  x^{\left(  n\right)  }%
\mid\tau^{1-\left(  j-\mu_{j}\right)  }a\right)  \right)  _{i,j\in\left[
m\right]  }%
\end{align*}
(since $\mu_{j}-j+1=1-\left(  j-\mu_{j}\right)  $ for each $j\in\left[
m\right]  $). Taking the limit as $n\rightarrow\infty$ on both sides of this
equality, we obtain\footnote{All the limits here are well-defined, due to the
equality (\ref{pf.lem.p21.6.19skew.limh}) below.}%
\begin{align}
&  \lim\limits_{n\rightarrow\infty}s_{\lambda/\mu}\left(  x^{\left(  n\right)
}\mid a\right) \nonumber\\
&  =\lim\limits_{n\rightarrow\infty}\det\left(  h_{\lambda_{i}-\mu_{j}%
-i+j}\left(  x^{\left(  n\right)  }\mid\tau^{1-\left(  j-\mu_{j}\right)
}a\right)  \right)  _{i,j\in\left[  m\right]  }\nonumber\\
&  =\det\left(  \lim\limits_{n\rightarrow\infty}h_{\lambda_{i}-\mu_{j}%
-i+j}\left(  x^{\left(  n\right)  }\mid\tau^{1-\left(  j-\mu_{j}\right)
}a\right)  \right)  _{i,j\in\left[  m\right]  }. \label{pf.lem.p21.6.19skew.3}%
\end{align}
However, for all $i,j\in\left[  m\right]  $, we have%
\begin{align}
&  h_{\lambda_{i}-\mu_{j}-i+j}\left(  x\parallel\tau^{1-\left(  j-\mu
_{j}\right)  }a\right) \nonumber\\
&  =\lim\limits_{n\rightarrow\infty}h_{\lambda_{i}-\mu_{j}-i+j}\left(
x^{\left(  n\right)  }\mid\tau^{1-\left(  j-\mu_{j}\right)  }a\right)
\label{pf.lem.p21.6.19skew.limh}%
\end{align}
(by Lemma \ref{lem.p21.TA}, applied to $\tau^{1-\left(  j-\mu_{j}\right)  }a$
and $\lambda_{i}-\mu_{j}-i+j$ instead of $a$ and $r$). Thus,
\begin{align*}
&  \det\left(  h_{\lambda_{i}-\mu_{j}-i+j}\left(  x\parallel\tau^{1-\left(
j-\mu_{j}\right)  }a\right)  \right)  _{i,j\in\left[  m\right]  }\\
&  =\det\left(  \lim\limits_{n\rightarrow\infty}h_{\lambda_{i}-\mu_{j}%
-i+j}\left(  x^{\left(  n\right)  }\mid\tau^{1-\left(  j-\mu_{j}\right)
}a\right)  \right)  _{i,j\in\left[  m\right]  }.
\end{align*}
Comparing this with (\ref{pf.lem.p21.6.19skew.3}), we obtain%
\begin{align}
&  \lim\limits_{n\rightarrow\infty}s_{\lambda/\mu}\left(  x^{\left(  n\right)
}\mid a\right) \nonumber\\
&  =\det\left(  h_{\lambda_{i}-\mu_{j}-i+j}\left(  x\parallel\tau^{1-\left(
j-\mu_{j}\right)  }a\right)  \right)  _{i,j\in\left[  m\right]  }.
\label{pf.lem.p21.6.19skew.5}%
\end{align}


Now, let us set all the variables $a_{0},a_{-1},a_{-2},\ldots$ to $0$. Then,
for any $i,j\in\left[  m\right]  $, the polynomial $h_{\lambda_{i}-\mu
_{j}-i+j}\left(  x\parallel\tau^{1-\left(  j-\mu_{j}\right)  }a\right)  $ on
the right hand side of (\ref{pf.lem.p21.6.19skew.5}) becomes $h_{\lambda
_{i}-\mu_{j}-i+j}\left(  x\parallel a\right)  $ (by Lemma \ref{lem.p21.TB},
applied to $r=\lambda_{i}-\mu_{j}-i+j$ and $k=j-\mu_{j}$\ \ \ \ \footnote{Here
we are using the fact that $j-\mu_{j}\geq1$. This is because
(\ref{eq.lem.p21.6.19skew.ass}) yields $\mu_{j}<j$, so that $\mu_{j}+1\leq
j$.}), whereas the left hand side of (\ref{pf.lem.p21.6.19skew.5}) stays
unchanged (since the polynomial $s_{\lambda/\mu}\left(  x^{\left(  n\right)
}\mid a\right)  $ does not involve the variables $a_{0},a_{-1},a_{-2},\ldots$
at all). Thus, the equality (\ref{pf.lem.p21.6.19skew.5}) turns into%
\[
\lim\limits_{n\rightarrow\infty}s_{\lambda/\mu}\left(  x^{\left(  n\right)
}\mid a\right)  =\det\left(  h_{\lambda_{i}-\mu_{j}-i+j}\left(  x\parallel
a\right)  \right)  _{i,j\in\left[  m\right]  }.
\]
Comparing this with (\ref{eq.lem.p21.6.19skew.def}), we obtain%
\[
\lim\limits_{n\rightarrow\infty}s_{\lambda/\mu}\left(  x^{\left(  n\right)
}\mid a\right)  =s_{\lambda/\mu}\left(  x\parallel a\right)  .
\]
Thus, Lemma \ref{lem.p21.6.19skew} is proved.
\end{proof}
\end{noncompile}

\item \textbf{page 22, (6.21):} Let me sketch the (somewhat slippery) proof of (6.21).

We will use some previously proved results such as the identity (6.16) and
Lemma \ref{lem.p21.slam/mu=det}, but we must account for the fact that
\textquotedblleft$x$\textquotedblright\ now has a different meaning than in
these latter results. Namely, the letter $x$ now denotes the doubly infinite
sequence $\left(  x_{n}\right)  _{n\in\mathbb{Z}}$, whereas it previously used
to denote either the sequence $\left(  x_{1},x_{2},x_{3},\ldots\right)  $ (in
Lemma \ref{lem.p21.slam/mu=det} and in (6.20)) or a finite tuple $\left(
x_{1},x_{2},\ldots,x_{n}\right)  $ (in (6.16)). Thus, let us introduce a new
notation. Namely, for any $k\in\mathbb{Z}$, we define two sequences $x_{\geq
k}$ and $a_{\geq k}$ of indeterminates by%
\begin{align*}
x_{\geq k}  &  :=\left(  x_{k},x_{k+1},x_{k+2},\ldots\right)
\ \ \ \ \ \ \ \ \ \ \text{and}\\
a_{\geq k}  &  :=\left(  a_{k},a_{k+1},a_{k+2},\ldots\right)  .
\end{align*}
Thus, $x_{\geq1}=\left(  x_{1},x_{2},x_{3},\ldots\right)  $ is the sequence
that was denoted by $x$ in Lemma \ref{lem.p21.slam/mu=det}. The finite tuple
$\left(  x_{1},x_{2},\ldots,x_{n}\right)  $ will be denoted by $x^{\left(
n\right)  }$, as before.

In these new notations, (6.16) rewrites as%
\begin{equation}
s_{\lambda/\mu}\left(  x^{\left(  n\right)  }\mid a\right)  =\sum
_{\substack{T:\lambda-\mu\rightarrow\left[  1,n\right]  \text{ is
a}\\\text{column-strict tableau}}}\left(  x^{\left(  n\right)  }\mid a\right)
^{T}. \label{p22.6.21.6.16}%
\end{equation}


Another notational discrepancy concerns the meaning of \textquotedblleft%
$h_{r}\left(  x\parallel a\right)  $\textquotedblright. We agree to use the
notation $h_{r}\left(  x\parallel a\right)  $ for the coefficient of $t^{r}$
in the power series $\prod_{i\in\mathbb{Z}}\left(  1-tx_{i}\right)  ^{-1}%
\prod_{j\in\mathbb{Z}}\left(  1+ta_{j}\right)  $, whereas the notation
$h_{r}\left(  x_{\geq p}\parallel a_{\geq q}\right)  $ (for given integers $p$
and $q$) shall denote the coefficient of $t^{r}$ in the power series
$\prod_{i\geq p}\left(  1-tx_{i}\right)  ^{-1}\prod_{j\geq q}\left(
1+ta_{j}\right)  $. Thus, what was denoted by $h_{r}\left(  x\parallel
a\right)  $ in (6.19) is actually $h_{r}\left(  x_{\geq1}\parallel a_{\geq
1}\right)  $ in our new notation. More generally, for any integers $r$ and
$q$, the power series that was denoted by $h_{r}\left(  x\parallel\tau
^{q}a\right)  $ in Lemma \ref{lem.p21.slam/mu=det} is actually $h_{r}\left(
x_{\geq1}\parallel a_{\geq q+1}\right)  $ in our new notation, since it is the
coefficient of $t^{r}$ in the power series
\begin{align*}
&  \prod_{i\geq1}\left(  1-tx_{i}\right)  ^{-1}\prod_{j\geq1}\left(
1+t\underbrace{\left(  \tau^{q}a\right)  _{j}}_{=a_{q+j}}\right) \\
&  =\prod_{i\geq1}\left(  1-tx_{i}\right)  ^{-1}\prod_{j\geq1}\left(
1+ta_{q+j}\right) \\
&  =\prod_{i\geq1}\left(  1-tx_{i}\right)  ^{-1}\prod_{j\geq q+1}\left(
1+ta_{j}\right)  \ \ \ \ \ \ \ \ \ \ \left(
\begin{array}
[c]{c}%
\text{here, we have substituted }j\\
\text{for }q+j\text{ in the second product}%
\end{array}
\right)  .
\end{align*}


Thus, Lemma \ref{lem.p21.slam/mu=det} now rewrites as follows:

\begin{lemmain}
\label{lem.p21.slam/mu=det.ver2}Let $\lambda=\left(  \lambda_{1},\lambda
_{2},\ldots,\lambda_{m}\right)  $ and $\mu=\left(  \mu_{1},\mu_{2},\ldots
,\mu_{m}\right)  $ be two partitions. Then,%
\[
\lim\limits_{n\rightarrow\infty}s_{\lambda/\mu}\left(  x^{\left(  n\right)
}\mid a\right)  =\det\left(  h_{\lambda_{i}-\mu_{j}-i+j}\left(  x_{\geq
1}\parallel a_{\geq\mu_{j}-j+1}\right)  \right)  _{i,j\in\left[  m\right]  }.
\]

\end{lemmain}

We shall use the notation%
\[
\left(  x\mid a\right)  ^{T}:=\prod_{s\in\lambda-\mu}\left(  x_{T\left(
s\right)  }+a_{T^{\ast}\left(  s\right)  }\right)
\]
whenever $T:\lambda-\mu\rightarrow\mathbb{Z}$ is a column-strict tableau with
integer entries.

We can now easily obtain the following:

\begin{lemmain}
\label{lem.p21.slam/mu=tab}Let $\lambda=\left(  \lambda_{1},\lambda_{2}%
,\ldots,\lambda_{m}\right)  $ and $\mu=\left(  \mu_{1},\mu_{2},\ldots,\mu
_{m}\right)  $ be two partitions. Then,%
\[
\det\left(  h_{\lambda_{i}-\mu_{j}-i+j}\left(  x_{\geq1}\parallel a_{\geq
\mu_{j}-j+1}\right)  \right)  _{i,j\in\left[  m\right]  }=\sum
_{\substack{T:\lambda-\mu\rightarrow\left\{  1,2,3,\ldots\right\}  \text{ is
a}\\\text{column-strict tableau}}}\left(  x\mid a\right)  ^{T}.
\]

\end{lemmain}

\begin{proof}
[Proof of Lemma \ref{lem.p21.slam/mu=tab}.]Lemma
\ref{lem.p21.slam/mu=det.ver2} yields%
\begin{align*}
&  \det\left(  h_{\lambda_{i}-\mu_{j}-i+j}\left(  x_{\geq1}\parallel
a_{\geq\mu_{j}-j+1}\right)  \right)  _{i,j\in\left[  m\right]  }\\
&  =\lim\limits_{n\rightarrow\infty}s_{\lambda/\mu}\left(  x^{\left(
n\right)  }\mid a\right) \\
&  =\lim\limits_{n\rightarrow\infty}\sum_{\substack{T:\lambda-\mu
\rightarrow\left[  1,n\right]  \text{ is a}\\\text{column-strict tableau}%
}}\underbrace{\left(  x^{\left(  n\right)  }\mid a\right)  ^{T}}%
_{\substack{=\left(  x\mid a\right)  ^{T}\\\text{(since both }\left(
x^{\left(  n\right)  }\mid a\right)  ^{T}\text{ and }\left(  x\mid a\right)
^{T}\text{ are}\\\text{defined as }\prod_{s\in\lambda-\mu}\left(  x_{T\left(
s\right)  }+a_{T^{\ast}\left(  s\right)  }\right)  \text{)}}%
}\ \ \ \ \ \ \ \ \ \ \left(  \text{by (\ref{p22.6.21.6.16})}\right) \\
&  =\lim\limits_{n\rightarrow\infty}\sum_{\substack{T:\lambda-\mu
\rightarrow\left[  1,n\right]  \text{ is a}\\\text{column-strict tableau}%
}}\left(  x\mid a\right)  ^{T}=\sum_{\substack{T:\lambda-\mu\rightarrow
\left\{  1,2,3,\ldots\right\}  \text{ is a}\\\text{column-strict tableau}%
}}\left(  x\mid a\right)  ^{T}%
\end{align*}
(since each column-strict tableau $T:\lambda-\mu\rightarrow\left\{
1,2,3,\ldots\right\}  $ has the property that $T:\lambda-\mu\rightarrow\left[
1,n\right]  $ for all sufficiently large $n$). This proves Lemma
\ref{lem.p21.slam/mu=tab}.
\end{proof}

We can easily modify Lemma \ref{lem.p21.slam/mu=tab} to use \textquotedblleft
shifted variables\textquotedblright:

\begin{lemmain}
\label{lem.p21.slam/mu=tab.shift}Let $p\in\mathbb{Z}$. Let $\lambda=\left(
\lambda_{1},\lambda_{2},\ldots,\lambda_{m}\right)  $ and $\mu=\left(  \mu
_{1},\mu_{2},\ldots,\mu_{m}\right)  $ be two partitions. Then,%
\[
\det\left(  h_{\lambda_{i}-\mu_{j}-i+j}\left(  x_{\geq p}\parallel a_{\geq
\mu_{j}-j+p}\right)  \right)  _{i,j\in\left[  m\right]  }=\sum
_{\substack{T:\lambda-\mu\rightarrow\left\{  p,p+1,p+2,\ldots\right\}
\\\text{is a column-strict tableau}}}\left(  x\mid a\right)  ^{T}.
\]

\end{lemmain}

\begin{proof}
[Proof of Lemma \ref{lem.p21.slam/mu=tab.shift}.]In a nutshell: Lemma
\ref{lem.p21.slam/mu=tab.shift} is simply Lemma \ref{lem.p21.slam/mu=tab},
after all the variables $x_{i}$ and $a_{i}$ have been renamed as $x_{i+p-1}$
and $a_{i+p-1}$, respectively.

In more detail: Recall that $\tau^{p-1}a$ is the doubly infinite sequence of
variables $\left(  a_{n+p-1}\right)  _{n\in\mathbb{Z}}$ (that is, the original
sequence $a$ with its indices shifted by $p-1$). Likewise, let $\omega^{p-1}x$
be the doubly infinite sequence of variables $\left(  x_{n+p-1}\right)
_{n\in\mathbb{Z}}$ (that is, the original sequence $x$ with its indices
shifted by $p-1$). Applying Lemma \ref{lem.p21.slam/mu=tab} to $\omega^{p-1}x$
and $\tau^{p-1}a$ instead of $x$ and $a$, we obtain%
\begin{align}
&  \det\left(  h_{\lambda_{i}-\mu_{j}-i+j}\left(  \left(  \omega
^{p-1}x\right)  _{\geq1}\parallel\left(  \tau^{p-1}a\right)  _{\geq\mu
_{j}-j+1}\right)  \right)  _{i,j\in\left[  m\right]  }\nonumber\\
&  =\sum_{\substack{T:\lambda-\mu\rightarrow\left\{  1,2,3,\ldots\right\}
\text{ is a}\\\text{column-strict tableau}}}\left(  \omega^{p-1}x\mid
\tau^{p-1}a\right)  ^{T}, \label{pf.lem.p21.slam/mu=tab.shift.1}%
\end{align}
where we set%
\[
\left(  \omega^{p-1}x\mid\tau^{p-1}a\right)  ^{T}:=\prod_{s\in\lambda-\mu
}\left(  \left(  \omega^{p-1}x\right)  _{T\left(  s\right)  }+\left(
\tau^{p-1}a\right)  _{T^{\ast}\left(  s\right)  }\right)  .
\]


However, it is clear that
\begin{align*}
\left(  \omega^{p-1}x\right)  _{\geq1}  &  =\left(  \left(  \omega
^{p-1}x\right)  _{1},\left(  \omega^{p-1}x\right)  _{2},\left(  \omega
^{p-1}x\right)  _{3},\ldots\right) \\
&  =\left(  x_{p},x_{p+1},x_{p+2},\ldots\right)  =x_{\geq p}%
\end{align*}
and
\begin{align*}
\left(  \tau^{p-1}a\right)  _{\geq\mu_{j}-j+1}  &  =\left(  \left(  \tau
^{p-1}a\right)  _{\mu_{j}-j+1},\left(  \tau^{p-1}a\right)  _{\mu_{j}%
-j+2},\left(  \tau^{p-1}a\right)  _{\mu_{j}-j+3},\ldots\right) \\
&  =\left(  a_{\mu_{j}-j+1+\left(  p-1\right)  },a_{\mu_{j}-j+p+2+\left(
p-1\right)  },a_{\mu_{j}-j+p+3+\left(  p-1\right)  },\ldots\right) \\
&  =\left(  a_{\mu_{j}-j+p},a_{\mu_{j}-j+p+1},a_{\mu_{j}-j+p+2},\ldots\right)
=a_{\geq\mu_{j}-j+p}%
\end{align*}
for each $j\in\left[  m\right]  $. Thus, we can rewrite
(\ref{pf.lem.p21.slam/mu=tab.shift.1}) as%
\begin{align}
&  \det\left(  h_{\lambda_{i}-\mu_{j}-i+j}\left(  x_{\geq p}\parallel
a_{\geq\mu_{j}-j+p}\right)  \right)  _{i,j\in\left[  m\right]  }\nonumber\\
&  =\sum_{\substack{T:\lambda-\mu\rightarrow\left\{  1,2,3,\ldots\right\}
\text{ is a}\\\text{column-strict tableau}}}\left(  \omega^{p-1}x\mid
\tau^{p-1}a\right)  ^{T}. \label{pf.lem.p21.slam/mu=tab.shift.2}%
\end{align}


Furthermore, if $T:\lambda-\mu\rightarrow\left\{  1,2,3,\ldots\right\}  $ is a
column-strict tableau, then we define $\Omega^{p-1}T:\lambda-\mu
\rightarrow\left\{  p,p+1,p+2,\ldots\right\}  $ to be the tableau obtained
from $T$ by adding $p-1$ to each entry. This new tableau $\Omega^{p-1}T$ is
again column-strict, and we have thus found a bijection
\begin{align*}
&  \text{from }\left\{  \text{column-strict tableaux }\lambda-\mu
\rightarrow\left\{  1,2,3,\ldots\right\}  \right\} \\
&  \text{to }\left\{  \text{column-strict tableaux }\lambda-\mu\rightarrow
\left\{  p,p+1,p+2,\ldots\right\}  \right\}
\end{align*}
that sends each $T$ to $\Omega^{p-1}T$. We can use this bijection to
substitute $\Omega^{p-1}T$ for $T$ in the sum $\sum_{\substack{T:\lambda
-\mu\rightarrow\left\{  p,p+1,p+2,\ldots\right\}  \\\text{is a column-strict
tableau}}}\left(  x\mid a\right)  ^{T}$, thus obtaining%
\begin{align}
&  \sum_{\substack{T:\lambda-\mu\rightarrow\left\{  p,p+1,p+2,\ldots\right\}
\\\text{is a column-strict tableau}}}\left(  x\mid a\right)  ^{T}\nonumber\\
&  =\sum_{\substack{T:\lambda-\mu\rightarrow\left\{  1,2,3,\ldots\right\}
\text{ is a}\\\text{column-strict tableau}}}\left(  x\mid a\right)
^{\Omega^{p-1}T}. \label{pf.lem.p21.slam/mu=tab.shift.sum=sum}%
\end{align}


For any column-strict tableau $T:\lambda-\mu\rightarrow\left\{  1,2,3,\ldots
\right\}  $ and any cell $s\in\lambda-\mu$, we have
\begin{equation}
\left(  \Omega^{p-1}T\right)  \left(  s\right)  =T\left(  s\right)  +\left(
p-1\right)  \label{pf.lem.p21.slam/mu=tab.shift.3a}%
\end{equation}
(by the definition of $\Omega^{p-1}T$) and thus also
\begin{equation}
\left(  \Omega^{p-1}T\right)  ^{\ast}\left(  s\right)  =T^{\ast}\left(
s\right)  +\left(  p-1\right)  \label{pf.lem.p21.slam/mu=tab.shift.3b}%
\end{equation}
(because if we write $s$ as $\left(  i,j\right)  $, then $T^{\ast}\left(
s\right)  =T\left(  s\right)  +j-i$ and likewise
\begin{align*}
\left(  \Omega^{p-1}T\right)  ^{\ast}\left(  s\right)   &
=\underbrace{\left(  \Omega^{p-1}T\right)  \left(  s\right)  }%
_{\substack{=T\left(  s\right)  +\left(  p-1\right)  \\\text{(by
(\ref{pf.lem.p21.slam/mu=tab.shift.3a}))}}}+\,j-i=T\left(  s\right)  +\left(
p-1\right)  +j-i\\
&  =\underbrace{T\left(  s\right)  +j-i}_{=T^{\ast}\left(  s\right)  }+\left(
p-1\right)  =T^{\ast}\left(  s\right)  +\left(  p-1\right)
\end{align*}
) and therefore%
\begin{align}
x_{\left(  \Omega^{p-1}T\right)  \left(  s\right)  }  &  =x_{T\left(
s\right)  +\left(  p-1\right)  }\ \ \ \ \ \ \ \ \ \ \left(  \text{by
(\ref{pf.lem.p21.slam/mu=tab.shift.3a})}\right) \nonumber\\
&  =\left(  \omega^{p-1}x\right)  _{T\left(  s\right)  }
\label{pf.lem.p21.slam/mu=tab.shift.3c}%
\end{align}
and%
\begin{align}
a_{\left(  \Omega^{p-1}T\right)  ^{\ast}\left(  s\right)  }  &  =a_{T^{\ast
}\left(  s\right)  +\left(  p-1\right)  }\ \ \ \ \ \ \ \ \ \ \left(  \text{by
(\ref{pf.lem.p21.slam/mu=tab.shift.3b})}\right) \nonumber\\
&  =\left(  \tau^{p-1}a\right)  _{T^{\ast}\left(  s\right)  }.
\label{pf.lem.p21.slam/mu=tab.shift.3d}%
\end{align}
Hence, for any column-strict tableau $T:\lambda-\mu\rightarrow\left\{
1,2,3,\ldots\right\}  $, we have
\begin{align}
\left(  x\mid a\right)  ^{\Omega^{p-1}T}  &  =\prod_{s\in\lambda-\mu}\left(
\underbrace{x_{\left(  \Omega^{p-1}T\right)  \left(  s\right)  }%
}_{\substack{=\left(  \omega^{p-1}x\right)  _{T\left(  s\right)  }\\\text{(by
(\ref{pf.lem.p21.slam/mu=tab.shift.3c}))}}}+\underbrace{a_{\left(
\Omega^{p-1}T\right)  ^{\ast}\left(  s\right)  }}_{\substack{=\left(
\tau^{p-1}a\right)  _{T^{\ast}\left(  s\right)  }\\\text{(by
(\ref{pf.lem.p21.slam/mu=tab.shift.3d}))}}}\right) \nonumber\\
&  \ \ \ \ \ \ \ \ \ \ \ \ \ \ \ \ \ \ \ \ \left(  \text{by the definition of
}\left(  x\mid a\right)  ^{\Omega^{p-1}T}\right) \nonumber\\
&  =\prod_{s\in\lambda-\mu}\left(  \left(  \omega^{p-1}x\right)  _{T\left(
s\right)  }+\left(  \tau^{p-1}a\right)  _{T^{\ast}\left(  s\right)  }\right)
\nonumber\\
&  =\left(  \omega^{p-1}x\mid\tau^{p-1}a\right)  ^{T}
\label{pf.lem.p21.slam/mu=tab.shift.4}%
\end{align}
(since $\left(  \omega^{p-1}x\mid\tau^{p-1}a\right)  ^{T}$ is defined to be
$\prod_{s\in\lambda-\mu}\left(  \left(  \omega^{p-1}x\right)  _{T\left(
s\right)  }+\left(  \tau^{p-1}a\right)  _{T^{\ast}\left(  s\right)  }\right)
$). Using this equality, we can rewrite
(\ref{pf.lem.p21.slam/mu=tab.shift.sum=sum}) as%
\begin{align*}
&  \sum_{\substack{T:\lambda-\mu\rightarrow\left\{  p,p+1,p+2,\ldots\right\}
\\\text{is a column-strict tableau}}}\left(  x\mid a\right)  ^{T}\\
&  =\sum_{\substack{T:\lambda-\mu\rightarrow\left\{  1,2,3,\ldots\right\}
\text{ is a}\\\text{column-strict tableau}}}\left(  \omega^{p-1}x\mid
\tau^{p-1}a\right)  ^{T}\\
&  =\det\left(  h_{\lambda_{i}-\mu_{j}-i+j}\left(  x_{\geq p}\parallel
a_{\geq\mu_{j}-j+p}\right)  \right)  _{i,j\in\left[  m\right]  }%
\ \ \ \ \ \ \ \ \ \ \left(  \text{by (\ref{pf.lem.p21.slam/mu=tab.shift.2}%
)}\right)  .
\end{align*}
This proves Lemma \ref{lem.p21.slam/mu=tab.shift}.
\end{proof}

\begin{lemmain}
\label{lem.p21.lim-h-down}Let $r$ and $m$ be two integers. Then,%
\[
\lim\limits_{p\rightarrow-\infty}h_{r}\left(  x_{\geq p}\parallel a_{\geq
m+p}\right)  =h_{r}\left(  x\parallel a\right)  .
\]

\end{lemmain}

\begin{proof}
[Proof of Lemma \ref{lem.p21.lim-h-down}.]For any $p\in\mathbb{Z}$, the power
series $h_{r}\left(  x_{\geq p}\parallel a_{\geq m+p}\right)  $ is defined as
the coefficient of $t^{r}$ in the power series
\begin{equation}
\prod_{i\geq p}\left(  1-tx_{i}\right)  ^{-1}\prod_{j\geq m+p}\left(
1+ta_{j}\right)  . \label{pf.lem.p21.lim-h-down.1}%
\end{equation}
As $p$ goes to $-\infty$, the products $\prod_{i\geq p}\left(  1-tx_{i}%
\right)  ^{-1}$ and $\prod_{j\geq m+p}\left(  1+ta_{j}\right)  $ approach
$\prod_{i\in\mathbb{Z}}\left(  1-tx_{i}\right)  ^{-1}$ and $\prod
_{j\in\mathbb{Z}}\left(  1+ta_{j}\right)  $, respectively, and thus the power
series (\ref{pf.lem.p21.lim-h-down.1}) approaches $\prod_{i\in\mathbb{Z}%
}\left(  1-tx_{i}\right)  ^{-1}\prod_{j\in\mathbb{Z}}\left(  1+ta_{j}\right)
$. Hence, the coefficient of $t^{r}$ in the former power series approaches the
coefficient of $t^{r}$ in the latter. In other words, $h_{r}\left(  x_{\geq
p}\parallel a_{\geq m+p}\right)  $ approaches $h_{r}\left(  x\parallel
a\right)  $ (since $h_{r}\left(  x_{\geq p}\parallel a_{\geq m+p}\right)  $ is
the coefficient of $t^{r}$ in the former power series, whereas $h_{r}\left(
x\parallel a\right)  $ is the coefficient of $t^{r}$ in the latter). This
proves Lemma \ref{lem.p21.lim-h-down}.
\end{proof}

Finally, we can prove (6.21):

\begin{proof}
[Proof of (6.21).]Write the partitions $\lambda$ and $\mu$ as $\lambda=\left(
\lambda_{1},\lambda_{2},\ldots,\lambda_{m}\right)  $ and $\mu=\left(  \mu
_{1},\mu_{2},\ldots,\mu_{m}\right)  $. Then, Lemma
\ref{lem.p21.slam/mu=tab.shift} yields%
\begin{align}
&  \det\left(  h_{\lambda_{i}-\mu_{j}-i+j}\left(  x_{\geq p}\parallel
a_{\geq\mu_{j}-j+p}\right)  \right)  _{i,j\in\left[  m\right]  }\nonumber\\
&  =\sum_{\substack{T:\lambda-\mu\rightarrow\left\{  p,p+1,p+2,\ldots\right\}
\\\text{is a column-strict tableau}}}\left(  x\mid a\right)  ^{T}
\label{p22.6.21.pf.1}%
\end{align}
for each $p\in\mathbb{Z}$.

Now, let us take the limit $p\rightarrow-\infty$ in this equality. As $p$ goes
to $-\infty$, each of the power series $h_{\lambda_{i}-\mu_{j}-i+j}\left(
x_{\geq p}\parallel a_{\geq\mu_{j}-j+p}\right)  $ (for any given
$i,j\in\left[  m\right]  $) approaches $h_{\lambda_{i}-\mu_{j}-i+j}\left(
x\parallel a\right)  $, because Lemma \ref{lem.p21.lim-h-down} (applied to
$m=\mu_{j}-j$ and $r=\lambda_{i}-\mu_{j}-i+j$) yields $\lim
\limits_{p\rightarrow-\infty}h_{\lambda_{i}-\mu_{j}-i+j}\left(  x_{\geq
p}\parallel a_{\geq\mu_{j}-j+p}\right)  =h_{\lambda_{i}-\mu_{j}-i+j}\left(
x\parallel a\right)  $. Meanwhile, the $\sum_{\substack{T:\lambda
-\mu\rightarrow\left\{  p,p+1,p+2,\ldots\right\}  \\\text{is a column-strict
tableau}}}\left(  x\mid a\right)  ^{T}$ approaches $\sum_{\substack{T:\lambda
-\mu\rightarrow\mathbb{Z}\\\text{is a column-strict tableau}}}\left(  x\mid
a\right)  ^{T}$ (as $p$ goes to $-\infty$), since each column-strict tableau
$T:\lambda-\mu\rightarrow\mathbb{Z}$ is of the form $T:\lambda-\mu
\rightarrow\left\{  p,p+1,p+2,\ldots\right\}  $ for all sufficiently small
integers $p$. Hence, taking the limit $p\rightarrow-\infty$ in the equality
(\ref{p22.6.21.pf.1}), we obtain%
\[
\det\left(  h_{\lambda_{i}-\mu_{j}-i+j}\left(  x\parallel a\right)  \right)
_{i,j\in\left[  m\right]  }=\sum_{\substack{T:\lambda-\mu\rightarrow
\mathbb{Z}\\\text{is a column-strict tableau}}}\left(  x\mid a\right)  ^{T}.
\]
Since $s_{\lambda/\mu}\left(  x\parallel a\right)  $ is defined as
$\det\left(  h_{\lambda_{i}-\mu_{j}-i+j}\left(  x\parallel a\right)  \right)
_{i,j\in\left[  m\right]  }$, we can rewrite this as%
\[
s_{\lambda/\mu}\left(  x\parallel a\right)  =\sum_{\substack{T:\lambda
-\mu\rightarrow\mathbb{Z}\\\text{is a column-strict tableau}}}\left(  x\mid
a\right)  ^{T}.
\]
Thus, (6.21) is proved.
\end{proof}
\end{enumerate}

\subsection{Corrections to the 7th Variation}

\begin{enumerate}
\item \textbf{page 23, (7.4):} It is worth adding that $A_{\delta}%
=P=\prod_{v\in V_{0}}v$ is known as the \emph{Moore determinant} of
$x_{1},x_{2},\ldots,x_{n}$, and is denoted by $\Delta\left(  x_{1}%
,x_{2},\ldots,x_{n}\right)  $ in \cite[Definition 1.3.2]{Goss98}. The equality
$A_{\delta}=P$ itself is \cite[Corollary 1.3.7]{Goss98}, and goes back to
Moore \cite[\S III]{Moore96}. This determinant is fundamental to the theory of
Carlitz modules \cite[Chapter 3]{Goss98}.

\item \textbf{page 24, definition of }$H_{r}\left(  V\right)  $\textbf{:} It
should also be said that $H_{r}\left(  V\right)  $ is defined to be $0$ when
$n=0$ and $r>0$. (In this case, $S_{\left(  r\right)  }\left(  V\right)  $ is
undefined, since $\left(  r\right)  $ is not a partition of length $\leq n$,
although it would make perfect sense to just define $S_{\lambda}\left(
V\right)  $ to be $0$ when $\lambda$ is a partition of length $>n$.)

\item \textbf{page 24:} \textquotedblleft A well-known theorem of Dickson
states that the subalgebra of $GL\left(  V\right)  $-invariant elements of
$\mathbf{S}\left(  V\right)  $ is a polynomial algebra over $F$, generated by
the $E_{r}\left(  V\right)  $ ($1\leq r\leq n$)\textquotedblright: Let me add
that this theorem is nicely proved in \cite[\S I.2, Example 27 parts (a) and
(d)]{Macdon95}.

\item \textbf{page 24:} In \textquotedblleft form an $F$-basis of $S\left(
V\right)  ^{GL\left(  V\right)  }$\textquotedblright, the \textquotedblleft%
$S$\textquotedblright\ should be a boldface \textquotedblleft$\mathbf{S}%
$\textquotedblright.

\item \textbf{page 25, proof of (7.9):} It is worth pointing out that the
expressions \textquotedblleft$S_{\lambda}\left(  V\right)  $\textquotedblright%
, \textquotedblleft$H_{r}\left(  V\right)  $\textquotedblright\ and
\textquotedblleft$E_{r}\left(  V\right)  $\textquotedblright\ are abbreviated
as \textquotedblleft$S_{\lambda}$\textquotedblright, \textquotedblleft$H_{r}%
$\textquotedblright\ and \textquotedblleft$E_{r}$\textquotedblright\ when the
argument $V$ is clear from the context.

\item \textbf{page 25, proof of (7.9):} Just as in the proof of (6.6), this
computation only shows that $\mathbf{H}\left(  V\right)  \cdot\mathbf{E}%
\left(  V\right)  =I_{\mathbb{Z}}$ (where $I_{\mathbb{Z}}$ is the identity
matrix with rows and columns indexed by all integers). In order to prove that
the two matrices $\mathbf{H}\left(  V\right)  $ and $\mathbf{E}\left(
V\right)  $ are inverse to each other (i.e., in order to prove (7.9)), it must
also be shown that $\mathbf{E}\left(  V\right)  \cdot\mathbf{H}\left(
V\right)  =I_{\mathbb{Z}}$.

Again, this can be done by observing that both matrices belong to the group
$\operatorname*{UT}\nolimits_{\mathbb{Z}}$ of all upper unitriangular matrices
with rows and columns indexed by all integers.

\item \textbf{page 25, proof of (7.9):} Starting with \textquotedblleft%
$f_{V}\left(  x_{i}\right)  =0$\textquotedblright\ and ending with the
equality (2), the \textquotedblleft$i$\textquotedblright\ stands for an
arbitrary element of $\left[  n\right]  $ rather than the $i$ from the
beginning of this proof.

\item \textbf{page 25, proof of (7.9):} It is crucial to point out that (2)
holds for all integers $r>-n$. (Merely stating it for $r\geq0$ is
insufficient, since it is later used for all $r>-n$. On the other hand, it
usually does not hold for $r\leq-n$.)

Let me explain why the proof of (2) does work for all $r>-n$. Indeed, this
proof relies on the fact that $H_{r}=A_{\left(  r\right)  +\delta}/A_{\delta}%
$. For $r\geq0$, this is clear from the definition of $H_{r}$ (because this
definition says that $H_{r}=S_{\left(  r\right)  }=A_{\left(  r\right)
+\delta}/A_{\delta}$). It remains to prove this fact in the case when
$0>r>-n$. In this case, the $n$-tuple $\left(  r\right)  +\delta=\left(
r+n-1,\ n-2,\ n-3,\ \ldots,\ 0\right)  $ has two equal entries (namely, its
first entry equals its $\left(  1-r\right)  $-th entry), and thus the matrix
$\left(  x_{i}^{q^{\left(  \left(  r\right)  +\delta\right)  _{j}}}\right)
_{1\leq i,j\leq n}$ has two equal columns. Therefore, the determinant of this
matrix is $0$. Since $A_{\left(  r\right)  +\delta}$ is defined to be this
determinant, we thus have shown that $A_{\left(  r\right)  +\delta}=0$.
Comparing this with $H_{r}=0$ (which is immediate from the definition of
$H_{r}$, since $r<0$), we obtain $H_{r}=A_{\left(  r\right)  +\delta
}/A_{\delta}$.

\item \textbf{page 26, proof of (7.9):} After (3), replace \textquotedblleft
for each $r\geq0$\textquotedblright\ by \textquotedblleft for each
$r>0$\textquotedblright.

\item \textbf{page 26, proof of (7.9):} Let me give a few details about how
\begin{equation}
\sum_{i\leq j\leq k}\left(  -1\right)  ^{k-j}\varphi^{k}\left(  E_{k-j}%
\right)  \varphi^{i+1}\left(  H_{j-i}\right)  =0 \label{p26.7.9.a}%
\end{equation}
is proved. Indeed, we rewrite the equality (3) as
\[
\sum_{p=0}^{n}\left(  -1\right)  ^{p}\varphi^{r-1}\left(  E_{p}\right)
H_{r-p}=0.
\]
Applying this to $r=k-i$ (which is allowed since $i<k$ and thus $k-i>0$), we
obtain%
\[
\sum_{p=0}^{n}\left(  -1\right)  ^{p}\varphi^{k-i-1}\left(  E_{p}\right)
H_{k-i-p}=0.
\]
Applying the algebra morphism $\varphi^{i+1}$ to both sides of this equality,
we obtain%
\[
\sum_{p=0}^{n}\left(  -1\right)  ^{p}\varphi^{k}\left(  E_{p}\right)
\varphi^{i+1}\left(  H_{k-i-p}\right)  =0.
\]
Substituting $k-j$ for $p$ in the sum, we can rewrite this as%
\[
\sum_{j=k-n}^{k}\left(  -1\right)  ^{k-j}\varphi^{k}\left(  E_{k-j}\right)
\varphi^{i+1}\left(  H_{k-i-\left(  k-j\right)  }\right)  =0.
\]
Since $k-i-\left(  k-j\right)  =j-i$, we can rewrite this as%
\[
\sum_{j=k-n}^{k}\left(  -1\right)  ^{k-j}\varphi^{k}\left(  E_{k-j}\right)
\varphi^{i+1}\left(  H_{j-i}\right)  =0.
\]
The lower limit $k-n$ of the sum can be replaced by $i$, because the addends
vanish whenever $j$ is smaller than \textbf{either} of $k-n$ and $i$ (indeed,
if $j<k-n$, then $k-j>n$ and thus $E_{k-j}=0$; on the other hand, if $j<i$,
then $j-i<0$ and thus $H_{j-i}=0$). Hence, the above equality takes the form%
\[
\sum_{j=i}^{k}\left(  -1\right)  ^{k-j}\varphi^{k}\left(  E_{k-j}\right)
\varphi^{i+1}\left(  H_{j-i}\right)  =0,
\]
which is precisely (\ref{p26.7.9.a}).

\item \textbf{page 26, (7.10):} To clarify:

\begin{itemize}
\item The expression \textquotedblleft$\det\left(  \varphi^{1-j}H_{\lambda
_{i}-i+j}\left(  V\right)  \right)  $\textquotedblright\ is understood to mean
the determinant%
\[
\det\left(  \varphi^{1-j}H_{\lambda_{i}-i+j}\left(  V\right)  \right)
_{i,j\in\left[  m\right]  },
\]
where $m$ is any integer such that $\ell\left(  \lambda\right)  \leq m$. (This
determinant does not depend on the choice of $m$, since one can easily see
that if $\lambda_{m}=0$, then the last row of the matrix $\left(
\varphi^{1-j}H_{\lambda_{i}-i+j}\left(  V\right)  \right)  _{i,j\in\left[
m\right]  }$ is $\left(  0,0,\ldots,0,1\right)  $, and therefore Laplace
expansion along this row reduces the determinant to the determinant of the
submatrix $\left(  \varphi^{1-j}H_{\lambda_{i}-i+j}\left(  V\right)  \right)
_{i,j\in\left[  m-1\right]  }$.)

\item The expression \textquotedblleft$\det\left(  \varphi^{j-1}E_{\lambda
_{i}^{\prime}-i+j}\left(  V\right)  \right)  $\textquotedblright\ is
understood to mean the determinant%
\[
\det\left(  \varphi^{j-1}E_{\lambda_{i}^{\prime}-i+j}\left(  V\right)
\right)  _{i,j\in\left[  p\right]  },
\]
where $p$ is any integer such that $\ell\left(  \lambda^{\prime}\right)  \leq
p$. (Again, this determinant does not depend on the choice of $p$, for a
reason similar to the one we just gave.)
\end{itemize}

\item \textbf{page 26, proof of (7.10):} Another clarification: All the
matrices appearing in this proof (more precisely, in the proof of the first
formula of (7.10)) are understood to be $n\times n$-matrices.

\item \textbf{page 26, proof of (7.10):} After \textquotedblleft the left-hand
side of (1) becomes equal to $1$\textquotedblright, I would add
\textquotedblleft(because the matrix $\left(  \varphi^{1-j}H_{\delta_{i}%
-n+j}\right)  _{i,j}=\left(  \varphi^{1-j}H_{j-i}\right)  _{i,j}$ is
unitriangular)\textquotedblright.

\item \textbf{page 26, proof of (7.10):} Let me explain how the second formula
in (7.10) can be obtained from the first.

Actually, I will prove the more general identity (7.11').

\begin{proof}
[Proof of (7.11').]Write the partitions $\lambda$, $\lambda^{\prime}$, $\mu$
and $\mu^{\prime}$ in the forms%
\begin{align*}
\lambda &  =\left(  \lambda_{1},\lambda_{2},\ldots,\lambda_{q}\right)
,\ \ \ \ \ \ \ \ \ \ \lambda^{\prime}=\left(  \lambda_{1}^{\prime},\lambda
_{2}^{\prime},\ldots,\lambda_{p}^{\prime}\right)  ,\\
\mu &  =\left(  \mu_{1},\mu_{2},\ldots,\mu_{q}\right)
,\ \ \ \ \ \ \ \ \ \ \mu^{\prime}=\left(  \mu_{1}^{\prime},\mu_{2}^{\prime
},\ldots,\mu_{p}^{\prime}\right)
\end{align*}
for some $p,q\in\mathbb{N}$ (where, of course, the entries $\lambda_{i}$,
$\lambda_{j}^{\prime}$, $\mu_{i}$ and $\mu_{j}^{\prime}$ are allowed to be $0$).

Recall that%
\begin{align*}
\mathbf{H}\left(  V\right)   &  =\left(  \varphi^{i+1}H_{j-i}\left(  V\right)
\right)  _{i,j\in\mathbb{Z}}\ \ \ \ \ \ \ \ \ \ \text{and}\\
\mathbf{E}\left(  V\right)   &  =\left(  \left(  -1\right)  ^{j-i}\varphi
^{j}E_{j-i}\left(  V\right)  \right)  _{i,j\in\mathbb{Z}}%
\end{align*}
are two upper unitriangular matrices in $\operatorname*{UT}%
\nolimits_{\mathbb{Z}}$ (since $H_{0}\left(  V\right)  =E_{0}\left(  V\right)
=1$ and $H_{k}\left(  V\right)  =E_{k}\left(  V\right)  =0$ whenever $k<0$).
From (7.9), we know that the matrix $\mathbf{E}\left(  V\right)  $ is the
inverse matrix of $\mathbf{H}\left(  V\right)  $. Hence, Lemma \ref{lem.A5}
(applied to $B=\mathbf{H}\left(  V\right)  $ and $b_{i,j}=\varphi^{i+1}%
H_{j-i}\left(  V\right)  $ and $c_{i,j}=\left(  -1\right)  ^{j-i}\varphi
^{j}E_{j-i}\left(  V\right)  $) yields%
\begin{align*}
&  \det\left(  \varphi^{\mu_{i}-i+1}H_{\left(  \lambda_{j}-j\right)  -\left(
\mu_{i}-i\right)  }\left(  V\right)  \right)  _{i,j\in\left[  q\right]  }\\
&  =\left(  -1\right)  ^{\left(  \lambda_{1}+\lambda_{2}+\cdots+\lambda
_{q}\right)  +\left(  \mu_{1}+\mu_{2}+\cdots+\mu_{q}\right)  }\\
&  \ \ \ \ \ \ \ \ \ \ \det\left(  \left(  -1\right)  ^{\left(  j-\mu
_{j}^{\prime}-1\right)  -\left(  i-\lambda_{i}^{\prime}-1\right)  }%
\varphi^{j-\mu_{j}^{\prime}-1}E_{\left(  j-\mu_{j}^{\prime}-1\right)  -\left(
i-\lambda_{i}^{\prime}-1\right)  }\left(  V\right)  \right)  _{i,j\in\left[
p\right]  }.
\end{align*}
In view of $\left(  \lambda_{j}-j\right)  -\left(  \mu_{i}-i\right)
=\lambda_{j}-\mu_{i}-j+i$ and $\left(  j-\mu_{j}^{\prime}-1\right)  -\left(
i-\lambda_{i}^{\prime}-1\right)  =\lambda_{i}^{\prime}-\mu_{j}^{\prime}-i+j$,
we can rewrite this as%
\begin{align*}
&  \det\left(  \varphi^{\mu_{i}-i+1}H_{\lambda_{j}-\mu_{i}-j+i}\left(
V\right)  \right)  _{i,j\in\left[  q\right]  }\\
&  =\left(  -1\right)  ^{\left(  \lambda_{1}+\lambda_{2}+\cdots+\lambda
_{q}\right)  +\left(  \mu_{1}+\mu_{2}+\cdots+\mu_{q}\right)  }\\
&  \ \ \ \ \ \ \ \ \ \ \det\left(  \left(  -1\right)  ^{\lambda_{i}^{\prime
}-\mu_{j}^{\prime}-i+j}\varphi^{j-\mu_{j}^{\prime}-1}E_{\lambda_{i}^{\prime
}-\mu_{j}^{\prime}-i+j}\left(  V\right)  \right)  _{i,j\in\left[  p\right]  }.
\end{align*}
In view of%
\[
\det\left(  \varphi^{\mu_{i}-i+1}H_{\lambda_{j}-\mu_{i}-j+i}\left(  V\right)
\right)  _{i,j\in\left[  q\right]  }=\det\left(  \varphi^{\mu_{j}%
-j+1}H_{\lambda_{i}-\mu_{j}-i+j}\left(  V\right)  \right)  _{i,j\in\left[
q\right]  }%
\]
(since the determinant of a matrix does not change when we transpose it) and%
\begin{align*}
&  \det\left(  \underbrace{\left(  -1\right)  ^{\lambda_{i}^{\prime}-\mu
_{j}^{\prime}-i+j}}_{=\left(  -1\right)  ^{\lambda_{i}^{\prime}-i}\left(
-1\right)  ^{\mu_{j}^{\prime}-j}}\varphi^{j-\mu_{j}^{\prime}-1}E_{\lambda
_{i}^{\prime}-\mu_{j}^{\prime}-i+j}\left(  V\right)  \right)  _{i,j\in\left[
p\right]  }\\
&  =\det\left(  \left(  -1\right)  ^{\lambda_{i}^{\prime}-i}\left(  -1\right)
^{\mu_{j}^{\prime}-j}\varphi^{j-\mu_{j}^{\prime}-1}E_{\lambda_{i}^{\prime}%
-\mu_{j}^{\prime}-i+j}\left(  V\right)  \right)  _{i,j\in\left[  p\right]  }\\
&  =\left(  \prod_{i=1}^{p}\left(  -1\right)  ^{\lambda_{i}^{\prime}%
-i}\right)  \left(  \prod_{j=1}^{p}\left(  -1\right)  ^{\mu_{j}^{\prime}%
-j}\right)  \cdot\det\left(  \varphi^{j-\mu_{j}^{\prime}-1}E_{\lambda
_{i}^{\prime}-\mu_{j}^{\prime}-i+j}\left(  V\right)  \right)  _{i,j\in\left[
p\right]  }%
\end{align*}
(by Lemma \ref{lem.A6}, applied to $\alpha_{i}=\left(  -1\right)
^{\lambda_{i}^{\prime}-i}$ and $\beta_{j}=\left(  -1\right)  ^{\mu_{j}%
^{\prime}-j}$ and $u_{i,j}=\varphi^{j-\mu_{j}^{\prime}-1}E_{\lambda
_{i}^{\prime}-\mu_{j}^{\prime}-i+j}\left(  V\right)  $), we can rewrite this
as%
\begin{align*}
&  \det\left(  \varphi^{\mu_{j}-j+1}H_{\lambda_{i}-\mu_{j}-i+j}\left(
V\right)  \right)  _{i,j\in\left[  q\right]  }\\
&  =\left(  -1\right)  ^{\left(  \lambda_{1}+\lambda_{2}+\cdots+\lambda
_{q}\right)  +\left(  \mu_{1}+\mu_{2}+\cdots+\mu_{q}\right)  }\\
&  \ \ \ \ \ \ \ \ \ \ \left(  \prod_{i=1}^{p}\left(  -1\right)  ^{\lambda
_{i}^{\prime}-i}\right)  \left(  \prod_{j=1}^{p}\left(  -1\right)  ^{\mu
_{j}^{\prime}-j}\right)  \cdot\det\left(  \varphi^{j-\mu_{j}^{\prime}%
-1}E_{\lambda_{i}^{\prime}-\mu_{j}^{\prime}-i+j}\left(  V\right)  \right)
_{i,j\in\left[  p\right]  }\\
&  =\underbrace{\left(  -1\right)  ^{\left(  \lambda_{1}+\lambda_{2}%
+\cdots+\lambda_{q}\right)  +\left(  \mu_{1}+\mu_{2}+\cdots+\mu_{q}\right)
}\left(  \prod_{i=1}^{p}\left(  -1\right)  ^{\lambda_{i}^{\prime}-i}\right)
\left(  \prod_{j=1}^{p}\left(  -1\right)  ^{\mu_{j}^{\prime}-j}\right)
}_{\substack{=1\\\text{(by Lemma \ref{lem.A6prods})}}}\\
&  \ \ \ \ \ \ \ \ \ \ \cdot\det\left(  \varphi^{j-\mu_{j}^{\prime}%
-1}E_{\lambda_{i}^{\prime}-\mu_{j}^{\prime}-i+j}\left(  V\right)  \right)
_{i,j\in\left[  p\right]  }\\
&  =\det\left(  \varphi^{j-\mu_{j}^{\prime}-1}E_{\lambda_{i}^{\prime}-\mu
_{j}^{\prime}-i+j}\left(  V\right)  \right)  _{i,j\in\left[  p\right]  }\\
&  =\det\left(  \varphi^{-\mu_{j}^{\prime}+j-1}E_{\lambda_{i}^{\prime}-\mu
_{j}^{\prime}-i+j}\left(  V\right)  \right)  _{i,j\in\left[  p\right]
}\ \ \ \ \ \ \ \ \ \ \left(  \text{since }j-\mu_{j}^{\prime}-1=-\mu
_{j}^{\prime}+j-1\right)  .
\end{align*}
In view of (7.11), this can be rewritten as%
\[
S_{\lambda/\mu}\left(  V\right)  =\det\left(  \varphi^{-\mu_{j}^{\prime}%
+j-1}E_{\lambda_{i}^{\prime}-\mu_{j}^{\prime}-i+j}\left(  V\right)  \right)
_{i,j\in\left[  p\right]  }.
\]
This proves (7.11').
\end{proof}

Applying (7.11') to $\mu=\varnothing$, we obtain the second equality in (7.10).

\item \textbf{page 26, (7.11):} The matrix on the right hand side of (7.11) is
understood to be an $m\times m$-matrix (with $i$ and $j$ ranging over $\left[
m\right]  $), where $m$ is any integer such that $\ell\left(  \lambda\right)
\leq m$ and $\ell\left(  \mu\right)  \leq m$. (This determinant of this matrix
does not depend on the choice of $m$, since one can easily see that if
$\lambda_{m}=0$ and $\mu_{m}=0$, then the last row of the matrix $\left(
\varphi^{\mu_{j}-j+1}H_{\lambda_{i}-\mu_{j}-i+j}\left(  V\right)  \right)
_{i,j\in\left[  m\right]  }$ is $\left(  0,0,\ldots,0,1\right)  $, and
therefore Laplace expansion along this row reduces the determinant to the
determinant of the submatrix $\left(  \varphi^{\mu_{j}-j+1}H_{\lambda_{i}%
-\mu_{j}-i+j}\left(  V\right)  \right)  _{i,j\in\left[  m-1\right]  }$.)

\item \textbf{page 26, (7.12):} The proof of (7.12) is similar to the above
proof of (6.10), except that it uses (7.11') instead of (6.9). (Of course, it
crucially relies on the fact that $E_{\ell}\left(  V\right)  =0$ whenever
$\ell<0$ or $\ell>n$.)

\item \textbf{page 27, Remark:} It should be said that \textquotedblleft
degree\textquotedblright\ means \textquotedblleft homogeneous
degree\textquotedblright\ here (i.e., the polynomials $H_{r}\left(  V\right)
$ and $\varphi^{\mu_{j}-j+1}H_{\lambda_{i}-\mu_{j}-i+j}\left(  V\right)  $ are
homogeneous of the respective degrees provided in this remark), and that the
zero polynomial is considered to be homogeneous of every degree (otherwise,
the claim that \textquotedblleft$H_{r}\left(  V\right)  $ has degree $\left(
q^{r}-1\right)  q^{n-1}$\textquotedblright\ would make no sense when $r$ is negative).

\item \textbf{page 27, Remark:} \textquotedblleft(and likewise
(7.12))\textquotedblright\ should be \textquotedblleft(and likewise
(7.11'))\textquotedblright.

\item \textbf{page 27, Remark:} The word \textquotedblleft
isobaric\textquotedblright\ is rather obscure, so let me recall its definition
here: A determinant $\det\left(  a_{i,j}\right)  _{i,j\in\left[  m\right]  }$
(or, better, the underlying matrix $\left(  a_{i,j}\right)  _{i,j\in\left[
m\right]  }$) over some graded ring $R$ is said to be \emph{isobaric} if each
entry $a_{i,j}$ is homogeneous of degree $d_{i}-e_{j}$, for certain integers
$d_{1},d_{2},\ldots,d_{m},e_{1},e_{2},\ldots,e_{m}$. In this case, each of the
products $a_{1,\sigma\left(  1\right)  }a_{2,\sigma\left(  2\right)  }\cdots
a_{m,\sigma\left(  m\right)  }$ in the Leibniz formula $\det\left(
a_{i,j}\right)  _{i,j\in\left[  m\right]  }=\sum_{\sigma\in S_{m}}\left(
\operatorname*{sign}\sigma\right)  a_{1,\sigma\left(  1\right)  }%
a_{2,\sigma\left(  2\right)  }\cdots a_{m,\sigma\left(  m\right)  }$ is
homogeneous of degree%
\begin{align*}
&  \left(  d_{1}-e_{\sigma\left(  1\right)  }\right)  +\left(  d_{2}%
-e_{\sigma\left(  2\right)  }\right)  +\cdots+\left(  d_{m}-e_{\sigma\left(
m\right)  }\right) \\
&  =\left(  d_{1}+d_{2}+\cdots+d_{m}\right)  -\underbrace{\left(
e_{\sigma\left(  1\right)  }+e_{\sigma\left(  2\right)  }+\cdots
+e_{\sigma\left(  m\right)  }\right)  }_{=e_{1}+e_{2}+\cdots+e_{m}}\\
&  =\left(  d_{1}+d_{2}+\cdots+d_{m}\right)  -\left(  e_{1}+e_{2}+\cdots
+e_{m}\right)  ,
\end{align*}
and therefore the whole determinant $\det\left(  a_{i,j}\right)
_{i,j\in\left[  m\right]  }$ is also homogeneous of this degree.

\item \textbf{page 27:} In the \textquotedblleft addition
formula\textquotedblright%
\[
s_{\lambda/\mu}\left(  x,y\right)  =\sum_{\nu}s_{\lambda/\mu}\left(  x\right)
s_{\nu/\mu}\left(  y\right)  ,
\]
the \textquotedblleft$s_{\lambda/\mu}\left(  x\right)  $\textquotedblright%
\ should be \textquotedblleft$s_{\lambda/\nu}\left(  x\right)  $%
\textquotedblright.

\item \textbf{page 27, definition of }$f_{U}$\textbf{, and onwards:} The
discussion of $f_{U}$, $f_{V/U}$ and other things related to subspaces $U$
suffers from an insufficiently general definition. Namely, all the main
objects of the \textquotedblleft7th Variation\textquotedblright\ section, such
as $S_{\lambda}\left(  V\right)  $ and $f_{V}$, have been defined only for the
case when $V$ is the span of the indeterminates $x_{1},x_{2},\ldots,x_{n}$ in
the polynomial ring $F\left[  x_{1},x_{2},\ldots,x_{n}\right]  $. In order to
make sense of things such as $f_{V/U}$ (and to make better sense of $f_{U}$ as
well), a more flexible definition is needed. Taking things literally,
$f_{V/U}$ would be a polynomial in $t$ over $\mathbf{S}\left(  V/U\right)  $,
but this is not what is intended; rather, a polynomial over $\mathbf{S}\left(
V\right)  $ is desired.

Let me thus introduce the definitions that (I believe) are needed to support
the claims in this section.

Recall that $F=\mathbb{F}_{q}$. Thus, each commutative $F$-algebra
$\mathbf{A}$ has an $F$-algebra endomorphism $\varphi:\mathbf{A}%
\rightarrow\mathbf{A},\ a\mapsto a^{q}$, which is known as the \emph{Frobenius
endomorphism}. We say that $\mathbf{A}$ is $\varphi$\emph{-bijective} if this
endomorphism $\varphi$ is bijective (i.e., an isomorphism). For instance,
$\widehat{\mathbf{S}}\left(  V\right)  $ is $\varphi$-bijective. Note that
every morphism $f:\mathbf{A}\rightarrow\mathbf{B}$ of commutative $F$-algebras
automatically respects the Frobenius endomorphisms (i.e., satisfies
$\varphi\circ f=f\circ\varphi$), since it clearly respects taking the $q$-th power.

If $V$ is any finite-dimensional vector space, then the element $S_{\lambda
}\left(  V\right)  \in\mathbf{S}\left(  V\right)  $ is defined for each
partition $\lambda$ of length $\leq n$ (since we can fix a basis $\left(
x_{1},x_{2},\ldots,x_{n}\right)  $ of $V$ and regard $\mathbf{S}\left(
V\right)  $ as the polynomial ring $F\left[  x_{1},x_{2},\ldots,x_{n}\right]
$, so that $S_{\lambda}\left(  V\right)  $ is defined as on page 24; note also
that the value of $S_{\lambda}\left(  V\right)  $ does not depend on this
basis $\left(  x_{1},x_{2},\ldots,x_{n}\right)  $, as was explained on page
24). Furthermore, if $V$ is any finite-dimensional vector space, then the
polynomial $f_{V}\left(  t\right)  :=\prod_{v\in V}\left(  t+v\right)
\in\left(  \mathbf{S}\left(  V\right)  \right)  \left[  t\right]  $ is
defined. So far, none of this is new.

But let us now extend these notions: Let $\mathbf{A}$ be a commutative
$F$-algebra, and let $V$ be a finite-dimensional $F$-vector subspace of
$\mathbf{A}$. Then, by the universal property of the symmetric algebra
$\mathbf{S}\left(  V\right)  $, there is a canonical $F$-algebra morphism
$\Phi_{V}^{\mathbf{A}}:\mathbf{S}\left(  V\right)  \rightarrow\mathbf{A}$ that
sends each $v\in V$ (viewed as a degree-$1$ element of $\mathbf{S}\left(
V\right)  $) to $v$ (but viewed as an element of $\mathbf{A}$). If we view
$\mathbf{S}\left(  V\right)  $ as a polynomial ring $F\left[  x_{1}%
,x_{2},\ldots,x_{n}\right]  $, then this is simply an evaluation morphism,
plugging each basis vector of $V$ for the appropriate indeterminate.

For any partition $\lambda$ of length $\leq n$, we let $S_{\lambda
}^{\mathbf{A}}\left(  V\right)  \in\mathbf{A}$ denote the image of
$S_{\lambda}\left(  V\right)  \in\mathbf{S}\left(  V\right)  $ under this
morphism $\Phi_{V}^{\mathbf{A}}:\mathbf{S}\left(  V\right)  \rightarrow
\mathbf{A}$. This image $S_{\lambda}^{\mathbf{A}}\left(  V\right)  $ is what
often is called $S_{\lambda}\left(  V\right)  $ in places such as (7.18),
although we shall consistently call it $S_{\lambda}^{\mathbf{A}}\left(
V\right)  $ here. Generally, we shall use \textquotedblleft$\mathbf{A}%
$\textquotedblright\ superscripts to signal that we are mapping things into
$\mathbf{A}$ via the morphism $\Phi_{V}^{\mathbf{A}}:\mathbf{S}\left(
V\right)  \rightarrow\mathbf{A}$ or its close relatives.

In particular, we let $H_{r}^{\mathbf{A}}\left(  V\right)  $ and
$E_{r}^{\mathbf{A}}\left(  V\right)  $ (for any $r\in\mathbb{Z}$) be the
images of $H_{r}\left(  V\right)  $ and $E_{r}\left(  V\right)  $ under the
morphism $\Phi_{V}^{\mathbf{A}}:\mathbf{S}\left(  V\right)  \rightarrow
\mathbf{A}$.

Furthermore, we let $f_{V}^{\mathbf{A}}\in\mathbf{A}\left[  t\right]  $ be the
image of the polynomial $f_{V}=\prod_{v\in V}\left(  t+v\right)  \in\left(
\mathbf{S}\left(  V\right)  \right)  \left[  t\right]  $ under the $F\left[
t\right]  $-algebra morphism $\Phi_{V}^{\mathbf{A}}\left[  t\right]  :\left(
\mathbf{S}\left(  V\right)  \right)  \left[  t\right]  \rightarrow
\mathbf{A}\left[  t\right]  $ (which applies $\Phi_{V}^{\mathbf{A}}$ to each
coefficient of the polynomial). Explicitly, this just means that%
\begin{equation}
f_{V}^{\mathbf{A}}=\prod_{v\in V}\left(  t+v\right)
\ \ \ \ \ \ \ \ \ \ \left(  \text{product taken in }\mathbf{A}\left[
t\right]  \right)  . \label{p26.fAV}%
\end{equation}


Since $\Phi_{V}^{\mathbf{A}}$ is an $F$-algebra morphism, many of the
properties of the original elements $S_{\lambda}\left(  V\right)  $,
$H_{r}\left(  V\right)  $, $E_{r}\left(  V\right)  $ and $f_{V}$ can be
transferred mechanically to their images $S_{\lambda}^{\mathbf{A}}\left(
V\right)  $, $H_{r}^{\mathbf{A}}\left(  V\right)  $, $E_{r}^{\mathbf{A}%
}\left(  V\right)  $ and $f_{V}^{\mathbf{A}}$ under $\Phi_{V}^{\mathbf{A}}$.
For example:

\begin{itemize}
\item Applying the $F\left[  t\right]  $-algebra morphism $\Phi_{V}%
^{\mathbf{A}}\left[  t\right]  :\left(  \mathbf{S}\left(  V\right)  \right)
\left[  t\right]  \rightarrow\mathbf{A}\left[  t\right]  $ (which applies
$\Phi_{V}^{\mathbf{A}}$ to each coefficient of the polynomial) to the equality
(7.7), we obtain%
\begin{align}
f_{V}^{\mathbf{A}}\left(  t\right)   &  =t^{q^{n}}-E_{1}^{\mathbf{A}}\left(
V\right)  t^{q^{n-1}}+\cdots+\left(  -1\right)  ^{n}E_{n}^{\mathbf{A}}\left(
V\right)  t\nonumber\\
&  =\sum_{i=0}^{n}\left(  -1\right)  ^{i}E_{i}^{\mathbf{A}}\left(  V\right)
t^{q^{n-i}}. \label{p26.7.7}%
\end{align}


\item Applying the $F\left[  t\right]  $-algebra morphism $\Phi_{V}%
^{\mathbf{A}}\left[  t\right]  :\left(  \mathbf{S}\left(  V\right)  \right)
\left[  t\right]  \rightarrow\mathbf{A}\left[  t\right]  $ to the equality
(7.8), we obtain%
\begin{equation}
f_{V}^{\mathbf{A}}\left(  at+bu\right)  =af_{V}^{\mathbf{A}}\left(  t\right)
+bf_{V}^{\mathbf{A}}\left(  u\right)  . \label{p26.7.8}%
\end{equation}
That is, $f_{V}^{\mathbf{A}}$ is an additive (and $F$-linear) polynomial. An
alternative proof of this can be found in \cite[Theorem 1.6]{ppoly} (where
$F$-linear polynomials are called \textquotedblleft$q$%
-polynomials\textquotedblright).

\item If $\mathbf{A}$ is $\varphi$-bijective, then negative powers
$\varphi^{-j}$ of the Frobenius endomorphism $\varphi:\mathbf{A}%
\rightarrow\mathbf{A}$ are defined, and the $F$-algebra morphism $\Phi
_{V}^{\mathbf{A}}:\mathbf{S}\left(  V\right)  \rightarrow\mathbf{A}$ can be
canonically extended to an $F$-algebra morphism $\widehat{\Phi}_{V}%
^{\mathbf{A}}:\widehat{\mathbf{S}}\left(  V\right)  \rightarrow\mathbf{A}$ (by
requiring it to respect the Frobenius endomorphism). Applying this latter
morphism $\widehat{\Phi}_{V}^{\mathbf{A}}$ to the two matrices $\mathbf{H}%
\left(  V\right)  $ and $\mathbf{E}\left(  V\right)  $ (more precisely: to
each entry of these two matrices), we obtain the two infinite matrices%
\begin{align*}
\mathbf{H}^{\mathbf{A}}\left(  V\right)   &  :=\left(  \varphi^{i+1}%
H_{j-i}^{\mathbf{A}}\left(  V\right)  \right)  _{i,j\in\mathbb{Z}%
}\ \ \ \ \ \ \ \ \ \ \text{and}\\
\mathbf{E}^{\mathbf{A}}\left(  V\right)   &  :=\left(  \left(  -1\right)
^{j-i}\varphi^{j}E_{j-i}^{\mathbf{A}}\left(  V\right)  \right)  _{i,j\in
\mathbb{Z}}%
\end{align*}
over $\mathbf{A}$. Hence, by applying this morphism to (7.9), we obtain%
\begin{equation}
\mathbf{E}^{\mathbf{A}}\left(  V\right)  =\mathbf{H}^{\mathbf{A}}\left(
V\right)  ^{-1}. \label{p26.7.9}%
\end{equation}


\item Applying the morphism $\widehat{\Phi}_{V}^{\mathbf{A}}$ to the equality
(7.10), we conclude that%
\begin{align}
S_{\lambda}^{\mathbf{A}}\left(  V\right)   &  =\det\left(  \varphi
^{1-j}H_{\lambda_{i}-i+j}^{\mathbf{A}}\left(  V\right)  \right)
\label{p26.7.10a}\\
&  =\det\left(  \varphi^{j-1}E_{\lambda_{i}^{\prime}-i+j}^{\mathbf{A}}\left(
V\right)  \right)  \label{p26.7.10b}%
\end{align}
whenever $\mathbf{A}$ is $\varphi$-bijective and $\lambda$ is a partition of
length $\leq n=\dim V$.

\item If $\mathbf{A}$ is $\varphi$-bijective, then, for any two partitions
$\lambda$ and $\mu$, we define $S_{\lambda/\mu}^{\mathbf{A}}\left(  V\right)
$ to be the image of $S_{\lambda/\mu}\left(  V\right)  $ under the $F$-algebra
morphism $\widehat{\Phi}_{V}^{\mathbf{A}}:\widehat{\mathbf{S}}\left(
V\right)  \rightarrow\mathbf{A}$. Hence, the formula (7.11) yields that%
\begin{equation}
S_{\lambda/\mu}^{\mathbf{A}}\left(  V\right)  =\det\left(  \varphi^{\mu
_{j}-j+1}H_{\lambda_{i}-\mu_{j}-i+j}^{\mathbf{A}}\left(  V\right)  \right)  ;
\label{p26.7.11}%
\end{equation}
the formula (7.11') yields%
\begin{equation}
S_{\lambda/\mu}^{\mathbf{A}}\left(  V\right)  =\det\left(  \varphi^{-\mu
_{j}^{\prime}+j-1}E_{\lambda_{i}^{\prime}-\mu_{j}^{\prime}-i+j}^{\mathbf{A}%
}\left(  V\right)  \right)  ; \label{p26.7.11'}%
\end{equation}
and the formula (7.12) leads to%
\begin{equation}
S_{\lambda/\mu}^{\mathbf{A}}\left(  V\right)
=0\ \ \ \ \ \ \ \ \ \ \text{unless }0\leq\lambda_{i}^{\prime}-\mu_{i}^{\prime
}\leq n\text{ for all }i\geq1. \label{p26.7.12}%
\end{equation}

\end{itemize}

We can now properly define the $f_{U}$ that is used at the bottom of page 27:
It is simply the polynomial $f_{U}^{\mathbf{S}\left(  V\right)  }\in\left(
\mathbf{S}\left(  V\right)  \right)  \left[  t\right]  $ (not the polynomial
$f_{U}^{\mathbf{S}\left(  U\right)  }\in\left(  \mathbf{S}\left(  U\right)
\right)  \left[  t\right]  $). Of course, when writing \textquotedblleft%
$f_{U}\left(  V\right)  $\textquotedblright, the intended meaning is the set
of all values of this polynomial $f_{U}=f_{U}^{\mathbf{S}\left(  V\right)  }$
taken on the values of $V$, that is, the set $\left\{  f_{U}\left(  v\right)
\ \mid\ v\in V\right\}  $.

More generally, all the \textquotedblleft$f_{W}$\textquotedblright s in what
follows are to be understood as either $f_{W}^{\mathbf{S}\left(  V\right)  }$
or $f_{W}^{\widehat{\mathbf{S}}\left(  V\right)  }$ (depending on whether
there is need for negative powers of $\varphi$).

Finally, let $\mathbf{A}$ be a commutative $F$-algebra, and let $U$ and $V$ be
two finite-dimensional $F$-vector subspaces of $\mathbf{A}$ satisfying
$U\subseteq V$. Then, we consider the set%
\[
f_{U}^{\mathbf{A}}\left(  V\right)  =\left\{  f_{U}^{\mathbf{A}}\left(
v\right)  \ \mid\ v\in V\right\}  =\left\{  \prod_{u\in U}\left(  v+u\right)
\ \mid\ v\in V\right\}  .
\]
This is simply the image of the map%
\begin{align*}
\psi_{U,V}^{\mathbf{A}}:V  &  \rightarrow\mathbf{A},\\
v  &  \mapsto f_{U}^{\mathbf{A}}\left(  v\right)  =\prod_{u\in U}\left(
v+u\right)  ,
\end{align*}
which is $F$-linear (by (\ref{p26.7.8})); thus, $f_{U}^{\mathbf{A}}\left(
V\right)  $ is an $F$-vector subspace of $\mathbf{A}$. Also, $U$ is always
contained in the kernel of $\psi_{U,V}^{\mathbf{A}}$ (since the product
$\prod_{u\in U}\left(  v+u\right)  $ has a zero factor when $v\in U$). If
$\mathbf{A}$ is an integral domain, then $U$ is precisely the kernel of
$\psi_{U,V}^{\mathbf{A}}$, and thus the isomorphism theorem yields that the
image $f_{U}^{\mathbf{A}}\left(  V\right)  $ of $\psi_{U,V}^{\mathbf{A}}$ is
isomorphic to the quotient space $V/\operatorname*{Ker}\psi_{U,V}^{\mathbf{A}%
}=V/U$. Hence, in this case, we \textbf{define} the notation%
\begin{equation}
V/U:=f_{U}^{\mathbf{A}}\left(  V\right)  , \label{p26.7.14}%
\end{equation}
generalizing (7.14). However, if $\mathbf{A}$ is not an integral domain, then
$f_{U}^{\mathbf{A}}\left(  V\right)  $ might have smaller dimension than
$V/U$, and so this notation would be misleading.

It is worth remarking that every finite-dimensional $F$-vector subspace $V$ of
$\mathbf{A}$ satisfies the equalities (not just isomorphisms!) $V/V=\left\{
0\right\}  $ and $V/\left\{  0\right\}  =V$.

\item \textbf{page 28, (7.15):} The \textquotedblleft$f_{U}$\textquotedblright%
\ here is probably meant to be $f_{U}^{\mathbf{S}\left(  V\right)  }$ (and
likewise, the \textquotedblleft$f_{V/U}$\textquotedblright\ probably means
$f_{V/U}^{\mathbf{S}\left(  V\right)  }$).

However, it is actually better to generalize (7.15) to the following setup: We
let $\mathbf{A}$ be a commutative $F$-algebra that is an integral domain (in
order for the notation $V/U$ to make sense), and we let $U$ and $V$ be two
finite-dimensional $F$-vector subspaces of $\mathbf{A}$ satisfying $U\subseteq
V$. In this case, (7.15) holds, where we understand each \textquotedblleft%
$f$\textquotedblright\ in (7.15) to have an \textquotedblleft$\mathbf{A}%
$\textquotedblright\ superscript -- i.e., the equality \textquotedblleft%
$f_{V}=f_{V/U}\circ f_{U}$\textquotedblright\ should be read as
\textquotedblleft$f_{V}^{\mathbf{A}}=f_{V/U}^{\mathbf{A}}\circ f_{U}%
^{\mathbf{A}}$\textquotedblright.

The proof of this generalization is just as easy as the original proof of
(7.15), since the polynomials $f_{V}^{\mathbf{A}}$ are given by (\ref{p26.fAV}%
). It is this general version that I believe is necessary to prove (7.19) by
induction further below.

\item \textbf{page 28, (7.16):} Again, the proper level of generality at which
this claim holds is the following: We let $\mathbf{A}$ be a commutative
$F$-algebra that is an integral domain (in order for the notations $V/U$,
$V/T$ and $U/T$ to make sense), and we let $T$, $U$ and $V$ be three
finite-dimensional $F$-vector subspaces of $\mathbf{A}$ satisfying $T\subseteq
U\subseteq V$. Then, we have $U/T\subseteq V/T$ and $V/U=\left(  V/T\right)
/\left(  U/T\right)  $ (as an equality between subspaces of $\mathbf{A}$).

The proof given here applies verbatim to this generalization (as long as we
understand each \textquotedblleft$f$\textquotedblright\ to have an
\textquotedblleft$\mathbf{A}$\textquotedblright\ superscript).

\item \textbf{page 28, (7.17):} Again, this can and should be generalized: We
let $\mathbf{A}$ be a commutative $F$-algebra that is $\varphi$-bijective and
is an integral domain, and we let $U$ and $V$ be two finite-dimensional
$F$-vector subspaces of $\mathbf{A}$ satisfying $U\subseteq V$. Then, we have%
\begin{align*}
\text{(i)}\ \ \ \ \ \ \ \ \ \  &  \mathbf{E}^{\mathbf{A}}\left(  V\right)
=\varphi^{\dim\left(  V/U\right)  }\left(  \mathbf{E}^{\mathbf{A}}\left(
U\right)  \right)  \cdot\mathbf{E}^{\mathbf{A}}\left(  V/U\right)  ;\\
\text{(ii)}\ \ \ \ \ \ \ \ \ \  &  \mathbf{H}^{\mathbf{A}}\left(  V\right)
=\mathbf{H}^{\mathbf{A}}\left(  V/U\right)  \cdot\varphi^{\dim\left(
V/U\right)  }\left(  \mathbf{H}^{\mathbf{A}}\left(  U\right)  \right)  .
\end{align*}


Again, the proof applies verbatim to this generalization, once the obvious
changes are made (e.g., when the equalities (7.9) and (7.11) are cited here,
what is actually being used are the equalities (\ref{p26.7.9}) and
(\ref{p26.7.11})).

\item \textbf{page 29, (7.18):} Replace \textquotedblleft$S_{\lambda/\mu
}\left(  U\right)  $\textquotedblright\ by \textquotedblleft$S_{\lambda/\nu
}\left(  U\right)  $\textquotedblright.

Moreover, again, the claim should be generalized: We let $\mathbf{A}$ be a
commutative $F$-algebra that is $\varphi$-bijective and is an integral domain,
and we let $U$ and $V$ be two finite-dimensional $F$-vector subspaces of
$\mathbf{A}$ satisfying $U\subseteq V$. Then, the claim of (7.18) holds (once
the \textquotedblleft$S_{\lambda/\mu}\left(  U\right)  $\textquotedblright%
\ has been corrected to \textquotedblleft$S_{\lambda/\nu}\left(  U\right)
$\textquotedblright, and once each of the three \textquotedblleft%
$S$\textquotedblright s has received an \textquotedblleft$\mathbf{A}%
$\textquotedblright\ superscript atop of it\footnote{In other words,
\textquotedblleft$S_{\lambda/\mu}\left(  V\right)  $\textquotedblright,
\textquotedblleft$S_{\nu/\mu}\left(  V/U\right)  $\textquotedblright\ and
\textquotedblleft$S_{\lambda/\nu}\left(  U\right)  $\textquotedblright\ should
be read as \textquotedblleft$S_{\lambda/\mu}^{\mathbf{A}}\left(  V\right)
$\textquotedblright, \textquotedblleft$S_{\nu/\mu}^{\mathbf{A}}\left(
V/U\right)  $\textquotedblright\ and \textquotedblleft$S_{\lambda/\nu
}^{\mathbf{A}}\left(  U\right)  $\textquotedblright, respectively.}).

\item \textbf{page 29, proof of (7.18):} Again, this proof is an application
of the Cauchy--Binet theorem, analogous to the proof of (6.13) which I have
detailed above.

\item \textbf{page 29, (7.19):} Replace \textquotedblleft subspaces of
$V$\textquotedblright\ by \textquotedblleft finite-dimensional subspaces of a
$\varphi$-bijective commutative $\mathbf{F}$-algebra $\mathbf{A}$ that is an
integral domain\textquotedblright. And again, put an \textquotedblleft%
$\mathbf{A}$\textquotedblright\ superscript atop each \textquotedblleft%
$S$\textquotedblright.

\item \textbf{page 29, (7.19):} It is worth pointing out that the base case
($r=0$) in the induction proof of (7.19) relies on the fact that the
$0$-dimensional subspace $\left\{  0\right\}  $ of $\mathbf{A}$ satisfies
$S_{\lambda/\mu}\left(  \left\{  0\right\}  \right)  =%
\begin{cases}
1, & \text{if }\lambda=\mu;\\
0, & \text{if }\lambda\neq\mu
\end{cases}
$\ \ for all $\lambda$ and $\mu$. (This is easiest to see from (\ref{p26.7.11}%
) and (\ref{p26.7.12}).)

\item \textbf{page 30:} In \textquotedblleft If $U$ is any finite-dimensional
subspace of $F\left(  V\right)  $\textquotedblright, replace \textquotedblleft%
$F\left(  V\right)  $\textquotedblright\ by \textquotedblleft$\mathbf{S}%
\left(  V\right)  $\textquotedblright\ -- or, better, generalize it to
\textquotedblleft any commutative $F$-algebra $\mathbf{A}$\textquotedblright.

\item \textbf{page 30, (7.22):} Here, too, \textquotedblleft$\mathbf{S}\left(
V\right)  $\textquotedblright\ can be replaced by any $\varphi$-bijective
commutative $F$-algebra $\mathbf{A}$. Again, this level of generality can be
obtained either by adapting the given proof (putting \textquotedblleft%
$\mathbf{A}$\textquotedblright\ superscripts over many of the symbols), or by
applying the $F$-algebra morphism $\widehat{\Phi}_{V}^{\mathbf{A}%
}:\widehat{\mathbf{S}}\left(  V\right)  \rightarrow\mathbf{A}$ to the original (7.22).

\item \textbf{page 30, proof of (7.22):} \textquotedblleft by
(7.10)\textquotedblright\ could better be replaced by \textquotedblleft by the
second formula of (7.10)\textquotedblright. Indeed, the second formula of
(7.10) (or the equality (\ref{p26.7.10b}), if you are working in the general
setting of an $F$-algebra $\mathbf{A}$) yields%
\begin{align*}
S_{\left(  r\right)  }\left(  U\right)   &  =\det\left(  \varphi
^{j-1}E_{1-i+j}\left(  U\right)  \right)  _{i,j\in\left[  r\right]
}\ \ \ \ \ \ \ \ \ \ \left(  \text{since }\left(  r\right)  ^{\prime}=\left(
1,1,\ldots,1\right)  \in\mathbb{N}^{r}\right)  \\
&  =\det\left(
\begin{array}
[c]{ccccc}%
\varphi^{0}E_{1}\left(  U\right)   & 0 & 0 & \cdots & 0\\
\varphi^{0}E_{0}\left(  U\right)   & \varphi^{1}E_{1}\left(  U\right)   & 0 &
\cdots & 0\\
\varphi^{0}E_{-1}\left(  U\right)   & \varphi^{1}E_{0}\left(  U\right)   &
\varphi^{2}E_{1}\left(  U\right)   & \cdots & 0\\
\vdots & \vdots & \vdots & \ddots & \vdots\\
\varphi^{0}E_{2-r}\left(  U\right)   & \varphi^{1}E_{3-r}\left(  U\right)   &
\varphi^{2}E_{4-r}\left(  U\right)   & \cdots & \varphi^{r-1}E_{1}\left(
U\right)
\end{array}
\right)  \\
&  \ \ \ \ \ \ \ \ \ \ \ \ \ \ \ \ \ \ \ \ \left(  \text{since }E_{k}\left(
U\right)  =0\text{ for all }k\geq2\right)  \\
&  =\prod_{j=1}^{r}\varphi^{j-1}E_{1}\left(  U\right)
\end{align*}
(since the determinant of a triangular matrix equals the product of its
diagonal entries).

\item \textbf{page 31, proof of (7.22):} The first equality sign on this page
follows from (7.11) (or from (\ref{p26.7.11}), if you are working in the
general setting of an $F$-algebra $\mathbf{A}$). The second equality sign is
analogous to (\ref{p19.last.diag}).

\item \textbf{page 31, (7.23):} Again, this holds not only in
$\widehat{\mathbf{S}}\left(  V\right)  $ but also in any $\varphi$-bijective
commutative $F$-algebra $\mathbf{A}$ that is an integral domain.

\item \textbf{page 31:} \textquotedblleft\textit{column strict}%
\textquotedblright\ should be \textquotedblleft\textit{column-strict}%
\textquotedblright.

\item \textbf{page 32:} The claim that \textquotedblleft the number of flags
$\mathfrak{V}$ in $V$ is congruent to $1$ modulo $q$\textquotedblright\ is
worth giving a proof.

(The simplest proof that comes to my mind is by induction on $n$. In the
induction step, we can argue that\footnote{We pick an arbitrary nondegenerate
symmetric bilinear form on $V$. The word \textquotedblleft
normal\textquotedblright\ shall refer to orthogonality with respect to this
form.}
\begin{align*}
&  \left(  \text{number of flags in }V\right) \\
&  =\sum_{H\text{ is a hyperplane in }V}\underbrace{\left(  \text{number of
flags in }V\text{ starting with }V>H\right)  }_{\substack{=\left(
\text{number of flags in }H\right)  \\\equiv1\operatorname{mod}q\\\text{(by
the induction hypothesis)}}}\\
&  \equiv\sum_{H\text{ is a hyperplane in }V}1\\
&  =\left(  \text{number of hyperplanes in }V\right) \\
&  =\dfrac{q^{n}-1}{q-1}\ \ \ \ \ \ \ \ \ \ \left(
\begin{array}
[c]{c}%
\text{since each of the }q^{n}-1\text{ nonzero vectors in }V\\
\text{determines a unique hyperplane normal to it, while}\\
\text{each hyperplane has }q-1\text{ nonzero normal vectors}%
\end{array}
\right) \\
&  =1+q+q^{2}+\cdots+q^{n-1}\equiv1\operatorname{mod}q.
\end{align*}
)

\item \textbf{page 32, (7.23'):} The comma at the end of this equality should
be a period.

\item \textbf{page 32, display between (7.24 ?) and (7.25 ?):} Replace
\textquotedblleft$\psi^{j-1}$\textquotedblright\ by \textquotedblleft%
$\varphi^{j-1}$\textquotedblright.

\item \textbf{page 32:} It is worth saying that the validity of the formula
(7.24 ?) for $\lambda=\left(  1^{r}\right)  $ is proved in \cite[\S I.2,
Example 26 (d)]{Macdon95}.

\item \textbf{page 33, (7.26):} The right hand side of this equality (or
alternatively the left hand side) needs to be multiplied by $\left(
-1\right)  ^{nm}$.

\item \textbf{page 33, (7.26):} It is worth pointing out that the two sides of
this equality live in the polynomial ring $F\left[  x_{1},x_{2},\ldots
,x_{n},y_{1},y_{2},\ldots,y_{m}\right]  =\mathbf{S}\left(  V\oplus W\right)  $
(which, of course, contains $\mathbf{S}\left(  V\right)  $ and $\mathbf{S}%
\left(  W\right)  $ as subrings).

\item \textbf{page 33, proof of (7.26):} The product sign \textquotedblleft%
$\prod_{\substack{w\in V_{0}\\w\in W_{0}}}$\textquotedblright\ should be
\textquotedblleft$\prod_{\substack{v\in V_{0}\\w\in W_{0}}}$\textquotedblright.

\item \textbf{page 33, proof of (7.26):} Let me explain why the first equality
sign in the equality (1) of this proof holds, i.e., why we have%
\begin{equation}
\prod_{\substack{v\in V_{0};\\w\in W;\\w\neq0}}\left(  v+w\right)
=\prod_{\substack{v\in V_{0};\\w\in W_{0}}}\left(  v^{q-1}-w^{q-1}\right)  .
\label{p33.7.26.1}%
\end{equation}
Clearly, it suffices to show that each $v\in V_{0}$ satisfies%
\begin{equation}
\prod_{\substack{w\in W;\\w\neq0}}\left(  v+w\right)  =\prod_{w\in W_{0}%
}\left(  v^{q-1}-w^{q-1}\right)  . \label{p33.7.26.2}%
\end{equation}
Before we prove this, we need two lemmas:

\begin{lemmain}
\label{lem.t-q-1}In the polynomial ring $F\left[  t\right]  $, we have%
\begin{equation}
\prod_{\alpha\in F\setminus\left\{  0\right\}  }\left(  t-\alpha\right)
=t^{q-1}-1. \label{eq.lem.t-q-1.eq}%
\end{equation}

\end{lemmain}

\begin{proof}
[Proof of Lemma \ref{lem.t-q-1}.]Lagrange's theorem (applied to the
multiplicative group $\left(  F\setminus\left\{  0\right\}  ,\ \cdot
,\ 1\right)  $ of $F$) shows that $\alpha^{\left\vert F\right\vert -1}=1$ for
each $\alpha\in F\setminus\left\{  0\right\}  $. In other words, the
polynomial $t^{q-1}-1$ vanishes at each $\alpha\in F\setminus\left\{
0\right\}  $. Hence, it can be written as%
\begin{equation}
t^{q-1}-1=w\left(  t\right)  \cdot\prod_{\alpha\in F\setminus\left\{
0\right\}  }\left(  t-\alpha\right)  \label{p33.7.26.4}%
\end{equation}
for some polynomial $w\left(  t\right)  \in F\left[  t\right]  $. This latter
polynomial $w\left(  t\right)  $ must have degree $0$ (by comparing degrees in
(\ref{p33.7.26.4}), since both polynomials $t^{q-1}-1$ and $\prod_{\alpha\in
F\setminus\left\{  0\right\}  }\left(  t-\alpha\right)  $ have degree
$\left\vert F\setminus\left\{  0\right\}  \right\vert =q-1$), and thus is a
constant; furthermore, this constant must be $1$ (by comparing leading terms
in (\ref{p33.7.26.4}), since both polynomials $t^{q-1}-1$ and $\prod
_{\alpha\in F\setminus\left\{  0\right\}  }\left(  t-\alpha\right)  $ are
monic). That is, $w\left(  t\right)  =1$. Hence, (\ref{p33.7.26.4}) simplifies
to $t^{q-1}-1=\prod_{\alpha\in F\setminus\left\{  0\right\}  }\left(
t-\alpha\right)  $. This proves Lemma \ref{lem.t-q-1}.
\end{proof}

\begin{lemmain}
\label{lem.t-q-2}In the polynomial ring $F\left[  t,s\right]  $, we have%
\begin{equation}
\prod_{\alpha\in F\setminus\left\{  0\right\}  }\left(  t+\alpha s\right)
=t^{q-1}-s^{q-1}. \label{eq.lem.t-q-2.eq}%
\end{equation}

\end{lemmain}

\begin{proof}
[Proof of Lemma \ref{lem.t-q-2}.]Consider the field of rational functions
$F\left(  t,s\right)  $. Substituting $t/s\in F\left(  t,s\right)  $ for $t$
in (\ref{eq.lem.t-q-1.eq}), we obtain%
\[
\prod_{\alpha\in F\setminus\left\{  0\right\}  }\left(  t/s-\alpha\right)
=\left(  t/s\right)  ^{q-1}-1.
\]


However,%
\[
\prod_{\alpha\in F\setminus\left\{  0\right\}  }\left(  t+\alpha s\right)
=\prod_{\alpha\in F\setminus\left\{  0\right\}  }\left(  t+\left(
-\alpha\right)  s\right)
\]
(here, we have substituted $-\alpha$ for $\alpha$ in the product, since the
map $F\setminus\left\{  0\right\}  \rightarrow F\setminus\left\{  0\right\}
,\ \alpha\mapsto-\alpha$ is a bijection). Thus,%
\begin{align*}
\prod_{\alpha\in F\setminus\left\{  0\right\}  }\left(  t+\alpha s\right)   &
=\prod_{\alpha\in F\setminus\left\{  0\right\}  }\underbrace{\left(  t+\left(
-\alpha\right)  s\right)  }_{=t-\alpha s=s\left(  t/s-\alpha\right)  }\\
&  =\prod_{\alpha\in F\setminus\left\{  0\right\}  }\left(  s\left(
t/s-\alpha\right)  \right)  =\underbrace{s^{\left\vert F\setminus\left\{
0\right\}  \right\vert }}_{\substack{=s^{q-1}\\\text{(since }\left\vert
F\setminus\left\{  0\right\}  \right\vert =q-1\text{)}}}\ \ \underbrace{\prod
_{\alpha\in F\setminus\left\{  0\right\}  }\left(  t/s-\alpha\right)
}_{=\left(  t/s\right)  ^{q-1}-1}\\
&  =s^{q-1}\left(  \left(  t/s\right)  ^{q-1}-1\right)  =t^{q-1}-s^{q-1}.
\end{align*}
This proves Lemma \ref{lem.t-q-2}.
\end{proof}

\begin{proof}
[Proof of (\ref{p33.7.26.2}).]Fix $v\in V_{0}$. Each nonzero vector $w\in W$
can be written as $\alpha w^{\prime}$ for a unique pair $\left(
\alpha,w^{\prime}\right)  $ of a nonzero scalar $\alpha\in F\setminus\left\{
0\right\}  $ and a vector $w^{\prime}\in W_{0}$ (indeed, $\alpha$ is simply
the first nonzero coordinate of $w$, while $w^{\prime}$ is $\dfrac{1}{\alpha
}w$). Thus, in $\mathbf{S}\left(  V\oplus W\right)  $, we have
\begin{align*}
\prod_{\substack{w\in W;\\w\neq0}}\left(  v+w\right)   &  =\prod
_{\substack{\alpha\in F\setminus\left\{  0\right\}  ;\\w^{\prime}\in W_{0}%
}}\left(  v+\alpha w^{\prime}\right)  =\prod_{w^{\prime}\in W_{0}%
}\ \ \underbrace{\prod_{\alpha\in F\setminus\left\{  0\right\}  }\left(
v+\alpha w^{\prime}\right)  }_{\substack{=v^{q-1}-\left(  w^{\prime}\right)
^{q-1}\\\text{(by (\ref{eq.lem.t-q-2.eq}), with }v\text{ and }w^{\prime
}\\\text{substituted for }t\text{ and }s\text{)}}}\\
&  =\prod_{w^{\prime}\in W_{0}}\left(  v^{q-1}-\left(  w^{\prime}\right)
^{q-1}\right)  =\prod_{w\in W_{0}}\left(  v^{q-1}-w^{q-1}\right)  .
\end{align*}
This proves (\ref{p33.7.26.2}).
\end{proof}

\item \textbf{page 33, proof of (7.26):} On the right hand side of the
equality (1) of this proof, replace \textquotedblleft$\pi\left(  L\right)
-\pi\left(  M\right)  $\textquotedblright\ by \textquotedblleft$\pi\left(
M\right)  -\pi\left(  L\right)  $\textquotedblright.

\item \textbf{page 33, proof of (7.26):} The second equality sign in (1) --
that is, the identity%
\[
\prod_{\substack{v\in V_{0};\\w\in W_{0}}}\left(  v^{q-1}-w^{q-1}\right)
=\prod_{L,M}\left(  \pi\left(  M\right)  -\pi\left(  L\right)  \right)
\]
-- holds for the following reasons:

Firstly, each line $L$ in $V$ can be written as $\operatorname*{span}\left\{
v\right\}  $ for a unique $v\in V_{0}$. Thus, taking a product over $v\in
V_{0}$ is tantamount to taking a product over lines $L$ in $V$.

Secondly, each line $M$ in $W$ can be written as $\operatorname*{span}\left\{
w\right\}  $ for a unique $w\in W_{0}$. Thus, taking a product over $w\in
W_{0}$ is tantamount to taking a product over lines $M$ in $W$.

Thirdly, for any line $K=\operatorname*{span}\left\{  x\right\}  $ in $V\oplus
W$, we have $x^{q-1}=-\pi\left(  K\right)  $. In order to prove this, we need
to recall the following basic fact (Wilson's theorem for finite fields):

\begin{lemmain}
\label{lem.wilson-F}We have $\prod_{\alpha\in F\setminus\left\{  0\right\}
}\alpha=-1$.
\end{lemmain}

\begin{proof}
[Proof of Lemma \ref{lem.wilson-F}.]In the multiplicative group $\left(
F\setminus\left\{  0\right\}  ,\ \cdot,\ 1\right)  $, each element $\alpha$
has an inverse $\alpha^{-1}$. Thus, almost all factors of the product
$\prod_{\alpha\in F\setminus\left\{  0\right\}  }\alpha$ can be paired up into
pairs of the form $\left(  \alpha,\alpha^{-1}\right)  $ (since $\left(
\alpha^{-1}\right)  ^{-1}$ is again $\alpha$). The only factors that end up
unpaired are the factors $\alpha$ that are their own inverses; these are $1$
and $-1$ (since $\alpha=\alpha^{-1}$ entails $\alpha^{2}=1$, thus $\left(
\alpha-1\right)  \left(  \alpha+1\right)  =\alpha^{2}-1=0$, hence $\alpha=1$
or $\alpha=-1$). The paired-up factors $\alpha,\alpha^{-1}$ always contribute
$\alpha\alpha^{-1}=1$ to the product $\prod_{\alpha\in F\setminus\left\{
0\right\}  }\alpha$, and thus can be dropped without changing the product. The
product therefore simplifies to $\prod_{\substack{\alpha\in F\setminus\left\{
0\right\}  \\\text{is unpaired}}}\alpha$, which is either $1\cdot\left(
-1\right)  $ or $-1$ (depending on whether the two unpaired factors $1$ and
$-1$ are distinct or not), but the result is $-1$ in either case. Hence, we
have shown that $\prod_{\alpha\in F\setminus\left\{  0\right\}  }\alpha=-1$.
This proves Lemma \ref{lem.wilson-F}.
\end{proof}

Now, our claim that $x^{q-1}=-\pi\left(  K\right)  $ for
$K=\operatorname*{span}\left\{  x\right\}  $ is easy to see: Indeed,
$K=\operatorname*{span}\left\{  x\right\}  =\left\{  \alpha x\ \mid\ \alpha\in
F\right\}  $, so that%
\begin{align*}
\pi\left(  K\right)   &  =\left(  \text{product of the nonzero elements of
}K\right) \\
&  =\prod_{\alpha\in F\setminus\left\{  0\right\}  }\left(  \alpha x\right)
=\underbrace{\left(  \prod_{\alpha\in F\setminus\left\{  0\right\}  }%
\alpha\right)  }_{\substack{=-1\\\text{(by Lemma \ref{lem.wilson-F})}%
}}\underbrace{x^{\left\vert F\setminus\left\{  0\right\}  \right\vert }%
}_{\substack{=x^{q-1}\\\text{(since }\left\vert F\setminus\left\{  0\right\}
\right\vert =q-1\text{)}}}\\
&  =-x^{q-1},
\end{align*}
and therefore $x^{q-1}=-\pi\left(  K\right)  $, qed.

\item \textbf{page 33, proof of (7.26):} The equality (2) is analogous to
(\ref{p20.claim}), and so is its proof.

\item \textbf{page 33, proof of (7.26):} In the last displayed equality in
this proof, again replace \textquotedblleft$\pi\left(  L\right)  -\pi\left(
M\right)  $\textquotedblright\ by \textquotedblleft$\pi\left(  M\right)
-\pi\left(  L\right)  $\textquotedblright.

\item \textbf{page 33, proof of (7.26):} In the last sentence of the proof,
replace \textquotedblleft each $y_{j}$ by $\omega y_{j}$\textquotedblright\ by
\textquotedblleft each $x_{i}$ by $\omega x_{i}$\textquotedblright.

To give a bit more detail: Substituting each $x_{i}$ by $\omega x_{i}$ affects
the left hand side of the preceding equality by replacing each $\pi\left(
L\right)  $ by $\underbrace{\omega^{q-1}}_{=-1}\pi\left(  L\right)
=-\pi\left(  L\right)  $, which causes the difference $\pi\left(  M\right)
-\pi\left(  L\right)  $ to turn into the sum $\pi\left(  L\right)  +\pi\left(
M\right)  $. On the right hand side, each $S_{\lambda}\left(  V\right)  $
turns into $\omega^{\deg S_{\lambda}\left(  V\right)  }S_{\lambda}\left(
V\right)  $, which can be easily seen to equal $\left(  -1\right)
^{\left\vert \lambda\right\vert }S_{\lambda}\left(  V\right)  $ (because when
$q$ is odd, the polynomial $S_{\lambda}\left(  V\right)  $ is homogeneous of
degree%
\begin{align*}
&  \sum_{i=1}^{n}\underbrace{\left(  q^{\lambda_{i}}-1\right)  }_{=\left(
q-1\right)  \left(  1+q+q^{2}+\cdots+q^{\lambda_{i}-1}\right)  }q^{n-i}\\
&  =\left(  q-1\right)  \sum_{i=1}^{n}\underbrace{\left(  1+q+q^{2}%
+\cdots+q^{\lambda_{i}-1}\right)  }_{\equiv\lambda_{i}\operatorname{mod}%
2}\underbrace{q^{n-i}}_{\equiv1\operatorname{mod}2}\\
&  \equiv\left(  q-1\right)  \cdot\underbrace{\sum_{i=1}^{n}\lambda_{i}%
}_{=\left\vert \lambda\right\vert }=\left(  q-1\right)  \cdot\left\vert
\lambda\right\vert \operatorname{mod}2\left(  q-1\right)  ;
\end{align*}
whereas when $q$ is even, all signs are equal in $F$), and the $\left(
-1\right)  ^{\left\vert \lambda\right\vert }$ factor here combines with the
$\left(  -1\right)  ^{\left\vert \widehat{\lambda}\right\vert }$ sign in the
equality to form the product $\left(  -1\right)  ^{\left\vert \lambda
\right\vert }\left(  -1\right)  ^{\left\vert \widehat{\lambda}\right\vert
}=\left(  -1\right)  ^{\left\vert \lambda\right\vert +\left\vert
\widehat{\lambda}\right\vert }=\left(  -1\right)  ^{\left\vert \left(
m^{n}\right)  \right\vert }=\left(  -1\right)  ^{nm}$.
\end{enumerate}

\subsection{Corrections to the 8th Variation}

\begin{enumerate}
\item \textbf{page 34, (8.3):} Another sufficient criterion for when
row-flagged and column-flagged Schur functions agree can be found in
\cite[Theorem 1.4]{McDowe23}.
\end{enumerate}

\subsection{Corrections to the 9th Variation}

\begin{enumerate}
\item \textbf{page 35:} In the first paragraph of page 35, after
\textquotedblleft Define an automorphism\textquotedblright, add
\textquotedblleft$\varphi$\textquotedblright.

\item \textbf{page 35, (9.1):} Note that the determinant $\det\left(
\varphi^{\mu_{j}-j+1}h_{\lambda_{i}-\mu_{j}-i+j}\right)  _{1\leq i,j\leq n}$
does not depend on the choice of $n$, since one can easily see that if
$\lambda_{n}=\mu_{n}=0$, then the last row of the matrix $\left(  \varphi
^{\mu_{j}-j+1}h_{\lambda_{i}-\mu_{j}-i+j}\right)  _{1\leq i,j\leq n}$ is
$\left(  0,0,\ldots,0,1\right)  $, and therefore Laplace expansion along this
row reduces the determinant to the determinant of the submatrix $\left(
\varphi^{\mu_{j}-j+1}h_{\lambda_{i}-\mu_{j}-i+j}\right)  _{1\leq i,j\leq n-1}%
$. Thus, $s_{\lambda/\mu}$ does not depend on $n$.

\item \textbf{page 35:} Yet another specialization of the \textquotedblleft
Schur functions\textquotedblright\ defined by (9.1) are the \emph{loop Schur
functions} (with loop parameter $n$) introduced by Lam and Pylyavskyy in
\cite[\S 7.2]{LamPyl12}. Indeed, the determinantal formula in \cite[Theorem
7.4]{LamPyl12} shows that they are obtained by specializing $s_{\lambda/\mu}$
for $h_{rs}=h_{r}^{\left(  s-1\operatorname{mod}n\right)  }\left(
\mathbf{x}\right)  $, where $\left(  s-1\operatorname{mod}n\right)  $ denotes
the residue class of $s-1$ modulo $n$. Note that the $h_{r}^{\left(  k\right)
}\left(  \mathbf{x}\right)  $ for all $r>0$ and $k\in\mathbb{Z}/n$ are
algebraically independent (see \cite[Remark 7.1]{LamPyl12}). Thus, properties
of loop Schur functions proved for each sufficiently large $n$ can be
translated into properties of the $s_{\lambda/\mu}$ from (9.1).

The loop Schur functions have since been studied in other sources, e.g.,
\cite{Lam10}, \cite{Friede20}, \cite{Ross13}.

\item \textbf{page 35, proof of (9.3):} Just as in the proof of (6.6), this
computation only shows that $HE$ is the identity matrix. In order to prove
that the two matrices $H$ and $E$ are inverse to each other (i.e., in order to
prove (9.3)), it must also be shown that $EH$ is the identity matrix.

Again, this can be done by observing that both matrices belong to the group
$\operatorname*{UT}\nolimits_{I}$ of all upper unitriangular matrices with
rows and columns indexed by the elements of $I$.

\item \textbf{page 36, proof of (9.3):} Let me explain how \textquotedblleft
expansion of this determinant along the top row gives (1)\textquotedblright.
The fact that is used here is the following lemma:

\begin{lemmain}
\label{lem.36.det-hess}Let $A=\left(  a_{i,j}\right)  _{i,j\in\left[
n\right]  }$ be an $n\times n$-matrix. Assume that its entries satisfy%
\begin{equation}
a_{i,j}=1\ \ \ \ \ \ \ \ \ \ \text{whenever }i=j+1.
\label{eq.lem.36.det-hess.ass.1}%
\end{equation}
Assume furthermore that%
\begin{equation}
a_{i,j}=0\ \ \ \ \ \ \ \ \ \ \text{whenever }i>j+1.
\label{eq.lem.36.det-hess.ass.0}%
\end{equation}
Then,%
\[
\det A=\sum_{k=1}^{n}\left(  -1\right)  ^{k+1}a_{1,k}\det\left(
a_{i+k,j+k}\right)  _{i,j\in\left[  n-k\right]  }.
\]

\end{lemmain}

\begin{proof}
[Proof of Lemma \ref{lem.36.det-hess}.]We have $A=\left(  a_{i,j}\right)
_{i,j\in\left[  n\right]  }$. Thus, Laplace expansion along the first row of
$\det A$ yields%
\begin{equation}
\det A=\sum_{k=1}^{n}\left(  -1\right)  ^{1+k}a_{1,k}\det\left(
\operatorname*{sub}\nolimits_{\left[  n\right]  \setminus\left\{  1\right\}
}^{\left[  n\right]  \setminus\left\{  k\right\}  }A\right)  .
\label{pf.lem.36.det-hess.1}%
\end{equation}


Now, let us fix $k\in\left[  n\right]  $. We shall prove that%
\[
\det\left(  \operatorname*{sub}\nolimits_{\left[  n\right]  \setminus\left\{
1\right\}  }^{\left[  n\right]  \setminus\left\{  k\right\}  }A\right)
=\det\left(  a_{i+k,j+k}\right)  _{i,j\in\left[  n-k\right]  }.
\]


Indeed, from $A=\left(  a_{i,j}\right)  _{i,j\in\left[  n\right]  }$, we
obtain%
\begin{align*}
&  \operatorname*{sub}\nolimits_{\left[  n\right]  \setminus\left\{
1\right\}  }^{\left[  n\right]  \setminus\left\{  k\right\}  }A\\
&  =\left(
\begin{array}
[c]{cccccccc}%
a_{2,1} & a_{2,2} & \cdots & a_{2,k-1} & a_{2,k+1} & a_{2,k+2} & \cdots &
a_{2,n}\\
a_{3,1} & a_{3,2} & \cdots & a_{3,k-1} & a_{3,k+1} & a_{3,k+2} & \cdots &
a_{3,n}\\
\vdots & \vdots & \ddots & \vdots & \vdots & \vdots & \ddots & \vdots\\
a_{k,1} & a_{k,2} & \cdots & a_{k,k-1} & a_{k,k+1} & a_{k,k+2} & \cdots &
a_{k,n}\\
a_{k+1,1} & a_{k+1,2} & \cdots & a_{k+1,k-1} & a_{k+1,k+1} & a_{k+1,k+2} &
\cdots & a_{k+1,n}\\
a_{k+2,1} & a_{k+2,2} & \cdots & a_{k+2,k-1} & a_{k+2,k+1} & a_{k+2,k+2} &
\cdots & a_{k+2,n}\\
\vdots & \vdots & \ddots & \vdots & \vdots & \vdots & \ddots & \vdots\\
a_{n,1} & a_{n,2} & \cdots & a_{n,k-1} & a_{n,k+1} & a_{n,k+2} & \cdots &
a_{n,n}%
\end{array}
\right) \\
&  =\left(
\begin{array}
[c]{cccccccc}%
1 & a_{2,2} & \cdots & a_{2,k-1} & a_{2,k+1} & a_{2,k+2} & \cdots & a_{2,n}\\
0 & 1 & \cdots & a_{3,k-1} & a_{3,k+1} & a_{3,k+2} & \cdots & a_{3,n}\\
\vdots & \vdots & \ddots & \vdots & \vdots & \vdots & \ddots & \vdots\\
0 & 0 & \cdots & 1 & a_{k,k+1} & a_{k,k+2} & \cdots & a_{k,n}\\
0 & 0 & \cdots & 0 & a_{k+1,k+1} & a_{k+1,k+2} & \cdots & a_{k+1,n}\\
0 & 0 & \cdots & 0 & a_{k+2,k+1} & a_{k+2,k+2} & \cdots & a_{k+2,n}\\
\vdots & \vdots & \ddots & \vdots & \vdots & \vdots & \ddots & \vdots\\
0 & 0 & \cdots & 0 & a_{n,k+1} & a_{n,k+2} & \cdots & a_{n,n}%
\end{array}
\right)
\end{align*}
(here, we have used (\ref{eq.lem.36.det-hess.ass.0}) and
(\ref{eq.lem.36.det-hess.ass.1}) to identify several entries in the first
$k-1$ columns of our matrix as being $0$ or $1$). In other words,
$\operatorname*{sub}\nolimits_{\left[  n\right]  \setminus\left\{  1\right\}
}^{\left[  n\right]  \setminus\left\{  k\right\}  }A$ can be written as block
matrix%
\[
\operatorname*{sub}\nolimits_{\left[  n\right]  \setminus\left\{  1\right\}
}^{\left[  n\right]  \setminus\left\{  k\right\}  }A=\left(
\begin{array}
[c]{cc}%
P & Q\\
0_{\left(  n-k\right)  \times\left(  k-1\right)  } & R
\end{array}
\right)  ,
\]
where

\begin{itemize}
\item the matrix $P=\left(
\begin{array}
[c]{cccc}%
1 & a_{2,2} & \cdots & a_{2,k-1}\\
0 & 1 & \cdots & a_{3,k-1}\\
\vdots & \vdots & \ddots & \vdots\\
0 & 0 & \cdots & 1
\end{array}
\right)  $ is an upper unitriangular matrix and thus has determinant $\det
P=1$;

\item the matrix $Q=\left(
\begin{array}
[c]{cccc}%
a_{2,k+1} & a_{2,k+2} & \cdots & a_{2,n}\\
a_{3,k+1} & a_{3,k+2} & \cdots & a_{3,n}\\
\vdots & \vdots & \ddots & \vdots\\
a_{k,k+1} & a_{k,k+2} & \cdots & a_{k,n}%
\end{array}
\right)  $ is a $\left(  k-1\right)  \times\left(  n-k\right)  $-matrix whose
entries don't interest us;

\item the matrix $0_{\left(  n-k\right)  \times\left(  k-1\right)  }$ is a
zero matrix of size $\left(  n-k\right)  \times\left(  k-1\right)  $;

\item the matrix $R=\left(
\begin{array}
[c]{cccc}%
a_{k+1,k+1} & a_{k+1,k+2} & \cdots & a_{k+1,n}\\
a_{k+2,k+1} & a_{k+2,k+2} & \cdots & a_{k+2,n}\\
\vdots & \vdots & \ddots & \vdots\\
a_{n,k+1} & a_{n,k+2} & \cdots & a_{n,n}%
\end{array}
\right)  $ is precisely the $\left(  n-k\right)  \times\left(  n-k\right)
$-matrix $\left(  a_{i+k,j+k}\right)  _{i,j\in\left[  n-k\right]  }$.
\end{itemize}

Hence,%
\[
\det\left(  \operatorname*{sub}\nolimits_{\left[  n\right]  \setminus\left\{
1\right\}  }^{\left[  n\right]  \setminus\left\{  k\right\}  }A\right)
=\det\left(
\begin{array}
[c]{cc}%
P & Q\\
0_{\left(  n-k\right)  \times\left(  k-1\right)  } & R
\end{array}
\right)  =\det P\cdot\det R
\]
by the well-known formula for the determinant of a block-triangular matrix
(see, e.g., \cite[Exercise 6.29]{detnotes}). Thus,%
\begin{align}
\det\left(  \operatorname*{sub}\nolimits_{\left[  n\right]  \setminus\left\{
1\right\}  }^{\left[  n\right]  \setminus\left\{  k\right\}  }A\right)   &
=\underbrace{\det P}_{=1}\cdot\det\underbrace{R}_{=\left(  a_{i+k,j+k}\right)
_{i,j\in\left[  n-k\right]  }}\nonumber\\
&  =\det\left(  a_{i+k,j+k}\right)  _{i,j\in\left[  n-k\right]  }.
\label{pf.lem.36.det-hess.7}%
\end{align}


Now forget that we fixed $k$. We thus have proved (\ref{pf.lem.36.det-hess.7})
for each $k\in\left[  n\right]  $. Hence, (\ref{pf.lem.36.det-hess.1}) becomes%
\begin{align*}
\det A  &  =\sum_{k=1}^{n}\underbrace{\left(  -1\right)  ^{1+k}}_{=\left(
-1\right)  ^{k+1}}a_{1,k}\underbrace{\det\left(  \operatorname*{sub}%
\nolimits_{\left[  n\right]  \setminus\left\{  1\right\}  }^{\left[  n\right]
\setminus\left\{  k\right\}  }A\right)  }_{\substack{=\det\left(
a_{i+k,j+k}\right)  _{i,j\in\left[  n-k\right]  }\\\text{(by
(\ref{pf.lem.36.det-hess.7}))}}}\\
&  =\sum_{k=1}^{n}\left(  -1\right)  ^{k+1}a_{1,k}\det\left(  a_{i+k,j+k}%
\right)  _{i,j\in\left[  n-k\right]  }.
\end{align*}
This proves Lemma \ref{lem.36.det-hess}.
\end{proof}

Now, we can prove the equality (1) from the proof of (9.3):

Indeed, the $r\times r$-matrix $\left(  \varphi^{1-j}\left(  h_{1-i+j}\right)
\right)  _{i,j\in\left[  r\right]  }$ has the property that%
\[
\varphi^{1-j}\left(  h_{1-i+j}\right)  =1\ \ \ \ \ \ \ \ \ \ \text{whenever
}i=j+1
\]
(because $i=j+1$ entails $1-i+j=1-\left(  j+1\right)  +j=0$ and thus
$h_{1-i+j}=h_{0}=1$, so that $\varphi^{1-j}\left(  h_{1-i+j}\right)
=\varphi^{1-j}\left(  1\right)  =1$). Furthermore, it has the property that%
\[
\varphi^{1-j}\left(  h_{1-i+j}\right)  =0\ \ \ \ \ \ \ \ \ \ \text{whenever
}i>j+1
\]
(since $i>j+1$ entails $1-i+j<1-\left(  j+1\right)  +j=0$ and thus
$h_{1-i+j}=0$, so that $\varphi^{1-j}\left(  h_{1-i+j}\right)  =\varphi
^{1-j}\left(  0\right)  =0$). Hence, Lemma \ref{lem.36.det-hess} (applied to
$n=r$ and $a_{i,j}=\varphi^{1-j}\left(  h_{1-i+j}\right)  $ and $A=\left(
\varphi^{1-j}\left(  h_{1-i+j}\right)  \right)  _{i,j\in\left[  r\right]  }$)
yields%
\begin{align*}
&  \det\left(  \varphi^{1-j}\left(  h_{1-i+j}\right)  \right)  _{i,j\in\left[
r\right]  }\\
&  =\sum_{k=1}^{r}\left(  -1\right)  ^{k+1}\varphi^{1-k}\underbrace{\left(
h_{1-1+k}\right)  }_{=h_{k}}\det\left(  \underbrace{\varphi^{1-\left(
j+k\right)  }}_{\substack{=\varphi^{-k+\left(  1-j\right)  }\\=\varphi
^{-k}\varphi^{1-j}}}\left(  \underbrace{h_{1-\left(  i+k\right)  +\left(
j+k\right)  }}_{=h_{1-i+j}}\right)  \right)  _{i,j\in\left[  r-k\right]  }\\
&  =\sum_{k=1}^{r}\left(  -1\right)  ^{k+1}\varphi^{1-k}\left(  h_{k}\right)
\underbrace{\det\left(  \varphi^{-k}\varphi^{1-j}\left(  h_{1-i+j}\right)
\right)  _{i,j\in\left[  r-k\right]  }}_{\substack{=\varphi^{-k}\det\left(
\varphi^{1-j}\left(  h_{1-i+j}\right)  \right)  _{i,j\in\left[  r-k\right]
}\\\text{(since }\varphi^{-k}\text{ is an algebra morphism)}}}\\
&  =\sum_{k=1}^{r}\left(  -1\right)  ^{k+1}\varphi^{1-k}\left(  h_{k}\right)
\varphi^{-k}\underbrace{\det\left(  \varphi^{1-j}\left(  h_{1-i+j}\right)
\right)  _{i,j\in\left[  r-k\right]  }}_{\substack{=e_{r-k}\\\text{(by the
definition of }e_{r-k}=s_{\left(  1^{r-k}\right)  }\text{)}}}\\
&  =\sum_{k=1}^{r}\left(  -1\right)  ^{k+1}\varphi^{1-k}\left(  h_{k}\right)
\varphi^{-k}\left(  e_{r-k}\right) \\
&  =\sum_{j=1}^{r}\left(  -1\right)  ^{j+1}\varphi^{1-j}\left(  h_{j}\right)
\varphi^{-j}\left(  e_{r-j}\right)  .
\end{align*}
Comparing this with%
\[
\det\left(  \varphi^{1-j}\left(  h_{1-i+j}\right)  \right)  _{i,j\in\left[
r\right]  }=e_{r}\ \ \ \ \ \ \ \ \ \ \left(  \text{by the definition of }%
e_{r}=s_{\left(  1^{r}\right)  }\right)  ,
\]
we find%
\begin{equation}
e_{r}=\sum_{j=1}^{r}\left(  -1\right)  ^{j+1}\varphi^{1-j}\left(
h_{j}\right)  \varphi^{-j}\left(  e_{r-j}\right)  . \label{p36.9.3.5}%
\end{equation}
Thus,%
\begin{align*}
&  \sum_{j=0}^{r}\left(  -1\right)  ^{j}\varphi^{1-j}\left(  h_{j}\right)
\varphi^{-j}\left(  e_{r-j}\right) \\
&  =\underbrace{\left(  -1\right)  ^{0}}_{=1}\varphi^{1-0}\left(
\underbrace{h_{0}}_{=1}\right)  \underbrace{\varphi^{-0}}_{=\operatorname*{id}%
}\left(  \underbrace{e_{r-0}}_{=e_{r}}\right)  +\sum_{j=1}^{r}%
\underbrace{\left(  -1\right)  ^{j}}_{=-\left(  -1\right)  ^{j+1}}%
\varphi^{1-j}\left(  h_{j}\right)  \varphi^{-j}\left(  e_{r-j}\right) \\
&  =\underbrace{\varphi^{1-0}\left(  1\right)  }_{=1}e_{r}-\underbrace{\sum
_{j=1}^{r}\left(  -1\right)  ^{j+1}\varphi^{1-j}\left(  h_{j}\right)
\varphi^{-j}\left(  e_{r-j}\right)  }_{\substack{=e_{r}\\\text{(by
(\ref{p36.9.3.5}))}}}=e_{r}-e_{r}=0,
\end{align*}
and so the equality (1) is proved.

\item \textbf{page 36:} The claim \textquotedblleft The function $s_{\theta
}=s_{\lambda/\mu}$ depends not only on the skew shape $\theta$ but also on its
location in the lattice plane\textquotedblright\ is a bit confusing, even
though it is correct: In fact, it is already non-obvious that $s_{\theta
}=s_{\lambda/\mu}$ is uniquely determined by the skew shape $\theta$ rather
than the pair of partitions $\left(  \lambda,\mu\right)  $. After all,
$s_{\lambda/\mu}$ was defined in terms of the pair of partitions $\left(
\lambda,\mu\right)  $, but different pairs of partitions $\left(  \lambda
,\mu\right)  $ will often lead to the same skew shape $\lambda-\mu$. For
instance, the skew partitions $\left(  3,2,1\right)  /\left(  2,2\right)  $
and $\left(  3,1,1\right)  /\left(  2,1\right)  $ lead to the same skew shape
$\left(  3,2,1\right)  -\left(  2,2\right)  =\left(  3,1,1\right)  -\left(
2,1\right)  $. Why are the corresponding \textquotedblleft
functions\textquotedblright\ $s_{\left(  3,2,1\right)  /\left(  2,2\right)  }$
and $s_{\left(  3,1,1\right)  /\left(  2,1\right)  }$ equal?

Let me answer this question in detail. We want to prove the following lemma:

\begin{lemmain}
\label{lem.9v.shape}Let $\lambda,\mu,\alpha,\beta$ be four partitions such
that $\mu\subseteq\lambda$ and $\beta\subseteq\alpha$. Assume that
$\lambda-\mu=\alpha-\beta$ as sets of cells (i.e., the cells that belong to
$\lambda-\mu$ are precisely the cells that belong to $\alpha-\beta$). Then,%
\[
s_{\lambda/\mu}=s_{\alpha/\beta}.
\]

\end{lemmain}

Before we prove this, we need four other lemmas:

\begin{lemmain}
\label{lem.9v.shape1}Let $\lambda,\mu,\alpha,\beta$ be four partitions such
that $\mu\subseteq\lambda$ and $\beta\subseteq\alpha$. Let $n\in\mathbb{N}$ be
such that all four lengths $\ell\left(  \lambda\right)  ,\ell\left(
\mu\right)  ,\ell\left(  \alpha\right)  ,\ell\left(  \beta\right)  $ are $\leq
n$. Assume that $\lambda-\mu=\alpha-\beta$ as sets of cells (i.e., the cells
that belong to $\lambda-\mu$ are precisely the cells that belong to
$\alpha-\beta$), but we don't have $\lambda=\alpha$ and $\mu=\beta$. Then,
there exists some $k\in\left[  n\right]  $ satisfying $\lambda_{k}=\mu_{k}%
\neq\alpha_{k}=\beta_{k}$.
\end{lemmain}

\begin{proof}
[Proof of Lemma \ref{lem.9v.shape1}.]We don't have $\lambda=\alpha$ and
$\mu=\beta$. Hence, we have $\lambda\neq\alpha$ or $\mu\neq\beta$. In other
words, there exists some $i\geq1$ such that $\lambda_{i}\neq\alpha_{i}$ or
$\mu_{i}\neq\beta_{i}$. Consider this $i$. If we had $i>n$, then we would have
$\lambda_{i}=0$ (since we assumed that $\ell\left(  \lambda\right)  \leq n<i$)
and similarly $\alpha_{i}=0$ and $\mu_{i}=0$ and $\beta_{i}=0$; but this would
contradict the fact that $\lambda_{i}\neq\alpha_{i}$ or $\mu_{i}\neq\beta_{i}%
$. Hence, we cannot have $i>n$. Thus, $i\leq n$, so that $i\in\left[
n\right]  $.

Note that $\lambda_{i}\neq\alpha_{i}$ or $\mu_{i}\neq\beta_{i}$. In other
words, $\lambda_{i}\neq\alpha_{i}$ or $\mu_{i}+1\neq\beta_{i}+1$ (since
$\mu_{i}\neq\beta_{i}$ is equivalent to $\mu_{i}+1\neq\beta_{i}+1$).

The $i$-th row of the skew shape $\lambda-\mu$ has cells in the columns
$\mu_{i}+1,\mu_{i}+2,\ldots,\lambda_{i}$ but in no other columns. That is,%
\[
\left\{  j\in\mathbb{Z}\ \mid\ \left(  i,j\right)  \in\lambda-\mu\right\}
=\left[  \mu_{i}+1,\ \lambda_{i}\right]  .
\]
Similarly,%
\[
\left\{  j\in\mathbb{Z}\ \mid\ \left(  i,j\right)  \in\alpha-\beta\right\}
=\left[  \beta_{i}+1,\ \alpha_{i}\right]  .
\]
The left hand sides of these two equalities are equal (since $\lambda
-\mu=\alpha-\beta$). Hence, so are the right hand sides. That is, we have
$\left[  \mu_{i}+1,\ \lambda_{i}\right]  =\left[  \beta_{i}+1,\ \alpha
_{i}\right]  $. If this interval $\left[  \mu_{i}+1,\ \lambda_{i}\right]
=\left[  \beta_{i}+1,\ \alpha_{i}\right]  $ was nonempty, then it would have a
largest element and a smallest element, and thus we would obtain%
\[
\lambda_{i}=\max\underbrace{\left[  \mu_{i}+1,\ \lambda_{i}\right]
}_{=\left[  \beta_{i}+1,\ \alpha_{i}\right]  }=\max\left[  \beta
_{i}+1,\ \alpha_{i}\right]  =\alpha_{i}%
\]
and%
\[
\mu_{i}+1=\min\underbrace{\left[  \mu_{i}+1,\ \lambda_{i}\right]  }_{=\left[
\beta_{i}+1,\ \alpha_{i}\right]  }=\min\left[  \beta_{i}+1,\ \alpha
_{i}\right]  =\beta_{i}+1;
\]
but these two equalities would contradict the fact that $\lambda_{i}\neq
\alpha_{i}$ or $\mu_{i}+1\neq\beta_{i}+1$. Hence, this interval $\left[
\mu_{i}+1,\ \lambda_{i}\right]  =\left[  \beta_{i}+1,\ \alpha_{i}\right]  $
must be empty. But $\mu\subseteq\lambda$ entails $\mu_{i}\leq\lambda_{i}$.
Hence, $\mu_{i}=\lambda_{i}$ (since the interval $\left[  \mu_{i}%
+1,\ \lambda_{i}\right]  $ is empty), so that $\lambda_{i}=\mu_{i}$.
Furthermore, $\beta\subseteq\alpha$ entails $\beta_{i}\leq\alpha_{i}$. Hence,
$\beta_{i}=\alpha_{i}$ (since the interval $\left[  \beta_{i}+1,\ \alpha
_{i}\right]  $ is empty), so that $\alpha_{i}=\beta_{i}$. If we had $\mu
_{i}=\alpha_{i}$, then we would thus have $\lambda_{i}=\mu_{i}=\alpha
_{i}=\beta_{i}$, which would contradict the fact that $\lambda_{i}\neq
\alpha_{i}$ or $\mu_{i}\neq\beta_{i}$. Thus, we cannot have $\mu_{i}%
=\alpha_{i}$. Hence, we have $\mu_{i}\neq\alpha_{i}$. Thus, $\lambda_{i}%
=\mu_{i}\neq\alpha_{i}=\beta_{i}$. Recall also that $i\in\left[  n\right]  $.

Thus we conclude that there exists some $k\in\left[  n\right]  $ satisfying
$\lambda_{k}=\mu_{k}\neq\alpha_{k}=\beta_{k}$ (namely, $k=i$). This proves
Lemma \ref{lem.9v.shape1}.
\end{proof}

\begin{lemmain}
\label{lem.9v.shape2}Let $\lambda$ and $\mu$ be two partitions. Let
$n\in\mathbb{N}$ be such that $\ell\left(  \lambda\right)  \leq n$ and
$\ell\left(  \mu\right)  \leq n$. Let $k\in\left[  0,n\right]  $ be such that
$k=0$ or $\mu_{k}\geq\lambda_{k+1}$. Define the four partitions%
\begin{align*}
\lambda^{\leq k}  &  =\left(  \lambda_{1},\lambda_{2},\ldots,\lambda
_{k}\right)  ,\ \ \ \ \ \ \ \ \ \ \lambda^{>k}=\left(  \lambda_{k+1}%
,\lambda_{k+2},\ldots,\lambda_{n}\right)  ,\\
\mu^{\leq k}  &  =\left(  \mu_{1},\mu_{2},\ldots,\mu_{k}\right)
,\ \ \ \ \ \ \ \ \ \ \mu^{>k}=\left(  \mu_{k+1},\mu_{k+2},\ldots,\mu
_{n}\right)  .
\end{align*}
Then,%
\[
s_{\lambda/\mu}=s_{\lambda^{\leq k}/\mu^{\leq k}}\cdot\varphi^{-k}\left(
s_{\lambda^{>k}/\mu^{>k}}\right)  .
\]

\end{lemmain}

\begin{proof}
[Proof of Lemma \ref{lem.9v.shape2}.]The definition of $s_{\lambda/\mu}$ shows
that%
\begin{equation}
s_{\lambda/\mu}=\det\left(  \varphi^{\mu_{j}-j+1}h_{\lambda_{i}-\mu_{j}%
-i+j}\right)  _{1\leq i,j\leq n}. \label{pf.lem.9v.shape2.s=1}%
\end{equation}
However, all $i\in\left[  k+1,n\right]  $ and $j\in\left[  k\right]  $ satisfy%
\begin{equation}
\lambda_{i}-\mu_{j}-i+j<0. \label{pf.lem.9v.shape2.=0}%
\end{equation}


(\textit{Proof:} Let $i\in\left[  k+1,n\right]  $ and $j\in\left[  k\right]  $
be arbitrary. From $j\in\left[  k\right]  $, we obtain $1\leq j\leq k$, so
that $k\geq1$. Hence, $k=0$ cannot happen. Thus, we must have $\mu_{k}%
\geq\lambda_{k+1}$ (since we have assumed that $k=0$ or $\mu_{k}\geq
\lambda_{k+1}$).

From $\lambda_{1}\geq\lambda_{2}\geq\lambda_{3}\geq\cdots$ (since $\lambda$ is
a partition) and $k+1\leq i$ (since $i\in\left[  k+1,n\right]  $), we obtain
$\lambda_{k+1}\geq\lambda_{i}$. From $\mu_{1}\geq\mu_{2}\geq\mu_{3}\geq\cdots$
(since $\mu$ is a partition) and $j\leq k$ (since $j\in\left[  k\right]  $),
we obtain $\mu_{j}\geq\mu_{k}$. Thus, $\mu_{j}\geq\mu_{k}\geq\lambda_{k+1}%
\geq\lambda_{i}$. Hence, $\lambda_{i}\leq\mu_{j}$ and therefore $\lambda
_{i}-\mu_{j}\leq0$. Adding the strict inequality $-i+j<0$ (which follows from
$j\leq k<k+1\leq i$) to this, we obtain $\lambda_{i}-\mu_{j}-i+j<0$. This
proves (\ref{pf.lem.9v.shape2.=0}).)

Now, all $i\in\left[  k+1,n\right]  $ and $j\in\left[  k\right]  $ satisfy
$\lambda_{i}-\mu_{j}-i+j<0$ (by (\ref{pf.lem.9v.shape2.=0})) and therefore
$h_{\lambda_{i}-\mu_{j}-i+j}=0$, so that $\varphi^{\mu_{j}-j+1}h_{\lambda
_{i}-\mu_{j}-i+j}=\varphi^{\mu_{j}-j+1}0=0$. In other words, the matrix
$\left(  \varphi^{\mu_{j}-j+1}h_{\lambda_{i}-\mu_{j}-i+j}\right)  _{1\leq
i,j\leq n}$ has the property that its $\left(  i,j\right)  $-th entries are
$0$ for all $i\in\left[  k+1,n\right]  $ and $j\in\left[  k\right]  $. Hence,
this matrix can be written as a block matrix of the form%
\[
\left(  \varphi^{\mu_{j}-j+1}h_{\lambda_{i}-\mu_{j}-i+j}\right)  _{1\leq
i,j\leq n}=\left(
\begin{array}
[c]{cc}%
A & B\\
0_{\left(  n-k\right)  \times k} & D
\end{array}
\right)  ,
\]
where

\begin{itemize}
\item the matrix $A$ is the $k\times k$-matrix
\[
\left(  \varphi^{\mu_{j}-j+1}h_{\lambda_{i}-\mu_{j}-i+j}\right)  _{1\leq
i,j\leq k};
\]


\item the matrix $B$ is the $k\times\left(  n-k\right)  $-matrix
\[
\left(  \varphi^{\mu_{j+k}-\left(  j+k\right)  +1}h_{\lambda_{i}-\mu_{\left(
j+k\right)  }-i+\left(  j+k\right)  }\right)  _{i\in\left[  k\right]
,\ j\in\left[  n-k\right]  };
\]


\item the matrix $0_{\left(  n-k\right)  \times k}$ is a zero matrix of size
$\left(  n-k\right)  \times k$;

\item the matrix $D$ is the $\left(  n-k\right)  \times\left(  n-k\right)
$-matrix
\[
\left(  \varphi^{\mu_{j+k}-\left(  j+k\right)  +1}h_{\lambda_{i+k}-\mu
_{j+k}-\left(  i+k\right)  +\left(  j+k\right)  }\right)  _{1\leq i,j\leq
n-k}.
\]

\end{itemize}

Therefore,
\begin{align}
\det\left(  \varphi^{\mu_{j}-j+1}h_{\lambda_{i}-\mu_{j}-i+j}\right)  _{1\leq
i,j\leq n}  &  =\det\left(
\begin{array}
[c]{cc}%
A & B\\
0_{\left(  n-k\right)  \times k} & D
\end{array}
\right) \nonumber\\
&  =\det A\cdot\det D \label{pf.lem.9v.shape2.detfact}%
\end{align}
by the well-known formula for the determinant of a block-triangular matrix
(see, e.g., \cite[Exercise 6.29]{detnotes}). However, from $A=\left(
\varphi^{\mu_{j}-j+1}h_{\lambda_{i}-\mu_{j}-i+j}\right)  _{1\leq i,j\leq k}$,
we obtain%
\begin{align}
\det A  &  =\det\left(  \varphi^{\mu_{j}-j+1}h_{\lambda_{i}-\mu_{j}%
-i+j}\right)  _{1\leq i,j\leq k}\nonumber\\
&  =s_{\lambda^{\leq k}/\mu^{\leq k}} \label{pf.lem.9v.shape2.A}%
\end{align}
(by the definition of $s_{\lambda^{\leq k}/\mu^{\leq k}}$, since
$\lambda^{\leq k}=\left(  \lambda_{1},\lambda_{2},\ldots,\lambda_{k}\right)  $
and $\mu^{\leq k}=\left(  \mu_{1},\mu_{2},\ldots,\mu_{k}\right)  $).
Furthermore, for any $i,j\in\left[  n-k\right]  $, we have%
\begin{align}
&  \underbrace{\varphi^{\mu_{j+k}-\left(  j+k\right)  +1}}_{\substack{=\varphi
^{-k+\left(  \mu_{j+k}-j+1\right)  }\\\text{(since }\mu_{j+k}-\left(
j+k\right)  +1=-k+\left(  \mu_{j+k}-j+1\right)  \text{)}}%
}\underbrace{h_{\lambda_{i+k}-\mu_{j+k}-\left(  i+k\right)  +\left(
j+k\right)  }}_{=h_{\lambda_{i+k}-\mu_{j+k}-i+j}}\nonumber\\
&  =\underbrace{\varphi^{-k+\left(  \mu_{j+k}-j+1\right)  }}_{=\varphi
^{-k}\varphi^{\mu_{j+k}-j+1}}h_{\lambda_{i+k}-\mu_{j+k}-i+j}\nonumber\\
&  =\varphi^{-k}\left(  \varphi^{\mu_{j+k}-j+1}h_{\lambda_{i+k}-\mu_{j+k}%
-i+j}\right)  . \label{pf.lem.9v.shape2.4}%
\end{align}
Now,
\begin{align*}
D  &  =\left(  \varphi^{\mu_{j+k}-\left(  j+k\right)  +1}h_{\lambda_{i+k}%
-\mu_{j+k}-\left(  i+k\right)  +\left(  j+k\right)  }\right)  _{1\leq i,j\leq
n-k}\\
&  =\left(  \varphi^{-k}\left(  \varphi^{\mu_{j+k}-j+1}h_{\lambda_{i+k}%
-\mu_{j+k}-i+j}\right)  \right)  _{1\leq i,j\leq n-k}%
\ \ \ \ \ \ \ \ \ \ \left(  \text{by (\ref{pf.lem.9v.shape2.4})}\right)  ,
\end{align*}
and thus%
\begin{align}
\det D  &  =\det\left(  \varphi^{-k}\left(  \varphi^{\mu_{j+k}-j+1}%
h_{\lambda_{i+k}-\mu_{j+k}-i+j}\right)  \right)  _{1\leq i,j\leq
n-k}\nonumber\\
&  =\varphi^{-k}\left(  \det\left(  \varphi^{\mu_{j+k}-j+1}h_{\lambda
_{i+k}-\mu_{j+k}-i+j}\right)  _{1\leq i,j\leq n-k}\right)
\label{pf.lem.9v.shape2.6}%
\end{align}
(since $\varphi^{-k}$ is a ring morphism). However, the definition of
$s_{\lambda^{>k}/\mu^{>k}}$ shows that%
\[
s_{\lambda^{>k}/\mu^{>k}}=\det\left(  \varphi^{\mu_{j+k}-j+1}h_{\lambda
_{i+k}-\mu_{j+k}-i+j}\right)  _{1\leq i,j\leq n-k}%
\]
(since $\lambda^{>k}=\left(  \lambda_{k+1},\lambda_{k+2},\ldots,\lambda
_{n}\right)  $ and $\mu^{>k}=\left(  \mu_{k+1},\mu_{k+2},\ldots,\mu
_{n}\right)  $). In view of this, we can rewrite (\ref{pf.lem.9v.shape2.6}) as%
\begin{equation}
\det D=\varphi^{-k}\left(  s_{\lambda^{>k}/\mu^{>k}}\right)  .
\label{pf.lem.9v.shape2.D}%
\end{equation}


Now, (\ref{pf.lem.9v.shape2.s=1}) becomes%
\begin{align*}
s_{\lambda/\mu}  &  =\det\left(  \varphi^{\mu_{j}-j+1}h_{\lambda_{i}-\mu
_{j}-i+j}\right)  _{1\leq i,j\leq n}\\
&  =\underbrace{\det A}_{\substack{=s_{\lambda^{\leq k}/\mu^{\leq k}%
}\\\text{(by (\ref{pf.lem.9v.shape2.A}))}}}\cdot\underbrace{\det
D}_{\substack{=\varphi^{-k}\left(  s_{\lambda^{>k}/\mu^{>k}}\right)
\\\text{(by (\ref{pf.lem.9v.shape2.D}))}}}\ \ \ \ \ \ \ \ \ \ \left(  \text{by
(\ref{pf.lem.9v.shape2.detfact})}\right) \\
&  =s_{\lambda^{\leq k}/\mu^{\leq k}}\cdot\varphi^{-k}\left(  s_{\lambda
^{>k}/\mu^{>k}}\right)  .
\end{align*}
This proves Lemma \ref{lem.9v.shape2}.
\end{proof}

\begin{lemmain}
\label{lem.9v.shape3}Let $\lambda$ and $\mu$ be two partitions. Let $n$ be a
positive integer such that $\ell\left(  \lambda\right)  \leq n$ and
$\ell\left(  \mu\right)  \leq n$ and $\lambda_{n}=\mu_{n}$. Define the two
partitions%
\[
\overline{\lambda}=\left(  \lambda_{1},\lambda_{2},\ldots,\lambda
_{n-1}\right)  \ \ \ \ \ \ \ \ \ \ \text{and}\ \ \ \ \ \ \ \ \ \ \overline
{\mu}=\left(  \mu_{1},\mu_{2},\ldots,\mu_{n-1}\right)  .
\]
Then,%
\[
s_{\lambda/\mu}=s_{\overline{\lambda}/\overline{\mu}}.
\]

\end{lemmain}

\begin{proof}
[Proof of Lemma \ref{lem.9v.shape3}.]Define the four partitions%
\begin{align*}
\lambda^{\leq n-1}  &  =\left(  \lambda_{1},\lambda_{2},\ldots,\lambda
_{n-1}\right)  ,\ \ \ \ \ \ \ \ \ \ \lambda^{>n-1}=\left(  \lambda
_{n-1+1},\lambda_{n-1+2},\ldots,\lambda_{n}\right)  ,\\
\mu^{\leq n-1}  &  =\left(  \mu_{1},\mu_{2},\ldots,\mu_{n-1}\right)
,\ \ \ \ \ \ \ \ \ \ \mu^{>n-1}=\left(  \mu_{n-1+1},\mu_{n-1+2},\ldots,\mu
_{n}\right)  .
\end{align*}
Of course, $\lambda^{\leq n-1}=\left(  \lambda_{1},\lambda_{2},\ldots
,\lambda_{n-1}\right)  =\overline{\lambda}$ and similarly $\mu^{\leq
n-1}=\overline{\mu}$. Moreover, $\lambda^{>n-1}=\left(  \lambda_{n-1+1}%
,\lambda_{n-1+2},\ldots,\lambda_{n}\right)  =\left(  \lambda_{n},\lambda
_{n+1},\ldots,\lambda_{n}\right)  =\left(  \lambda_{n}\right)  $ and similarly
$\mu^{>n-1}=\left(  \mu_{n}\right)  $.

If $n-1>0$, then $\mu_{n-1}\geq\lambda_{\left(  n-1\right)  +1}$ (because
$\mu$ is a partition, so that $\mu_{1}\geq\mu_{2}\geq\mu_{3}\geq\cdots$, and
therefore $\mu_{n-1}\geq\mu_{n}=\lambda_{n}=\lambda_{\left(  n-1\right)  +1}%
$). In other words, we have $n-1=0$ or $\mu_{n-1}\geq\lambda_{\left(
n-1\right)  +1}$. Hence, Lemma \ref{lem.9v.shape2} (applied to $k=n-1$) yields%
\begin{equation}
s_{\lambda/\mu}=s_{\lambda^{\leq n-1}/\mu^{\leq n-1}}\cdot\varphi^{-\left(
n-1\right)  }\left(  s_{\lambda^{>n-1}/\mu^{>n-1}}\right)  .
\label{pf.lem.9v.shape3.4}%
\end{equation}


But the definition of $s_{\lambda^{>n-1}/\mu^{>n-1}}$ yields%
\begin{align*}
s_{\lambda^{>n-1}/\mu^{>n-1}} &  =\det\left(  \varphi^{\mu_{n}-j+1}%
h_{\lambda_{n}-\mu_{n}-i+j}\right)  _{1\leq i,j\leq1}\\
&  \ \ \ \ \ \ \ \ \ \ \ \ \ \ \ \ \ \ \ \ \left(  \text{since }\lambda
^{>n-1}=\left(  \lambda_{n}\right)  \text{ and }\mu^{>n-1}=\left(  \mu
_{n}\right)  \right)  \\
&  =\varphi^{\mu_{n}-1+1}h_{\lambda_{n}-\mu_{n}-1+1}\\
&  =\varphi^{\mu_{n}-1+1}\underbrace{h_{0}}_{=1}\ \ \ \ \ \ \ \ \ \ \left(
\text{since }\underbrace{\lambda_{n}}_{=\mu_{n}}-\,\mu_{n}-1+1=\mu_{n}-\mu
_{n}-1+1=0\right)  \\
&  =\varphi^{\mu_{n}-1+1}1=1.
\end{align*}
Thus, $\varphi^{-\left(  n-1\right)  }\left(  s_{\lambda^{>n-1}/\mu^{>n-1}%
}\right)  =\varphi^{-\left(  n-1\right)  }\left(  1\right)  =1$ as well.
Hence, (\ref{pf.lem.9v.shape3.4}) becomes%
\[
s_{\lambda/\mu}=s_{\lambda^{\leq n-1}/\mu^{\leq n-1}}\cdot\underbrace{\varphi
^{-\left(  n-1\right)  }\left(  s_{\lambda^{>n-1}/\mu^{>n-1}}\right)  }%
_{=1}=s_{\lambda^{\leq n-1}/\mu^{\leq n-1}}=s_{\overline{\lambda}%
/\overline{\mu}}%
\]
(since $\lambda^{\leq n-1}=\overline{\lambda}$ and $\mu^{\leq n-1}%
=\overline{\mu}$). This proves Lemma \ref{lem.9v.shape3}.
\end{proof}

\begin{lemmain}
\label{lem.9v.shape4}Let $\lambda$ and $\mu$ be two partitions. Let
$n\in\mathbb{N}$ be such that $\ell\left(  \lambda\right)  \leq n$ and
$\ell\left(  \mu\right)  \leq n$. Let $k\in\left[  n\right]  $ be such that
$\lambda_{k}=\mu_{k}$. Define the four partitions%
\begin{align*}
\lambda^{\leq k-1}  &  =\left(  \lambda_{1},\lambda_{2},\ldots,\lambda
_{k-1}\right)  ,\ \ \ \ \ \ \ \ \ \ \lambda^{>k}=\left(  \lambda_{k+1}%
,\lambda_{k+2},\ldots,\lambda_{n}\right)  ,\\
\mu^{\leq k-1}  &  =\left(  \mu_{1},\mu_{2},\ldots,\mu_{k-1}\right)
,\ \ \ \ \ \ \ \ \ \ \mu^{>k}=\left(  \mu_{k+1},\mu_{k+2},\ldots,\mu
_{n}\right)  .
\end{align*}
Then,%
\[
s_{\lambda/\mu}=s_{\lambda^{\leq k-1}/\mu^{\leq k-1}}\cdot\varphi^{-k}\left(
s_{\lambda^{>k}/\mu^{>k}}\right)  .
\]

\end{lemmain}

\begin{proof}
[Proof of Lemma \ref{lem.9v.shape4}.]Define the two partitions%
\[
\lambda^{\leq k}=\left(  \lambda_{1},\lambda_{2},\ldots,\lambda_{k}\right)
\ \ \ \ \ \ \ \ \ \ \text{and}\ \ \ \ \ \ \ \ \ \ \mu^{\leq k}=\left(  \mu
_{1},\mu_{2},\ldots,\mu_{k}\right)  .
\]
We have $\lambda_{k}=\mu_{k}$, hence $\mu_{k}=\lambda_{k}\geq\lambda_{k+1}$
(since $\lambda$ is a partition). Thus, Lemma \ref{lem.9v.shape2} yields%
\begin{equation}
s_{\lambda/\mu}=s_{\lambda^{\leq k}/\mu^{\leq k}}\cdot\varphi^{-k}\left(
s_{\lambda^{>k}/\mu^{>k}}\right)  . \label{pf.lem.9v.shape4.1}%
\end{equation}
However, the partitions $\lambda^{\leq k}=\left(  \lambda_{1},\lambda
_{2},\ldots,\lambda_{k}\right)  $ and $\mu^{\leq k}=\left(  \mu_{1},\mu
_{2},\ldots,\mu_{k}\right)  $ clearly have lengths $\ell\left(  \lambda^{\leq
k}\right)  \leq k$ and $\ell\left(  \mu^{\leq k}\right)  \leq k$ and satisfy
$\lambda_{k}=\mu_{k}$. Hence, if we define the two partitions
\[
\overline{\lambda^{\leq k}}=\left(  \lambda_{1},\lambda_{2},\ldots
,\lambda_{k-1}\right)  \ \ \ \ \ \ \ \ \ \ \text{and}%
\ \ \ \ \ \ \ \ \ \ \overline{\mu^{\leq k}}=\left(  \mu_{1},\mu_{2},\ldots
,\mu_{k-1}\right)
\]
(which consist of the first $k-1$ entries of $\lambda^{\leq k}$ and $\mu^{\leq
k}$, respectively), then Lemma \ref{lem.9v.shape3} (applied to $k$,
$\lambda^{\leq k}$, $\mu^{\leq k}$, $\overline{\lambda^{\leq k}}$ and
$\overline{\mu^{\leq k}}$ instead of $n$, $\lambda$, $\mu$, $\overline
{\lambda}$ and $\overline{\mu}$) yields%
\begin{equation}
s_{\lambda^{\leq k}/\mu^{\leq k}}=s_{\overline{\lambda^{\leq k}}/\overline
{\mu^{\leq k}}}. \label{pf.lem.9v.shape4.2}%
\end{equation}
However, $\overline{\lambda^{\leq k}}=\left(  \lambda_{1},\lambda_{2}%
,\ldots,\lambda_{k-1}\right)  =\lambda^{\leq k-1}$ and similarly
$\overline{\mu^{\leq k}}=\mu^{\leq k-1}$. In view of this, we can rewrite
(\ref{pf.lem.9v.shape4.2}) as%
\[
s_{\lambda^{\leq k}/\mu^{\leq k}}=s_{\lambda^{\leq k-1}/\mu^{\leq k-1}}.
\]
This lets us rewrite (\ref{pf.lem.9v.shape4.1}) as%
\[
s_{\lambda/\mu}=s_{\lambda^{\leq k-1}/\mu^{\leq k-1}}\cdot\varphi^{-k}\left(
s_{\lambda^{>k}/\mu^{>k}}\right)  .
\]
Hence, Lemma \ref{lem.9v.shape4} is proven.
\end{proof}

\begin{proof}
[Proof of Lemma \ref{lem.9v.shape}.]We proceed by strong induction on $n$,
where $n$ is an integer such that all four lengths $\ell\left(  \lambda
\right)  ,\ell\left(  \mu\right)  ,\ell\left(  \alpha\right)  ,\ell\left(
\beta\right)  $ are $\leq n$. The \textit{base case} ($n=0$) is trivial, since
both $s_{\lambda/\mu}$ and $s_{\alpha/\beta}$ are equal in this case.

Now to the \textit{induction step:} Fix a positive integer $n$. Assume (as
induction hypothesis) that Lemma \ref{lem.9v.shape} is already proved in the
case when all four lengths $\ell\left(  \lambda\right)  ,\ell\left(
\mu\right)  ,\ell\left(  \alpha\right)  ,\ell\left(  \beta\right)  $ are $<n$.
We must then prove Lemma \ref{lem.9v.shape} in the case when all four lengths
$\ell\left(  \lambda\right)  ,\ell\left(  \mu\right)  ,\ell\left(
\alpha\right)  ,\ell\left(  \beta\right)  $ are $\leq n$.

So let $\lambda,\mu,\alpha,\beta$ be four partitions as in Lemma
\ref{lem.9v.shape} (that is, satisfying $\mu\subseteq\lambda$ and
$\beta\subseteq\alpha$ and $\lambda-\mu=\alpha-\beta$), and assume that all
four lengths $\ell\left(  \lambda\right)  ,\ell\left(  \mu\right)
,\ell\left(  \alpha\right)  ,\ell\left(  \beta\right)  $ are $\leq n$. We must
show that $s_{\lambda/\mu}=s_{\alpha/\beta}$.

If $\lambda=\alpha$ and $\mu=\beta$, then this is obvious. Thus, we WLOG
assume that we don't have $\lambda=\alpha$ and $\mu=\beta$. Hence, Lemma
\ref{lem.9v.shape1} shows that there exists some $k\in\left[  n\right]  $
satisfying $\lambda_{k}=\mu_{k}\neq\alpha_{k}=\beta_{k}$. Consider this $k$.

Now, define the eight partitions
\begin{align*}
\lambda^{\leq k-1}  &  =\left(  \lambda_{1},\lambda_{2},\ldots,\lambda
_{k-1}\right)  ,\ \ \ \ \ \ \ \ \ \ \lambda^{>k}=\left(  \lambda_{k+1}%
,\lambda_{k+2},\ldots,\lambda_{n}\right)  ,\\
\mu^{\leq k-1}  &  =\left(  \mu_{1},\mu_{2},\ldots,\mu_{k-1}\right)
,\ \ \ \ \ \ \ \ \ \ \mu^{>k}=\left(  \mu_{k+1},\mu_{k+2},\ldots,\mu
_{n}\right)  ,\\
\alpha^{\leq k-1}  &  =\left(  \alpha_{1},\alpha_{2},\ldots,\alpha
_{k-1}\right)  ,\ \ \ \ \ \ \ \ \ \ \alpha^{>k}=\left(  \alpha_{k+1}%
,\alpha_{k+2},\ldots,\alpha_{n}\right)  ,\\
\beta^{\leq k-1}  &  =\left(  \beta_{1},\beta_{2},\ldots,\beta_{k-1}\right)
,\ \ \ \ \ \ \ \ \ \ \beta^{>k}=\left(  \beta_{k+1},\beta_{k+2},\ldots
,\beta_{n}\right)  .
\end{align*}
All these eight partitions have lengths $<n$ (indeed, the four partitions
$\lambda^{\leq k-1},\mu^{\leq k-1},\alpha^{\leq k-1},\beta^{\leq k-1}$ have
lengths $\leq k-1<k\leq n$, whereas the four partitions $\lambda^{>k},\mu
^{>k},\alpha^{>k},\beta^{>k}$ have lengths $\leq n-k<n$).

The definitions of $\lambda^{\leq k-1}$ and $\mu^{\leq k-1}$ show that the
Young diagrams of $\lambda^{\leq k-1}$ and $\mu^{\leq k-1}$ consist of the
first $k-1$ rows of the Young diagrams of $\lambda$ and $\mu$, respectively.
Hence, the skew shape $\lambda^{\leq k-1}-\mu^{\leq k-1}$ consists of the
first $k-1$ rows of the skew shape $\lambda-\mu$. Similarly, the skew shape
$\alpha^{\leq k-1}-\beta^{\leq k-1}$ consists of the first $k-1$ rows of the
skew shape $\alpha-\beta$. Hence, from $\lambda-\mu=\alpha-\beta$, we obtain
$\lambda^{\leq k-1}-\mu^{\leq k-1}=\alpha^{\leq k-1}-\beta^{\leq k-1}$
(because if two skew shapes are equal, then the skew shapes consisting of
their respective first $k-1$ rows are also equal). Since the partitions
$\lambda^{\leq k-1},\mu^{\leq k-1},\alpha^{\leq k-1},\beta^{\leq k-1}$ all
have lengths $<n$, we thus can apply the induction hypothesis to
$\lambda^{\leq k-1},\mu^{\leq k-1},\alpha^{\leq k-1},\beta^{\leq k-1}$ instead
of $\lambda,\mu,\alpha,\beta$, and conclude that%
\begin{equation}
s_{\lambda^{\leq k-1}/\mu^{\leq k-1}}=s_{\alpha^{\leq k-1}/\beta^{\leq k-1}}.
\label{pf.lem.9v.shape.s=sL}%
\end{equation}


The definitions of $\lambda^{>k}$ and $\mu^{>k}$ show that the Young diagrams
of $\lambda^{>k}$ and $\mu^{>k}$ consist of the $\left(  k+1\right)  $-st,
$\left(  k+2\right)  $-nd, $...$, $n$-th rows of the Young diagrams of
$\lambda$ and $\mu$, respectively. Thus, an argument (using the induction
hypothesis) analogous to the proof of (\ref{pf.lem.9v.shape.s=sL}) shows that%
\begin{equation}
s_{\lambda^{>k}/\mu^{>k}}=s_{\alpha^{>k}/\beta^{>k}}.
\label{pf.lem.9v.shape.s=sR}%
\end{equation}


However, we know that $\ell\left(  \lambda\right)  \leq n$ and $\ell\left(
\mu\right)  \leq n$ and $k\in\left[  n\right]  $ and $\lambda_{k}=\mu_{k}$.
Thus, Lemma \ref{lem.9v.shape4} yields%
\[
s_{\lambda/\mu}=s_{\lambda^{\leq k-1}/\mu^{\leq k-1}}\cdot\varphi^{-k}\left(
s_{\lambda^{>k}/\mu^{>k}}\right)  .
\]
The same argument (but made for the partitions $\alpha$ and $\beta$ instead of
$\lambda$ and $\mu$) shows that%
\[
s_{\alpha/\beta}=s_{\alpha^{\leq k-1}/\beta^{\leq k-1}}\cdot\varphi
^{-k}\left(  s_{\alpha^{>k}/\beta^{>k}}\right)  .
\]
Thus,%
\begin{align*}
s_{\lambda/\mu}  &  =\underbrace{s_{\lambda^{\leq k-1}/\mu^{\leq k-1}}%
}_{\substack{=s_{\alpha^{\leq k-1}/\beta^{\leq k-1}}\\\text{(by
(\ref{pf.lem.9v.shape.s=sL}))}}}\cdot\varphi^{-k}\left(
\underbrace{s_{\lambda^{>k}/\mu^{>k}}}_{\substack{=s_{\alpha^{>k}/\beta^{>k}%
}\\\text{(by (\ref{pf.lem.9v.shape.s=sR}))}}}\right) \\
&  =s_{\alpha^{\leq k-1}/\beta^{\leq k-1}}\cdot\varphi^{-k}\left(
s_{\alpha^{>k}/\beta^{>k}}\right)  =s_{\alpha/\beta}.
\end{align*}
This completes the induction step. Thus, Lemma \ref{lem.9v.shape} is proved by induction.
\end{proof}

\item \textbf{page 36:} Let me give some details on how to prove the two
identities%
\[
s_{\tau_{0,1}\left(  \theta\right)  }=\varphi s_{\theta}%
\ \ \ \ \ \ \ \ \ \ \text{and}\ \ \ \ \ \ \ \ \ \ s_{\tau_{1,0}\left(
\theta\right)  }=\varphi^{-1}s_{\theta}%
\]
that make up the bulk of the proof of (9.4).

\begin{lemmain}
\label{lem.9v.move}Let $\theta=\lambda-\mu$ be a skew shape. Then,%
\[
s_{\tau_{0,1}\left(  \theta\right)  }=\varphi s_{\theta}%
\ \ \ \ \ \ \ \ \ \ \text{and}\ \ \ \ \ \ \ \ \ \ s_{\tau_{1,0}\left(
\theta\right)  }=\varphi^{-1}s_{\theta}.
\]

\end{lemmain}

Before we prove this, we introduce one further lemma:

\begin{lemmain}
\label{lem.9v.shape4-1}Let $\lambda$ and $\mu$ be two partitions. Let
$n\in\mathbb{N}$ be such that $\ell\left(  \lambda\right)  \leq n$ and
$\ell\left(  \mu\right)  \leq n$. Assume that $\lambda_{1}=\mu_{1}$. Define
the two partitions%
\[
\lambda^{>1}=\left(  \lambda_{2},\lambda_{3},\ldots,\lambda_{n}\right)
\ \ \ \ \ \ \ \ \ \ \text{and}\ \ \ \ \ \ \ \ \ \ \mu^{>1}=\left(  \mu_{2}%
,\mu_{3},\ldots,\mu_{n}\right)  .
\]
Then,%
\[
s_{\lambda/\mu}=\varphi^{-1}\left(  s_{\lambda^{>1}/\mu^{>1}}\right)  .
\]

\end{lemmain}

\begin{proof}
[Proof of Lemma \ref{lem.9v.shape4-1}.]Our definitions of $\lambda^{>1}$ and
$\mu^{>1}$ agree with the definitions of $\lambda^{>k}$ and $\mu^{>k}$ in
Lemma \ref{lem.9v.shape4} for $k=1$. We have $\lambda_{1}=\mu_{1}$ by
assumption. Hence, applying Lemma \ref{lem.9v.shape4} to $k=1$, we obtain%
\begin{equation}
s_{\lambda/\mu}=s_{\lambda^{\leq1-1}/\mu^{\leq1-1}}\cdot\varphi^{-1}\left(
s_{\lambda^{>1}/\mu^{>1}}\right)  , \label{pf.lem.9v.shape4-1.1}%
\end{equation}
where
\[
\lambda^{\leq1-1}=\left(  \lambda_{1},\lambda_{2},\ldots,\lambda_{1-1}\right)
\ \ \ \ \ \ \ \ \ \ \text{and}\ \ \ \ \ \ \ \ \ \ \mu^{\leq1-1}=\left(
\mu_{1},\mu_{2},\ldots,\mu_{1-1}\right)  .
\]
However, both partitions $\lambda^{\leq1-1}$ and $\mu^{\leq1-1}$ are empty
(i.e., have no entries at all), i.e., equal the empty partition $\varnothing$.
Hence, $s_{\lambda^{\leq1-1}/\mu^{\leq1-1}}=s_{\varnothing/\varnothing}=1$
(since $s_{\varnothing/\varnothing}$ is defined as the determinant of a
$0\times0$-matrix, but such a determinant is always $1$). Therefore,
(\ref{pf.lem.9v.shape4-1.1}) rewrites as%
\[
s_{\lambda/\mu}=1\cdot\varphi^{-1}\left(  s_{\lambda^{>1}/\mu^{>1}}\right)
=\varphi^{-1}\left(  s_{\lambda^{>1}/\mu^{>1}}\right)  .
\]
This proves Lemma \ref{lem.9v.shape4-1}.
\end{proof}

\begin{proof}
[Proof of Lemma \ref{lem.9v.move}.]Let us pick an integer $n\in\mathbb{N}$
such that $\ell\left(  \lambda\right)  \leq n$ and $\ell\left(  \mu\right)
\leq n$. Thus, $\lambda=\left(  \lambda_{1},\lambda_{2},\ldots,\lambda
_{n}\right)  $ and $\mu=\left(  \mu_{1},\mu_{2},\ldots,\mu_{n}\right)  $. Let
$g$ be an integer such that $g\geq\lambda_{1}$ and $g\geq\mu_{1}$. Let us
furthermore define the four partitions%
\begin{align*}
\lambda^{+}  &  =\left(  \lambda_{1}+1,\ \lambda_{2}+1,\ \ldots,\ \lambda
_{n}+1\right)  ,\\
\mu^{+}  &  =\left(  \mu_{1}+1,\ \mu_{2}+1,\ \ldots,\ \mu_{n}+1\right)  ,\\
\widetilde{\lambda}  &  =\left(  g,\lambda_{1},\lambda_{2},\ldots,\lambda
_{n}\right)  ,\\
\widetilde{\mu}  &  =\left(  g,\mu_{1},\mu_{2},\ldots,\mu_{n}\right)
\end{align*}
(note that $\widetilde{\lambda}$ and $\widetilde{\mu}$ are partitions since
$g\geq\lambda_{1}$ and $g\geq\mu_{1}$). Thus, the Young diagrams of
$\lambda^{+}$ and $\mu^{+}$ are obtained from those of $\lambda$ and $\mu$ by
inserting an extra column of length $n$ on the left end and moving the
existing cells $1$ step to the right. Hence,
\[
\tau_{0,1}\left(  \lambda-\mu\right)  =\lambda^{+}-\mu^{+}.
\]
Furthermore, the Young diagrams of $\widetilde{\lambda}$ and $\widetilde{\mu}$
are obtained from those of $\lambda$ and $\mu$ by inserting an extra row of
length $g$ at the top and moving the existing cells $1$ step down. Hence,
\[
\tau_{1,0}\left(  \lambda-\mu\right)  =\widetilde{\lambda}-\widetilde{\mu}.
\]


We have $\theta=\lambda-\mu$ and thus%
\begin{equation}
s_{\theta}=s_{\lambda/\mu}=\det\left(  \varphi^{\mu_{j}-j+1}h_{\lambda_{i}%
-\mu_{j}-i+j}\right)  _{1\leq i,j\leq n} \label{pf.lem.9v.move.slm}%
\end{equation}
(by the definition of $s_{\lambda/\mu}$, since $\lambda=\left(  \lambda
_{1},\lambda_{2},\ldots,\lambda_{n}\right)  $ and $\mu=\left(  \mu_{1},\mu
_{2},\ldots,\mu_{n}\right)  $).

From $\lambda^{+}=\left(  \lambda_{1}+1,\ \lambda_{2}+1,\ \ldots,\ \lambda
_{n}+1\right)  $ and $\mu^{+}=\left(  \mu_{1}+1,\ \mu_{2}+1,\ \ldots,\ \mu
_{n}+1\right)  $, we obtain%
\begin{align*}
s_{\lambda^{+}/\mu^{+}}  &  =\det\left(  \underbrace{\varphi^{\left(  \mu
_{j}+1\right)  -j+1}}_{\substack{=\varphi^{1+\left(  \mu_{j}-j+1\right)
}\\=\varphi\varphi^{\mu_{j}-j+1}}}\underbrace{h_{\left(  \lambda_{i}+1\right)
-\left(  \mu_{j}+1\right)  -i+j}}_{=h_{\lambda_{i}-\mu_{j}-i+j}}\right)
_{1\leq i,j\leq n}\\
&  =\det\left(  \varphi\varphi^{\mu_{j}-j+1}h_{\lambda_{i}-\mu_{j}%
-i+j}\right)  _{1\leq i,j\leq n}\\
&  =\varphi\underbrace{\left(  \det\left(  \varphi^{\mu_{j}-j+1}h_{\lambda
_{i}-\mu_{j}-i+j}\right)  _{1\leq i,j\leq n}\right)  }_{\substack{=s_{\theta
}\\\text{(by (\ref{pf.lem.9v.move.slm}))}}}\ \ \ \ \ \ \ \ \ \ \left(
\text{since }\varphi\text{ is a ring morphism}\right) \\
&  =\varphi s_{\theta}.
\end{align*}
But $\theta=\lambda-\mu$ and thus $\tau_{0,1}\left(  \theta\right)
=\tau_{0,1}\left(  \lambda-\mu\right)  =\lambda^{+}-\mu^{+}$. Hence,%
\[
s_{\tau_{0,1}\left(  \theta\right)  }=s_{\lambda^{+}/\mu^{+}}=\varphi
s_{\theta}.
\]
This proves the first equality in Lemma \ref{lem.9v.move}. It remains to prove
the second one.

Write the $\left(  n+1\right)  $-tuple $\widetilde{\lambda}=\left(
g,\lambda_{1},\lambda_{2},\ldots,\lambda_{n}\right)  $ as $\widetilde{\lambda
}=\left(  \widetilde{\lambda}_{1},\widetilde{\lambda}_{2},\ldots
,\widetilde{\lambda}_{n+1}\right)  $. Thus,%
\begin{equation}
\widetilde{\lambda}_{1}=g\ \ \ \ \ \ \ \ \ \ \text{and}
\label{pf.lem.9v.move.lt1}%
\end{equation}%
\begin{equation}
\widetilde{\lambda}_{i}=\lambda_{i-1}\ \ \ \ \ \ \ \ \ \ \text{for each }%
i\in\left[  2,n+1\right]  . \label{pf.lem.9v.move.lti}%
\end{equation}


Write the $\left(  n+1\right)  $-tuple $\widetilde{\mu}=\left(  g,\mu_{1}%
,\mu_{2},\ldots,\mu_{n}\right)  $ as $\widetilde{\mu}=\left(  \widetilde{\mu
}_{1},\widetilde{\mu}_{2},\ldots,\widetilde{\mu}_{n+1}\right)  $. Thus,%
\begin{equation}
\widetilde{\mu}_{1}=g\ \ \ \ \ \ \ \ \ \ \text{and} \label{pf.lem.9v.move.mt1}%
\end{equation}%
\begin{equation}
\widetilde{\mu}_{i}=\mu_{i-1}\ \ \ \ \ \ \ \ \ \ \text{for each }i\in\left[
2,n+1\right]  . \label{pf.lem.9v.move.mti}%
\end{equation}


Clearly, $\ell\left(  \widetilde{\lambda}\right)  \leq n+1$ and $\ell\left(
\widetilde{\mu}\right)  \leq n+1$ and $\widetilde{\lambda}_{1}=\widetilde{\mu
}_{1}$ (by comparing (\ref{pf.lem.9v.move.lt1}) with (\ref{pf.lem.9v.move.mt1}%
)). Thus, Lemma \ref{lem.9v.shape4-1} (applied to $n+1$, $\widetilde{\lambda}%
$, $\widetilde{\mu}$, $\widetilde{\lambda}_{i}$ and $\widetilde{\mu}_{i}$
instead of $n$, $\lambda$, $\mu$, $\lambda_{i}$ and $\mu_{i}$) yields%
\begin{equation}
s_{\widetilde{\lambda}/\widetilde{\mu}}=\varphi^{-1}\left(
s_{\widetilde{\lambda}^{>1}/\widetilde{\mu}^{>1}}\right)  ,
\label{pf.lem.9v.move.slti}%
\end{equation}
where%
\[
\widetilde{\lambda}^{>1}=\left(  \widetilde{\lambda}_{2},\widetilde{\lambda
}_{3},\ldots,\widetilde{\lambda}_{n+1}\right)  \ \ \ \ \ \ \ \ \ \ \text{and}%
\ \ \ \ \ \ \ \ \ \ \widetilde{\mu}^{>1}=\left(  \widetilde{\mu}%
_{2},\widetilde{\mu}_{3},\ldots,\widetilde{\mu}_{n+1}\right)  .
\]


However,
\begin{align*}
\widetilde{\lambda}^{>1}  &  =\left(  \widetilde{\lambda}_{2}%
,\widetilde{\lambda}_{3},\ldots,\widetilde{\lambda}_{n+1}\right)  =\left(
\lambda_{1},\lambda_{2},\ldots,\lambda_{n}\right)  \ \ \ \ \ \ \ \ \ \ \left(
\text{by (\ref{pf.lem.9v.move.lti})}\right) \\
&  =\lambda
\end{align*}
and $\widetilde{\mu}^{>1}=\mu$ (similarly). Thus, we can rewrite
(\ref{pf.lem.9v.move.slti}) as%
\[
s_{\widetilde{\lambda}/\widetilde{\mu}}=\varphi^{-1}\left(  s_{\lambda/\mu
}\right)  .
\]
But $\theta=\lambda-\mu$ and thus $\tau_{1,0}\left(  \theta\right)
=\tau_{1,0}\left(  \lambda-\mu\right)  =\widetilde{\lambda}-\widetilde{\mu}$.
Hence,
\[
s_{\tau_{1,0}\left(  \theta\right)  }=s_{\widetilde{\lambda}/\widetilde{\mu}%
}=\varphi^{-1}\left(  s_{\lambda/\mu}\right)  =\varphi^{-1}\left(  s_{\theta
}\right)
\]
(since $s_{\lambda/\mu}=s_{\theta}$). This proves the second equality in Lemma
\ref{lem.9v.move}. The proof of the lemma is thus complete.
\end{proof}

\item \textbf{page 36, proof of (9.5):} In the displayed equation, replace the
subscript \textquotedblleft$\widehat{\mu}_{j}-\widehat{\lambda}_{i}%
+i+j$\textquotedblright\ by \textquotedblleft$\widehat{\mu}_{j}%
-\widehat{\lambda}_{i}+i-j$\textquotedblright.

\item \textbf{page 37, between (9.6) and (9.6'):} \textquotedblleft where
$\omega$ is the involution defined by $\varphi^{s}h_{r}\mapsto\varphi
^{-s}e_{r}$ for all $r,s$\textquotedblright\ presupposes the existence of such
an involution. This is indeed the case, but it needs to be proved. What
\textbf{is} immediately clear is that there exists a unique $\mathbb{Z}%
$\textbf{-algebra endomorphism} of $R$ that sends each $\varphi^{s}h_{r}$ with
$r\geq1$ and $s\in\mathbb{Z}$ to $\varphi^{-s}e_{r}$ (since the $\varphi
^{s}h_{r}=h_{r,s}$ for $r\geq1$ and $s\in\mathbb{Z}$ are the indeterminates of
the polynomial ring $R$). Denoting this endomorphism by $\omega$, we can then
easily see that it also sends $\varphi^{s}h_{r}$ to $\varphi^{-s}e_{r}$ when
$r$ is negative (because if $r=0$, then both $\varphi^{s}h_{r}$ and
$\varphi^{-s}e_{r}$ are $1$, whereas if $r<0$, then both $\varphi^{s}h_{r}$
and $\varphi^{-s}e_{r}$ are $0$). That is, it satisfies%
\begin{equation}
\omega\left(  \varphi^{s}h_{r}\right)  =\varphi^{-s}e_{r}%
\ \ \ \ \ \ \ \ \ \ \text{for all }r,s\in\mathbb{Z}. \label{p37.omg}%
\end{equation}
It remains to show that $\omega$ is an involution. As it happens, this is
easiest to do \textbf{after} proving (9.6') and (9.6).

Let me first expand on the proofs of (9.6') and (9.6), since they are not as
easy as Macdonald is making them look. I will need two lemmas:

\begin{lemmain}
\label{lem.matprod.-180}Let $A=\left(  a_{i,j}\right)  _{i,j\in\mathbb{Z}}$
and $B=\left(  b_{i,j}\right)  _{i,j\in\mathbb{Z}}$ be two upper-triangular
$\mathbb{Z}\times\mathbb{Z}$-matrices such that $AB=I_{\mathbb{Z}}$. Define
the $\mathbb{Z}\times\mathbb{Z}$-matrices $\widetilde{A}:=\left(
a_{-j,-i}\right)  _{i,j\in\mathbb{Z}}$ and $\widetilde{B}:=\left(
b_{-j,-i}\right)  _{i,j\in\mathbb{Z}}$. Then, $\widetilde{B}\widetilde{A}%
=I_{\mathbb{Z}}$ as well.
\end{lemmain}

\begin{proof}
[Proof of Lemma \ref{lem.matprod.-180}.]For each $i,k\in\mathbb{Z}$, we have%
\begin{equation}
\sum_{j\in\mathbb{Z}}a_{i,j}b_{j,k}=\delta_{i,k} \label{pf.lem.matprod.-180.1}%
\end{equation}
(by comparing the $\left(  i,k\right)  $-th entries of the matrices on both
sides of $AB=I_{\mathbb{Z}}$).

Now, for each $i,k\in\mathbb{Z}$, we have%
\begin{align*}
&  \left(  \text{the }\left(  i,k\right)  \text{-th entry of the matrix
}\widetilde{B}\widetilde{A}\right) \\
&  =\sum_{j\in\mathbb{Z}}b_{-j,-i}a_{-k,-j}\ \ \ \ \ \ \ \ \ \ \left(
\text{since }\widetilde{B}=\left(  b_{-j,-i}\right)  _{i,j\in\mathbb{Z}}\text{
and }\widetilde{A}=\left(  a_{-j,-i}\right)  _{i,j\in\mathbb{Z}}\right) \\
&  =\sum_{j\in\mathbb{Z}}\underbrace{b_{j,-i}a_{-k,j}}_{=a_{-k,j}b_{j,-i}%
}\ \ \ \ \ \ \ \ \ \ \left(
\begin{array}
[c]{c}%
\text{here, we have substituted }-j\\
\text{for }j\text{ in the sum}%
\end{array}
\right) \\
&  =\sum_{j\in\mathbb{Z}}a_{-k,j}b_{j,-i}=\delta_{-k,-i}\\
&  \ \ \ \ \ \ \ \ \ \ \ \ \ \ \ \ \ \ \ \ \left(  \text{by
(\ref{pf.lem.matprod.-180.1}), applied to }-k\text{ and }-i\text{ instead of
}i\text{ and }k\right) \\
&  =\delta_{i,k}\ \ \ \ \ \ \ \ \ \ \left(  \text{since }-k=-i\text{ holds if
and only if }i=k\right) \\
&  =\left(  \text{the }\left(  i,k\right)  \text{-th entry of the matrix
}I_{\mathbb{Z}}\right)  .
\end{align*}
In other words, $\widetilde{B}\widetilde{A}=I_{\mathbb{Z}}$. This proves Lemma
\ref{lem.matprod.-180}.
\end{proof}

Our next lemma is a variant of (9.3):

\begin{lemmain}
\label{lem.9.3.prime}The two $\mathbb{Z}\times\mathbb{Z}$-matrices%
\begin{align*}
\widetilde{H}  &  =\left(  \varphi^{1+i}h_{j-i}\right)  _{i,j\in\mathbb{Z}%
}\ \ \ \ \ \ \ \ \ \ \text{and}\\
\widetilde{E}  &  =\left(  \left(  -1\right)  ^{j-i}\varphi^{j}e_{j-i}\right)
_{i,j\in\mathbb{Z}}%
\end{align*}
are inverses of each other.
\end{lemmain}

\begin{proof}
[Proof of Lemma \ref{lem.9.3.prime}.]First, we observe that the matrices
$\widetilde{E}$ and $\widetilde{H}$ are upper unitriangular (since
$h_{0}=e_{0}=1$ and $h_{k}=e_{k}=0$ for all $k<0$). Hence, they belong to the
group $\operatorname*{UT}\nolimits_{\mathbb{Z}}$ of all upper unitriangular
matrices with rows and columns indexed by all integers.

Let $I=\mathbb{Z}$. Consider the $I\times I$-matrices (that is, $\mathbb{Z}%
\times\mathbb{Z}$-matrices) $H=\left(  \varphi^{1-j}h_{j-i}\right)
_{i,j\in\mathbb{Z}}$ and $E=\left(  \left(  -1\right)  ^{j-i}\varphi
^{-i}e_{j-i}\right)  _{i,j\in\mathbb{Z}}$ defined in (9.3). For comparison,
our two matrices $\widetilde{H}$ and $\widetilde{E}$ are%
\[
\widetilde{H}=\left(  \varphi^{1+i}h_{j-i}\right)  _{i,j\in\mathbb{Z}}=\left(
\varphi^{1-\left(  -i\right)  }h_{\left(  -i\right)  -\left(  -j\right)
}\right)  _{i,j\in\mathbb{Z}}%
\]
(since $1+i=1-\left(  -i\right)  $ and $j-i=\left(  -i\right)  -\left(
-j\right)  $) and%
\[
\widetilde{E}=\left(  \left(  -1\right)  ^{j-i}\varphi^{j}e_{j-i}\right)
_{i,j\in\mathbb{Z}}=\left(  \left(  -1\right)  ^{\left(  -i\right)  -\left(
-j\right)  }\varphi^{-\left(  -j\right)  }e_{\left(  -i\right)  -\left(
-j\right)  }\right)  _{i,j\in\mathbb{Z}}%
\]
(since $j-i=\left(  -i\right)  -\left(  -j\right)  $ and $j=-\left(
-j\right)  $). Hence, our two matrices $\widetilde{H}$ and $\widetilde{E}$ are
obtained from the two matrices $H$ and $E$ in the same way as the matrices
$\widetilde{A}$ and $\widetilde{B}$ in Lemma \ref{lem.matprod.-180} are
obtained from $A$ and $B$. Thus, Lemma \ref{lem.matprod.-180} (applied to $H$,
$E$, $\widetilde{H}$ and $\widetilde{E}$ instead of $A$, $B$, $\widetilde{A}$
and $\widetilde{B}$) shows that $\widetilde{E}\widetilde{H}=I_{\mathbb{Z}}$
(since (9.3) shows that $HE=I_{\mathbb{Z}}$). Since the matrices
$\widetilde{E}$ and $\widetilde{H}$ belong to the group $\operatorname*{UT}%
\nolimits_{\mathbb{Z}}$, this equality entails that they are inverses of each
other. In other words, $\widetilde{H}$ and $\widetilde{E}$ are inverses of
each other. This proves Lemma \ref{lem.9.3.prime}.
\end{proof}

\begin{proof}
[Proof of (9.6').]Write the partitions $\lambda$, $\lambda^{\prime}$, $\mu$
and $\mu^{\prime}$ in the forms%
\begin{align*}
\lambda &  =\left(  \lambda_{1},\lambda_{2},\ldots,\lambda_{q}\right)
,\ \ \ \ \ \ \ \ \ \ \lambda^{\prime}=\left(  \lambda_{1}^{\prime},\lambda
_{2}^{\prime},\ldots,\lambda_{p}^{\prime}\right)  ,\\
\mu &  =\left(  \mu_{1},\mu_{2},\ldots,\mu_{q}\right)
,\ \ \ \ \ \ \ \ \ \ \mu^{\prime}=\left(  \mu_{1}^{\prime},\mu_{2}^{\prime
},\ldots,\mu_{p}^{\prime}\right)
\end{align*}
for some $p,q\in\mathbb{N}$ (where, of course, the entries $\lambda_{i}$,
$\lambda_{j}^{\prime}$, $\mu_{i}$ and $\mu_{j}^{\prime}$ are allowed to be $0$).

Consider the two matrices%
\[
\widetilde{H}=\left(  \varphi^{1+i}h_{j-i}\right)  _{i,j\in\mathbb{Z}%
}\ \ \ \ \ \ \ \ \ \ \text{and}\ \ \ \ \ \ \ \ \ \ \widetilde{E}=\left(
\left(  -1\right)  ^{j-i}\varphi^{j}e_{j-i}\right)  _{i,j\in\mathbb{Z}}%
\]
from Lemma \ref{lem.9.3.prime}. These two matrices are upper unitriangular
(since $h_{0}=e_{0}=1$ and $h_{k}=e_{k}=0$ for all $k<0$). Moreover, the
matrix $\widetilde{E}$ is the inverse matrix of $\widetilde{H}$ (by Lemma
\ref{lem.9.3.prime}). Hence, Lemma \ref{lem.A5} (applied to $B=\widetilde{H}$
and $b_{i,j}=\varphi^{1+i}h_{j-i}$ and $c_{i,j}=\left(  -1\right)
^{j-i}\varphi^{j}e_{j-i}$) yields%
\begin{align*}
&  \det\left(  \varphi^{1+\left(  \mu_{i}-i\right)  }h_{\left(  \lambda
_{j}-j\right)  -\left(  \mu_{i}-i\right)  }\right)  _{i,j\in\left[  q\right]
}\\
&  =\left(  -1\right)  ^{\left(  \lambda_{1}+\lambda_{2}+\cdots+\lambda
_{q}\right)  +\left(  \mu_{1}+\mu_{2}+\cdots+\mu_{q}\right)  }\\
&  \ \ \ \ \ \ \ \ \ \ \det\left(  \left(  -1\right)  ^{\left(  j-\mu
_{j}^{\prime}-1\right)  -\left(  i-\lambda_{i}^{\prime}-1\right)  }%
\varphi^{j-\mu_{j}^{\prime}-1}e_{\left(  j-\mu_{j}^{\prime}-1\right)  -\left(
i-\lambda_{i}^{\prime}-1\right)  }\right)  _{i,j\in\left[  p\right]  }.
\end{align*}
In view of $\left(  \lambda_{j}-j\right)  -\left(  \mu_{i}-i\right)
=\lambda_{j}-\mu_{i}-j+i$ and $1+\left(  \mu_{i}-i\right)  =\mu_{i}-i+1$ and
$\left(  j-\mu_{j}^{\prime}-1\right)  -\left(  i-\lambda_{i}^{\prime
}-1\right)  =\lambda_{i}^{\prime}-\mu_{j}^{\prime}-i+j$, we can rewrite this
as%
\begin{align*}
&  \det\left(  \varphi^{\mu_{i}-i+1}h_{\lambda_{j}-\mu_{i}-j+i}\right)
_{i,j\in\left[  q\right]  }\\
&  =\left(  -1\right)  ^{\left(  \lambda_{1}+\lambda_{2}+\cdots+\lambda
_{q}\right)  +\left(  \mu_{1}+\mu_{2}+\cdots+\mu_{q}\right)  }\\
&  \ \ \ \ \ \ \ \ \ \ \det\left(  \left(  -1\right)  ^{\lambda_{i}^{\prime
}-\mu_{j}^{\prime}-i+j}\varphi^{j-\mu_{j}^{\prime}-1}e_{\lambda_{i}^{\prime
}-\mu_{j}^{\prime}-i+j}\right)  _{i,j\in\left[  p\right]  }.
\end{align*}
In view of%
\[
\det\left(  \varphi^{\mu_{i}-i+1}h_{\lambda_{j}-\mu_{i}-j+i}\right)
_{i,j\in\left[  q\right]  }=\det\left(  \varphi^{\mu_{j}-j+1}h_{\lambda
_{i}-\mu_{j}-i+j}\right)  _{i,j\in\left[  q\right]  }%
\]
(since the determinant of a matrix does not change when we transpose it) and%
\begin{align*}
&  \det\left(  \underbrace{\left(  -1\right)  ^{\lambda_{i}^{\prime}-\mu
_{j}^{\prime}-i+j}}_{=\left(  -1\right)  ^{\lambda_{i}^{\prime}-i}\left(
-1\right)  ^{\mu_{j}^{\prime}-j}}\varphi^{j-\mu_{j}^{\prime}-1}e_{\lambda
_{i}^{\prime}-\mu_{j}^{\prime}-i+j}\right)  _{i,j\in\left[  p\right]  }\\
&  =\det\left(  \left(  -1\right)  ^{\lambda_{i}^{\prime}-i}\left(  -1\right)
^{\mu_{j}^{\prime}-j}\varphi^{j-\mu_{j}^{\prime}-1}e_{\lambda_{i}^{\prime}%
-\mu_{j}^{\prime}-i+j}\right)  _{i,j\in\left[  p\right]  }\\
&  =\left(  \prod_{i=1}^{p}\left(  -1\right)  ^{\lambda_{i}^{\prime}%
-i}\right)  \left(  \prod_{j=1}^{p}\left(  -1\right)  ^{\mu_{j}^{\prime}%
-j}\right)  \cdot\det\left(  \varphi^{j-\mu_{j}^{\prime}-1}e_{\lambda
_{i}^{\prime}-\mu_{j}^{\prime}-i+j}\right)  _{i,j\in\left[  p\right]  }%
\end{align*}
(by Lemma \ref{lem.A6}, applied to $\alpha_{i}=\left(  -1\right)
^{\lambda_{i}^{\prime}-i}$ and $\beta_{j}=\left(  -1\right)  ^{\mu_{j}%
^{\prime}-j}$ and $u_{i,j}=\varphi^{j-\mu_{j}^{\prime}-1}e_{\lambda
_{i}^{\prime}-\mu_{j}^{\prime}-i+j}$), we can rewrite this as%
\begin{align*}
&  \det\left(  \varphi^{\mu_{j}-j+1}h_{\lambda_{i}-\mu_{j}-i+j}\right)
_{i,j\in\left[  q\right]  }\\
&  =\left(  -1\right)  ^{\left(  \lambda_{1}+\lambda_{2}+\cdots+\lambda
_{q}\right)  +\left(  \mu_{1}+\mu_{2}+\cdots+\mu_{q}\right)  }\\
&  \ \ \ \ \ \ \ \ \ \ \left(  \prod_{i=1}^{p}\left(  -1\right)  ^{\lambda
_{i}^{\prime}-i}\right)  \left(  \prod_{j=1}^{p}\left(  -1\right)  ^{\mu
_{j}^{\prime}-j}\right)  \cdot\det\left(  \varphi^{j-\mu_{j}^{\prime}%
-1}e_{\lambda_{i}^{\prime}-\mu_{j}^{\prime}-i+j}\right)  _{i,j\in\left[
p\right]  }\\
&  =\underbrace{\left(  -1\right)  ^{\left(  \lambda_{1}+\lambda_{2}%
+\cdots+\lambda_{q}\right)  +\left(  \mu_{1}+\mu_{2}+\cdots+\mu_{q}\right)
}\left(  \prod_{i=1}^{p}\left(  -1\right)  ^{\lambda_{i}^{\prime}-i}\right)
\left(  \prod_{j=1}^{p}\left(  -1\right)  ^{\mu_{j}^{\prime}-j}\right)
}_{\substack{=1\\\text{(by Lemma \ref{lem.A6prods})}}}\\
&  \ \ \ \ \ \ \ \ \ \ \cdot\det\left(  \varphi^{j-\mu_{j}^{\prime}%
-1}e_{\lambda_{i}^{\prime}-\mu_{j}^{\prime}-i+j}\right)  _{i,j\in\left[
p\right]  }\\
&  =\det\left(  \varphi^{j-\mu_{j}^{\prime}-1}e_{\lambda_{i}^{\prime}-\mu
_{j}^{\prime}-i+j}\right)  _{i,j\in\left[  p\right]  }\\
&  =\det\left(  \varphi^{-\mu_{j}^{\prime}+j-1}e_{\lambda_{i}^{\prime}-\mu
_{j}^{\prime}-i+j}\right)  _{i,j\in\left[  p\right]  }%
\ \ \ \ \ \ \ \ \ \ \left(  \text{since }j-\mu_{j}^{\prime}-1=-\mu_{j}%
^{\prime}+j-1\right)  .
\end{align*}
In view of (9.1), this can be rewritten as%
\[
s_{\lambda/\mu}=\det\left(  \varphi^{-\mu_{j}^{\prime}+j-1}e_{\lambda
_{i}^{\prime}-\mu_{j}^{\prime}-i+j}\right)  _{i,j\in\left[  p\right]  }.
\]
This proves (9.6').
\end{proof}

\begin{proof}
[Proof of (9.6).]From $\theta=\lambda-\mu$, we obtain $s_{\theta}%
=s_{\lambda/\mu}$. From $\theta^{\prime}=\lambda^{\prime}-\mu^{\prime}$, we
obtain $s_{\theta^{\prime}}=s_{\lambda^{\prime}/\mu^{\prime}}$. Now, (9.6')
(applied to $\lambda^{\prime}$ and $\mu^{\prime}$ instead of $\lambda$ and
$\mu$) shows that%
\[
s_{\lambda^{\prime}/\mu^{\prime}}=\det\left(  \varphi^{-\mu_{j}^{\prime\prime
}+j-1}e_{\lambda_{i}^{\prime\prime}-\mu_{j}^{\prime\prime}-i+j}\right)
_{i,j\in\left[  q\right]  },
\]
where $q\in\mathbb{N}$ is chosen such that $q\geq\ell\left(  \lambda
^{\prime\prime}\right)  $ and $q\geq\ell\left(  \mu^{\prime\prime}\right)  $.
Since $\lambda^{\prime\prime}=\lambda$ and $\mu^{\prime\prime}=\mu$, we can
rewrite this as%
\begin{equation}
s_{\lambda^{\prime}/\mu^{\prime}}=\det\left(  \varphi^{-\mu_{j}+j-1}%
e_{\lambda_{i}-\mu_{j}-i+j}\right)  _{i,j\in\left[  q\right]  }%
,\label{p37.9.6.pf.1}%
\end{equation}
where $q\in\mathbb{N}$ is chosen such that $q\geq\ell\left(  \lambda\right)  $
and $q\geq\ell\left(  \mu\right)  $. Consider this $q$.

However,%
\[
s_{\theta}=s_{\lambda/\mu}=\det\left(  \varphi^{\mu_{j}-j+1}h_{\lambda_{i}%
-\mu_{j}-i+j}\right)  _{i,j\in\left[  q\right]  }\ \ \ \ \ \ \ \ \ \ \left(
\text{by (9.1)}\right)  .
\]
Applying the map $\omega$ to both sides of this equality, we obtain%
\begin{align*}
\omega s_{\theta}  &  =\omega\left(  \det\left(  \varphi^{\mu_{j}%
-j+1}h_{\lambda_{i}-\mu_{j}-i+j}\right)  _{i,j\in\left[  q\right]  }\right) \\
&  =\det\left(  \omega\left(  \varphi^{\mu_{j}-j+1}h_{\lambda_{i}-\mu_{j}%
-i+j}\right)  \right)  _{i,j\in\left[  q\right]  }\ \ \ \ \ \ \ \ \ \ \left(
\text{since }\omega\text{ is a ring morphism}\right)  .
\end{align*}
Since each $i,j\in\left[  q\right]  $ satisfy%
\begin{align*}
&  \omega\left(  \varphi^{\mu_{j}-j+1}h_{\lambda_{i}-\mu_{j}-i+j}\right) \\
&  =\varphi^{-\left(  \mu_{j}-j+1\right)  }e_{\lambda_{i}-\mu_{j}%
-i+j}\ \ \ \ \ \ \ \ \ \ \left(  \text{by (\ref{p37.omg})}\right) \\
&  =\varphi^{-\mu_{j}+j-1}e_{\lambda_{i}-\mu_{j}-i+j}%
\ \ \ \ \ \ \ \ \ \ \left(  \text{since }-\left(  \mu_{j}-j+1\right)
=-\mu_{j}+j-1\right)  ,
\end{align*}
we can rewrite this as%
\[
\omega s_{\theta}=\det\left(  \varphi^{-\mu_{j}+j-1}e_{\lambda_{i}-\mu
_{j}-i+j}\right)  _{i,j\in\left[  q\right]  }.
\]
Comparing this with (\ref{p37.9.6.pf.1}), we obtain%
\[
\omega s_{\theta}=s_{\lambda^{\prime}/\mu^{\prime}}=s_{\theta^{\prime}%
}\ \ \ \ \ \ \ \ \ \ \left(  \text{since }s_{\theta^{\prime}}=s_{\lambda
^{\prime}/\mu^{\prime}}\right)  .
\]
This proves (9.6).
\end{proof}

Next, we shall show that%
\begin{equation}
\omega\circ\varphi^{t}=\varphi^{-t}\circ\omega\ \ \ \ \ \ \ \ \ \ \text{for
each }t\in\mathbb{Z}. \label{p37.omg-commute}%
\end{equation}


\begin{proof}
[Proof of (\ref{p37.omg-commute}).]Both $\varphi$ and $\omega$ are
$\mathbb{Z}$-algebra morphisms (by their definitions), with $\varphi$ also
being invertible. Thus, $\omega\circ\varphi$ and $\varphi^{-1}\circ\omega$ are
$\mathbb{Z}$-algebra morphisms. These two $\mathbb{Z}$-algebra morphisms
$\omega\circ\varphi$ and $\varphi^{-1}\circ\omega$ agree on each of the
generators $h_{r,s}=\varphi^{s}h_{r}$ of the $\mathbb{Z}$-algebra $R$, since
each $r,s\in\mathbb{Z}$ satisfy%
\begin{align*}
\left(  \omega\circ\varphi\right)  \left(  \varphi^{s}h_{r}\right)   &
=\omega\left(  \underbrace{\varphi\left(  \varphi^{s}h_{r}\right)  }%
_{=\varphi^{s+1}h_{r}}\right)  =\omega\left(  \varphi^{s+1}h_{r}\right)
=\underbrace{\varphi^{-\left(  s+1\right)  }}_{=\varphi^{-1}\varphi^{-s}}%
e_{r}\ \ \ \ \ \ \ \ \ \ \left(  \text{by (\ref{p37.omg})}\right) \\
&  =\varphi^{-1}\underbrace{\left(  \varphi^{-s}e_{r}\right)  }%
_{\substack{=\omega\left(  \varphi^{s}h_{r}\right)  \\\text{(by (\ref{p37.omg}%
))}}}=\varphi^{-1}\left(  \omega\left(  \varphi^{s}h_{r}\right)  \right)
=\left(  \varphi^{-1}\circ\omega\right)  \left(  \varphi^{s}h_{r}\right)  .
\end{align*}
Hence, these two $\mathbb{Z}$-algebra morphisms $\omega\circ\varphi$ and
$\varphi^{-1}\circ\omega$ must be equal (since two $\mathbb{Z}$-algebra
morphisms that agree on each of the generators of their domain must be equal).
In other words,%
\begin{equation}
\omega\circ\varphi=\varphi^{-1}\circ\omega. \label{p37.omg-commute.pf.4}%
\end{equation}
Thus we can easily show that (\ref{p37.omg-commute}) holds for all nonnegative
integers $t$ (by induction on $t$; the base case is obvious, whereas the
induction step proceeds from the assumption $\omega\circ\varphi^{t}%
=\varphi^{-t}\circ\omega$ to argue that $\omega\circ\underbrace{\varphi^{t+1}%
}_{=\varphi\circ\varphi^{t}}=\underbrace{\omega\circ\varphi}_{=\varphi
^{-1}\circ\omega}\circ\,\varphi^{t}=\varphi^{-1}\circ\underbrace{\omega
\circ\varphi^{t}}_{=\varphi^{-t}\circ\omega}=\underbrace{\varphi^{-1}%
\circ\varphi^{-t}}_{=\varphi^{-\left(  t+1\right)  }}\circ\,\omega
=\varphi^{-\left(  t+1\right)  }\circ\omega$). But this easily entails that
(\ref{p37.omg-commute}) holds for all negative integers $t$ as well (because
if $t$ is negative, then $-t$ is nonnegative, and therefore
(\ref{p37.omg-commute}) holds for $-t$ instead of $t$ according to the
preceding sentence; but this means $\omega\circ\varphi^{-t}=\varphi^{-\left(
-t\right)  }\circ\omega=\varphi^{t}\circ\omega$, and therefore $\varphi
^{-t}\circ\underbrace{\omega\circ\varphi^{-t}}_{=\varphi^{t}\circ\omega}%
\circ\,\varphi^{t}=\underbrace{\varphi^{-t}\circ\varphi^{t}}%
_{=\operatorname*{id}}\circ\,\omega\circ\varphi^{t}=\omega\circ\varphi^{t}$,
so that $\omega\circ\varphi^{t}=\varphi^{-t}\circ\omega\circ
\underbrace{\varphi^{-t}\circ\varphi^{t}}_{=\operatorname*{id}}=\varphi
^{-t}\circ\omega$, which shows that (\ref{p37.omg-commute}) holds for $t$ as
well). Thus, (\ref{p37.omg-commute}) holds for all integers $t$ (nonnegative
and negative both), and the proof is complete.
\end{proof}

We can now finally prove that%
\begin{equation}
\omega\text{ is an involution.} \label{p37.omg-invol}%
\end{equation}


\begin{proof}
[Proof of (\ref{p37.omg-invol}).]We know that $\omega$ is a $\mathbb{Z}%
$-algebra morphism. Hence, $\omega\circ\omega$ is a $\mathbb{Z}$-algebra
morphism as well. We shall now show that $\omega\circ\omega$ fixes each of the
generators $h_{r,s}=\varphi^{s}h_{r}$ of the $\mathbb{Z}$-algebra $R$ (with
$r\geq1$).

Indeed, let $r,s\in\mathbb{Z}$ with $r\geq1$. Then, $h_{r}=s_{\left(
r\right)  }=s_{\left(  r\right)  /\varnothing}$ and $e_{r}=s_{\left(
1^{r}\right)  }=s_{\left(  1^{r}\right)  /\varnothing}$. But (9.6) (applied to
$\theta=\left(  1^{r}\right)  $) yields $s_{\left(  1^{r}\right)  ^{\prime}%
}=\omega s_{\left(  1^{r}\right)  }$. Hence, $\omega s_{\left(  1^{r}\right)
}=s_{\left(  1^{r}\right)  ^{\prime}}=s_{\left(  r\right)  }$ (since $\left(
1^{r}\right)  ^{\prime}=\left(  r\right)  $). In view of $h_{r}=s_{\left(
r\right)  }$ and $e_{r}=s_{\left(  1^{r}\right)  }$, we can rewrite this as
$\omega e_{r}=h_{r}$.

The definition of $\omega$ yields $\omega\left(  \varphi^{s}h_{r}\right)
=\varphi^{-s}e_{r}$. But (\ref{p37.omg-commute}) (applied to $t=-s$) yields
$\omega\circ\varphi^{-s}=\varphi^{-\left(  -s\right)  }\circ\omega=\varphi
^{s}\circ\omega$. Now,
\begin{align*}
\left(  \omega\circ\omega\right)  \left(  \varphi^{s}h_{r}\right)   &
=\omega\underbrace{\left(  \omega\left(  \varphi^{s}h_{r}\right)  \right)
}_{=\varphi^{-s}e_{r}}=\omega\left(  \varphi^{-s}e_{r}\right) \\
&  =\underbrace{\left(  \omega\circ\varphi^{-s}\right)  }_{=\varphi^{s}%
\circ\omega}\left(  e_{r}\right)  =\left(  \varphi^{s}\circ\omega\right)
\left(  e_{r}\right)  =\varphi^{s}\underbrace{\left(  \omega e_{r}\right)
}_{=h_{r}}=\varphi^{s}h_{r}.
\end{align*}
In other words, $\omega\circ\omega$ fixes $\varphi^{s}h_{r}$.

Forget that we fixed $r,s$. We thus have shown that the $\mathbb{Z}$-algebra
morphism $\omega\circ\omega$ fixes each of the generators $\varphi^{s}h_{r}$
of the $\mathbb{Z}$-algebra $R$. This shows that this morphism $\omega
\circ\omega$ must be the identity (because a $\mathbb{Z}$-algebra morphism
that fixes each of the generators of its domain must be the identity). In
other words, $\omega$ is an involution. This proves (\ref{p37.omg-invol}).
\end{proof}

\item \textbf{page 37, proof of (9.7):} \textquotedblleft in
which\textquotedblright\ should be \textquotedblleft in which
case\textquotedblright.

\item \textbf{page 37, proof of (9.7):} Again, let me add some more details.

The \textquotedblleft expansion of the determinant along the top
row\textquotedblright\ that proves the first displayed equation in this proof
is actually an application of Lemma \ref{lem.36.det-hess} to the $\left(
b+1\right)  \times\left(  b+1\right)  $-matrix%
\[
\left(
\begin{array}
[c]{cccccc}%
\varphi^{-0}\left(  h_{a+1}\right)  & \varphi^{-1}\left(  h_{a+2}\right)  &
\varphi^{-2}\left(  h_{a+3}\right)  & \varphi^{-3}\left(  h_{a+4}\right)  &
\cdots & \varphi^{-b}\left(  h_{a+b+1}\right) \\
1 & \varphi^{-1}\left(  h_{1}\right)  & \varphi^{-2}\left(  h_{2}\right)  &
\varphi^{-3}\left(  h_{3}\right)  & \cdots & \varphi^{-b}\left(  h_{b}\right)
\\
0 & 1 & \varphi^{-2}\left(  h_{1}\right)  & \varphi^{-3}\left(  h_{2}\right)
& \cdots & \varphi^{-b}\left(  h_{b-1}\right) \\
0 & 0 & 1 & \varphi^{-3}\left(  h_{1}\right)  & \cdots & \varphi^{-b}\left(
h_{b-2}\right) \\
\vdots & \vdots & \vdots & \vdots & \ddots & \vdots\\
0 & 0 & 0 & 0 & \cdots & \varphi^{-b}\left(  h_{1}\right)
\end{array}
\right)  ,
\]
whose determinant is $s_{\left(  a+1,1^{b}\right)  }=s_{\left(  a\mid
b\right)  }$.

Moreover, the same argument shows that%
\begin{equation}
s_{\left(  a\mid b\right)  }=0\ \ \ \ \ \ \ \ \ \ \text{for all }a\in\left[
-b,-1\right]  \label{p37.9.7.hook-greater}%
\end{equation}
(because in this case, $s_{\left(  a\mid b\right)  }$ is still the determinant
of the above $\left(  b+1\right)  \times\left(  b+1\right)  $-matrix, but now
the above $\left(  b+1\right)  \times\left(  b+1\right)  $-matrix has two
equal rows\footnote{Namely: its first row equals its $\left(  1-a\right)  $-th
row.} and therefore has determinant $0$). Combining this with the obvious fact
that%
\begin{equation}
s_{\left(  a\mid b\right)  }=0\ \ \ \ \ \ \ \ \ \ \text{for all }a<-1-b
\label{p37.9.7.hook-smaller}%
\end{equation}
(because $s_{\left(  a\mid b\right)  }$ is homogeneous of degree $a+b+1$,
which in this case is negative; or, alternatively, because all addends in the
sum defining $s_{\left(  a\mid b\right)  }$ are $0$), we conclude that%
\begin{equation}
s_{\left(  a\mid b\right)  }=0\ \ \ \ \ \ \ \ \ \ \text{for all negative
}a\neq-1-b. \label{p37.9.7.hook-neq}%
\end{equation}
Now, combining (\ref{p37.9.7.hook-neq}) with the easy fact that%
\begin{equation}
s_{\left(  a\mid b\right)  }=\left(  -1\right)  ^{b}%
\ \ \ \ \ \ \ \ \ \ \text{for }a=-1-b \label{p37.9.7.hook-equal}%
\end{equation}
(because all addends in the sum defining $s_{\left(  a\mid b\right)  }$ are
$0$ except for the addend for $i=b$, which is $\left(  -1\right)  ^{b}$), we
obtain a proof of the claim that $s_{\left(  a\mid b\right)  }=0$ for all
$a<0$ and $b\geq0$ except for the case $a=-1-b$, in which case $s_{\left(
a\mid b\right)  }=\left(  -1\right)  ^{b}$. In other words, we have proved
that%
\begin{equation}
s_{\left(  a\mid b\right)  }=\left(  -1\right)  ^{b}\delta_{a,-1-b}%
\ \ \ \ \ \ \ \ \ \ \text{for all }a<0\text{ and }b\geq0
\label{p37.9.7.hook-all}%
\end{equation}
(where the \textquotedblleft$\delta$\textquotedblright\ is a Kronecker delta).

\item \textbf{page 37, proof of (9.7):} Let me detail the argument in the last
sentence of this proof, specifically the part that reduces $\det\left(
s_{\left(  \lambda_{i}-i\mid n-k\right)  }\right)  _{1\leq i,k\leq n}$ to
$\det\left(  s_{\left(  \alpha_{i}\mid\beta_{j}\right)  }\right)  _{1\leq
i,j\leq r}$. That is, let me give a proof of
\begin{equation}
\det\left(  s_{\left(  \lambda_{i}-i\mid n-k\right)  }\right)  _{1\leq i,k\leq
n}=\det\left(  s_{\left(  \alpha_{i}\mid\beta_{j}\right)  }\right)  _{1\leq
i,j\leq r}. \label{p37.9.7.det=det}%
\end{equation}


First, let me show two lemmas about determinants:

\begin{lemmain}
\label{lem.p37.det1}Let $n\in\mathbb{N}$. Let $\left(  a_{i,j}\right)  _{1\leq
i,j\leq n}$ be an $n\times n$-matrix. Assume that its entries satisfy%
\begin{equation}
a_{i,j}=0\ \ \ \ \ \ \ \ \ \ \text{whenever }i+j\neq n+1.
\label{eq.lem.p37.det1.ass}%
\end{equation}
Then,%
\[
\det\left(  a_{i,j}\right)  _{1\leq i,j\leq n}=\left(  -1\right)  ^{\tbinom
{n}{2}}\prod_{i=1}^{n}a_{i,n+1-i}.
\]

\end{lemmain}

\begin{proof}
[Proof of Lemma \ref{lem.p37.det1}.]Let $w_{0}\in S_{n}$ be the permutation
that sends each $i\in\left[  n\right]  $ to $n+1-i$. Then, each of the
$\dbinom{n}{2}$ many pairs $\left(  i,j\right)  $ of elements of $\left[
n\right]  $ satisfying $i<j$ is an inversion of $w_{0}$ (since $i<j$ entails
$n+1-i>n+1-j$, that is, $w_{0}\left(  i\right)  >w_{0}\left(  j\right)  $).
Thus, the sign of this permutation $w_{0}$ is $\operatorname*{sign}\left(
w_{0}\right)  =\left(  -1\right)  ^{\tbinom{n}{2}}$.

It is well-known that, if we permute the columns of a matrix, then its
determinant is multiplied by the sign of the permutation used. Thus, in
particular, if we permute the columns of a matrix using the permutation
$w_{0}$, then its determinant is multiplied by $\operatorname*{sign}\left(
w_{0}\right)  =\left(  -1\right)  ^{\tbinom{n}{2}}$. Hence,%
\begin{equation}
\det\left(  a_{i,w_{0}\left(  j\right)  }\right)  _{1\leq i,j\leq n}=\left(
-1\right)  ^{\tbinom{n}{2}}\det\left(  a_{i,j}\right)  _{1\leq i,j\leq n}
\label{pf.lem.p37.det1.1}%
\end{equation}
(since the matrix $\left(  a_{i,w_{0}\left(  j\right)  }\right)  _{1\leq
i,j\leq n}$ is obtained from $\left(  a_{i,j}\right)  _{1\leq i,j\leq n}$ by
permuting the columns using the permutation $w_{0}$). However, the definition
of $w_{0}$ shows that $w_{0}\left(  j\right)  =n+1-j$ for each $j\in\left[
n\right]  $. Hence, $\left(  a_{i,w_{0}\left(  j\right)  }\right)  _{1\leq
i,j\leq n}=\left(  a_{i,n+1-j}\right)  _{1\leq i,j\leq n}$. Moreover,
(\ref{eq.lem.p37.det1.ass}) shows that the matrix $\left(  a_{i,n+1-j}\right)
_{1\leq i,j\leq n}$ is diagonal\footnote{\textit{Proof.} We must show that
$a_{i,n+1-j}=0$ for all $i,j\in\left[  n\right]  $ satisfying $i\neq j$.
\par
So let $i,j\in\left[  n\right]  $ satisfy $i\neq j$. Then, $i-j\neq0$, and now
$i+\left(  n+1-j\right)  =n+1+\underbrace{i-j}_{\neq0}\neq n+1$. Hence,
(\ref{eq.lem.p37.det1.ass}) (applied to $n+1-j$ instead of $j$) yields
$a_{i,n+1-j}=0$, qed.}; hence, its determinant equals the product of its
diagonal entries. That is,%
\[
\det\left(  a_{i,n+1-j}\right)  _{1\leq i,j\leq n}=\prod_{i=1}^{n}%
a_{i,n+1-i}.
\]
In view of $\left(  a_{i,w_{0}\left(  j\right)  }\right)  _{1\leq i,j\leq
n}=\left(  a_{i,n+1-j}\right)  _{1\leq i,j\leq n}$, we can rewrite this as%
\[
\det\left(  a_{i,w_{0}\left(  j\right)  }\right)  _{1\leq i,j\leq n}%
=\prod_{i=1}^{n}a_{i,n+1-i}.
\]
Comparing this with (\ref{pf.lem.p37.det1.1}), we obtain%
\[
\left(  -1\right)  ^{\tbinom{n}{2}}\det\left(  a_{i,j}\right)  _{1\leq i,j\leq
n}=\prod_{i=1}^{n}a_{i,n+1-i}.
\]
Solving this for $\det\left(  a_{i,j}\right)  _{1\leq i,j\leq n}$, we find%
\[
\det\left(  a_{i,j}\right)  _{1\leq i,j\leq n}=\underbrace{\dfrac{1}{\left(
-1\right)  ^{\tbinom{n}{2}}}}_{\substack{=\left(  -1\right)  ^{\tbinom{n}{2}%
}\\\text{(since }\left(  -1\right)  ^{\tbinom{n}{2}}\in\left\{  1,-1\right\}
\text{)}}}\prod_{i=1}^{n}a_{i,n+1-i}=\left(  -1\right)  ^{\tbinom{n}{2}}%
\prod_{i=1}^{n}a_{i,n+1-i}.
\]
This proves Lemma \ref{lem.p37.det1}.
\end{proof}

\begin{lemmain}
\label{lem.p37.det2}Let $n\in\mathbb{N}$ and $r\in\left[  0,n\right]  $. Let
$C=\left(  c_{i,j}\right)  _{1\leq i,j\leq n}$ be an $n\times n$-matrix. Let
$\omega_{r+1},\omega_{r+2},\ldots,\omega_{n}$ be $n-r$ elements of $\left[
n\right]  $ such that $\omega_{r+1}>\omega_{r+2}>\cdots>\omega_{n}$. Assume
that the entries of $C$ satisfy%
\begin{equation}
c_{i,j}=0\ \ \ \ \ \ \ \ \ \ \text{for all }i>r\text{ and }j\neq\omega_{i}.
\label{eq.lem.p37.det2.ass}%
\end{equation}
Set $Q_{1}:=\left\{  \omega_{r+1},\omega_{r+2},\ldots,\omega_{n}\right\}  $.
Then,%
\[
\det C=\left(  -1\right)  ^{\left(  n-r\right)  n+\left(  \omega_{r+1}%
+\omega_{r+2}+\cdots+\omega_{n}\right)  }\left(  \prod_{i=r+1}^{n}%
c_{i,\omega_{i}}\right)  \det\left(  \operatorname*{sub}\nolimits_{\left[
r\right]  }^{\left[  n\right]  \setminus Q_{1}}C\right)  .
\]

\end{lemmain}

\begin{statement}
\textit{Example:} For $n=5$ and $r=3$ and $\omega_{4}=2$ and $\omega_{5}=4$,
Lemma \ref{lem.p37.det2} says that%
\begin{align*}
&  \det\left(
\begin{array}
[c]{ccccc}%
c_{1,1} & c_{1,2} & c_{1,3} & c_{1,4} & c_{1,5}\\
c_{2,1} & c_{2,2} & c_{2,3} & c_{2,4} & c_{2,5}\\
c_{3,1} & c_{3,2} & c_{3,3} & c_{3,4} & c_{3,5}\\
0 & c_{4,2} & 0 & 0 & 0\\
0 & 0 & 0 & c_{5,4} & 0
\end{array}
\right) \\
&  =\left(  -1\right)  ^{\left(  5-3\right)  \cdot5+\left(  2+4\right)
}\left(  c_{4,2}c_{5,4}\right)  \det\left(
\begin{array}
[c]{ccc}%
c_{1,1} & c_{1,3} & c_{1,5}\\
c_{2,1} & c_{2,3} & c_{2,5}\\
c_{3,1} & c_{3,3} & c_{3,5}%
\end{array}
\right)  .
\end{align*}

\end{statement}

\begin{proof}
[Proof of Lemma \ref{lem.p37.det2}.]This is easy to see by Laplace expansion
along the last $n-r$ rows of the matrix $C$. What makes the proof laborious
and painful is the bookkeeping involved in getting the signs right. Here are
the details, for the sake of completeness:

Let $P=\left[  r\right]  $. This is a subset of $\left[  n\right]  $ (since
$r\in\left[  0,n\right]  $ entails $r\leq n$), and has size $\left\vert
P\right\vert =\left\vert \left[  r\right]  \right\vert =r$.

For any subset $I$ of $\left[  n\right]  $, we let $\sum I$ denote the sum of
all elements of $I$, and we let $\widetilde{I}$ denote the complement $\left[
n\right]  \setminus I$ of $I$. (These are the same notations we used in Lemma
\ref{lem.lap-rows}, but now relative to $n$ rather than to $k$.) Then, Lemma
\ref{lem.lap-rows} (applied to $n$ and $C$ instead of $k$ and $A$) yields%
\begin{equation}
\det C=\sum_{\substack{Q\subseteq\left[  n\right]  ;\\\left\vert Q\right\vert
=\left\vert P\right\vert }}\left(  -1\right)  ^{\sum P+\sum Q}\det\left(
\operatorname*{sub}\nolimits_{P}^{Q}C\right)  \det\left(  \operatorname*{sub}%
\nolimits_{\widetilde{P}}^{\widetilde{Q}}C\right)
\ \ \ \ \ \ \ \ \ \ \label{pf.lem.p37.det2.detC=}%
\end{equation}
(since $P$ is a subset of $\left[  n\right]  $). Note that the definition of
$\widetilde{P}$ yields%
\begin{align*}
\widetilde{P}  &  =\left[  n\right]  \setminus\underbrace{P}_{=\left[
r\right]  }=\left[  n\right]  \setminus\left[  r\right]  =\left[
r+1,n\right]  =\left\{  r+1,r+2,\ldots,n\right\} \\
&  =\left\{  r+1<r+2<\cdots<n\right\}  .
\end{align*}


Recall that%
\begin{align*}
Q_{1}  &  =\left\{  \omega_{r+1},\omega_{r+2},\ldots,\omega_{n}\right\} \\
&  =\left\{  \omega_{r+1}>\omega_{r+2}>\cdots>\omega_{n}\right\}
\ \ \ \ \ \ \ \ \ \ \left(  \text{since }\omega_{r+1}>\omega_{r+2}%
>\cdots>\omega_{n}\right) \\
&  =\left\{  \omega_{n}<\omega_{n-1}<\cdots<\omega_{r+1}\right\}  ,
\end{align*}
so that $\left\vert Q_{1}\right\vert =n-r$. Moreover, $Q_{1}=\left\{
\omega_{r+1},\omega_{r+2},\ldots,\omega_{n}\right\}  \subseteq\left[
n\right]  $ (since $\omega_{r+1},\omega_{r+2},\ldots,\omega_{n}$ are elements
of $\left[  n\right]  $).

Set%
\[
Q_{0}:=\widetilde{Q_{1}}.
\]
Thus, $Q_{0}=\widetilde{Q_{1}}=\left[  n\right]  \setminus Q_{1}$ (by the
definition of $\widetilde{Q_{1}}$), whence
\begin{align*}
\left\vert Q_{0}\right\vert  &  =\left\vert \left[  n\right]  \setminus
Q_{1}\right\vert =\underbrace{\left\vert \left[  n\right]  \right\vert }%
_{=n}-\underbrace{\left\vert Q_{1}\right\vert }_{=n-r}%
\ \ \ \ \ \ \ \ \ \ \left(  \text{since }Q_{1}\subseteq\left[  n\right]
\right) \\
&  =n-\left(  n-r\right)  =r=\left\vert P\right\vert .
\end{align*}


Furthermore, from $Q_{0}=\left[  n\right]  \setminus Q_{1}$, we obtain%
\begin{align*}
\sum Q_{0}  &  =\sum\left(  \left[  n\right]  \setminus Q_{1}\right) \\
&  =\sum\left[  n\right]  -\sum Q_{1}\ \ \ \ \ \ \ \ \ \ \left(  \text{since
}Q_{1}\subseteq\left[  n\right]  \right) \\
&  \equiv\underbrace{\sum\left[  n\right]  }_{=1+2+\cdots+n}+\underbrace{\sum
Q_{1}}_{\substack{=\omega_{r+1}+\omega_{r+2}+\cdots+\omega_{n}\\\text{(since
}Q_{1}=\left\{  \omega_{r+1}>\omega_{r+2}>\cdots>\omega_{n}\right\}  \text{)}%
}}\\
&  \ \ \ \ \ \ \ \ \ \ \ \ \ \ \ \ \ \ \ \ \left(  \text{since }x-y\equiv
x+y\operatorname{mod}2\text{ for all }x,y\in\mathbb{Z}\right) \\
&  =\left(  1+2+\cdots+n\right)  +\left(  \omega_{r+1}+\omega_{r+2}%
+\cdots+\omega_{n}\right)  \operatorname{mod}2.
\end{align*}
Also, from $P=\left[  r\right]  $, we obtain $\sum P=\sum\left[  r\right]
=1+2+\cdots+r$.

Recall that $Q_{0}=\left[  n\right]  \setminus Q_{1}\subseteq\left[  n\right]
$ and $\left\vert Q_{0}\right\vert =\left\vert P\right\vert $. Hence, the sum
on the right hand side of (\ref{pf.lem.p37.det2.detC=}) has an addend for
$Q=Q_{0}$. We shall now show that all its other addends are $0$. Slightly
better, we will show that each subset $Q$ of $\left[  n\right]  $ satisfying
$\left\vert Q\right\vert =\left\vert P\right\vert $ and $Q\neq Q_{0}$
satisfies%
\begin{equation}
\det\left(  \operatorname*{sub}\nolimits_{\widetilde{P}}^{\widetilde{Q}%
}C\right)  =0. \label{pf.lem.p37.det2.=0}%
\end{equation}


\begin{proof}
[Proof of (\ref{pf.lem.p37.det2.=0}).]Let $Q$ be a subset of $\left[
n\right]  $ satisfying $\left\vert Q\right\vert =\left\vert P\right\vert $ and
$Q\neq Q_{0}$.

We have $\left\vert Q\right\vert =\left\vert P\right\vert =r=\left\vert
Q_{0}\right\vert $ (since $\left\vert Q_{0}\right\vert =r$). That is, the sets
$Q$ and $Q_{0}$ have the same size. If we had $Q_{0}\subseteq Q$, then we
would thus have $Q_{0}=Q$, which would contradict $Q\neq Q_{0}$. Hence, we
don't have $Q_{0}\subseteq Q$. Thus, there exists some $q\in Q_{0}$ such that
$q\notin Q$. Consider this $q$. From $q\in Q_{0}\subseteq\left[  n\right]  $
and $q\notin Q$, we obtain $q\in\left[  n\right]  \setminus Q=\widetilde{Q}$
(by the definition of $\widetilde{Q}$). From $Q\subseteq\left[  n\right]  $,
we also obtain $\left\vert \left[  n\right]  \setminus Q\right\vert
=\underbrace{\left\vert \left[  n\right]  \right\vert }_{=n}%
-\underbrace{\left\vert Q\right\vert }_{=r}=n-r$. In other words, $\left\vert
\widetilde{Q}\right\vert =n-r$ (since $\widetilde{Q}=\left[  n\right]
\setminus Q$).

We have $\widetilde{P}=\left\{  r+1<r+2<\cdots<n\right\}  $. Write
$\widetilde{Q}$ as $\widetilde{Q}=\left\{  q_{1}<q_{2}<\cdots<q_{n-r}\right\}
$ (we can do this, since $\left\vert \widetilde{Q}\right\vert =n-r$). Using
these two equalities, we obtain%
\begin{align}
\operatorname*{sub}\nolimits_{\widetilde{P}}^{\widetilde{Q}}C  &
=\operatorname*{sub}\nolimits_{\left\{  r+1<r+2<\cdots<n\right\}  }^{\left\{
q_{1}<q_{2}<\cdots<q_{n-r}\right\}  }C\nonumber\\
&  =\left(  c_{r+i,q_{j}}\right)  _{1\leq i,j\leq n-r}
\label{pf.lem.p37.det2.=0.pf.3}%
\end{align}
(since $C=\left(  c_{i,j}\right)  _{1\leq i,j\leq n}$). Also, we have
$q\in\widetilde{Q}=\left\{  q_{1}<q_{2}<\cdots<q_{n-r}\right\}  $, so that
$q=q_{k}$ for some $k\in\left[  n-r\right]  $. Consider this $k$. Thus,
$q_{k}=q\in Q_{0}=\left[  n\right]  \setminus Q_{1}$, so that $q_{k}\in\left[
n\right]  $ and $q_{k}\notin Q_{1}$.

Let $i\in\left[  n-r\right]  $ be arbitrary. Then, $q_{k}\neq\omega_{r+i}$
(because $q_{k}=\omega_{r+i}$ would entail $q_{k}=\omega_{r+i}\in\left\{
\omega_{r+1},\omega_{r+2},\ldots,\omega_{n}\right\}  =Q_{1}$, contradicting
$q_{k}\notin Q_{1}$). Therefore, (\ref{eq.lem.p37.det2.ass}) (applied to $r+i$
and $q_{k}$ instead of $i$ and $j$) yields $c_{r+i,q_{k}}=0$ (since $r+i>r$).

Forget that we fixed $i$. We thus have shown that $c_{r+i,q_{k}}=0$ for each
$i\in\left[  n-r\right]  $. In other words, the whole $k$-th column of the
matrix $\left(  c_{r+i,q_{j}}\right)  _{1\leq i,j\leq n-r}$ is $0$. Hence,
this matrix has a zero column, and therefore has determinant $0$. In other
words, $\det\left(  c_{r+i,q_{j}}\right)  _{1\leq i,j\leq n-r}=0$. In view of
(\ref{pf.lem.p37.det2.=0.pf.3}), we can rewrite this as $\det\left(
\operatorname*{sub}\nolimits_{\widetilde{P}}^{\widetilde{Q}}C\right)  =0$.
This proves (\ref{pf.lem.p37.det2.=0}).
\end{proof}

Now, (\ref{pf.lem.p37.det2.detC=}) becomes%
\begin{align}
\det C  &  =\sum_{\substack{Q\subseteq\left[  n\right]  ;\\\left\vert
Q\right\vert =\left\vert P\right\vert }}\left(  -1\right)  ^{\sum P+\sum
Q}\det\left(  \operatorname*{sub}\nolimits_{P}^{Q}C\right)  \det\left(
\operatorname*{sub}\nolimits_{\widetilde{P}}^{\widetilde{Q}}C\right)
\nonumber\\
&  =\left(  -1\right)  ^{\sum P+\sum Q_{0}}\det\left(  \operatorname*{sub}%
\nolimits_{P}^{Q_{0}}C\right)  \det\left(  \operatorname*{sub}%
\nolimits_{\widetilde{P}}^{\widetilde{Q_{0}}}C\right) \nonumber\\
&  \ \ \ \ \ \ \ \ \ \ +\sum_{\substack{Q\subseteq\left[  n\right]
;\\\left\vert Q\right\vert =\left\vert P\right\vert ;\\Q\neq Q_{0}}}\left(
-1\right)  ^{\sum P+\sum Q}\det\left(  \operatorname*{sub}\nolimits_{P}%
^{Q}C\right)  \underbrace{\det\left(  \operatorname*{sub}%
\nolimits_{\widetilde{P}}^{\widetilde{Q}}C\right)  }_{\substack{=0\\\text{(by
(\ref{pf.lem.p37.det2.=0}))}}}\nonumber\\
&  \ \ \ \ \ \ \ \ \ \ \ \ \ \ \ \ \ \ \ \ \left(
\begin{array}
[c]{c}%
\text{here, we have split off the addend for }Q=Q_{0}\\
\text{from the sum, since we know that this addend exists}%
\end{array}
\right) \nonumber\\
&  =\left(  -1\right)  ^{\sum P+\sum Q_{0}}\det\left(  \operatorname*{sub}%
\nolimits_{P}^{Q_{0}}C\right)  \det\left(  \operatorname*{sub}%
\nolimits_{\widetilde{P}}^{\widetilde{Q_{0}}}C\right) \nonumber\\
&  \ \ \ \ \ \ \ \ \ \ +\underbrace{\sum_{\substack{Q\subseteq\left[
n\right]  ;\\\left\vert Q\right\vert =\left\vert P\right\vert ;\\Q\neq Q_{0}%
}}\left(  -1\right)  ^{\sum P+\sum Q}\det\left(  \operatorname*{sub}%
\nolimits_{P}^{Q}C\right)  0}_{=0}\nonumber\\
&  =\left(  -1\right)  ^{\sum P+\sum Q_{0}}\det\left(  \operatorname*{sub}%
\nolimits_{P}^{Q_{0}}C\right)  \det\left(  \operatorname*{sub}%
\nolimits_{\widetilde{P}}^{\widetilde{Q_{0}}}C\right)  .
\label{pf.lem.p37.det2.detC=2}%
\end{align}


Next, we claim that%
\begin{equation}
\det\left(  \operatorname*{sub}\nolimits_{\widetilde{P}}^{\widetilde{Q_{0}}%
}C\right)  =\left(  -1\right)  ^{\tbinom{n-r}{2}}\prod_{i=r+1}^{n}%
c_{i,\omega_{i}}. \label{pf.lem.p37.det2.5}%
\end{equation}


\begin{proof}
[Proof of (\ref{pf.lem.p37.det2.5}).]The definition of $\widetilde{Q_{0}}$
yields%
\begin{align*}
\widetilde{Q_{0}}  &  =\left[  n\right]  \setminus\underbrace{Q_{0}}_{=\left[
n\right]  \setminus Q_{1}}=\left[  n\right]  \setminus\left(  \left[
n\right]  \setminus Q_{1}\right)  =Q_{1}\ \ \ \ \ \ \ \ \ \ \left(
\text{since }Q_{1}\subseteq\left[  n\right]  \right) \\
&  =\left\{  \omega_{n}<\omega_{n-1}<\cdots<\omega_{r+1}\right\}  .
\end{align*}
Using this equality as well as $\widetilde{P}=\left\{  r+1<r+2<\cdots
<n\right\}  $, we obtain%
\begin{align}
\operatorname*{sub}\nolimits_{\widetilde{P}}^{\widetilde{Q_{0}}}C  &
=\operatorname*{sub}\nolimits_{\left\{  r+1<r+2<\cdots<n\right\}  }^{\left\{
\omega_{n}<\omega_{n-1}<\cdots<\omega_{r+1}\right\}  }C\nonumber\\
&  =\left(  c_{r+i,\omega_{n+1-j}}\right)  _{1\leq i,j\leq n-r}%
\ \ \ \ \ \ \ \ \ \ \left(  \text{since }C=\left(  c_{i,j}\right)  _{1\leq
i,j\leq n}\right)  .\nonumber
\end{align}


However, using (\ref{eq.lem.p37.det2.ass}), it is easy to see that the entries
of this matrix satisfy%
\[
c_{r+i,\omega_{n+1-j}}=0\ \ \ \ \ \ \ \ \ \ \text{whenever }i,j\in\left[
n-r\right]  \text{ satisfy }i+j\neq n-r+1
\]
\footnote{\textit{Proof.} Let $i,j\in\left[  n-r\right]  $ be such that
$i+j\neq n-r+1$. We must prove that $c_{r+i,\omega_{n+1-j}}=0$.
\par
From $i+j\neq n-r+1$, we obtain $r+i+j\neq n+1$, thus $r+i\neq n+1-j$. In
other words, $n+1-j\neq r+i$.
\par
We have $\omega_{r+1}>\omega_{r+2}>\cdots>\omega_{n}$. Thus, the numbers
$\omega_{r+1},\omega_{r+2},\ldots,\omega_{n}$ are distinct. Hence, from
$n+1-j\neq r+i$, we obtain $\omega_{n+1-j}\neq\omega_{r+i}$.
\par
Thus, we have $r+i>r$ and $\omega_{n+1-j}\neq\omega_{r+i}$. Hence,
(\ref{eq.lem.p37.det2.ass}) (applied to $r+i$ and $\omega_{n+1-j}$ instead of
$i$ and $j$) yields $c_{r+i,\omega_{n+1-j}}=0$, qed.}. Thus, Lemma
\ref{lem.p37.det1} (applied to $n-r$ and $c_{r+i,\omega_{n+1-j}}$ instead of
$n$ and $a_{i,j}$) yields
\begin{align*}
\det\left(  c_{r+i,\omega_{n+1-j}}\right)  _{1\leq i,j\leq n-r}  &  =\left(
-1\right)  ^{\tbinom{n-r}{2}}\prod_{i=1}^{n-r}\underbrace{c_{r+i,\omega
_{n+1-\left(  n-r+1-i\right)  }}}_{\substack{=c_{r+i,\omega_{r+i}%
}\\\text{(since }n+1-\left(  n-r+1-i\right)  =r+i\text{)}}}\\
&  =\left(  -1\right)  ^{\tbinom{n-r}{2}}\prod_{i=1}^{n-r}c_{r+i,\omega_{r+i}%
}=\left(  -1\right)  ^{\tbinom{n-r}{2}}\prod_{i=r+1}^{n}c_{i,\omega_{i}}%
\end{align*}
(here, we have substituted $i$ for $r+i$ in the product). Since
$\operatorname*{sub}\nolimits_{\widetilde{P}}^{\widetilde{Q_{0}}}C=\left(
c_{r+i,\omega_{n+1-j}}\right)  _{1\leq i,j\leq n-r}$, we can rewrite this as%
\[
\det\left(  \operatorname*{sub}\nolimits_{\widetilde{P}}^{\widetilde{Q_{0}}%
}C\right)  =\left(  -1\right)  ^{\tbinom{n-r}{2}}\prod_{i=r+1}^{n}%
c_{i,\omega_{i}}.
\]
This proves (\ref{pf.lem.p37.det2.5}).
\end{proof}

Now, (\ref{pf.lem.p37.det2.detC=2}) becomes%
\begin{align}
\det C  &  =\left(  -1\right)  ^{\sum P+\sum Q_{0}}\det\left(
\operatorname*{sub}\nolimits_{P}^{Q_{0}}C\right)  \underbrace{\det\left(
\operatorname*{sub}\nolimits_{\widetilde{P}}^{\widetilde{Q_{0}}}C\right)
}_{\substack{=\left(  -1\right)  ^{\tbinom{n-r}{2}}\prod_{i=r+1}%
^{n}c_{i,\omega_{i}}\\\text{(by (\ref{pf.lem.p37.det2.5}))}}}\nonumber\\
&  =\left(  -1\right)  ^{\sum P+\sum Q_{0}}\det\left(  \operatorname*{sub}%
\nolimits_{P}^{Q_{0}}C\right)  \left(  -1\right)  ^{\tbinom{n-r}{2}}%
\prod_{i=r+1}^{n}c_{i,\omega_{i}}\nonumber\\
&  =\underbrace{\left(  -1\right)  ^{\sum P+\sum Q_{0}}\left(  -1\right)
^{\tbinom{n-r}{2}}}_{=\left(  -1\right)  ^{\sum P+\sum Q_{0}+\tbinom{n-r}{2}}%
}\left(  \prod_{i=r+1}^{n}c_{i,\omega_{i}}\right)  \det\left(
\operatorname*{sub}\nolimits_{P}^{Q_{0}}C\right) \nonumber\\
&  =\left(  -1\right)  ^{\sum P+\sum Q_{0}+\tbinom{n-r}{2}}\left(
\prod_{i=r+1}^{n}c_{i,\omega_{i}}\right)  \det\left(  \operatorname*{sub}%
\nolimits_{P}^{Q_{0}}C\right)
.\ \ \ \ \ \ \ \ \ \ \ \ \ \label{pf.lem.p37.det2.detC=7}%
\end{align}


Now, let us simplify the sign on the right hand side. Indeed,%
\begin{align*}
&  \sum P+\sum Q_{0}+\dbinom{n-r}{2}\\
&  =\underbrace{\dbinom{n-r}{2}}_{\substack{=0+1+\cdots+\left(  n-r-1\right)
\\\text{(since }\dbinom{u}{2}=0+1+\cdots+\left(  u-1\right)  \\\text{for each
}u\in\mathbb{N}\text{)}}}+\underbrace{\sum P}_{\substack{=1+2+\cdots
+r\\\equiv-\left(  1+2+\cdots+r\right)  \operatorname{mod}2\\\text{(since
}u\equiv-u\operatorname{mod}2\text{ for each }u\in\mathbb{Z}\text{)}}}\\
&  \ \ \ \ \ \ \ \ \ \ +\underbrace{\sum Q_{0}}_{\equiv\left(  1+2+\cdots
+n\right)  +\left(  \omega_{r+1}+\omega_{r+2}+\cdots+\omega_{n}\right)
\operatorname{mod}2}\\
&  \equiv\underbrace{\left(  0+1+\cdots+\left(  n-r-1\right)  \right)
}_{=\sum_{k=0}^{n-r-1}k}+\underbrace{\left(  -\left(  1+2+\cdots+r\right)
\right)  +\left(  1+2+\cdots+n\right)  }_{\substack{=\left(  1+2+\cdots
+n\right)  -\left(  1+2+\cdots+r\right)  \\=\left(  r+1\right)  +\left(
r+2\right)  +\cdots+n\\=n+\left(  n-1\right)  +\left(  n-2\right)
+\cdots+\left(  r+1\right)  \\=\sum_{k=0}^{n-r-1}\left(  n-k\right)  }}\\
&  \ \ \ \ \ \ \ \ \ \ +\left(  \omega_{r+1}+\omega_{r+2}+\cdots+\omega
_{n}\right) \\
&  =\underbrace{\sum_{k=0}^{n-r-1}k+\sum_{k=0}^{n-r-1}\left(  n-k\right)
}_{=\sum_{k=0}^{n-r-1}\left(  k+\left(  n-k\right)  \right)  }+\left(
\omega_{r+1}+\omega_{r+2}+\cdots+\omega_{n}\right) \\
&  =\sum_{k=0}^{n-r-1}\underbrace{\left(  k+\left(  n-k\right)  \right)
}_{=n}+\left(  \omega_{r+1}+\omega_{r+2}+\cdots+\omega_{n}\right) \\
&  =\underbrace{\sum_{k=0}^{n-r-1}n}_{=\left(  n-r\right)  n}+\left(
\omega_{r+1}+\omega_{r+2}+\cdots+\omega_{n}\right) \\
&  =\left(  n-r\right)  n+\left(  \omega_{r+1}+\omega_{r+2}+\cdots+\omega
_{n}\right)  \operatorname{mod}2,
\end{align*}
and therefore%
\[
\left(  -1\right)  ^{\sum P+\sum Q_{0}+\tbinom{n-r}{2}}=\left(  -1\right)
^{\left(  n-r\right)  n+\left(  \omega_{r+1}+\omega_{r+2}+\cdots+\omega
_{n}\right)  }.
\]
Hence, (\ref{pf.lem.p37.det2.detC=7}) rewrites as%
\begin{align*}
\det C  &  =\left(  -1\right)  ^{\left(  n-r\right)  n+\left(  \omega
_{r+1}+\omega_{r+2}+\cdots+\omega_{n}\right)  }\left(  \prod_{i=r+1}%
^{n}c_{i,\omega_{i}}\right)  \det\left(  \operatorname*{sub}\nolimits_{P}%
^{Q_{0}}C\right) \\
&  =\left(  -1\right)  ^{\left(  n-r\right)  n+\left(  \omega_{r+1}%
+\omega_{r+2}+\cdots+\omega_{n}\right)  }\left(  \prod_{i=r+1}^{n}%
c_{i,\omega_{i}}\right)  \det\left(  \operatorname*{sub}\nolimits_{\left[
r\right]  }^{\left[  n\right]  \setminus Q_{1}}C\right)
\end{align*}
(since $P=\left[  r\right]  $ and $Q_{0}=\left[  n\right]  \setminus Q_{1}$).
This proves Lemma \ref{lem.p37.det2}.
\end{proof}

\begin{proof}
[Proof of (\ref{p37.9.7.det=det}).]Recall that $r$ is defined to be the
largest integer $i\geq1$ such that the cell $\left(  i,i\right)  $ lies in the
Young diagram of $\lambda$. Thus, the cells $\left(  1,1\right)  ,\ \left(
2,2\right)  ,\ \ldots,$ $\left(  r,r\right)  $ lie in the Young diagram of
$\lambda$, but the cells \newline$\left(  r+1,r+1\right)  ,\ \left(
r+2,r+2\right)  ,\ \left(  r+3,r+3\right)  ,\ \ldots$ do not.

Thus, it is easy to see that $r\leq n$\ \ \ \ \footnote{\textit{Proof.} Assume
the contrary. Thus, $r>n$, so that $r>n\geq\ell\left(  \lambda\right)  $.
Hence, $\lambda_{r}=0$. But $r$ is defined to be the largest integer $i\geq1$
such that the cell $\left(  i,i\right)  $ lies in the Young diagram of
$\lambda$. Thus, the cell $\left(  r,r\right)  $ lies in the Young diagram of
$\lambda$ (since $r>n\geq0$). In other words, $r\leq\lambda_{r}$. Hence,
$r\leq\lambda_{r}=0$, which contradicts $r>n\geq0$. This contradiction shows
that our assumption was wrong; thus, $r\leq n$ is proved.}, so that
$r\in\left[  0,n\right]  $.

The numbers $\alpha_{i}$ and $\beta_{i}$ are defined for all $i\in\left[
r\right]  $, but we shall now extend their definition to all positive integers
$i$. Namely, for each positive integer $i$, we set%
\[
\alpha_{i}:=\lambda_{i}-i\ \ \ \ \ \ \ \ \ \ \text{and}%
\ \ \ \ \ \ \ \ \ \ \beta_{i}:=\lambda_{i}^{\prime}-i
\]
(just as we previously did for $i\in\left[  r\right]  $). Let us furthermore
set%
\[
\omega_{i}:=n+1+\alpha_{i}\ \ \ \ \ \ \ \ \ \ \text{for each positive integer
}i.
\]


Since $\lambda$ is a partition, we have $\lambda_{1}\geq\lambda_{2}\geq
\lambda_{3}\geq\cdots$, and thus $\lambda_{1}-1>\lambda_{2}-2>\lambda
_{3}-3>\cdots$. The latter chain of inequalities rewrites as%
\begin{equation}
\alpha_{1}>\alpha_{2}>\alpha_{3}>\cdots\label{p37.9.7.det=det.pf.alpha-ineq}%
\end{equation}
(since each $i\geq1$ satisfies $\alpha_{i}=\lambda_{i}-i$). Similarly,%
\begin{equation}
\beta_{1}>\beta_{2}>\beta_{3}>\cdots\label{p37.9.7.det=det.pf.beta-ineq}%
\end{equation}
(since $\lambda^{\prime}$ is a partition). Furthermore, note that%
\begin{equation}
\alpha_{i}<0\ \ \ \ \ \ \ \ \ \ \text{for each }i>r
\label{p37.9.7.det=det.pf.alpha-neg}%
\end{equation}
\footnote{\textit{Proof.} Let $i>r$ be an integer. Thus, $i\in\left\{
r+1,r+2,r+3,\ldots\right\}  $. Hence, the cell $\left(  i,i\right)  $ does not
lie in the Young diagram of $\lambda$ (since the cells $\left(
r+1,r+1\right)  ,\ \left(  r+2,r+2\right)  ,\ \left(  r+3,r+3\right)
,\ \ldots$ do not lie in the Young diagram of $\lambda$). In other words,
$i>\lambda_{i}$. But the definition of $\alpha_{i}$ says $\alpha_{i}%
=\lambda_{i}-i<0$ (since $i>\lambda_{i}$). This proves
(\ref{p37.9.7.det=det.pf.alpha-neg}).}.

For any two positive integers $i$ and $j$, we have $\lambda_{i}+\lambda
_{j}^{\prime}-i-j\neq1$ (by \cite[Proposition 3.18 \textbf{(c)}]{petriesym},
applied to $\mu=\lambda^{\prime}$) and thus%
\begin{align}
\underbrace{\alpha_{i}}_{\substack{=\lambda_{i}-i\\\text{(by the definition of
}\alpha_{i}\text{)}}}+\underbrace{\beta_{j}}_{\substack{=\lambda_{j}^{\prime
}-j\\\text{(by the definition of }\beta_{j}\text{)}}}  &  =\lambda
_{i}-i+\lambda_{j}^{\prime}-j=\lambda_{i}+\lambda_{j}^{\prime}-i-j\nonumber\\
&  \neq1. \label{p37.9.7.det=det.pf.not1}%
\end{align}


We have $\beta_{1}>\beta_{2}>\cdots>\beta_{r}$ (by
(\ref{p37.9.7.det=det.pf.beta-ineq})) and thus
\[
n-\beta_{1}<n-\beta_{2}<\cdots<n-\beta_{r}.
\]
Thus, let us define a set%
\begin{align*}
Q_{0}:=  &  \ \left\{  n-\beta_{1}<n-\beta_{2}<\cdots<n-\beta_{r}\right\} \\
=  &  \ \left\{  n-\beta_{j}\ \mid\ j\in\left[  r\right]  \right\}  .
\end{align*}
Clearly, $\left\vert Q_{0}\right\vert =r$, since $Q_{0}=\left\{  n-\beta
_{1}<n-\beta_{2}<\cdots<n-\beta_{r}\right\}  $.

Furthermore, we have $\alpha_{r+1}>\alpha_{r+2}>\cdots>\alpha_{n}$ (by
(\ref{p37.9.7.det=det.pf.alpha-ineq})) and thus%
\[
n+1+\alpha_{r+1}>n+1+\alpha_{r+2}>\cdots>n+1+\alpha_{n}.
\]
This chain of inequalities rewrites as%
\[
\omega_{r+1}>\omega_{r+2}>\cdots>\omega_{n}%
\]
(since $\omega_{i}=n+1+\alpha_{i}$ for each positive integer $i$). Let us
define a set%
\[
Q_{1}:=\left\{  \omega_{r+1},\omega_{r+2},\ldots,\omega_{n}\right\}  =\left\{
\omega_{i}\ \mid\ i\in\left[  r+1,n\right]  \right\}  .
\]
Then,
\[
Q_{1}=\left\{  \omega_{r+1},\omega_{r+2},\ldots,\omega_{n}\right\}  =\left\{
\omega_{r+1}>\omega_{r+2}>\cdots>\omega_{n}\right\}
\]
(since we know that $\omega_{r+1}>\omega_{r+2}>\cdots>\omega_{n}$), and thus
$\left\vert Q_{1}\right\vert =n-r$.

From (\ref{p37.9.7.det=det.pf.not1}), we see that the sets $Q_{0}$ and $Q_{1}$
are disjoint\footnote{\textit{Proof.} Assume the contrary. Thus, there exists
some $k\in Q_{0}\cap Q_{1}$. Consider this $k$. Then,%
\[
k\in Q_{0}\cap Q_{1}\subseteq Q_{1}=\left\{  \omega_{i}\ \mid\ i\in\left[
r+1,n\right]  \right\}  .
\]
That is, $k=\omega_{i}$ for some $i\in\left[  r+1,n\right]  $. Consider this
$i$. Furthermore,%
\[
k\in Q_{0}\cap Q_{1}\subseteq Q_{0}=\left\{  n-\beta_{j}\ \mid\ j\in\left[
r\right]  \right\}  .
\]
In other words, $k=n-\beta_{j}$ for some $j\in\left[  r\right]  $. Consider
this $j$. Hence, $n-\beta_{j}=k=\omega_{i}=n+1+\alpha_{i}$ (by the definition
of $\omega_{i}$), so that $\beta_{j}=n-\left(  n+1+\alpha_{i}\right)
=-\alpha_{i}-1$ and therefore $\alpha_{i}+\beta_{j}=1$. But this contradicts
(\ref{p37.9.7.det=det.pf.not1}). This contradiction shows that our assumption
was false, qed.}. Hence,
\[
\left\vert Q_{0}\cup Q_{1}\right\vert =\underbrace{\left\vert Q_{0}\right\vert
}_{=r}+\underbrace{\left\vert Q_{1}\right\vert }_{=n-r}=r+\left(  n-r\right)
=n.
\]


Next, we can easily see that $Q_{0}\subseteq\left[  n\right]  $%
\ \ \ \ \footnote{\textit{Proof.} Let $j\in\left[  r\right]  $. We shall prove
that $n-\beta_{j}\in\left[  n\right]  $.
\par
We have $j\in\left[  r\right]  $. Thus, the cell $\left(  j,j\right)  $ lies
in the Young diagram of $\lambda$ (since the cells $\left(  1,1\right)
,\ \left(  2,2\right)  ,\ \ldots,$ $\left(  r,r\right)  $ lie in the Young
diagram of $\lambda$). In other words, the $j$-th column of $\lambda$ has at
least $j$ cells. In other words, $\lambda_{j}^{\prime}\geq j$ (since
$\lambda_{j}^{\prime}$ is the height of the $j$-th column of $\lambda$). But
the definition of $\beta_{j}$ yields $\beta_{j}=\lambda_{j}^{\prime}-j\geq0$
(since $\lambda_{j}^{\prime}\geq j$). Hence, $n-\underbrace{\beta_{j}}_{\geq
0}\leq n$.
\par
On the other hand, $\lambda_{1}^{\prime}\geq\lambda_{2}^{\prime}\geq
\lambda_{3}^{\prime}\geq\cdots$ (since $\lambda^{\prime}$ is a partition) and
thus $\lambda_{1}^{\prime}\geq\lambda_{j}^{\prime}$, so that $\lambda
_{j}^{\prime}\leq\lambda_{1}^{\prime}=\ell\left(  \lambda\right)  \leq n$.
Now, $\beta_{j}=\lambda_{j}^{\prime}-j$, so that $n-\beta_{j}=n-\left(
\lambda_{j}^{\prime}-j\right)  =n-\underbrace{\lambda_{j}^{\prime}}_{\leq
n}+\underbrace{j}_{\geq1}\geq n-n+1=1$. Combining this with $n-\beta_{j}\leq
n$, we obtain $n-\beta_{j}\in\left[  1,n\right]  =\left[  n\right]  $.
\par
Forget that we fixed $j$. We thus have shown that $n-\beta_{j}\in\left[
n\right]  $ for each $j\in\left[  r\right]  $. In other words, $\left\{
n-\beta_{j}\ \mid\ j\in\left[  r\right]  \right\}  \subseteq\left[  n\right]
$. Thus, $Q_{0}=\left\{  n-\beta_{j}\ \mid\ j\in\left[  r\right]  \right\}
\subseteq\left[  n\right]  $.} and $Q_{1}\subseteq\left[  n\right]
$\ \ \ \ \footnote{\textit{Proof.} Let $i\in\left[  r+1,n\right]  $. We shall
show that $\omega_{i}\in\left[  n\right]  $.
\par
Indeed, $i\in\left[  r+1,n\right]  $, so that $i\geq r+1>r$ and thus
$\alpha_{i}<0$ (by (\ref{p37.9.7.det=det.pf.alpha-neg})). Hence, $\alpha
_{i}\leq-1$. But the definition of $\omega_{i}$ shows that $\omega
_{i}=n+1+\underbrace{\alpha_{i}}_{\leq-1}\leq n+1+\left(  -1\right)  =n$.
\par
On the other hand, $i\in\left[  r+1,n\right]  $ shows that $i\leq n$ and thus
$\alpha_{i}=\underbrace{\lambda_{i}}_{\geq0}-\underbrace{i}_{\leq n}%
\geq0-n=-n$. Now, $\omega_{i}=n+1+\underbrace{\alpha_{i}}_{\geq-n}\geq
n+1+\left(  -n\right)  =1$. Combining this with $\omega_{i}\leq n$, we obtain
$\omega_{i}\in\left[  1,n\right]  =\left[  n\right]  $.
\par
Forget that we fixed $i$. We thus have proved that $\omega_{i}\in\left[
n\right]  $ for each $i\in\left[  r+1,n\right]  $. In other words, $\left\{
\omega_{i}\ \mid\ i\in\left[  r+1,n\right]  \right\}  \subseteq\left[
n\right]  $. Hence, $Q_{1}=\left\{  \omega_{i}\ \mid\ i\in\left[
r+1,n\right]  \right\}  \subseteq\left[  n\right]  $.}. In other words, the
sets $Q_{0}$ and $Q_{1}$ are subsets of $\left[  n\right]  $. Hence, their
union $Q_{0}\cup Q_{1}$ is a subset of $\left[  n\right]  $ as well. Since
this union has size $\left\vert Q_{0}\cup Q_{1}\right\vert =n$, we thus
conclude that this union $Q_{0}\cup Q_{1}$ must be $\left[  n\right]  $ itself
(since the only subset of $\left[  n\right]  $ that has size $n$ is $\left[
n\right]  $ itself). In other words, $Q_{0}\cup Q_{1}=\left[  n\right]  $.
Since the sets $Q_{0}$ and $Q_{1}$ are disjoint, we thus conclude that they
are each other's complement in $\left[  n\right]  $. In other words,
$Q_{1}=\left[  n\right]  \setminus Q_{0}$ and $Q_{0}=\left[  n\right]
\setminus Q_{1}$.

Note that $\omega_{r+1},\omega_{r+2},\ldots,\omega_{n}$ are elements of
$\left[  n\right]  $ (since $\left\{  \omega_{r+1},\omega_{r+2},\ldots
,\omega_{n}\right\}  =Q_{1}\subseteq\left[  n\right]  $).

Now, define the $n\times n$-matrix%
\[
C:=\left(  s_{\left(  \alpha_{i}\mid n-j\right)  }\right)  _{1\leq i,j\leq
n}=\left(  s_{\left(  \lambda_{i}-i\mid n-j\right)  }\right)  _{1\leq i,j\leq
n}%
\]
(the latter equality sign is due to the fact that $\alpha_{i}=\lambda_{i}-i$
for each $i\geq1$). Then, using (\ref{p37.9.7.hook-neq}), we easily see that
the entries of $C$ satisfy%
\[
s_{\left(  \alpha_{i}\mid n-j\right)  }=0\ \ \ \ \ \ \ \ \ \ \text{for all
}i>r\text{ and }j\neq\omega_{i}%
\]
\footnote{\textit{Proof.} Let $i,j\in\left[  n\right]  $ be such that $i>r$
and $j\neq\omega_{i}$. Then, $\alpha_{i}<0$ by
(\ref{p37.9.7.det=det.pf.alpha-neg}). That is, $\alpha_{i}$ is negative.
Moreover, $j\neq\omega_{i}=n+1+\alpha_{i}$ (by the definition of $\omega_{i}$)
and thus $\alpha_{i}\neq j-n-1=-1-\left(  n-j\right)  $. Hence,
(\ref{p37.9.7.hook-neq}) (applied to $a=\alpha_{i}$ and $b=n-j$) yields
$s_{\left(  \alpha_{i}\mid n-j\right)  }=0$, qed.}. Hence, Lemma
\ref{lem.p37.det2} (applied to $c_{i,j}=s_{\left(  \alpha_{i}\mid n-j\right)
}$) yields
\begin{equation}
\det C=\left(  -1\right)  ^{\left(  n-r\right)  n+\left(  \omega_{r+1}%
+\omega_{r+2}+\cdots+\omega_{n}\right)  }\left(  \prod_{i=r+1}^{n}s_{\left(
\alpha_{i}\mid n-\omega_{i}\right)  }\right)  \det\left(  \operatorname*{sub}%
\nolimits_{\left[  r\right]  }^{\left[  n\right]  \setminus Q_{1}}C\right)  .
\label{p37.9.7.det=det.pf.4}%
\end{equation}


However, each $i\in\left[  r+1,n\right]  $ satisfies $\alpha_{i}<0$ (by
(\ref{p37.9.7.det=det.pf.alpha-neg}), since $i\geq r+1>r$) and $\alpha
_{i}=-1-\left(  n-\omega_{i}\right)  $ (since the definition of $\omega_{i}$
yields $\omega_{i}=n+1+\alpha_{i}$, so that $\alpha_{i}=\omega_{i}%
-n-1=-1-\left(  n-\omega_{i}\right)  $) and therefore
\[
s_{\left(  \alpha_{i}\mid n-\omega_{i}\right)  }=\left(  -1\right)
^{n-\omega_{i}}\ \ \ \ \ \ \ \ \ \ \left(  \text{by (\ref{p37.9.7.hook-equal}%
), applied to }a=\alpha_{i}\text{ and }b=n-\omega_{i}\right)  .
\]
Multiplying these equalities over all $i\in\left[  r+1,n\right]  $, we obtain%
\begin{align*}
\prod_{i=r+1}^{n}s_{\left(  \alpha_{i}\mid n-\omega_{i}\right)  }  &
=\prod_{i=r+1}^{n}\left(  -1\right)  ^{n-\omega_{i}}=\left(  -1\right)
^{\left(  n-\omega_{r+1}\right)  +\left(  n-\omega_{r+2}\right)
+\cdots+\left(  n-\omega_{n}\right)  }\\
&  =\left(  -1\right)  ^{\left(  n-r\right)  n-\left(  \omega_{r+1}%
+\omega_{r+2}+\cdots+\omega_{n}\right)  },
\end{align*}
since%
\begin{align*}
&  \left(  n-\omega_{r+1}\right)  +\left(  n-\omega_{r+2}\right)
+\cdots+\left(  n-\omega_{n}\right) \\
&  =\underbrace{\underbrace{\left(  n+n+\cdots+n\right)  }_{n-r\text{ times}}%
}_{=\left(  n-r\right)  n}-\left(  \omega_{r+1}+\omega_{r+2}+\cdots+\omega
_{n}\right) \\
&  =\left(  n-r\right)  n-\left(  \omega_{r+1}+\omega_{r+2}+\cdots+\omega
_{n}\right)  .
\end{align*}
Thus,
\begin{align*}
&  \left(  -1\right)  ^{\left(  n-r\right)  n+\left(  \omega_{r+1}%
+\omega_{r+2}+\cdots+\omega_{n}\right)  }\underbrace{\left(  \prod_{i=r+1}%
^{n}s_{\left(  \alpha_{i}\mid n-\omega_{i}\right)  }\right)  }_{=\left(
-1\right)  ^{\left(  n-r\right)  n-\left(  \omega_{r+1}+\omega_{r+2}%
+\cdots+\omega_{n}\right)  }}\\
&  =\left(  -1\right)  ^{\left(  n-r\right)  n+\left(  \omega_{r+1}%
+\omega_{r+2}+\cdots+\omega_{n}\right)  }\left(  -1\right)  ^{\left(
n-r\right)  n-\left(  \omega_{r+1}+\omega_{r+2}+\cdots+\omega_{n}\right)  }\\
&  =\left(  -1\right)  ^{\left(  n-r\right)  n+\left(  \omega_{r+1}%
+\omega_{r+2}+\cdots+\omega_{n}\right)  +\left(  n-r\right)  n-\left(
\omega_{r+1}+\omega_{r+2}+\cdots+\omega_{n}\right)  }\\
&  =1
\end{align*}
(since $\left(  n-r\right)  n+\left(  \omega_{r+1}+\omega_{r+2}+\cdots
+\omega_{n}\right)  +\left(  n-r\right)  n-\left(  \omega_{r+1}+\omega
_{r+2}+\cdots+\omega_{n}\right)  =2\left(  n-r\right)  n$ is even). Hence,
(\ref{p37.9.7.det=det.pf.4}) becomes%
\begin{align}
\det C  &  =\underbrace{\left(  -1\right)  ^{\left(  n-r\right)  n+\left(
\omega_{r+1}+\omega_{r+2}+\cdots+\omega_{n}\right)  }\left(  \prod_{i=r+1}%
^{n}s_{\left(  \alpha_{i}\mid n-\omega_{i}\right)  }\right)  }_{=1}\det\left(
\operatorname*{sub}\nolimits_{\left[  r\right]  }^{\left[  n\right]  \setminus
Q_{1}}C\right) \nonumber\\
&  =\det\left(  \operatorname*{sub}\nolimits_{\left[  r\right]  }^{\left[
n\right]  \setminus Q_{1}}C\right)  . \label{p37.9.7.det=det.pf.6}%
\end{align}


Finally, since $\left[  n\right]  \setminus Q_{1}=Q_{0}=\left\{  n-\beta
_{1}<n-\beta_{2}<\cdots<n-\beta_{r}\right\}  $ and $\left[  r\right]
=\left\{  1<2<\cdots<r\right\}  $, we have%
\begin{align*}
\operatorname*{sub}\nolimits_{\left[  r\right]  }^{\left[  n\right]  \setminus
Q_{1}}C  &  =\operatorname*{sub}\nolimits_{\left\{  1<2<\cdots<r\right\}
}^{\left\{  n-\beta_{1}<n-\beta_{2}<\cdots<n-\beta_{r}\right\}  }C\\
&  =\left(  s_{\left(  \alpha_{i}\mid n-\left(  n-\beta_{j}\right)  \right)
}\right)  _{1\leq i,j\leq r}\ \ \ \ \ \ \ \ \ \ \left(  \text{since }C=\left(
s_{\left(  \alpha_{i}\mid n-j\right)  }\right)  _{1\leq i,j\leq n}\right) \\
&  =\left(  s_{\left(  \alpha_{i}\mid\beta_{j}\right)  }\right)  _{1\leq
i,j\leq r}\ \ \ \ \ \ \ \ \ \ \left(  \text{since }n-\left(  n-\beta
_{j}\right)  =\beta_{j}\text{ for each }j\right)  .
\end{align*}
Hence, we can rewrite (\ref{p37.9.7.det=det.pf.6}) as%
\[
\det C=\det\left(  s_{\left(  \alpha_{i}\mid\beta_{j}\right)  }\right)
_{1\leq i,j\leq r}.
\]
In view of $C=\left(  s_{\left(  \lambda_{i}-i\mid n-j\right)  }\right)
_{1\leq i,j\leq n}=\left(  s_{\left(  \lambda_{i}-i\mid n-k\right)  }\right)
_{1\leq i,k\leq n}$, we can rewrite this as%
\[
\det\left(  s_{\left(  \lambda_{i}-i\mid n-k\right)  }\right)  _{1\leq i,k\leq
n}=\det\left(  s_{\left(  \alpha_{i}\mid\beta_{j}\right)  }\right)  _{1\leq
i,j\leq r}.
\]
Thus, (\ref{p37.9.7.det=det}) is proven.
\end{proof}

\item \textbf{page 37:} Two lines above (9.8), a comma is missing in front of
\textquotedblleft$\varepsilon_{s}$\textquotedblright\ in \textquotedblleft let
$\mu=\left(  \gamma_{1},\ldots,\gamma_{s}\mid\varepsilon_{1},\ldots
\varepsilon_{s}\right)  $\textquotedblright.

\item \textbf{page 38:} In \textquotedblleft where $s_{\varphi_{ij}}=\left(
0\mid\beta_{j}\right)  -\left(  0\mid\varepsilon_{i}\right)  $%
\textquotedblright, replace \textquotedblleft$s_{\varphi_{ij}}$%
\textquotedblright\ by \textquotedblleft$\varphi_{ij}$\textquotedblright.
\end{enumerate}

\begin{thebibliography}{99999999}                                                                                         %


\bibitem[Friede20]{Friede20}%
\href{https://doi.org/10.4310/joc.2020.v11.n2.a7}{Gabriel Frieden, \textit{A
ratio of alternants formula for loop Schur functions}, Journal of
Combinatorics \textbf{11}(2) (2020), pp. 359--376.}

\bibitem[Goss98]{Goss98}\href{https://doi.org/10.1007/978-3-642-61480-4}{David
Goss, \textit{Basic Structures of Function Field Arithmetic}, Springer 1998.}

\bibitem[Grinbe16]{ppoly}Darij Grinberg, \textit{On }$p$\textit{-polynomials
and $\mathbb{F}_{p}$-vector subspaces of fields}, 9 December 2016.\newline\url{https://www.cip.ifi.lmu.de/~grinberg/algebra/ppoly-prob.pdf}

\bibitem[Grinbe18]{regpol}%
\href{https://www.cip.ifi.lmu.de/~grinberg/algebra/regpol.pdf}{Darij Grinberg,
\textit{Regular elements of a ring, monic polynomials and \textquotedblleft
lcm-coprimality\textquotedblright}, 19 September 2025}.

\bibitem[Grinbe19]{petriesym}\href{https://arxiv.org/abs/2004.11194v3}{Darij
Grinberg, \textit{Petrie symmetric functions}, arXiv:2004.11194v3}.

\bibitem[Grinbe20]{detnotes}\href{https://arxiv.org/abs/2008.09862v3}{Darij
Grinberg, \textit{Notes on the combinatorial fundamentals of algebra},
arXiv:2008.09862v3}.

\bibitem[Lam10]{Lam10}\href{https://arxiv.org/abs/1012.1262v2}{Thomas Lam,
\textit{Loop symmetric functions and factorizing matrix polynomials},
arXiv:1012.1262v2.}

\bibitem[LamPyl12]{LamPyl12}%
\href{https://doi.org/10.1016/j.aim.2012.03.012}{Thomas Lam and Pavlo
Pylyavskyy, \textit{Total positivity in loop groups, I: Whirls and curls},
Advances in Mathematics \textbf{230}(3) (2012), pp. 1222--1271.}

\bibitem[LLPT95]{LLPT95}D. Laksov, A. Lascoux, P. Pragacz, and A. Thorup,
\textit{The LLPT Notes}, edited by A. Thorup, updated 2018,\newline%
\url{https://noter.math.ku.dk/sympol.pdf} .

\bibitem[Macdon95]{Macdon95}Ian Grant Macdonald, \textit{Symmetric Functions
and Hall Polynomials}, 2nd edition, Oxford Science Publications 1995.

\bibitem[McDowe23]{McDowe23}\href{https://doi.org/10.37236/11200}{Eoghan
McDowell, \textit{Flagged Schur Polynomial Duality via a Lattice Path
Bijection}, Electronic Journal of Combinatorics \textbf{30} (2023), Issue 1,
\#P1.5.}

\bibitem[Moore96]{Moore96}%
\href{https://projecteuclid.org/journals/bulletin-of-the-american-mathematical-society/volume-2/issue-7/A-two-fold-generalization-of-Fermats-theorem/bams/1183414614.full}{Eliakim
Hastings Moore, \textit{A two-fold generalization of Fermat's theorem}, Bull.
Amer. Math. Soc. \textbf{2}(7) (1896), pp. 189--199.}

\bibitem[Ross13]{Ross13}%
\href{https://doi.org/10.1007/s10801-013-0436-y}{Dustin Ross, \textit{The loop
Murnaghan--Nakayama rule}, Journal of Algebraic Combinatorics \textbf{39}(1)
(2013), pp. 3--15.}
\end{thebibliography}


\end{document}